% concatenated from draft_jrnl_v12_revision.tex
\documentclass[journal,draftcls,onecolumn]{IEEEtran}
% \documentclass[journal]{IEEEtran}

% \documentclass[10pt, journal]{IEEEtran}
%% depending on your installation, you may wish to adjust the top margin:
%\addtolength{\topmargin}{9mm}

%%%%%%
%% Packages:
%% Some useful packages (and compatibility issues with the IEEE format)
%% are pointed out at the very end of this template source file (they are 
%% taken verbatim out of bare_conf.tex by Michael Shell).
%
% *** Do not adjust lengths that control margins, column widths, etc. ***
% *** Do not use packages that alter fonts (such as pslatex).         ***
%
\usepackage[utf8]{inputenc} 
\usepackage[T1]{fontenc}
\usepackage{url}
\usepackage{ifthen}
\usepackage{cite}
\usepackage[cmex10]{amsmath} % Use the [cmex10] option to ensure complicance
                             % with IEEE Xplore (see bare_conf.tex)
\usepackage{enumitem}

%% Please note that the amsthm package must not be loaded with
%% IEEEtran.cls because IEEEtran provides its own versions of
%% theorems. Also note that IEEEXplore does not accepts submissions
%% with hyperlinks, i.e., hyperref cannot be used.

%%added packages

\usepackage{ amsthm}
\newtheoremstyle{mystyle}
  {\topsep} % Space above
  {\topsep} % Space below
  {\normalfont} % Body font
  {\parindent} % Indent amount
  {\itshape} % Theorem head font
  {:} % Punctuation after theorem head
  { } % Space after theorem head (\newline = linebreak)
  {} % Theorem head spec (can be left empty, meaning 'normal')


\theoremstyle{mystyle}


%\theoremstyle{remark}
\newtheorem{prop}{Proposition}
\newtheorem{defn}{Definition}
\newtheorem{thm}{Theorem}
\newtheorem{lemma}{Lemma}
\newtheorem{cons}{Construction}
\newtheorem{claim}{Claim}
\newtheorem{prf}{}
\newtheorem{cor}{Corollary}
\newtheorem{remark}{Remark}
\newtheorem{example}{Example}
% \usepackage{graphicx}
% \usepackage{subcaption}
% \captionsetup[subfigure]{font=footnotesize}
% \captionsetup[figure]{justification=raggedright,singlelinecheck=false}











\usepackage{booktabs}
\usepackage{tabularx}
\usepackage{array}
\usepackage{multirow}







\usepackage{soul}









\usepackage{cite}
\usepackage{amsmath,amssymb,amsfonts}




\usepackage{dsfont}
\usepackage{algorithmic}
\usepackage{graphicx}
\usepackage{textcomp}
\usepackage{xcolor}
\usepackage{braket}
\usepackage{hyperref}
\usepackage{amsthm}
\usepackage{url}


\usepackage{amsmath}

\usepackage{graphicx}
\usepackage{subcaption}

%\usepackage[caption=false,font=footnotesize]{subfig}

\newcommand{\bfU}{\mathbf{U}}
\newcommand{\bfC}{\mathbf{C}}
\newcommand{\bfb}{\mathbf{b}}
\newcommand{\bfy}{\mathbf{y}}
\newcommand{\bfr}{\mathbf{r}}
\newcommand{\bfZ}{\mathbf{Z}}
\newcommand{\bfD}{\mathbf{D}}
\newcommand{\bfQ}{\mathbf{Q}}
\newcommand{\bfE}{\mathbf{E}}
\newcommand{\bfF}{\mathbf{F}}
\newcommand{\bfH}{\mathbf{H}}
\newcommand{\bfV}{\mathbf{V}}
\newcommand{\bfX}{\mathbf{X}}
\newcommand{\bfY}{\mathbf{Y}}
\newcommand{\bfW}{\mathbf{W}}
\newcommand{\bfA}{\mathbf{A}}
\newcommand{\bfT}{\mathbf{T}}
\newcommand{\bfG}{\mathbf{G}}
\newcommand{\bfS}{\mathbf{S}}
\newcommand{\bfg}{\mathbf{g}}
\newcommand{\bfu}{\mathbf{u}}
\newcommand{\bfm}{\mathbf{m}}
\newcommand{\bfx}{\mathbf{x}}
\newcommand{\bfw}{\mathbf{w}}
\newcommand{\bfc}{\mathbf{c}}
\newcommand{\bfB}{\mathbf{B}}
\newcommand{\bfI}{\mathbf{I}}
\newcommand{\bfz}{\mathbf{z}}
\newcommand{\bfM}{\mathbf{M}}
\newcommand{\bfR}{\mathbf{R}}
\newcommand{\bfK}{\mathbf{K}}
\newcommand{\bfL}{\mathbf{L}}
\newcommand{\bfv}{\mathbf{v}}
\newcommand{\bfP}{\mathbf{P}}


\newcommand{\bbC}{\mathbb{C}}
\newcommand{\bbZ}{\mathbb{Z}}
\newcommand{\bbR}{\mathbb{R}}
\newcommand{\bbF}{\mathbb{F}}
\newcommand{\cF}{\mathcal{F}}
\newcommand{\calD}{\mathcal{D}}

\DeclareMathOperator*{\argmin}{arg\,min}



\newcommand{\floor}[1]{\left\lfloor #1\right\rfloor}




\usepackage{algorithm}
%\usepackage{algorithmic}

% Customize the keywords
\renewcommand{\algorithmicrequire}{\textbf{Input:}}
\renewcommand{\algorithmicensure}{\textbf{Output:}}



\newcommand{\aditya}[1]{\marginpar{+}{\bf Aditya's remark}: {\em #1}}
\newcommand{\sifat}[1]{\marginpar{+}{\bf Sifat's remark}: {\em #1}}


%% added packages end

\interdisplaylinepenalty=2500 % As explained in bare_conf.tex


%%%%%%
% correct bad hyphenation here
\hyphenation{op-tical net-works semi-conduc-tor}

\pagestyle{plain}
% ------------------------------------------------------------

\title{Communication-Efficient Approximate Gradient Coding} 

% %%% Single author, or several authors with same affiliation:
% \author{%
%  \IEEEauthorblockN{Author 1 and Author 2}
% \IEEEauthorblockA{Department of Statistics and Data Science\\
%                    University 1\\
 %                   City 1\\
  %                  Email: author1@university1.edu}% }


%%% Several authors with up to three affiliations:
% \author{%
%   \IEEEauthorblockN{Sifat Munim and Aditya Ramamoorthy}
%   \IEEEauthorblockA{Department of Electrical and Computer Engineering \\
%                     Iowa State University\\
%                     Ames, IA 50011, U.S.A\\
%                     Email: \{smunim, adityar\}@iastate.edu
%                     }
% }

\author{Sifat Munim and Aditya Ramamoorthy%
\thanks{This paper was presented in part at the IEEE International Symposium on Information Theory (ISIT), Ann Arbor, MI, USA, 2025. This work was supported in part by the National Science Foundation (NSF) under grants CCF-2523473 and CCSS-2503640.\\
The authors are with the Department of Electrical and Computer Engineering, Iowa State University, Ames, IA 50011 USA (e-mail: \{smunim, adityar\}@iastate.edu).}
}


\begin{document}
\maketitle
\begin{abstract}
%In distributed large scale ML training, the problem is typically one of minimizing the empirical risk, a loss function $L$ that depends on the training dataset with respect to a $d$-length parameter vector. Workers are assigned portions of the original dataset and at each iteration of the algorithm, they compute gradients on the data points assigned to them. These (partial) gradients are sent to a parameter server (PS) which aggregates them to compute either an exact or approximate version of $\nabla L$ (gradient of $L$); the PS generates a new parameter vector and the iterations continue thereafter. Within distributed training, in addition to worker computation, communication is well-recognized to be a significant part of the overall job execution time. Unfortunately, in many large scale clusters (especially in the cloud), many workers are slower than their promised speed or even failure-prone. A 
Large-scale distributed learning aims at minimizing a loss function $L$ that depends on a training dataset with respect to a $d$-length parameter vector. The distributed cluster typically consists of a parameter server (PS) and multiple workers. Gradient coding is a technique that makes the learning process resilient to straggling workers. It introduces redundancy within the assignment of data points to the workers and uses coding theoretic ideas so that the PS can recover $\nabla L$ exactly or approximately, even in the presence of stragglers. Communication-efficient gradient coding allows the workers to communicate vectors of length smaller than $d$ to the PS, thus reducing the communication time. While there have been schemes that address the exact recovery of $\nabla L$ within communication-efficient gradient coding, to the best of our knowledge the approximate variant has not been considered in a systematic manner. In this work we present constructions of communication-efficient approximate gradient coding schemes. Our schemes use structured matrices that arise from bipartite graphs, combinatorial designs and strongly regular graphs, along with randomization, and algebraic techniques and constraints. We derive analytical upper bounds on the approximation error of our schemes that are tight in certain cases. Moreover, we derive a corresponding worst-case lower bound on the approximation error of any scheme. For a large class of our methods, under reasonable probabilistic worker failure models, we show that, after appropriate rescaling, the expected value of the computed gradient equals the true gradient. This in turn allows us to establish convergence of the training to a stationary point. Numerical experiments corroborate our theoretical findings.
\end{abstract}
\begin{IEEEkeywords}
Distributed computing, gradient coding, straggler, communication efficiency, structured matrices. 
\end{IEEEkeywords}

%    To our best knowledge, systematic approaches for the construction of communication-efficient approximate gradient coding schemes have not been considered in literature thus far. 
   
% Gradient coding has been a popular technique to address the issues that arise from slow or failed workers, i.e., stragglers in a distributed learning environment. Several techniques for approximating the gradient reliably in the presence of a high number of stragglers have also been proposed. With the increasing parameter size, the communication cost between the workers and the parameter server has also become a performance metric. Several works propose communication-efficient techniques to alleviate this issue while computing the exact gradient. However, these techniques are applicable to restricted problem parameters and are susceptible to a high number of stragglers. In this work, we consider approximating the gradient in a communication-efficient setting. We propose two communication-efficient constructions for which the approximation error can be upper-bounded reliably. This allows a trade-off between the approximation error and the communication cost. We also provide a lower bound on the approximation error for a fixed number of stragglers in the communication-efficient setting. 


\section{introduction}
Large-scale distributed learning is the workhorse of modern-day machine learning (ML) algorithms. The sheer size of the data and the corresponding computational needs necessitate the usage of huge clusters for the purpose of parameter fitting in most ML problems of practical interest. Such problems are essentially a guided search over a very large space of parameters; the number of parameters can even be in the billions. Examples of such problems include the training of deep neural networks \cite{GoodBengCour16} that have made huge advances in speech and image recognition, and the training of large language models (LLMs) \cite{zhao2024surveylargelanguagemodels}.

% Large scale machine learning (ML) is at the heart of several AI-based ideas that are revolutionizing several modern technologies. For instance, speech recognition, image recognition and large language models \cite{ 9419963, 10038754, 8978543,9276486, 3641289, NEURIPS2023_949f0f8f, 3534862} all require (\sifat{These citations are not there in the .bib files}) the training of large scale ML models with billions of parameters and a huge amount of data. Such processing is inevitably performed on large distributed clusters. 

At the top level, distributed machine learning operates by partitioning the relevant dataset into data subsets. The workers are assigned a subset of the data subsets and each worker is only responsible for processing its own data subsets. A prototypical example of this scenario is the training of deep neural networks. In this case each worker is responsible for computing the gradient (with respect to the parameter at the current iteration) of the appropriately defined loss function over its assigned data subsets. These gradients are then communicated to the central parameter server (PS) that coordinates the training, i.e., it aggregates the received gradients and determines the parameter for the next iteration. The PS broadcasts the new parameter and the iterations continue thereafter.

Distributed training is a key enabler of several ML technologies. Nevertheless, distributed clusters come with attendant challenges. A major issue is that large-scale clusters, especially those deployed within cloud platforms (e.g., Amazon Web Services (AWS)) are often heterogeneous in nature. These clusters often suffer from the problem of stragglers (slow or failed workers). Depending upon how the job is distributed amongst the workers, it is possible that the overall job execution time is limited by the speed of the slowest worker; this is clearly undesirable. For instance, the work of \cite{tandon_gradient} showed that stragglers can be up to $5\times$ slower than average workers on Amazon EC2. 
% \aditya{CHECK}
\begin{figure}[t!]
        \centering
        \includegraphics[scale = 1.0]{Journal/Fig_1.pdf}
        \caption{\label{fig:original_GC} {\small Worker $W_3$ has failed. $\mathbf{g}_i$ refers to the partial gradient on data subset $\calD_i$. Based on work completed by $W_1$ and $W_2$,  $\mathbf{g}_{1}+ \mathbf{g}_{2}+\mathbf{g}_{3}$ can be computed.}}
\end{figure}%
% \aditya{Fix the subscript in the figures} 

Distributing the dataset reduces the per-worker computational load. However, another relevant issue is the communication cost from the workers to the parameter server. In particular, the training process requires to-and-fro communication of high-dimensional vectors between the PS and the workers. When the number of parameters is very high, the communication cost of training can also be prohibitive.

% \aditya{DISCUSS}Large scale models often contain billions of parameters \cite{li2014scaling, 10.5555/2968826.2968829}, which in turn requires the workers to send large vectors to the parameter server, introducing a significant communication time \cite{NIPS2017_6c340f25, 10.5555/3618408.3619905}. The communication time associated with this can nullify the savings in computation time \cite{li2014scaling, 10.5555/2968826.2968829}.



Gradient coding, introduced in the work of \cite{tandon_gradient}, addresses worker slowdowns/failures by introducing redundancy within the assignment of data points to the workers (see Fig. \ref{fig:original_GC} where each data subset is replicated twice in the cluster). The workers transmit linear combinations of the gradients that are computed by them to the PS. As shown in Fig. \ref{fig:original_GC} an exact gradient coding solution allows the PS to precisely recover $\nabla L$ even in the presence of worker failures. It is straightforward to verify that if we insist on exact gradient recovery in the presence of any $s$ worker failures, then each data subset must be replicated at least $(s+1)$ times across the cluster. Thus, exact gradient recovery can be expensive. Moreover, in many parameter training scenarios, the exact gradient may not be needed for the convergence to the appropriate set of parameters, e.g., when the data distribution across the workers is i.i.d.

Accordingly, the approximate gradient coding variant \cite{dimakis_cyclic_mds} considers a setting where the gradient is recovered only approximately. The quality of the gradient is measured by its distance from the true gradient. The approximate gradient with rigorous guarantees on the gradient quality can typically be recovered even if the number of worker failures is much higher.

Communication-efficient gradient coding \cite{YeA18} considers an interesting point in the underlying design space. Specifically, it allows for trading off the redundancy in the data subset assignment for protecting against worker failures and also for communicating shorter length vectors from the workers to the PS. An example of communication-efficient gradient coding is illustrated in Fig. \ref{fig:comm_eff_GC}. However, we point out that most known constructions of communication-efficient gradient coding only consider the case of exact gradient recovery. We elaborate on these different variants of the gradient coding problem more formally in Section \ref{sec:bground}.


\begin{figure}[t!]        
        \centering
        \includegraphics[scale = 1.0]{Journal/Fig_2.pdf}
        \caption{\label{fig:comm_eff_GC} {\small Gradient vectors are split into half \cite{YeA18}, i.e., $\mathbf{g}_{i} = [\mathbf{g}_{i}[0]^T~ \mathbf{g}_{i}[1]^T]^T$.  If all workers return their results, i.e., if there are no failures and worker $W_i$ sends $\tilde{\mathbf{g}}_i$ (labeled on the outgoing edge from $W_i$), the PS can calculate $\sum_{i=1}^3 \mathbf{g}_{i}[0] = \tilde{\mathbf{g}}_1 + \tilde{\mathbf{g}}_2$, and $\sum_{i=1}^3 \mathbf{g}_{i}[1] = \tilde{\mathbf{g}}_2 + \tilde{\mathbf{g}}_3$, i.e., the exact gradient. }}
\end{figure}
% \sifat{Fig. 2 caption needs fixing}
    %\caption{{\small System with $N=m=3$ and $\delta = \gamma = 2$. Workers process the chunks in a top-to-bottom order. Here and in subsequent figures, red means that the worker did not process a task, whereas green means that it did. }}
%\vs\vs    
%\end{figure}

% The widespread usage of such clusters presents novel challenges, where ideas from coding theory have recently made significant contributions. 


% In machine learning, being able to train large models with large datasets can dramatically improve performance. Distributed machine learning is a way to facilitate the training of large models \cite{10.5555/2999134.2999271, 10.1145/2488388.2488393}. In this case, a dataset is divided among serval worker machines that collectively perform the task of training the model. Each worker computes the gradient of the data subset it is assigned, and then sends it back to the parameter server which aggregates the results from the workers to obtain the gradient of the whole dataset.


% \cite{10.5555/3294771.3294915, NIPS2017_6c340f25, aji-heafield-2017-sparse} introduced gradient sparsification and gradient quantization to address the issue of communication cost. \cite{10.1109/ALLERTON.2015.7447112, 7541478, tandon_gradient} introduced coding theory based techniques in distributed machine learning. In particular, \cite{tandon_gradient} proposes gradient coding to address the effect of stragglers, where each worker sends a linear combination of the gradients of the assigned data subsets and the parameter server linearly combines the results from the non-straggling workers. \cite{pmlr-v80-ye18a, DBLP:journals/corr/abs-2005-07184} proposed techniques that are both communication efficient and straggler resilient. However, these methods need high redundancy across workers for the parameter server to compute exact gradient. This motivated researchers to consider approximate gradient coding, where the parameter server computes an approximate of the gradient, for any number of straggling workers \cite{dimakis_cyclic_mds, kadheKR19, soft_BIBD_GC22,  charles2017approximate, glasgowW21}. However, these approaches do not consider the issue of communication cost. 



In this work, we propose approximate gradient coding techniques that are also communication-efficient. For our techniques, we provide upper bounds on the approximation error for a given communication reduction factor. Moreover, we provide a lower bound on the approximation error for any scheme with a fixed number of stragglers and communication reduction factor. Furthermore, under appropriate worker failure models, we establish a stationary-point convergence guarantee for the gradient descent algorithm. Our results are supported by numerical experiments. Table~\ref{tab:main_results} summarizes the proposed constructions and the main theoretical results.


% Added table



% Table~\ref{tab:main_results} summarizes the proposed constructions and
% their main theoretical guarantees.



\begin{table*}[t]
\centering
\caption{Summary of the proposed constructions and main theoretical results.}
\label{tab:main_results}
\renewcommand{\arraystretch}{1.25}
\setlength{\tabcolsep}{4pt}

\begin{tabularx}{\textwidth}{
    >{\raggedright\arraybackslash}p{3.0cm}
    >{\raggedright\arraybackslash}X
    >{\centering\arraybackslash}p{2.3cm}
    >{\centering\arraybackslash}p{2.1cm}
    >{\centering\arraybackslash}p{2.2cm}
}
\toprule
\textbf{Construction/result}
&
\textbf{Underlying structure or scope}
&
\textbf{Exact recovery when no workers straggle}
&
\textbf{Convergence guarantee}
&
\textbf{Analytical error guarantee}
\\
\midrule

\multicolumn{5}{l}{\textit{Proposed constructions}}
\\
\midrule

\multirow{3}{*}{%
\begin{tabular}[c]{@{}l@{}}
Post multiplication\\ by diagonal
matrices
\end{tabular}}
&
Balanced incomplete block design
&
No
&
Yes
&
Upper bound
\\
\cline{2-5}

&
Strongly regular graph
&
No
&
Yes, for the vertex-transitive case
&
Upper bound
\\
\cline{2-5}

&
Coset bipartite graph
&
Yes
&
Yes
&
Upper bound
\\
\midrule

Hadamard product with null-space constraints
&
Bi-regular bipartite graph or balanced incomplete block design
&
Yes
&
No
&
Upper bound
\\
\midrule

\multicolumn{5}{l}{\textit{General result}}
\\
\midrule

Worst-case lower bound
&
All communication-efficient approximate gradient coding schemes
&
---
&
---
&
Lower bound
\\
\bottomrule
\end{tabularx}
\end{table*}
The remainder of this paper is organized as follows. In Section \ref{sec:bground}, we present the relevant background on gradient coding, discuss related work, and summarize our main contributions. Section \ref{sec:prob_form} introduces the communication-efficient approximate gradient coding model and defines the approximation error metric. Section \ref{sec:structured_matrices} presents the structured assignment matrices used throughout the paper. Section \ref{sec:construction_RDM} presents our first method based on post-multiplication by diagonal matrices and derives corresponding upper bounds on the approximation error for several structured assignment matrices. Section \ref{sec:construction_RHPNSC} develops a second method based on Hadamard products with null-space constraints, which achieves exact recovery in the absence of stragglers and admits tractable error upper bounds. Section \ref{sec:lowerbound_proof} establishes a lower bound on the approximation error that applies to any communication-efficient approximate gradient coding scheme. Section \ref{sec:convergence} analyzes convergence of gradient descent under our proposed schemes and suitable straggler models. Finally, Section \ref{sec:numerical_exp} provides numerical experiments validating our theoretical bounds and comparing our constructions to baseline approaches.

% When the number of stragglers is high, exact computation requires a large computation load. In stochastic gradient descent (SGD), an unbiased approximation of the gradient guarantees convergence \cite{9081964}. This motivated several works that consider approximate gradient coding \cite{dimakis_cyclic_mds, kadheKR19, soft_BIBD_GC22,  charles2017approximate, glasgowW21}. In particular, \cite{dimakis_cyclic_mds, glasgowW21, charles2017approximate} uses graph-based analysis to propose approximate gradient coding schemes, which achieves reduced computation load at the cost of higher approximation error.  \cite{kadheKR19, soft_BIBD_GC22} uses combinatorial design based techniques to do so. In \cite{10.1145/3366700}, a trade-off among the computation load, accuracy of approximation and the number of stragglers is proposed. In addition, the authors proposed two schemes that achieve the trade-off. The main goal in approximate gradient coding is to guarantee the quality of the gradient even in the presence of (a potentially large number of) stragglers. This analysis can often be challenging (see upcoming Remark \ref{remark:opt_vs_fixed}). In this work, we propose approximate gradient coding schemes that are communication efficient. 

\section{Background, Related Work and Summary of Contributions}
\label{sec:bground}
 %Large-scale machine learning involves the following prototypical problem.
Let $\mathcal{D} = \{ (\mathbf{x}_i , y_i )\}_{i = 1}^N$ be a dataset where $(\mathbf{x}_i, y_i)$ are feature-label pairs. Let $L(\mathbf{w}) = \frac{1}{N} \sum_{i = 1}^N l(\mathbf{x}_i, y_i; \mathbf{w})$ be a loss function. Within learning, we wish to minimize $L$ with respect to the parameter $\mathbf{w} \in \mathbb{R}^d$ by gradient descent. Here, $l$ is a function that measures the prediction error for each point. Let $\mathbf{w}^{(t)}$ be the state of the parameter at iteration $t$. The parameter is updated as
\[ \mathbf{w}^{(t+1)} = \mathbf{w}^{(t)} - \eta^{(t)}\mathbf{g}^{(t)}, \]
where $\mathbf{g}^{(t)} := \frac{1}{N} \sum_{i = 1}^N \nabla l(\mathbf{x}_i, y_i, \mathbf{w}^{(t)})$ is the gradient of the loss function $L$ and $\eta^{(t)} \geq 0$ is the learning rate at iteration $t$.

%Typically, the size of $\mathcal{D}$ is large in most modern-day ML problems. 
%Thus, the computation of $\mathbf{g}^{(t)}$ needs to be performed in a distributed manner.
The central goal is the computation of $\mathbf{g}^{(t)}$ in a distributed manner.
A distributed learning setup involves a PS and $n$ workers denoted $W_1, W_2, \dots, W_n$. The dataset $\mathcal{D}$ is divided into $k$ disjoint data subsets of equal size denoted by $\mathcal{D}_1, \mathcal{D}_2, \dots, \mathcal{D}_k$ which are distributed among the workers; each worker computes the gradients on its assigned data subsets. The PS receives the gradients calculated by the workers and aggregates them to find the overall gradient and hence the updated parameter.


Let $[a] \triangleq \{1,2, \dots, a\}$ for any $a \in \mathbb{N}$, and for a matrix $\bfM$, let $\text{supp}(\bfM)$ denote the set of indices of the nonzero entries of $\bfM$. %\sifat{Set of indices of the nonzero entries?} 
For $i \in [n]$, let the computation load at worker $W_i$ be $\delta_i$, i.e., $W_i$ is assigned $\delta_i$ data subsets. For $j \in [k]$, the number of workers to which $\mathcal{D}_j$ is assigned is called its replication factor, denoted by $\gamma_j$. Throughout our work, we assume a regular assignment, i.e., $\delta_i = \delta$ for all $i \in [n]$, and $\gamma_j = \gamma$  for all $j \in [k]$. %Let $m$ be the communication reduction factor. For each $i \in [n]$, the worker $W_i$ is assigned $\delta$ data subsets $\mathcal{D}_{i_1}, \dots ,\mathcal{D}_{i_{\delta}}$. 


%Commenting discussion below as it is not used afterwards
% It then computes the gradients of the assigned data subsets and sends  $f_i(\mathbf{g}_{i_1}, \dots,\mathbf{g}_{i_{\delta}}) \in \mathbb{R}^{\frac{d}{m}}$ to the PS, where $f_i(\mathbf{g}_{i_1}, \dots,\mathbf{g}_{i_{\delta}}): \mathbb{R}^{d\delta} \rightarrow \mathbb{R}^{\frac{d}{m}}$ is a linear function of the gradients. The PS then computes linear combinations of the results from the non-straggling workers. In case of exact gradient computation, the PS computes $\mathbf{g}$. 
% % In a communication efficient scheme, if the communication reduction factor is $m$, then $f_i \in \mathbb{R}^{\frac{d}{m}}$ for each $i \in [n]$.

%Next we discuss the encoding and decoding techniques that we use throughout this work. 

For a given scheme we define the $(n,k,\delta,\gamma)$ assignment matrix $\mathbf{A} \in \mathbb{R}^{k \times n}$ which is such that $\mathbf{A}(i,j) \ne 0$ \footnote{We use MATLAB notation at various places in this paper.} if and only if worker $W_j$ is assigned the data subset $\calD_i$. Also, note that the assignment matrix stays the same over the iterations. Moreover, our schemes for recovering $\nabla L$ will not be time dependent. Accordingly, henceforth we drop the dependence of the gradient on $t$. %Note that for a regular assignment, each column of $\mathbf{A}$ will 

Let $\mathbf{g}_ i := \frac{1}{N} \sum_{(\mathbf{x}_j, y_j) \in \calD_i} \nabla l(\mathbf{x}_j, y_j; \mathbf{w})$ so that $\nabla L = \mathbf{g} = \sum_{i = 1}^k \mathbf{g}_i$. In the basic gradient coding protocol, worker $W_j$ calculates $\mathbf{g}_{i}$ for all $i \in \text{supp}(\mathbf{A}(:,j))$ (nonzero entries in the $j$-th column of $\mathbf{A}$) and linearly combines them to obtain %$g_j = \sum_{i =1}^N A_{i,j} g_{\calD_i}.$
 \begin{align}
     \tilde{\mathbf{g}}_j &= \sum_{i=1}^k \mathbf{A}(i,j) \mathbf{g}_{i}.\label{eq:linearly_comb_grad}
 \end{align}
It subsequently transmits $\tilde{\mathbf{g}}_j$ to the PS. %We emphasize that it is possible that worker $j$ is straggling. 
The PS wants to decode $\nabla L$, so it picks a decoding vector $\mathbf{r}$ which is such that $\mathbf{r}(j) =0$ if worker $j$ is straggling (i.e., if $\tilde{\mathbf{g}}_j$ is unavailable) and calculates %$\sum_{j=1}^m r_j g_j = \sum_{i=1}^N  \left( \sum_{j=1}^m   A_{i,j} r_j \right) g_{\calD_i}$ 
\begin{align}
\sum_{j=1}^n \mathbf{r}(j) \tilde{\mathbf{g}}_j &= \sum_{i=1}^k  \left( \sum_{j=1}^n   \mathbf{A}(i,j) \mathbf{r}(j) \right) {\mathbf{g}}_{i}, \label{eq:grad_cod_eq}
\end{align}
where we emphasize that the decoding vector $\mathbf{r}$ depends on the set of non-straggling workers. We now discuss three main threads of work within the gradient coding area.

\noindent {\bf Exact gradient coding:} 
Let $\mathbf{1}_a$ and $\mathbf{0}_a$ denote the all-ones and all-zeros column vector of size $a$ respectively. In exact gradient coding, we want to ensure that $\mathbf{A}\mathbf{r} = \mathbf{1}_k$ under any choice of at most $s$ stragglers, so that the PS can obtain $\mathbf{g}$. For exact gradient coding, each $\calD_i$ needs to be replicated at least $s+1$ times across the cluster.
The exact setting was the one that was discussed in the original work \cite{tandon_gradient} and some of the initial works in the area \cite{dimakis_cyclic_mds,reedsolomon_GC18,10.5555/3327345.3327412,tieredGC20,dynamicClusteringGC23,numerically_stableGC20}. %It is especially relevant in heterogeneous data distribution settings \cite{merugu2005distributed}, where working with approximate gradients may not be sufficient, as different parts of the data contribute non-trivially to the overall loss function.

\noindent {\bf Approximate gradient coding:} There are other scenarios in which the full gradient is not required, e.g., i.i.d. data distribution or too costly to work with because of the high replication factor needed. Indeed, many ML algorithms work with mini-batch SGD \cite{bottou_optimization} where the gradients are only calculated on a random subset of the data points. %Furthermore, for regular assignments and with adversarial stragglers it can be shown that in the presence of $s$ straggling nodes, we need the replication factor $\gamma(A) \geq s+1$. 
%Requiring the recovery of the full gradient can be quite expensive from a resource standpoint since $\gamma(A) \geq s+1$, i.e., each data point (or chunk) needs to appear in at least $s+1$ workers. %Thus, within approximate GD, we aim to guarantee that $\beta$ is close to $\mathds{1}$, e.g., $||\beta - \mathds{1}||_2 $ can be bounded as a function of the number of straggling workers.
Thus, the fundamental problem within approximate gradient coding \cite{dimakis_cyclic_mds, charles2017approximate,wang2019fundamental} is to design the assignment matrix $\mathbf{A}$ such that $|| \mathbf{A} \mathbf{r}- \mathbf{1}_k||_2$ ($\ell_2$-norm) can be bounded as a function of the number of straggling workers by an appropriate choice of $\mathbf{r}$. Furthermore, one needs to understand if and under what conditions the iterative algorithm converges to a stationary point of the loss function $L(\cdot)$ in this setting. %\aditya{Should we include this as part of proposed work?}

%Within approximate gradient coding, it is natural to consider two scenarios - random or an adversarial choice of stragglers. 
%In addition, under approximate gradient updates, one needs to understand the conditions under which the iterative algorithm converges to the minimum point (local or global) of the loss function $L(\cdot)$. Note that these considerations do not exist in the exact setting, as the full gradient computation is exact and thus the convergence of the iterative algorithm has no dependence on the gradient coding algorithm. 


%To further simplify the problem, we can assume that the $\calD$ dataset is partitioned into $m$ disjoint subsets denoted $\calD_i , i = 1, \dots, m$ each of size $N/m$ \footnote{discuss partition size}. Each worker node is assigned at most $d$ such subsets. Thus as compared to the previous formulation, the assignment is more coarse-grained. It will reduce to the previous formulation if 

% \begin{remark}
% \label{remark:naive_scheme}
% We point out that a simple scheme that ignores the unavailable gradients can be put in the formalism above ({\it cf.} \eqref{eq:linearly_comb_grad} \& \eqref{eq:grad_cod_eq}), and GC ideas provide clear benefits over them \cite{dimakis_cyclic_mds}, in the sense that $||Ar - \mathds{1}||_2$ is strictly smaller.   
% \end{remark}


\noindent {\bf Communication-efficient gradient coding:} %An important aspect of gradient coding has to do with communication efficiency. 
Within distributed training, a significant time cost is associated with the transmission of the computed gradients (vectors of length-$d$) by the workers to the PS, e.g., deep learning usually operates in the highly over-parameterized regime \cite{Belkin_2019} (many more parameters than data points). In real-world experiments, the work of \cite{alistarh2017qsgd} demonstrates that with increasing number of workers, the proportion of time spent on communication actually increases within distributed deep learning. %More recently, in 2023, work on large language model (LLM) training \cite{wang2023cocktailsgd} demonstrates similar findings.

For a specified assignment of data subsets to workers, if the number of failures is small, the extra redundancy can be judiciously used to transmit shorter vectors from the workers to the PS; this is known as communication-efficient gradient coding. 
For exact gradient coding, in the regular assignment setting if the replication factor is $s + m$, then the dimension of the transmitted vectors from the workers to the PS can be lowered to $d/ m$ \cite{YeA18} (see also \cite{tayyebehM21}) while still protecting against $s$ failures, i.e., one can trade off communication for computation. In this case, we call $m$ the communication reduction factor. Thus, if we know in advance that fewer failures are expected, the system can save communication time by operating in a communication-efficient mode.

% It follows that reducing the communication cost within GC is an important problem. The intuition of communication-efficient GC is as follows. If two copies of each $\calD_i$ are available across the cluster (even after $s$ failures), then it is conceivable that two equal halves of $\mathbf{g}_{i}, i = 1, \dots, N$ can be obtained from the different workers while ensuring that each worker only transmits a vector of length $d/2$. For exact GC, in the regular assignment setting if the 
% replication factor is $s + m$, then the dimension of the transmitted vectors from the workers can be lowered to $d/ m$ \cite{YeA18,tayyebehM21}, i.e., one can trade off communication for computation. Henceforth, we refer to $m$ as the communication reduction factor. %Depending on the problem instance, this can be a huge saving on the overall job execution time. %Depending on the time cost of communication, this can be a significant saving.

%The main result of \cite{YeA18} (see Fig. \ref{fig:comm_eff_GD}) in the regular assignment setting states that if the replication factor is $s + \ell$, then the dimension of the transmitted vectors from the workers can be lowered to $d/ \ell$ (a result for general assignments appears in \cite{tayyebehM21}). 

%Within GC, it is possible to trade-off computational redundancy with communication cost (dimension of the transmitted vectors from the workers) as well \cite{YeA18,tayyebehM21}. 

%Thus, while protecting the training process against stragglers is important, it is possible to further reduce the overall training time by saving on the amount of transmission from the workers to the PS.  Once again, within the realm of communication-efficient GC, we can consider both the exact and the approximate versions of the problem. Most of the work in the literature thus far has only considered the exact version \aditya{CHECK!}




\vspace{-0.1in}
\subsection{Discussion of Related Work} %\aditya{Edit as reqd.}
Ideas from coding theory have been the topic of intense investigation within the broad area of large-scale distributed computing within the past several years (see, e.g.,  \cite{yu2020straggler,saurav_g_analytics,DBLP:journals/corr/abs-1806-00939, LiMA16, dutta2016short, dutta_et_al20,tandon_gradient,dimakis_cyclic_mds}). %In particular, there has been a significant body of work that has used approaches from coding theory for problems within distributed matrix computation \cite{yu2020straggler,saurav_g_analytics,lagrange_cdc, LiMA16, dutta2016short, dutta_et_al20}. 

Exact gradient coding was introduced in the work of \cite{tandon_gradient}, which showed that for regular assignments, the replication factor $\gamma \geq s+1$ and demonstrated constructions that met this bound. These include a scheme based on fractional repetition \cite{el2010fractional,olmezR16} that requires the number of workers $n$ to be a multiple of $(s+1)$, and a more general cyclic assignment-based scheme. Several variants of the exact gradient coding problem have been examined \cite{10.5555/3327345.3327412,tieredGC20,dynamicClusteringGC23,numerically_stableGC20,byzantine_GC23}. The work of \cite{dimakis_cyclic_mds} demonstrated the intimate link of exact gradient coding with coding theory (see also \cite{reedsolomon_GC18}) and introduced the formulation of approximate gradient coding. We note that the performance of the exact gradient coding schemes proposed in \cite{tandon_gradient} deteriorates in the approximate setting. For approximate gradient coding, \cite{dimakis_cyclic_mds} showed connections with the spectral properties of graphs and constructed schemes from expander graphs.  The work of \cite{charles2017approximate} (see also \cite{charles2018gradient}) observed that Ramanujan graphs (expanders with the largest spectral gap) only exist in restrictive settings and considered the usage of sparse random graphs instead. The work of \cite{glasgowW21} also presented graph based schemes in this setting. Connections of approximate gradient coding with block designs were considered in \cite{kadheKR19} and subsequently in \cite{soft_BIBD_GC22}. Using convex optimization-based techniques, \cite{11195507} derived an improved upper bound on the approximation error compared with \cite{dimakis_cyclic_mds}. Fundamental limits of approximate gradient coding were examined in \cite{wang2019fundamental}. Furthermore, several variants of the original gradient coding algorithm have been examined (see \cite{BitarSGC19, treeGC19, partial_recoveryGC23,10.1109/TIT.2024.3420222, krishnan2023sequential} among others).% CITE MORE??
%Krishnan et al. \cite{krishnan2023sequential} consider the usage of GC ideas across iterations  \aditya{How to compare with this?} 

Communication-efficient gradient coding was introduced in \cite{YeA18} (in the exact setting). The work of \cite{tayyebehM21} considered arbitrary assignment matrices and provided converse results. Crucially, both these works use polynomial interpolation in their solution. This leads to significant numerical instability \cite{Pan16} (see discussion in \cite{ramamoorthyMG24}) to the extent that their solution is essentially unusable for systems with twenty or more workers. This issue was considered in part in the exact setting by \cite{kadheKR20}; however, their solution works only for a restrictive assignment setting, i.e., it requires certain divisibility conditions on the problem parameters to work (similar to the fractional repetition approach of \cite{tandon_gradient}). While the work of \cite{tayyebehM21} does discuss extensions of their Lagrange Interpolation based idea to the approximate case, there are no numerical results reported in their work. %The numerical instability associated with Lagrange interpolation continues to be an issue here. 
Broadly speaking, the design of ``communication-efficient approximate gradient coding schemes'' is an open problem.

The work of \cite{ramamoorthyMG24} presented a different class of gradient coding protocols that leverage a small amount of feedback between the workers and the parameter server. This allows for more effective usage of slow (as against failed) workers while remaining numerically stable.

% \aditya{Integrate above with below}
% Several prior works addressed the issue of stragglers. \cite{tandon_gradient} introduced gradient coding, which explores the role of redundant subset assignments within the workers. The authors of \cite{tandon_gradient} considered a setting with $n$ workers, $k$ data subsets where each worker is assigned $\delta > 1$ subsets, referred to as the computation load. They showed that if the number of stragglers is $s$, then $\frac{\delta}{k} \ge \frac{s+1}{n}$ is necessary and sufficient for computing the exact gradient under any choice of $s$ stragglers. Several prior work consider the same setting \cite{dimakis_cyclic_mds, 10.1109/ISIT.2018.8437467, NEURIPS2018_feecee9f, 9088154, 9755943, 10.1109/ISIT44484.2020.9174512}.


% The work of \cite{pmlr-v80-ye18a} introduced communication-efficient gradient coding. They showed that redundancy can also be traded off for shorter transmissions from the workers to the PS. In particular, if $\frac{\delta}{k} \ge \frac{s+m}{n}$, then the workers can transmit vectors of length $d/m$ to the PS, where $d$ is the length of the gradient vector and $m$ is referred to as the communication reduction factor. The work of \cite{tayyebehM21} generalized these results. However, both schemes operate via Lagrange polynomial interpolation and are numerically unstable for twenty or more workers \cite{doi:10.1137/15M1030170}. \cite{DBLP:journals/corr/abs-2005-07184} proposed a scheme that uses MDS encoding schemes to address this issue. However, this work applies only under certain restrictive problem parameters. All these works consider the exact scenario. While \cite{tayyebehM21} discusses extensions of their approach to the approximate case, stability issues exist there as well (no simulation results are presented in their work). In a different context, a modified gradient coding protocol that leverages feedback from the workers to the PS addresses these issues in part \cite{ramamoorthyMG24}.


% When the number of stragglers is high, exact computation requires a large computation load. In stochastic gradient descent (SGD), an unbiased approximation of the gradient guarantees convergence \cite{9081964}. This motivated several works that consider approximate gradient coding \cite{dimakis_cyclic_mds, kadheKR19, soft_BIBD_GC22,  charles2017approximate, glasgowW21}. In particular, \cite{dimakis_cyclic_mds, glasgowW21, charles2017approximate} uses expander and sparse random graph-based analysis to propose approximate gradient coding schemes, which achieves reduced computation load at the cost of higher approximation error.  \cite{kadheKR19, soft_BIBD_GC22} uses combinatorial design based techniques to do so. \cite{11195507} uses convex optimization based technique to get an improved upper bound of approximation error over \cite{dimakis_cyclic_mds}. In \cite{10.1145/3366700}, a trade-off among the computation load, accuracy of approximation and the number of stragglers is proposed. In addition, the authors proposed two schemes that achieve the trade-off. The main goal in approximate gradient coding is to guarantee the quality of the gradient even in the presence of (a potentially large number of) stragglers. This analysis can often be challenging (see upcoming Remark \ref{remark:opt_vs_fixed}). 
% % In this work, we propose approximate gradient coding schemes that are communication efficient.

\subsection{Main contributions}
%\aditya{Need to revisit}
In this work, we present the first systematic approaches for the construction of communication-efficient approximate gradient coding schemes.


Our first method uses structured matrices that arise from (i) incidence matrices of combinatorial designs, (ii) adjacency matrices of strongly regular graphs, and (iii) bi-adjacency matrices of bipartite graphs. For the bipartite graph setting, we consider an algebraic construction that we refer to as coset bipartite graphs. These matrices are used to define the underlying assignment matrix. The first construction is obtained by vertically stacking $m$ copies of the assignment matrix after post-multiplication by diagonal matrices with random entries. For coset bipartite graphs, we also consider a deterministic construction based on Hadamard matrices. Our second method uses the Hadamard product with appropriate vectors that additionally satisfy certain null-space conditions.
% Our first method uses structured matrices that arise from (i) incidence matrices of combinatorial designs, (ii) adjacency matrices of strongly regular graphs, and (iii) bi-adjacency matrices of bipartite graphs. Here, we consider random bipartite graphs and an algebraic construction that we call coset bipartite graphs. These matrices are used as a base for the assignment matrix. The schemes are obtained by vertically stacking $m$ of these matrices and post-multiplying them by diagonal matrices. We mostly use random diagonal matrices but also consider deterministic variants stemming from Hadamard matrices and orthogonal matrices as well. Our second scheme uses the Hadamard product with appropriate vectors that in addition have to satisfy certain null-space conditions. 
Both methods allow the workers to transmit gradients of length $d/m$. The PS solves a least-squares problem to decode the approximate gradient from the non-straggling workers.

It is well recognized in the literature \cite{charles2017approximate} that analyzing approximate gradient coding schemes is challenging as one needs to analyze the corresponding least-squares solution over all possible straggler sets. For our constructions, we provide analytical upper bounds on the approximation error of our schemes that are tight in certain cases. In addition, for several of our constructions, we demonstrate that the condition numbers of the recovery matrices are low; this implies numerically stable decoding of the gradient. Furthermore, we derive a corresponding worst-case lower bound on the approximation error of any scheme. We also demonstrate that under reasonable worker failure models, after appropriate rescaling, the expected value of the computed gradient equals the true gradient. Together with a bounded-variance result, this allows us to show that the expected squared gradient norm at a randomly selected iterate converges to zero. We also present numerical experiments that corroborate our theoretical findings.




% Our first scheme operates by post-multiplying the corresponding matrices by random diagonal matrices. The second scheme uses the Hadamard product with appropriate vectors that in addition have to satisfy certain null-space conditions. Moreover, we derive a corresponding lower bound on the approximation error of any scheme. In several cases, we are able to provide analytical upper bounds on the approximation error of our schemes. Furthermore, we provide proof of convergence for some specific constructions. Numerical experiments demonstrate that our schemes have low approximation error. Also, we verify convergence of schemes with fully connected neural networks. 


\section{Problem Formulation}
\label{sec:prob_form}
 
As we consider the communication-efficient variant in this work, we express it in terms of the formalism developed thus far.  We consider a scenario where the assignment of subsets to workers is fixed. Depending on the operating conditions, the cluster can decide to operate in a communication-efficient mode with communication reduction factor $m$. Thus, the communication-efficient schemes we consider in this work operate with a given assignment of subsets to workers. In Table~\ref{tab:notation}, we summarize the notations used in this paper. 

\begin{table}[!h]
\centering
\caption{Summary of the main notation.}
\label{tab:notation}
\renewcommand{\arraystretch}{1.12}
\begin{tabular}{@{}c p{0.73\columnwidth}@{}}
\hline
\textbf{Symbol} & \textbf{Description} \\
\hline
$n$ & Number of workers. \\

$k$ & Number of data subsets. \\

$s$ & Number of straggling workers. \\

$m$ & Communication reduction factor; each worker transmits a vector of length $d/m$. \\

$\gamma$ & Replication factor of each data subset. \\

$\delta$ & Computation load at each worker. \\

$\mathcal{F}$ & Index set of non-straggling workers, where $|\mathcal{F}|=n-s$. \\

$d$ & Dimension of the gradient vector. \\


$\mathbf{A}\in\mathbb{R}^{k\times n}$ &
Assignment matrix specifying the data subsets assigned to each worker. \\

$\mathbf{B}\in\mathbb{R}^{mk\times n}$ &
Encoding matrix used to generate the worker transmissions. \\

$\mathbf{D}_i\in\mathbb{R}^{n\times n}$ &
Random diagonal matrix used in the construction in~\eqref{eq:effAD1AD2}. \\

$\mathbf{g}_i\in\mathbb{R}^{d}$ &
Gradient corresponding to data subset $\mathcal{D}_i$. \\
$\mathbf{Z}\in\mathbb{R}^{\frac{d}{m}\times mk}$ &
Gradient block matrix formed from $\mathbf{g}_1,\ldots,\mathbf{g}_k$. \\

$\mathbf{R}\in\mathbb{R}^{n\times m}$ &
Decoding matrix used by the PS; its nonzero rows correspond to workers in $\mathcal{F}$. \\

$\mathbf{F}\in\mathbb{R}^{mk\times m}$ &
Target matrix corresponding to the $m$ blocks of the exact gradient. \\
\hline
\end{tabular}
\end{table}



Let $\bfA \in \mathbb{R}^{k \times n}$ be a $(n,k,\delta, \gamma)$ assignment matrix. For designing a scheme with communication reduction factor $m$, we define the encoding matrix denoted by $\mathbf{B} \in \mathbb{R}^{mk \times n}$ as follows.
\begin{align*}
\mathbf{B}^T = [\bfA_1^T ~|~ \dots ~|~ \bfA_m^T],
%\begin{bmatrix} \mathbf{A}_1 \\ \vdots \\ \mathbf{A}_m\end{bmatrix},
\end{align*}
where $\text{supp}(\bfA_i) \subseteq  \text{supp}(\bfA)$ for $i \in [m]$. This constraint ensures that the communication-efficient scheme operates with the same assignment of subsets to workers.

For $i \in [k]$, let $\bfg_i$ be partitioned into $m$ equal sized blocks (with zero-padding if required) so that $\bfg_i^T = [\bfg_i[1]^T ~|~ \dots ~|~\bfg_i[m]^T]$. Also, for $u \in [m]$, let $\mathbf{G}[u] \in \mathbb{R}^{\frac{d}{m} \times k}$ be such that $\mathbf{G}[u] := [\mathbf{g}_1[u] ~|~\mathbf{g}_2[u] ~|~ \dots ~|~\mathbf{g}_k[u]].$
% \[ \mathbf{G}[u] := \begin{bmatrix} \mathbf{g}_1[u] &\mathbf{g}_2[u] & \dots \mathbf{g}_k[u]
% \end{bmatrix}.\] 
Now define a matrix $\mathbf{Z} \in \mathbb{R}^{\frac{d}{m} \times mk}$ as
%\[\mathbf{Z} := \begin{bmatrix} \mathbf{G}[1] & \mathbf{G}[2] &\dots & \mathbf{G}[m]  \end{bmatrix}.\]
\[\mathbf{Z} := [ \mathbf{G}[1] ~|~ \mathbf{G}[2] ~|~\dots~|~ \mathbf{G}[m] ].\]


% Let $\bfZ \in \mathbb{R}^{\frac{d}{m} \times mk}$ be a matrix whose $i^{th}$ row $\mathbf{Z}(i,:)$ is defined as:

% \[ \mathbf{Z}(i, :) :=  \begin{bmatrix} \mathbf{z}_i^{(1)} & \mathbf{z}_i^{(2)} & \cdots & \mathbf{z}_i^{(m)} \end{bmatrix} \] 
% Where 
% $\mathbf{z}_i^{(u)} := \begin{bmatrix} \mathbf{g}_1((i-1)m +u) \dots \mathbf{g}_{k}((i-1)m +u) \end{bmatrix} \in \mathbb{R}^k$ for $u \in [m],i \in [\frac{d}{m}]$.
%Worker $W_i$ then sends $f_i(\mathbf{g}_{i_1}, \dots,\mathbf{g}_{i_{\delta}}) = \mathbf{Z}\mathbf{B}(:,i) \in \mathbb{R}^{\frac{d}{m}}$. 
Worker $W_i$ then sends $\mathbf{Z}\mathbf{B}(:,i) \in \mathbb{R}^{\frac{d}{m}}$, once it has finished processing all its assigned subsets. Otherwise, it does not send anything. For $i \in [m]$, let $\mathbf{f}_i^T := [\underbrace{\mathbf{0}_k^T ~|~ \dots ~|~ \mathbf{0}_k^T}_{(i-1) \text{blocks}} ~|~ \mathbf{1}_k^T ~|~\underbrace{\mathbf{0}_k^T~|~ \dots ~|~ \mathbf{0}_k^T}_{(m-i) \text{blocks}}]$ of length $mk$. Next, let 
$\bfF := [\mathbf{f}_1 ~|~ \dots ~|~ \mathbf{f}_m]$. It can be observed that the exact gradient can be obtained from $\mathbf{ZF}$, since
\begin{align*}
%\mathbf{ZF} &= \begin{bmatrix} \sum_{i = 1}^k \mathbf{g}_i[1] & \sum_{i = 1}^k \mathbf{g}_i[2] & \dots & \sum_{i = 1}^k \mathbf{g}_i[m] \end{bmatrix}.
\mathbf{ZF} &= [ \sum_{i = 1}^k \mathbf{g}_i[1] ~|~ \sum_{i = 1}^k \mathbf{g}_i[2] ~|~ \dots ~|~ \sum_{i = 1}^k \mathbf{g}_i[m]].
 \end{align*}

% Consider $\mathbf{F}  \in \mathbb{R}^{mk \times m}$, where 
% \begin{equation}\label{eq:F}
% \bfF :=
% \begin{bmatrix}
% \mathbf{1}_{k} & \mathbf{0}_{k} & \cdots & \mathbf{0}_{k} \\
% \mathbf{0}_{k} & \mathbf{1}_{k} & \cdots & \mathbf{0}_{k} \\
% \vdots & \vdots & \ddots & \vdots \\
% \mathbf{0}_{k}  &\mathbf{0}_{k}  & \cdots & \mathbf{1}_{k}
% \end{bmatrix}.
% \end{equation}
% The exact gradient can be obtained from $\mathbf{ZF}$, since  

% \begin{align*}
% % \mathbf{ZF} &= \begin{bmatrix} \mathbf{G}[1] & \mathbf{G}[2] &\dots & \mathbf{G}[m]  \end{bmatrix}\begin{bmatrix}
% % \mathbf{1}_{k} & \mathbf{0}_{k} & \cdots & \mathbf{0}_{k} \\
% % \mathbf{0}_{k} & \mathbf{1}_{k} & \cdots & \mathbf{0}_{k} \\
% % \vdots & \vdots & \ddots & \vdots \\
% % \mathbf{0}_{k}  &\mathbf{0}_{k}  & \cdots & \mathbf{1}_{k}
% % \end{bmatrix}\\
% \mathbf{ZF} &= \begin{bmatrix} \mathbf{G}[1]\mathbf{1}_k & \mathbf{G}[2]\mathbf{1}_k & \dots & \mathbf{G}[m]\mathbf{1}_k \end{bmatrix} \\
% &= \begin{bmatrix} \sum_{i = 1}^k \mathbf{g}_i[1] & \sum_{i = 1}^k \mathbf{g}_i[2] & \dots & \sum_{i = 1}^k \mathbf{g}_i[m] \end{bmatrix}.
%  \end{align*}

% \begin{align*}
% \mathbf{ZF} &= \begin{bmatrix} \mathbf{z}_1^{(1)}\mathbf{1}_k & \mathbf{z}_1^{(2)}\mathbf{1}_k & \dots & \mathbf{z}_1^{(m)}\mathbf{1}_k\\ 
% \vdots & \vdots & \ddots & \vdots \\
% \mathbf{z}_{\frac{d}{m}}^{(1)}\mathbf{1}_k & \mathbf{z}_{\frac{d}{m}}^{(2)}\mathbf{1}_k & \dots & \mathbf{z}_{\frac{d}{m}}^{(m)}\mathbf{1}_k
% \end{bmatrix}\\
% & = \begin{bmatrix}\bfg(1) & \bfg(2) & \dots & \bfg(m)\\
% \vdots & \vdots & \ddots & \vdots \\
% \bfg(d-m+ 1) & \bfg(d - m + 2) & \dots & \bfg(d)
% \end{bmatrix}
% \end{align*}
% The second equality holds since for $u \in [m]$ and $i \in [\frac{d}{m}]$,\[ 
%  \mathbf{z}_i^{(u)}\mathbf{1}_k = \sum_{ j =1}^k \bfg_j((i-1)m + u) = \bfg((i-1)m + u) \] 

Suppose that there are $s$ stragglers, and let $\mathcal{F} \subseteq [n]$ be a set such that $i \in \mathcal{F}$ if and only if $W_i$ is not a straggler. Let $\mathbf{R} \in \mathbb{R}^{n \times m}$ be the decoding matrix such that $\mathbf{R}(:,i) = \mathbf{r}_i$ and $\text{supp}(\mathbf{r}_i) \subseteq \mathcal{F}$ for each $i \in [m]$. The PS then computes  $\mathbf{ZBR}$ and it follows that if $\mathbf{BR} = \mathbf{F}$, the PS computes the gradient exactly. For a matrix $\bfM$, let $\|\bfM\|_F$ and $\|\bfM\|_2$ be the Frobenius norm and the spectral norm of $\bfM$, respectively. %For a vector $\mathbf{v}$, let $\|\mathbf{v}\|_2$ be the $l_2$ norm of $\mathbf{v}$. Let $\mathbf{f}_i = \mathbf{F}(:,i)$ and $\mathbf{r}_i = \mathbf{R}(:,i)$ for each $i \in [m]$. 
It is not too hard to see that
\begin{align*}
\|\mathbf{ZBR} - \mathbf{ZF}\|_F^2 
& \leq \|\mathbf{Z}\|_2^2 \|\mathbf{BR-F}\|_F^2.
\end{align*}

% We have, 
% \begin{align*}
% \|\mathbf{ZBR} - \mathbf{ZF}\|_F^2 &= \|\mathbf{ZB}\mathbf{r}_1 - \mathbf{Z}\mathbf{f}_1\|_2^2 + \dots + \|\mathbf{ZB}\mathbf{r}_m - \mathbf{Z}\mathbf{f}_m\|_2^2\\
% & \le \|\mathbf{Z}\|_2^2 (\|\mathbf{B}\mathbf{r}_1 - \mathbf{f}_1\|_2^2 + \dots + \|\mathbf{B}\mathbf{r}_m - \mathbf{f}_m\|_2^2)\\
% & = \|\mathbf{Z}\|_2^2 \|\mathbf{BR-F}\|_F^2
% \end{align*}

%So, the Frobenius norm of the difference between the computed gradient $\mathbf{ZBR}$ and the exact gradient $\mathbf{ZF}$ is upper bounded by $\|\mathbf{BR-F}\|_F$. 
Note that $\|\mathbf{Z}\|_2^2$ depends on the actual gradients and we do not have any control over them.
Thus, for a given set of non-stragglers corresponding to $\mathcal{F}$, we seek to minimize $\|\mathbf{BR-F}\|^2_F$. This motivates the following definition. 

\begin{defn}
For a given set of non-stragglers corresponding to $\cF \subseteq [n]$ of size $n-s$, and a given encoding matrix $\bfB$, the approximation error $\text{Err}_{\cF}(\bfB)$ is defined as
\begin{equation} \label{eq:ApproxError}
\text{Err}_{\cF}(\bfB) := \min\limits_{\substack{\mathbf{R}\in \mathbb{R}^{n \times m} \\ \text{supp}(\mathbf{R}(:,i)) \subseteq \mathcal{F},  \forall i \in [m]}}  {\|\mathbf{B} \mathbf{R}- \mathbf{F}\|_F^2}.
\end{equation}
\end{defn}

For a matrix $\mathbf{M} \in \mathbb{R}^{a \times b}$, let $\mathcal{H} \subseteq [a]$ and $\mathcal{K} \subseteq [b]$ be sets corresponding to the rows and columns of $\mathbf{M}$ respectively. We denote the submatrix of $\mathbf{M}$ with rows and columns corresponding to $\mathcal{H}$ and $\mathcal{K}$ by $\mathbf{M}(\mathcal{H}, \mathcal{K})$. Also, denote $\mathbf{M}_{\mathcal{H}} := \mathbf{M}(:, \mathcal{H})$. Now, for a given set of non-stragglers corresponding to $\mathcal{F}$, if $\mathbf{B}_{\mathcal{F}}^T\mathbf{B}_{\mathcal{F}}$ is invertible, then the RHS of \eqref{eq:ApproxError} can be expressed as
\begin{align}
% \text{Err}_{\cF}(\bfB) &= \min\limits_{\substack{\mathbf{R}\in \mathbb{R}^{n \times m} \\ \text{supp}(\mathbf{R}(:,i)) \subseteq \mathcal{F} \\ \forall i \in [m]}}  {\|\mathbf{B} \mathbf{R}- \mathbf{F}\|_F^2} \nonumber\\
%&\text{Err}_{\cF}(\bfB)\\ 
&\sum_{i =1}^m \min\limits_{\substack{\mathbf{r}_i\in \mathbb{R}^{n} \\ \text{supp}(\mathbf{r}_i) \subseteq \mathcal{F} }}   \|\bfB\mathbf{r}_i - \mathbf{f}_i\|_2^2 %\nonumber
 = \sum_{i =1}^m \min\limits_{\substack{\mathbf{r}_i\in \mathbb{R}^{n-s} }}   \|\bfB_{\mathcal{F}}\mathbf{r}_i - \mathbf{f}_i\|_2^2 \nonumber\\
& = \sum_{i =1}^m\mathbf{f}_i^T \mathbf{f}_i - \mathbf{f}_i^T \bfB_{\cF}(\bfB_{\cF}^T\bfB_{\cF})^{-1}\bfB_{\cF}^T\mathbf{f}_i.\label{eq:opt_dec_error}
\end{align}
Here, the first equality follows from observation, and the second follows from an analysis of the least-squares error \cite{ekf}. The expression in \eqref{eq:opt_dec_error} characterizes the optimal decoding error. In practice, the PS can compute the decoding matrix by solving the corresponding least-squares problem using a standard QR or SVD based solver, without explicitly forming 
$(\mathbf{B}_{\mathcal{F}}^{T}\mathbf{B}_{\mathcal{F}})^{-1}$. 
Using a QR based solver, the decoding complexity is at most 
$\mathcal{O}\bigl(mk(n-s)^2\bigr)$ \cite{golub2013matrix}. 
The numerical stability of this decoding step depends on the condition number of 
$\mathbf{B}_{\mathcal{F}}^{T}\mathbf{B}_{\mathcal{F}}$. For a positive definite matrix $\bfM$, its condition number is $\kappa(\bfM) = \frac{\lambda_{max}(\bfM)} {\lambda_{min}(\bfM)}$ \cite{horn2012matrix}. Condition number bounds on the system of equations encountered in the decoding process for several of the proposed schemes are provided in the forthcoming sections.

% Appendix \ref{sec:appendix_cond_number} provides explicit condition number bounds for some of our proposed constructions. 

The following example illustrates the encoding, worker transmissions, and least-squares recovery for a small communication-efficient scheme.
%by setting the gradient of $\|\bfB_{\mathcal{F}}\mathbf{r}_i - \mathbf{f}_i\|_2^2$ with respect to $\mathbf{r}_i$ to zero and substituting in the solution for $\mathbf{r}_i$ to $\|\bfB_{\mathcal{F}}\mathbf{r}_i - \mathbf{f}_i\|_2^2$, for each $i \in [m]$.\\ 
% \section{Main Contributions}
% Several previous works have proposed approximate gradient coding schemes for the case $m = 1$ {CITE}. 




%  added example 
\begin{example}
Consider $k=2$ data subsets and $n=4$ workers with
\[
\mathbf{A}=
\begin{bmatrix}
1&0&1&0\\
0&1&0&1
\end{bmatrix},
\qquad m=2.
\]
Thus, $W_1,W_3$ are assigned $\mathcal{D}_1$, while $W_2,W_4$ are
assigned $\mathcal{D}_2$. Partition
$\mathbf{g}_i^T=[\mathbf{g}_i[1]^T\mid\mathbf{g}_i[2]^T]$, $i\in[2]$,
so that
\[
\mathbf{Z}
=
[\mathbf{g}_1[1]\mid\mathbf{g}_2[1]\mid
 \mathbf{g}_1[2]\mid\mathbf{g}_2[2]].
\]
Choose
\[
\mathbf{A}_1=\mathbf{A},
\qquad
\mathbf{A}_2=
\begin{bmatrix}
1&0&-1&0\\
0&1&0&-1
\end{bmatrix}.
\]
Since $\operatorname{supp}(\mathbf{A}_1)
=\operatorname{supp}(\mathbf{A}_2)
=\operatorname{supp}(\mathbf{A})$, the encoding matrix is
\[
\mathbf{B}
=
\begin{bmatrix}
\mathbf{A}_1\\
\mathbf{A}_2
\end{bmatrix}
=
\begin{bmatrix}
1&0&1&0\\
0&1&0&1\\
1&0&-1&0\\
0&1&0&-1
\end{bmatrix}.
\]
The workers transmit, respectively,
\[
\begin{aligned}
W_1&:\mathbf{g}_1[1]+\mathbf{g}_1[2],&
W_2&:\mathbf{g}_2[1]+\mathbf{g}_2[2],\\
W_3&:\mathbf{g}_1[1]-\mathbf{g}_1[2],&
W_4&:\mathbf{g}_2[1]-\mathbf{g}_2[2],
\end{aligned}
\]
where each transmitted vector has length $d/2$. Suppose that $W_4$ is a straggler, so
$\mathcal{F}=\{1,2,3\}$. Then,
\[
\mathbf{F}=
\begin{bmatrix}
1&0\\
1&0\\
0&1\\
0&1
\end{bmatrix},
\qquad
\mathbf{R}^{\star}
=
\argmin_{\mathbf{R}\in\mathbb{R}^{3\times2}}
\|\mathbf{B}_{\mathcal{F}}\mathbf{R}-\mathbf{F}\|_F^2
=
\frac{1}{2}
\begin{bmatrix}
1&1\\
1&1\\
1&-1
\end{bmatrix}.
\]
Consequently, the recovered blocks are
\[
\begin{aligned}
\widehat{\mathbf{g}}[1]
&=\mathbf{g}_1[1]
+\frac{1}{2}\bigl(\mathbf{g}_2[1]+\mathbf{g}_2[2]\bigr),\\
\widehat{\mathbf{g}}[2]
&=\mathbf{g}_1[2]
+\frac{1}{2}\bigl(\mathbf{g}_2[1]+\mathbf{g}_2[2]\bigr).
\end{aligned}
\]
The PS concatenates these blocks to obtain the approximate gradient.
\end{example}



\begin{remark}
\label{remark:opt_vs_fixed}
    As pointed out in \cite{charles2017approximate}, the analysis of the approximate gradient coding error is challenging, since the error expression involves the inverse of $\bfB_{\cF}^T\bfB_{\cF}$. This quantity needs to be bounded over all possible $\cF \subset [n], |\cF| = n-s$. For the case of $m=1$, prior work \cite{dimakis_cyclic_mds} has aimed at upper bounding this error by using ``fixed-decoding'' approaches, where the decoding vector is determined in advance, rather than by solving a least-squares problem. This typically results in loose upper bounds. The work of \cite{glasgowW21} proposed schemes that use optimal decoding (again for $m=1$). In this work, our goal is not only to design schemes for $m > 1$ that have low approximation error, but also to provide techniques to analytically or numerically upper bound the least-squares error ({\it cf.} \eqref{eq:opt_dec_error}). %\aditya{need to revisit}
\end{remark}
%\textbf{Main Contributions:} In this work, we propose schemes for which the approximation error $\text{Err}_{\cF}(\bfB)$ can be upper bounded reliably for any $m \ge 1$. This allows a trade-off between the communication reduction factor and the approximation error. We also provide a lower bound on the worst case error over all possible choices of $\mathcal{F}$ for a given number of stragglers. 

%As discussed in Remark \ref{remark:opt_vs_fixed}, analysis of $\text{Err}_{\cF}(\bfB)$ over a range of $\cF$ is a non-trivial task. Accordingly, 


%To help facilitate this, we discuss types of structured matrices below.


%where each $\mathbf{A}_i, i\in [m]$ has the same indices of nonzero entries as the assignment matrix $\mathbf{A}$. Specifically, $\text{supp}(\mathbf{A}_i(:,j)) = \text{supp}(\mathbf{A}(:,j))$ for all $i \in [m], j \in [n]$, where $\text{supp}(\mathbf{v}) \subseteq [a]$ denotes the set of indices that correspond to the nonzero entries of the vector $\mathbf{v} \in \mathbb{R}^a$.  
% Let $\mathbf{B} \in \mathbb{R}^{mk \times n}$ be a matrix of  which each column corresponds to a worker and the nonzero entries of each column is determined by the assigned data subsets of the corresponding worker.
%Consequently, $\mathbf{B}(k(u-1)+i,j) \ne 0$ for all $u \in [m]$ if and only if worker $W_j$ is assigned the subset $\mathcal{D}_i$. 
% $\{\mathcal{D}_{i_1}, \dots ,\mathcal{D}_{i_{\delta}} \}$, then $\text{supp}(\mathbf{B}(:, i)) = \{ i_1, \dots, i_{\delta}, k+ i_i, \dots, k+ i_{\delta}, \dots,  k(m-1)+ i_i, \dots, km+ i_{\delta} \}$. 
% We refer to $\mathbf{B}$ as the encoding matrix. 
% How does the worker $W_i$ compute $f_i$?

%For each $i \in [k]$, we arrange $\mathbf{g}_i$ into $m$ groups, each denoted $\mathbf{g}_i[u] \in \mathbb{R}^{\frac{d}{m}}$, for $u\in [m]$ such that,
%\[\mathbf{g}_i[u] := \begin{bmatrix} \mathbf{g}_i((u-1)\frac{d}{m} + 1)\\  \vdots \\  \mathbf{g}_i((u-1)\frac{d}{m} + \frac{d}{m}) \end{bmatrix}.\]


%\section{Preliminaries}
%\label{sec:prem}



\section{Structured Assignment Matrices}

\label{sec:structured_matrices}


To facilitate the analysis of the approximate decoding error, we will focus our attention on certain classes of assignment matrices. In particular, we will need to use structured matrices that arise from different areas such as graph theory and combinatorial design theory (see \cite{godsil-royle-algebraic},\cite{Brouwer_SRG}, \cite{StinsonBook}). In this section, we briefly review the relevant notions. In what follows, let $\mathbf{I}_a$ and $\mathbf{J}_a$ be the $a \times a$ identity matrix and all ones matrix, respectively. Also, let $\mathbf{0}_{a \times a}$ be the $a \times a$ all zeros matrix. 



\begin{defn} {[{\it Balanced Incomplete Block Design (BIBD) and its Incidence Matrix}]} Let $(X, \mathcal{A})$ be a pair where $X$ is a set of elements called points and $\mathcal{A}$ is a collection of nonempty subsets of $X$ called blocks. A $(n,k, \gamma, \delta, \lambda)$- balanced incomplete block design (BIBD) is a pair $(X, \mathcal{A})$ that satisfies the following properties: {\it (i)} $X$ has size $n$, $\mathcal{A}$ has size $k$, each block in $\mathcal{A}$ has $\gamma$ points, and every point in $X$ is contained in exactly $\delta$ blocks. {\it (ii)} Any pair of distinct points is contained in exactly $\lambda$ blocks. The incidence matrix of a BIBD is a $n \times k$ binary matrix $\mathbf{E}$ such that $\mathbf{E}(i,j) = 1$ if $x_i \in A_j$ and $\mathbf{E}(i,j) = 0$ otherwise. 
     
\end{defn}

Let $\bfE$ be the incidence matrix of a $(n,k,\gamma, \delta, \lambda)$ BIBD and $\bfM = \bfE^T$. Then,  $\mathbf{M} \mathbf{1}_n = \gamma \mathbf{1}_k$ and $\mathbf{M}^T \mathbf{1}_k = \delta \mathbf{1}_n$. In addition, the inner product, $\langle \mathbf{M}(:,i), \mathbf{M}(:,j)\rangle = \lambda$, for $i \ne j$ and $i, j \in [n]$ (see  \cite{StinsonBook}). Consequently, $\bfM^T\bfM = (\delta - \lambda)\bfI_n + \lambda \mathbf{J}_n$ and, 
\begin{equation} \label{eq: bibd_prop}
\bfM_{\cF}^T\bfM_{\cF} = (\delta - \lambda)\mathbf{I}_{n-s} + \lambda \mathbf{J}_{n-s},
\end{equation}
since $|\mathcal{F}| = n-s$.

\begin{defn} {[{\it Strongly Regular Graph (SRG)}]} Let $G$ be a graph that is neither complete nor empty. Then, $G$ is a strongly regular graph (SRG) with parameters $(n, \delta, \lambda, \mu)$ if it has $n$ vertices where each vertex has degree $\delta$, any pair of adjacent vertices has $\lambda$ common neighbors, and any pair of 
% distinct
nonadjacent vertices has $\mu$ common neighbors. 
\end{defn}


If $\bfM$ is the adjacency matrix of a $(n, \delta, \lambda, \mu)$ SRG, then
$\mathbf{M} = \mathbf{M}^T$ and $\mathbf{M} \mathbf{1}_n = \delta \mathbf{1}_n$  and furthermore, $\mathbf{M}^2 = \delta \mathbf{I}_n + \lambda \mathbf{M} + \mu (\mathbf{J}_n - \mathbf{I}_n - \mathbf{M})$ (see \cite{Brouwer_SRG}). Consequently, %\aditya{Fix subscript issue here}
\begin{equation}\label{eq: srg_prop}
\mathbf{M}_{\cF}^T\mathbf{M}_{\cF}  = \delta \mathbf{I}_{n-s} + \lambda \mathbf{M}(\cF, \cF) + \mu (\mathbf{J}_{n-s} - \mathbf{I}_{n-s} - \mathbf{M}(\cF, \cF)).
\end{equation}
\begin{defn} {[{\it Bi-regular Bipartite Graph and Bi-adjacency Matrix}]}
Let $G = (L \cup R, E)$ be a graph with $L = [k]$, $R = [n]$, $L \cap R = \emptyset$, where each vertex in $L$ has degree $\gamma$ and each vertex in $R$ has degree $\delta$ such that the edges only exist between vertices in $L$ and $R$. Then, $G$ is a $(n,k,\delta, \gamma)$ bi-regular bipartite graph. The bi-adjacency matrix of a bi-regular bipartite graph is a matrix $\mathbf{M}$ such that for $i \in [k]$, $j \in [n]$, $\mathbf{M}(i,j) = 1$ if $i$ and $j$ are adjacent and $\mathbf{M}(i,j) = 0$ otherwise.   
\end{defn}

If $\bfM$ is the bi-adjacency matrix of a $(n,k,\delta, \gamma)$ bi-regular bipartite graph, then $\mathbf{M} \mathbf{1}_n = \gamma \mathbf{1}_k$ and $\mathbf{M}^T \mathbf{1}_k = \delta \mathbf{1}_n$. Next, we discuss the construction of a special class of bi-regular bipartite graphs that we refer to as a coset bipartite graph.  

\subsection{Coset Bipartite Graph}
Let $k,m,\delta$ be positive integers. A $(k, m, \delta)$ coset bipartite graph is constructed as follows. Set $n = mk$. 
We construct a bi-regular bipartite graph $G = (L \cup R, E)$ with $|L| = k$ and $|R| = n = mk$.\\
\textit{Vertex Sets:}
Let $H$ be an order $m$ subgroup of $\mathbb{Z}_{mk}$ (group operation is addition modulo $mk$) such that $H = \{0, k, 2k, \dots, (m-1)k\}$. For $i \in \mathbb{Z}_{mk}$, the coset $i+H$ of $H$ is defined as $i + H = \{ (i + a) \bmod mk \ | \ a \in H\}$. Note that there are $k$ distinct cosets of $H$.
We define the vertex sets corresponding to the distinct cosets of $H$ as $L = \{\, i+H : i=0,1,\dots,k-1 \,\}$ and  
$R = \mathbb{Z}_{mk}$. So, $|L| = k$ and $|R| = mk$.\\
\textit{Edges:} Let $B$ be a subset of the set $\{0,1, \dots, k-1\}$ of size $\delta$. Let $S = \bigcup_{b \in B} (b+H)$. For $i \in \{0,1,\dots,k-1\}$, we place an edge between $i+H \in L$ and $x \in R$ if and only if
$x \in i+S$. We refer to the set $B$ as the generating set.\\
\textit{Bi-regularity:} Each right vertex $x$ lies in exactly one $H$-coset, and since $|B|=\delta$, it is adjacent
to exactly $\delta$ left vertices. Hence, each right vertex has degree $\delta$.
Because $S$ consists of $\delta$ full cosets of $H$, each left vertex has degree
$
\delta \cdot |H|=\delta \cdot m = m\delta.
$
Thus $G$ is bi-regular where each vertex in $L$ has degree $m\delta$ and each vertex in $R$ has degree $\delta$.

If $\bfM$ is the bi-adjacency matrix of a $(k,m,\delta)$ coset bipartite graph, then with $n=mk$, $\mathbf{M} \mathbf{1}_n = m \delta \mathbf{1}_k$ and $\mathbf{M}^T \mathbf{1}_k = \delta \mathbf{1}_n$.




\subsection{Existence and Parameter Feasibility}

The structured assignment matrices discussed above are not available for
arbitrary system parameters. For a $(n,k,\gamma,\delta,\lambda)$ BIBD,
the necessary parameter relations are $n\delta=k\gamma$ and
$\delta(\gamma-1)=\lambda(n-1)$. These conditions are not
sufficient. 

Similarly, the parameters of a $(n,\delta,\lambda,\mu)$ SRG must satisfy
$(n-\delta-1)\mu=\delta(\delta-\lambda-1)$, and the multiplicities of the
two eigenvalues other than $\delta$ must be nonnegative integers. Consequently, SRGs exist only for specific parameter tuples.

In contrast, a $(k,m,\delta)$ coset bipartite graph can be constructed for
any positive integers $k$ and $m$ and any $1\leq\delta\leq k$. However, its
system parameters necessarily satisfy $n=mk$ and $\gamma=m\delta$. In the sequel we need additional assumptions: $k=p^a$, where $p$ is prime and
$a\in\mathbb{Z}_{\geq 1}$, and that $p\nmid\delta$. They are used later to establish
the almost-sure invertibility of the corresponding encoding matrix. 

In summary, the
choice of the structured assignment matrix is constrained by the desired
system parameters, and all three structures need not be simultaneously
available under a common parameter setting.




















% \begin{defn}
% \begin{defn}


% \aditya{Define BIBD, SRG in a formal way}
% \subsection{Structured Matrices} 
% \label{sec:struc_matrices}
% Throughout this work, we will need to use structured matrices that arise from different areas such as graph theory and combinatorial design theory. Let $\mathbf{M}$ denote a $k \times n$ matrix that is binary (consists only of 0's and 1's). Let $\mathbf{I}_a$ and $\mathbf{J}_a$ be the $a \times a$ identity matrix and all ones matrix, respectively.
% \begin{itemize}[leftmargin=*]
%     %\item {\it Incidence matrix of bipartite graph:} 
%     \item {\it Bi-adjacency matrix of bipartite graph:} Consider the case when $\mathbf{M} \mathbf{1}_n = \gamma \mathbf{1}_k$ and $\mathbf{M}^T \mathbf{1}_k = \delta \mathbf{1}_n$. Then, $\mathbf{M}$ corresponds to the bi-adjacency matrix of a bi-regular bipartite graph, i.e., $\mathbf{M}(i,j) = 1$ if $i$ and $j$ are connected and zero otherwise. The left and right degrees of the graph are $\gamma$ and $\delta$ respectively. We refer to this graph as a $(n, k, \delta, \gamma)$ bi-regular bipartite graph.
%     \item {\it Transposed incidence matrix of balanced incomplete block design (BIBD):} Suppose $\mathbf{M}$ satisfies the properties of the bi-adjacency matrix of a bi-regular bipartite graph and in addition, is such that the inner product, $\langle \mathbf{M}(:,i), \mathbf{M}(:,j)\rangle = \lambda$, for $i \ne j$ and $i, j \in [n]$. Then, $\mathbf{M}$ corresponds to the transpose of the incidence matrix of a BIBD \cite{StinsonBook} with parameters $(n,k,\gamma, \delta, \lambda)$. 
% %\aditya{complete in terms of N, K and dl, dr}
%     %\sifat{done. I will change the notations once everything is uniform.}
%     \item {\it Adjacency matrix of strongly-regular graph (SRG):} Suppose that $n = k$, $\mathbf{M} = \mathbf{M}^T$ and $\mathbf{M} \mathbf{1}_n = \delta \mathbf{1}_n$ and $\mathbf{M}$ satisfies%Furthermore, suppose that 
%     \begin{align*}
%         \mathbf{M}^2 &= \delta \mathbf{I}_n + \lambda \mathbf{M} + \mu (\mathbf{J}_n - \mathbf{I}_n - \mathbf{M}),
%     \end{align*}
%     where $\lambda,\mu$ are non-negative integers.
%     Then, $\mathbf{M}$ is the adjacency matrix of a $(n, \delta, \lambda, \mu)$ SRG \cite{Brouwer_SRG}. 
%     %\sifat{changed a line here}
    
% \end{itemize}
% We emphasize that such matrices have restrictions on their parameters; these are well-studied. We discuss these in the relevant subsections where these matrices are used.

%\section{First Construction}
\section{Construction based on Post Multiplication by diagonal matrices}
\label{sec:construction_RDM}


% Let $\mathbf{A}$ be the incidence matrix of a $(v,b,p,r, \lambda)$ BIBD.

%Our first scheme provides an upper bound on $\text{Err}_{\mathcal{F}}(\mathbf{B})$ for the general case $m \ge 1$ while keeping the data subset assignments the same as a BIBD.
%Our first scheme operates by post multiplying the baseline scheme by random diagonal matrices.

Let $\mathbf{A}$ be a $(n,k, \delta, \gamma)$ assignment matrix.  Let $\epsilon \in [0,1)$ and define the intervals $I_1 := [ 1-\epsilon  ,  1+\epsilon ]$, and 
$I_2 := [ -1-\epsilon , -1+\epsilon ]$. Let $I := I_1 \cup I_2$. %\aditya{Here we have repeated the usage of the symbol S}
%\sifat{fixed}
% Since $0 \le a <1$, $S_1 \cap S_2 = \emptyset$. So, $|S| = |S_1| + |S_2| = 2p$.
% Let $S$ be the union of a set of $p$ equidistant points from $-1-a$ to $-1 + a$ and a set of $p$ equidistant points from $1-a$ to $1 + a$, where $a \in \mathbb{R}$. So, $|S| = 2p$.\\
For $i \in [m]$, let $\mathbf{D}_i \in \mathbb{R}^{n \times n}$ be diagonal matrices where the diagonal entries are distributed i.i.d. uniformly over $I$. We construct the encoding matrix $\mathbf{B}$ such that
\begin{equation}\label{eq:effAD1AD2}
%\mathbf{B} = \begin{bmatrix} \bfA\bfD_1 \\ \bfA\bfD_2 \\ \vdots \\ \bfA\bfD_m \end{bmatrix}
\mathbf{B}^T = [ \bfD_1 \bfA^T ~|~ \bfD_2 \bfA^T ~|~ \dots ~|~ \bfD_m \bfA^T].
\end{equation}
% The resultant scheme will have number of workers $n = v$,  number of data subsets $k = b$, computation load $\delta = r$, replication factor $\gamma = p$ and a communication reduction factor of $m$. 
For this construction, we upper-bound the expected approximation error for a given non-straggler set $\mathcal{F} \subseteq [n]$; throughout this section the expectation is taken over the randomness in the diagonal matrices $\mathbf{D}_i$ for $i \in [m]$.




% added example

\begin{example}
Consider the $(3,3,2,2,1)$ BIBD whose incidence matrix is
\[
\mathbf{M}
=
\begin{bmatrix}
1 & 1 & 0\\
1 & 0 & 1\\
0 & 1 & 1
\end{bmatrix}.
\]
We choose the assignment matrix $\bfA = \bfM^T$. Let $m=2$ and $\epsilon=0$. Consider the realization
\[
\mathbf{D}_1=\mathbf{I}_3,
\qquad
\mathbf{D}_2= \begin{bmatrix}
1 & 0 & 0\\
0 & -1 & 0\\
0 & 0 & 1
\end{bmatrix}.
\]
Using the construction in~\eqref{eq:effAD1AD2}, we obtain %\aditya{Write equation below in transpose}
% \[
% \mathbf{B}
% =
% \begin{bmatrix}
% \mathbf{A}\mathbf{D}_1\\
% \mathbf{A}\mathbf{D}_2
% \end{bmatrix}
% =
% \begin{bmatrix}
% 1 & 1  & 0\\
% 1 & 0  & 1\\
% 0 & 1  & 1\\
% 1 & -1 & 0\\
% 1 & 0  & 1\\
% 0 & -1 & 1
% \end{bmatrix}.
% \]

\[
\mathbf{B}^T
=
\begin{bmatrix}
(\mathbf{A}\mathbf{D}_1)^T & (\mathbf{A}\mathbf{D}_2)^T
\end{bmatrix}
=
\begin{bmatrix}
1 & 1 & 0 & 1 & 1 & 0\\
1 & 0 & 1 & -1 & 0 & -1\\
0 & 1 & 1 & 0 & 1 & 1
\end{bmatrix}.
\]

Since the diagonal entries of $\mathbf{D}_1$ and $\mathbf{D}_2$ are nonzero,
\[
\operatorname{supp}(\mathbf{A}\mathbf{D}_1)
=
\operatorname{supp}(\mathbf{A}\mathbf{D}_2)
=
\operatorname{supp}(\mathbf{A}),
\]
and hence the encoding matrix preserves the original data-subset assignment.
\end{example}




%\aditya{There is a confusing point. You are using the superscript tilde over Delta to represent one thing in this section and a different thing in the other section. Use $\Delta^{(\text{diag})}$ here instead.}
%\sifat{done}
% \aditya{You are using c before it is defined.}

\begin{thm}  
\label{thm:random_diag}
For the construction as described above, let $\mathcal{F} \subseteq [n]$ be a set of size $n-s$ corresponding to non-stragglers and suppose that $\mathbf{B}_{\mathcal{F}}^T\mathbf{B}_{\mathcal{F}}  \in \mathbb{R}^{(n-s) \times (n-s)}$ is invertible. Furthermore, let $c := \mathbb{E}[X^2]\mathbb{E}[\frac{1}{X^2}]$, where $X$ is a random variable that is distributed uniformly  over $I$. Then, %the expected approximation error can be upper bounded as follows. 
% \begin{equation}\label{eq:AD1AD2err}
%  \mathbb{E}[\text{Err}_{\mathcal{F}}(\mathbf{B})] \le mk - \frac{m\delta^2(n-s)}{m\delta + \lambda (n-s-1)}   
% \end{equation} 
\begin{equation}
\label{eq:AD1AD2err}
 \mathbb{E}[\text{Err}_{\mathcal{F}}(\mathbf{B})] \le mk - m\mathbf{1}_k^T \bfA_{\cF} \left(\mathbf{K}\right)^{-1} \bfA_{\cF}^T\mathbf{1}_k, 
\end{equation} 
where $\mathbf{K}: = \mathbf{\Delta}_{\mathcal{F}} + c(m-1)\mathbf{\Delta}^{(\text{diag})}_{\mathcal{F}}$, $\mathbf{\Delta}_{\mathcal{F}} := \mathbf{A}_{\mathcal{F}}^T\mathbf{A}_{\mathcal{F}} \in \mathbb{R}^{(n-s) \times (n-s)}$ and  $\mathbf{\Delta}^{(\text{diag})}_{\mathcal{F}}$ is a diagonal matrix such that $\mathbf{\Delta}^{(\text{diag})}_{\mathcal{F}}(i,i) := \mathbf{\Delta}_{\mathcal{F}}(i,i)$ 
% and $\tilde{\mathbf{\Delta}}_{\mathcal{F}}(i,j) = 0$ for $i \ne j$ and 
for $i \in [n-s]$. 
\end{thm}

\begin{proof}%[Proof of Lemma \ref{lemma:random_diag}]
 % Let $\mathbf{B}_{\mathcal{F}}$ be a matrix consisting of the columns of $\mathbf{B}$ that correspond to non-straggling workers and 
For $i \in [m]$, let $\tilde{\mathbf{D}}_i$ be a submatrix of $\mathbf{D}_i$ such that $\tilde{\mathbf{D}}_i :=\mathbf{D}_i(\mathcal{F}, \mathcal{F})$. We have, %Also, $\mathbf{D}_i^T = \mathbf{D}_i$ and $\tilde{\mathbf{D}}^T_i = \tilde{\mathbf{D}}_i$, since $\mathbf{D}_i$ is diagonal. We have, %\aditya{Fix the transpose by apply tilde only on D, as we discussed yesterday. Do this throughout}
%\sifat{done}
\begin{align*}
\bfB_{\cF}^T &= [\tilde{\bfD}_1 \bfA_{\cF}^T ~|~ \tilde{\bfD}_2 \bfA_{\cF}^T ~|~ \dots ~|~\tilde{\bfD}_m\bfA_{\cF}^T], \text{~so that}\\
\bfB_{\cF}^T\bfB_{\cF} & =  \sum_{j = 1}^m \tilde{\bfD}_j\bfA_{\cF}^T\bfA_{\cF}\tilde{\bfD}_j = \sum_{j = 1}^m \tilde{\bfD}_j \mathbf{\Delta}_{\mathcal{F}} \tilde{\bfD}_j .
\end{align*}
% so that 
% \begin{align*}
% \bfB_{\cF}^T\bfB_{\cF} 
% & =  \sum_{j = 1}^m \tilde{\bfD}_j\bfA_{\cF}^T\bfA_{\cF}\tilde{\bfD}_j.
% \end{align*}
% % \[ \bfB_{\cF} = \begin{bmatrix} \bfA_{\cF}\tilde{\bfD}_1 \\ \bfA_{\cF}\tilde{\bfD}_2 \\ \vdots \\\bfA_{\cF}\tilde{\bfD}_m \end{bmatrix} \] 
% Note that, 
% \begin{align*}
% \bfB_{\cF}^T\bfB_{\cF} & = \begin{bmatrix} \tilde{\bfD}_1\bfA_{\cF}^T & \tilde{\bfD}_2\bfA_{\cF}^T & \hdots &\tilde{\bfD}_m\bfA_{\cF}^T \end{bmatrix}\begin{bmatrix} \bfA_{\cF}\tilde{\bfD}_1 \\ \bfA_{\cF}\tilde{\bfD}_2 \\ \vdots \\\bfA_{\cF}\tilde{\bfD}_m \end{bmatrix}\\
% & =  \sum_{j = 1}^m \tilde{\bfD}_j\bfA_{\cF}^T\bfA_{\cF}\tilde{\bfD}_j
% % \\
% % &= \sum_{i = 1}^m \tilde{\bfD_i}\tilde{\bfD_i}
% \end{align*}
% Let $\mathbf{\Delta}_{\mathcal{F}} := \bfA_{\cF}^T\bfA_{\cF}$. Let $\mathbf{I}_a$ be the $a \times a$ identity matrix and $\mathbf{J}_a$ be the $a \times a$ all ones matrix for any $a \in \mathbb{N}$. By \cite{kadheKR19}, $\mathbf{\Delta}_{\mathcal{F}} = (\delta - \lambda) \mathbf{I}_{n-s} + \lambda \mathbf{J}_{n-s} $. 
For $i \in [m]$, observe that $\mathbf{f}_i^T\mathbf{f}_i = k$ and $\mathbf{B}_{\mathcal{F}}^T \mathbf{f}_i = \tilde{\mathbf{D}_i}\mathbf{A}_{\mathcal{F}}^T\mathbf{1}_k$. Therefore, 
\begin{align*}
% &  \min\limits_{\substack{\mathbf{r}_i\in \mathbb{R}^{n} \\ supp(\mathbf{r}_i) \subseteq \mathcal{F} }}   \|\bfB\mathbf{r}_i - \mathbf{f}_i\|_2^2\\
% &= \min\limits_{\substack{\mathbf{r}_i\in \mathbb{R}^{n-s} }}   \|\bfB_{\mathcal{F}}\mathbf{r}_i - \mathbf{f}_i\|_2^2\\
&\mathbf{f}_i^T \mathbf{f}_i - \mathbf{f}_i^T \bfB_{\cF}(\bfB_{\cF}^T\bfB_{\cF})^{-1}\bfB_{\cF}^T\mathbf{f}_i\\
& = k - \mathbf{1}_k^T \bfA_{\cF} \tilde{\bfD}_i\Bigg(\sum_{j = 1}^m \tilde{\bfD}_j\mathbf{\Delta}_{\mathcal{F}}\tilde{\bfD}_j\Bigg)^{-1}\tilde{\bfD}_i \bfA_{\cF}^T \mathbf{1}_k\\
& = k - \mathbf{1}_k^T \bfA_{\cF} \Bigg(\tilde{\bfD}^{-1}_i\Bigg(\sum_{j = 1}^m \tilde{\bfD}_j\mathbf{\Delta}_{\mathcal{F}}\tilde{\bfD}_j\Bigg)\tilde{\bfD}^{-1}_i\Bigg)^{-1} \bfA_{\cF}^T \mathbf{1}_k\\
% & = k -  \mathbf{1}^T \bfA_{\cF} \Bigg(\tilde{\bfD_i}^{-1}\tilde{\bfD_i} \mathbf{\Delta}_{\mathcal{F}} \tilde{\bfD_i} \tilde{\bfD_i}^{-1}  +  \\
% &  \quad \quad \quad \quad \quad \quad  \quad \sum_{j \ne i}\tilde{\bfD_i}^{-1}\tilde{\bfD_j}\mathbf{\Delta}_{\mathcal{F}}\tilde{\bfD_j}\tilde{\bfD_i}^{-1} 
%   \Bigg)^{-1}\bfA_{\cF}^T\mathbf{1}\\
& = k -  \mathbf{1}_k^T \bfA_{\cF} \left( \mathbf{\Delta}_{\mathcal{F}}  + \sum_{\substack{j = 1 \\ j \neq i}}^m\tilde{\bfD}^{-1}_i\tilde{\bfD}_j\mathbf{\Delta}_{\mathcal{F}}\tilde{\bfD}_j\tilde{\bfD}^{-1}_i\right)^{-1}\bfA_{\cF}^T\mathbf{1}_k.
% & = k -  \mathbf{1}_k^T \bfA_{\cF} (\tilde{\bfD_i}^{-1}\tilde{\bfD_1} \mathbf{\Delta} \tilde{\bfD_1} \tilde{\bfD_i}^{-1} + \dots +  \mathbf{\Delta} \\
% & \quad \quad \quad \quad \quad \quad  \quad \quad \quad \dots +  \tilde{\bfD_i}^{-1}\tilde{\bfD_m} \mathbf{\Delta} \tilde{\bfD_m} \tilde{\bfD_i}^{-1})^{-1}\bfA_{\cF}^T\mathbf{1}_k
\end{align*}
% Note that for any positive definite matrix $\mathbf{X}$, the operation $\mathbf{X} \mapsto \mathbf{X}^{-1}$ is matrix convex and thus, $\mathbb{E}[\mathbf{X}^{-1}] \ge (\mathbb{E}[\mathbf{X}])^{-1}$ \CITE. (For matrices $\mathbf{X}, \mathbf{Y}$, $\mathbf{X} \ge \mathbf{Y}$ means $\mathbf{X}- \mathbf{Y}$ is positive semi-definite). Consequently, 

% \[ \mathbb{E}\left[ \bigg( \mathbf{\Delta}_{\mathcal{F}}  + \sum_{\substack{j = 1 \\ j \neq i}}^m\tilde{\bfD}^{-1}_i\tilde{\bfD}_j\mathbf{\Delta}_{\mathcal{F}}\tilde{\bfD}_j\tilde{\bfD}^{-1}_i\bigg)^{-1}\right] \ge \mathbb{E}\left[ \mathbf{\Delta}_{\mathcal{F}}  + \sum_{\substack{j = 1 \\ j \neq i}}^m\tilde{\bfD}^{-1}_i\tilde{\bfD}_j\mathbf{\Delta}_{\mathcal{F}}\tilde{\bfD}_j\tilde{\bfD}^{-1}_i\right] \]


Let
\begin{equation}\label{K_i}
\mathbf{K}_i := \mathbf{\Delta}_{\mathcal{F}}  +  \sum_{\substack{j = 1 \\ j \neq i}}^m\tilde{\bfD}^{-1}_i\tilde{\bfD}_j\mathbf{\Delta}_{\mathcal{F}}\tilde{\bfD}_j\tilde{\bfD}^{-1}_i.\end{equation}

Note that
\[
\mathbf{K}_i
= \tilde{\bf D}_i^{-1}
\left(\sum_{j=1}^m \tilde{\bf D}_j \mathbf{\Delta}_{\mathcal{F}} \tilde{\bf D}_j\right)
\tilde{\bf D}_i^{-1}
= \tilde{\bf D}_i^{-1}
(\mathbf{B}_{\mathcal{F}}^{\top}\mathbf{B}_{\mathcal{F}})
\tilde{\bf D}_i^{-1}.
\]

Since $\mathbf{B}_{\mathcal{F}}^{T}\mathbf{B}_{\mathcal{F}}$ is positive semidefinite and invertible (by assumption), it is positive definite. Also, since $\tilde{\bf D}_i$ is diagonal with nonzero entries, it is also invertible. Consequently,
$\mathbf{K}_i$ is positive definite. 





%\aditya{You can format this better by just focusing on the matrix term within the expectation. Just analyze it and then provide the overall result at the end.}
%\sifat{done}
% \begin{align*}
% % &\mathbb{E}\left[\min\limits_{\substack{\mathbf{r}_i\in \mathbb{R}^{n} \\ supp(\mathbf{r}_i) \subseteq \mathcal{F} }}   \|\bfB\mathbf{r}_i - \mathbf{f}_i\|_2^2\right] \\
% &\mathbb{E}\left[ \mathbf{f}_i^T \mathbf{f}_i - \mathbf{f}_i^T \bfB_{\cF}(\bfB_{\cF}^T\bfB_{\cF})^{-1}\bfB_{\cF}^T\mathbf{f}_i\right] \\
% % &= \mathbb{E} \bigg[ k -  \mathbf{1}^T \bfA_{\cF} \bigg( \mathbf{\Delta}_{\mathcal{F}}  + \\
% % & \quad \quad \quad \quad \quad \quad  \quad \sum_{j \ne i}\tilde{\bfD_i}^{-1}\tilde{\bfD_j}\mathbf{\Delta}_{\mathcal{F}}\tilde{\bfD_j}\tilde{\bfD_i}^{-1}\bigg)^{-1}\bfA_{\cF}^T\mathbf{1} \bigg]\\ 
% % &= k - \mathbf{1}_k^T \mathbf{A}_{\mathcal{F}}\mathbb{E} \left[\left( \mathbf{\Delta}_{\mathcal{F}}  +  \sum_{\substack{j = 1 \\ j \neq i}}^m\tilde{\bfD_i}^{-1}\tilde{\bfD_j}\mathbf{\Delta}_{\mathcal{F}}\tilde{\bfD_j}\tilde{\bfD_i}^{-1}\right)^{-1} \right] \mathbf{A}_{\mathcal{F}}^T \mathbf{1}_k\\
% & = k - \mathbf{1}_k^T \mathbf{A}_{\mathcal{F}} \mathbb{E} \Bigg[\bigg( 
% \mathbf{\Delta}_{\mathcal{F}} \\
% &  \quad \quad \quad \quad \quad \quad \quad + 
% \sum_{\substack{j = 1 \\ j \neq i}}^m 
% \tilde{\bfD}^{-1}_i \tilde{\bfD}_j \mathbf{\Delta}_{\mathcal{F}} 
% \tilde{\bfD}_j \tilde{\bfD}^{-1}_i
% \bigg)^{-1} \Bigg] \mathbf{A}_{\mathcal{F}}^T \mathbf{1}_k\\
% &\le k - \mathbf{1}_k^T \mathbf{A}_{\mathcal{F}}\Bigg(\mathbb{E}\bigg[ \mathbf{\Delta}_{\mathcal{F}}  \\
% & \quad \quad \quad \quad \quad \quad \quad+ \sum_{\substack{j = 1 \\ j \neq i}}^m\tilde{\bfD}^{-1}_i\tilde{\bfD}_j\mathbf{\Delta}_{\mathcal{F}}\tilde{\bfD}_j\tilde{\bfD}^{-1}_i \bigg] \Bigg)^{-1}\mathbf{A}_{\mathcal{F}}^T \mathbf{1}_k\\
% & = k - \mathbf{1}_k^T \bfA_{\cF} (\mathbf{\Delta}_{\mathcal{F}} + (m-1)\mathbf{\Delta}^{(\text{diag})}_{\mathcal{F}})^{-1} \bfA_{\cF}^T\mathbf{1}_k
% % & =  k -   \frac{\delta^2(n-s)} {m\delta + (n-s-1)\lambda}
% \end{align*}

% Here, the inequality holds since 

% The second equality holds because of the following:

Let $r_1, r_2 \in [m]$ and $r_1 \ne r_2$. Let $u,v \in [n-s]$.  Note that, 
\begin{align*}
&\mathbb{E}[\tilde{\bfD}^{-1}_{r_1}\tilde{\bfD}_{r_2} \mathbf{\Delta}_{\mathcal{F}} \tilde{\bfD}_{r_2} \tilde{\bfD}^{-1}_{r_1}](u,v) \\
&= \mathbb{E}[\mathbf{\Delta}_{\mathcal{F}}(u,v)  \frac{\tilde{\bfD}_{r_2}(u,u)\tilde{\bfD}_{r_2}(v,v)} {\tilde{\bfD}_{r_1}(u,u)\tilde{\bfD}_{r_1}(v,v)}].
\end{align*}
If $u = v$, then 
\begin{align*}
&\mathbb{E}[\mathbf{\Delta}_{\mathcal{F}}(u,v)  \frac{\tilde{\bfD}_{r_2}(u,u)\tilde{\bfD}_{r_2}(v,v)} {\tilde{\bfD}_{r_1}(u,u)\tilde{\bfD}_{r_1}(v,v)}]\\ 
&= \mathbb{E}[\mathbf{\Delta}_{\mathcal{F}}(u,u)  \frac{(\tilde{\bfD}_{r_2}(u,u))^2} {(\tilde{\bfD}_{r_1}(u,u))^2}]\\
&= {\mathbf{\Delta}_{\mathcal{F}}(u,u)}\mathbb{E}[(\tilde{\bfD}_{r_2}(u,u))^2] 
\mathbb{E}\left[\frac{1}{(\tilde{\bfD}_{r_1}(u,u))^2}\right]\\
&= c\mathbf{\Delta}_{\mathcal{F}}(u,u)  = c\mathbf{\Delta}^{(\text{diag})}_{\mathcal{F}}(u,u).
\end{align*}

The third equality holds since the diagonal entries of $\tilde{\bfD}_{r_2}$ and $\tilde{\bfD}_{r_1}$ are i.i.d. and uniformly distributed over the set $I$. Now, if $u \ne v$, then 
\begin{align*}
&\mathbb{E}[\mathbf{\Delta}_{\mathcal{F}}(u,v)  \frac{\tilde{\bfD}_{r_2}(u,u)\tilde{\bfD}_{r_2}(v,v)} {\tilde{\bfD}_{r_1}(u,u)\tilde{\bfD}_{r_1}(v,v)}]\\
& \overset{(a)}= \mathbf{\Delta}_{\mathcal{F}}(u,v)\mathbb{E}[\tilde{\bfD}_{r_2}(u,u)]\mathbb{E}[\tilde{\bfD}_{r_2}(v,v)]\\
& \quad \quad \quad \quad \quad \quad \mathbb{E}\left[\frac{1}{\tilde{\bfD}_{r_1}(u,u)}\right] 
\mathbb{E}\left[\frac{1}{\tilde{\bfD}_{r_1}(v,v)}\right]\\
& \overset{(b)}= \mathbf{\Delta}_{\mathcal{F}}(u,v) \cdot 0 = 0 =c\mathbf{\Delta}^{(\text{diag})}_{\mathcal{F}}(u,v).
\end{align*}
Here, $(a)$ holds since the diagonal entries of $\tilde{\bfD}_{r_1}$ and $\tilde{\bfD}_{r_2}$ are i.i.d. and $(b)$ holds since the expected value of each diagonal entry is zero, and the expected value of the reciprocal of the diagonal entry is also zero. Thus, 
\begin{align*}
\mathbb{E}[\tilde{\bfD}^{-1}_{r_1}\tilde{\bfD}_{r_2} \mathbf{\Delta}_{\cF} \tilde{\bfD}_{r_2} \tilde{\bfD}^{-1}_{r_1}] 
= c\mathbf{\Delta}^{(\text{diag})}_{\mathcal{F}}.
\end{align*}








% Let $a,b \in [m]$ and $a \ne b$. Let $u,v \in [n-s]$.  Note that, 
% \begin{align*}
% &\mathbb{E}[\tilde{\bfD}^{-1}_a\tilde{\bfD}_b \mathbf{\Delta}_{\mathcal{F}} \tilde{\bfD}_b \tilde{\bfD}^{-1}_a](u,v) \\
% &= \mathbb{E}[\mathbf{\Delta}_{\mathcal{F}}(u,v)  \frac{\tilde{\bfD}_b(u,u)\tilde{\bfD}_b(v,v)} {\tilde{\bfD}_a(u,u)\tilde{\bfD}_a(v,v)}].
% \end{align*}
% If $u = v$, then \begin{align*}
% &\mathbb{E}[\mathbf{\Delta}_{\mathcal{F}}(u,v)  \frac{\tilde{\bfD}_b(u,u)\tilde{\bfD}_b(v,v)} {\tilde{\bfD}_a(u,u)\tilde{\bfD}_a(v,v)}]\\ &= \mathbb{E}[\mathbf{\Delta}_{\mathcal{F}}(u,u)  \frac{(\tilde{\bfD}_b(u,u))^2} {(\tilde{\bfD}_a(u,u))^2}]\\
% &= {\mathbf{\Delta}_{\mathcal{F}}(u,u)}\mathbb{E}[\tilde{\bfD}_b(u,u))^2] \mathbb{E}[\frac{1}{\tilde{\bfD}_a(u,u))^2}]\\
% &= c\mathbf{\Delta}_{\mathcal{F}}(u,u)  = c\mathbf{\Delta}^{(\text{diag})}_{\mathcal{F}}(u,u).
% \end{align*}
% % If $i \ne j$, then 
% % \begin{align*}
% % &\mathbb{E}[\mathbf{\Delta}_{\mathcal{F}}(i,j)  \frac{\tilde{\bfD_b}(i,i)\tilde{\bfD_b}(j,j)} {\tilde{\bfD_a}(i,i)\tilde{\bfD_a}(j,j)}]\\
% % & = \mathbf{\Delta}_{\mathcal{F}}(i,j)\mathbb{E}[\tilde{\bfD_b}(i,i)]\mathbb{E}[\tilde{\bfD_b}(j,j)]\mathbb{E}\left[\frac{1}{\tilde{\bfD_a}(i,i)}\right] \mathbb{E}\left[\frac{1}{\tilde{\bfD_a}(j,j)}\right]\\
% % & = \mathbf{\Delta}_{\mathcal{F}}(i,j) . 0 = 0 =\tilde{\mathbf{\Delta}}_{\mathcal{F}}(i,j)
% % \end{align*}
% The third equality holds since the diagonal entries of $\tilde{\bfD}_b$ and $\tilde{\bfD}_a$ are independently and uniformly distributed over the set $S$. Now, if $u \ne v$, then 
% \begin{align*}
% &\mathbb{E}[\mathbf{\Delta}_{\mathcal{F}}(u,v)  \frac{\tilde{\bfD}_b(u,u)\tilde{\bfD}_b(v,v)} {\tilde{\bfD}_a(u,u)\tilde{\bfD}_a(v,v)}]\\
% & \overset{(a)}= \mathbf{\Delta}_{\mathcal{F}}(u,v)\mathbb{E}[\tilde{\bfD}_b(u,u)]\mathbb{E}[\tilde{\bfD}_b(v,v)]\\
% & \quad \quad \quad \quad \quad \quad \mathbb{E}\left[\frac{1}{\tilde{\bfD}_a(u,u)}\right] \mathbb{E}\left[\frac{1}{\tilde{\bfD}_a(v,v)}\right]\\
% & \overset{(b)}= \mathbf{\Delta}_{\mathcal{F}}(u,v) . 0 = 0 =c\mathbf{\Delta}^{(\text{diag})}_{\mathcal{F}}(u,v).
% \end{align*}
% Here, $(a)$ holds since the diagonal entries of $\tilde{\bfD}_a$ and $\tilde{\bfD}_b$ are i.i.d. and $(b)$ holds since the expected value of each diagonal entry is zero. Thus, 
% \begin{align*}\mathbb{E}[\tilde{\bfD}^{-1}_a\tilde{\bfD}_b \mathbf{\Delta}_{\cF} \tilde{\bfD}_b \tilde{\bfD}^{-1}_a] = c\mathbf{\Delta}^{(\text{diag})}_{\mathcal{F}}
% \end{align*}
Using linearity of expectation, 
% \begin{align*}
% \mathbf{K}_i= \mathbf{\Delta} + (m-1) \delta \mathbf{I}_{n-s} = (m\delta - \lambda) \mathbf{I}_{n-s} + \lambda \mathbf{J}_{n-s}
% \end{align*}
\begin{align*}\label{E[k_i]}
\mathbb{E}[\mathbf{K}_i] &= \mathbb{E}\left[ \mathbf{\Delta}_{\mathcal{F}}  +  \sum_{\substack{j = 1 \\ j \neq i}}^m\tilde{\bfD}^{-1}_i\tilde{\bfD}_j\mathbf{\Delta}_{\mathcal{F}}\tilde{\bfD}_j\tilde{\bfD}^{-1}_i \right]\\
& = \mathbf{\Delta}_{\mathcal{F}} + c(m-1) \mathbf{\Delta}^{(\text{diag})}_{\mathcal{F}} = \mathbf{K}.
\end{align*}
Finally,  \nopagebreak[4]
% Observe that $\mathbf{K}_i$ has eigenvalue $m\delta - \lambda$ with multiplicity $n-s-1$ and eigenvalue $m\delta - \lambda + \lambda(n-s)$ with multiplicity 1 and the corresponding eigenvector being $\mathbf{1}_{n-s}$. So, this matrix is invertible. Also, since the matrix is symmetric, the eigenvectors are orthonormal. Let $\mathbf{K}_i$ have the following eigen-decomposition, 
% $\mathbf{K}_i = \sum_{i = 1}^{n-s} \lambda_i \mathbf{v}_i\mathbf{v}_i^T$. So, $\mathbf{K}_i^{-1} = \sum_{i = 1}^{n-s}\frac{1}{\lambda_i}\mathbf{v}_i\mathbf{v}_i^T = \frac{1}{m\delta + \lambda(n-s-1)} \mathbf{1}\mathbf{1}^T + \sum_{i \ne 1}\frac{1}{\lambda_i}\mathbf{v}_i\mathbf{v}_i^T$
% % The final equality follows by applying matrix inversion lemma \cite{44d50e5e-b5a2-3887-ba7d-b3766ee95319}.

% Since $ \mathbf{A}_{\mathcal{F}}^T\mathbf{1}_k = \delta \mathbf{1}_{n-s}$, it follows that $\mathbf{1}_k^T \mathbf{A}_{\mathcal{F}} \mathbf{K}_i^{-1} \mathbf{A}_{\mathcal{F}}^T\mathbf{1}_k =  \frac{\delta^2(n-s)} {m\delta + (n-s-1)\lambda}$
% since $\mathbf{v}_i^T \mathbf{1}_{n-s} = 0$, for $i \ne 1, i \in [n-s]$. 

\begin{align*}
&\mathbb{E}[\text{Err}_{\mathcal{F}}(\mathbf{B})] \\
& \overset{(c)}{=} \mathbb{E} \left[ \sum_{i =1}^m\mathbf{f}_i^T \mathbf{f}_i - \mathbf{f}_i^T \bfB_{\cF}(\bfB_{\cF}^T\bfB_{\cF})^{-1}\bfB_{\cF}^T\mathbf{f}_i \right]\\
& \overset{(d)}{=} \sum_{i =1}^m \mathbb{E} \left[\mathbf{f}_i^T \mathbf{f}_i - \mathbf{f}_i^T \bfB_{\cF}(\bfB_{\cF}^T\bfB_{\cF})^{-1}\bfB_{\cF}^T\mathbf{f}_i \right]\\
& \overset{(e)}{=}  \sum_{i =1}^m \mathbb{E} \left[k -  \mathbf{1}_k^T \bfA_{\cF} \mathbf{K}^{-1}_i\bfA_{\cF}^T\mathbf{1}_k\right]\\
& \overset{(f)}{=}  \sum_{i =1}^m k -  \mathbf{1}_k^T \bfA_{\cF} \mathbb{E} \left[\mathbf{K}^{-1}_i\right]\bfA_{\cF}^T\mathbf{1}_k\\
&\overset{(g)}{\le}  \sum_{i =1}^m k -  \mathbf{1}_k^T \bfA_{\cF} \left(\mathbb{E} \left[\mathbf{K}_i\right]\right)^{-1}\bfA_{\cF}^T\mathbf{1}_k\\
& \overset{(h)}{=} \sum_{i = 1}^m  k - \mathbf{1}_k^T \bfA_{\cF} \left(\mathbf{K}\right)^{-1} \bfA_{\cF}^T\mathbf{1}_k\\
%& \overset{(i)}{=}  
& = mk - m\mathbf{1}_k^T \bfA_{\cF} \left(\mathbf{K}\right)^{-1} \bfA_{\cF}^T\mathbf{1}_k.
\end{align*}
Here, $(c)$ follows from $\eqref{eq:opt_dec_error}$, $(d)$ follows by linearity of expectation. 
$(e)$ follows from $\eqref{K_i}$. $(f)$ follows since the expectation is over the random diagonal 
matrices $\mathbf{D}_i, i \in [m]$. $(g)$ follows since for any positive definite matrix 
$\mathbf{X}$, the operation $\mathbf{X} \mapsto \mathbf{X}^{-1}$ is matrix convex and thus, 
$\mathbb{E}[\mathbf{X}^{-1}] \succeq (\mathbb{E}[\mathbf{X}])^{-1}$ (see \cite{bhatia2009positive}). 
(For matrices $\mathbf{X}, \mathbf{Y}$, $\mathbf{X} \succeq \mathbf{Y}$ means $\mathbf{X}- \mathbf{Y}$ 
is positive semidefinite). $(h)$ follows since $\mathbb{E}\left[\mathbf{K}_i\right] = \mathbf{K}$.
\end{proof}







 
% \begin{align*}
% &\mathbb{E}[\text{Err}_{\mathcal{F}}(\mathbf{B})] \\
% & \overset{(a)}{=} \mathbb{E} \left[ \sum_{i =1}^m\mathbf{f}_i^T \mathbf{f}_i - \mathbf{f}_i^T \bfB_{\cF}(\bfB_{\cF}^T\bfB_{\cF})^{-1}\bfB_{\cF}^T\mathbf{f}_i \right]\\
% % &= \mathbb{E}\left[\min\limits_{\substack{\mathbf{R}\in \mathbb{R}^{n \times m} \\ supp(\mathbf{R}(:,i)) \subseteq \mathcal{F} \\ \forall i \in [m]}}  {\|\mathbf{B} \mathbf{R}- \mathbf{F}\|_F^2} \right]\\
% % & = \mathbb{E} \left[\sum_{i = 1}^m  \min\limits_{\substack{\mathbf{r}_i\in \mathbb{R}^{n} \\ supp(\mathbf{r}_i) \subseteq \mathcal{F} }}   \|\bfB\mathbf{r}_i - \mathbf{f}_i\|_2^2\right]\\
% % & = \sum_{i = 1}^m \mathbb{E} \left[\min\limits_{\substack{\mathbf{r}_i\in \mathbb{R}^{n} \\ supp(\mathbf{r}_i) \subseteq \mathcal{F} }}   \|\bfB\mathbf{r}_i - \mathbf{f}_i\|_2^2\right]\\
% & \overset{(b)}{=} \sum_{i =1}^m \mathbb{E} \left[\mathbf{f}_i^T \mathbf{f}_i - \mathbf{f}_i^T \bfB_{\cF}(\bfB_{\cF}^T\bfB_{\cF})^{-1}\bfB_{\cF}^T\mathbf{f}_i \right]\\
% & \overset{(c)}{=}  \sum_{i =1}^m \mathbb{E} \left[k -  \mathbf{1}_k^T \bfA_{\cF} \mathbf{K}^{-1}_i\bfA_{\cF}^T\mathbf{1}_k\right]\\
% & \overset{(d)}{=}  \sum_{i =1}^m k -  \mathbf{1}_k^T \bfA_{\cF} \mathbb{E} \left[\mathbf{K}^{-1}_i\right]\bfA_{\cF}^T\mathbf{1}_k\\
% &\overset{(e)}{\le}  \sum_{i =1}^m k -  \mathbf{1}_k^T \bfA_{\cF} \left(\mathbb{E} \left[\mathbf{K}_i\right]\right)^{-1}\bfA_{\cF}^T\mathbf{1}_k\\
% & \overset{(f)}{=} \sum_{i = 1}^m  k - \mathbf{1}_k^T \bfA_{\cF} \left(\mathbf{K}\right)^{-1} \bfA_{\cF}^T\mathbf{1}_k\\
% & \overset{(g)}{=}  mk - m\mathbf{1}_k^T \bfA_{\cF} \left(\mathbf{K}\right)^{-1} \bfA_{\cF}^T\mathbf{1}_k.
% % & \le \sum_{i = 1}^m  k -   \frac{\delta^2(n-s)} {m\delta + (n-s-1)\lambda} \\
% % & = mk -   \frac{m\delta^2(n-s)} {m\delta + (n-s-1)\lambda}
% \end{align*}
% % \aditya{Rewrite below by using supercript on the operators}
% Here, $(a)$ follows from $\eqref{eq:opt_dec_error}$, $(b)$ follows by linearity of expectation. $(c)$ follows from $\eqref{K_i}$. $(d)$ follows since the expectation is over the random diagonal matrices $\mathbf{D}_i, i \in [m]$. $(e)$ follows since for any positive definite matrix $\mathbf{X}$, the operation $\mathbf{X} \mapsto \mathbf{X}^{-1}$ is matrix convex and thus, $\mathbb{E}[\mathbf{X}^{-1}] \succeq (\mathbb{E}[\mathbf{X}])^{-1}$ (see \cite{bhatia2009positive}). (For matrices $\mathbf{X}, \mathbf{Y}$, $\mathbf{X} \succeq \mathbf{Y}$ means $\mathbf{X}- \mathbf{Y}$ is positive semi-definite). $(f)$ follows since $\mathbb{E}\left[\mathbf{K}_i\right] = \mathbf{K}$.








% \noindent {\it Sketch of Proof:} Note that $\bfB_{\cF}^T\bfB_{\cF} =  \sum_{j = 1}^m \tilde{\bfD}_j\bfA_{\cF}^T\bfA_{\cF}\tilde{\bfD}_j$ and for $i \in [m]$,
% \begin{align*}
%     &\mathbf{f}_i^T \bfB_{\cF}(\bfB_{\cF}^T\bfB_{\cF})^{-1}\bfB_{\cF}^T\mathbf{f}_i \\
%     &=\mathbf{1}_k^T \bfA_{\cF} \Bigg(\tilde{\bfD}^{-1}_i\Bigg(\sum_{j = 1}^m \tilde{\bfD}_j\mathbf{\Delta}_{\mathcal{F}}\tilde{\bfD}_j\Bigg)\tilde{\bfD}^{-1}_i\Bigg)^{-1} \bfA_{\cF}^T \mathbf{1}_k\\
%     &=\mathbf{1}_k^T \bfA_{\cF} \Bigg( \mathbf{\Delta}_{\mathcal{F}}  + \sum_{\substack{j = 1, j \neq i}}^m\tilde{\bfD}^{-1}_i\tilde{\bfD}_j\mathbf{\Delta}_{\mathcal{F}}\tilde{\bfD}_j\tilde{\bfD}^{-1}_i\Bigg)^{-1}\bfA_{\cF}^T\mathbf{1}_k.
% \end{align*}
% Let, 
% $\mathbf{K}_i := \mathbf{\Delta}_{\mathcal{F}}  +  \sum_{\substack{j = 1 \\ j \neq i}}^m\tilde{\bfD}^{-1}_i\tilde{\bfD}_j\mathbf{\Delta}_{\mathcal{F}}\tilde{\bfD}_j\tilde{\bfD}^{-1}_i.$
% Then, by the matrix-convexity of inverse map over positive semi-definite matrices \cite{bhatia2009positive}, we have the $\mathbb{E} \left[\mathbf{K}^{-1}_i\right] \geq \left(\mathbb{E} \left[\mathbf{K}_i\right]\right)^{-1}$. The result follows by calculating $\mathbb{E} \left[\mathbf{K}_i\right]$, using the linearity of expectation, and calculating the relevant quantities.
We now consider constructions where $\bfA$ is chosen as an instance of a structured matrix presented in Section \ref{sec:structured_matrices}.

\begin{remark}\label{rem:c_equals_one}
Note that, by the Cauchy--Schwarz inequality, $
c=\mathbb{E}[X^2]\mathbb{E}\left[\frac{1}{X^2}\right]
\geq \left(\mathbb{E}\left[X\left(\frac{1}{X}\right)\right]\right)^2=1
$. When $\epsilon=0$, $X$ takes values in $\{-1,1\}$ with equal probability, and hence $c=1$, which is the smallest possible value of $c$. 

% For the BIBD and SRG constructions considered below, $\mathbf{B}_{\mathcal{F}}^T\mathbf{B}_{\mathcal{F}}$ is invertible for every $\epsilon\in[0,1)$ and every nonempty $\mathcal{F}$, as shown in the proofs of Corollaries~1 and~2. Thus, the choice of $\epsilon$ does not affect the applicability of Lemma~1 in these cases. We next explain why choosing $\epsilon=0$ yields the smallest upper bound in Lemma~1

Let $c_2\geq c_1\geq 1$, and let $\mathbf{K}_1$ and $\mathbf{K}_2$ denote the matrices obtained by substituting $c_1$ and $c_2$, respectively, for the parameter $c$ in the definition of $\mathbf{K}$ in Theorem \ref{thm:random_diag}. Since $\mathbf{\Delta}_{\mathcal{F}}=\mathbf{A}_{\mathcal{F}}^T\mathbf{A}_{\mathcal{F}}$ is positive semidefinite, its diagonal entries are nonnegative. Hence, $\mathbf{\Delta}_{\mathcal{F}}^{(\text{diag})}\succeq\mathbf{0}$. Then,
\[
\mathbf{K}_2-\mathbf{K}_1
=(c_2-c_1)(m-1)\mathbf{\Delta}_{\mathcal{F}}^{(\mathrm{diag})}
\succeq \mathbf{0}.
\]
Suppose $c_1$ and $c_2$ satisfy the invertibility assumption on $\mathbf{K}$ in Theorem ~\ref{thm:random_diag}. Then, $\mathbf{K}_1$ and $\mathbf{K}_2$ are positive definite. Therefore, the upper bound in \eqref{eq:AD1AD2err} increases with $c_i$.  Since $c\geq 1$, with equality when $\epsilon=0$, choosing $\epsilon=0$ yields the smallest upper bound. 



We note that the applicability of Theorem \ref{thm:random_diag} requires $\mathbf{B}_{\mathcal{F}}^T 
\mathbf{B}_{\mathcal{F}}$ to be invertible. This condition may not hold for all values of $\epsilon$ depending on the assignment matrix $\mathbf{A}$. For the BIBD and SRG constructions considered below, 
$\mathbf{B}_{\mathcal{F}}^T \mathbf{B}_{\mathcal{F}}$ is invertible for $\epsilon = 0$ and every 
nonempty $\mathcal{F}$, as established in the proofs of Corollaries~1
and~2. Therefore, for these constructions, the choice $\epsilon = 0$ both satisfies 
the assumptions of Theorem \ref{thm:random_diag} and yields the smallest upper bound.

\end{remark}



% \begin{remark} \label{rem:c_equals_one}
%  Note that by Cauchy-Schwarz inequality, $c = \mathbb{E}[X^2] \mathbb{E}[\frac{1}{X^2}] \ge (\mathbb{E} [X (\frac{1}{X})])^2 = 1$. Therefore, in the sequel, we often consider $S$ with $\epsilon =0$, as this achieves $c=1$ and results in the lowest upper bound. % For the following two constructions using BIBD and SRG, we consider the set $S$ with $a = 0$, as this achieves $c=1$ and consequently results in the lowest upper bound.    
% \end{remark}
% \aditya{Refer to this remark when using it}
 
% \aditya{include argument that the required matrix is non-singular.}
% \sifat{done in the proofs}
% \aditya{explain why the case considered is best }
\begin{cor} \label{cor:BIBD}
Suppose that the assignment matrix $\mathbf{A}$ is the transpose of the incidence matrix of a $(n,k,\gamma,\delta,\lambda)$ BIBD. %and let $\mathbf{A} = \mathbf{E}$ so that $\mathbf{A}$ is an $(n,k,\delta, \gamma)$ assignment matrix. 
Then, for the construction as in $\eqref{eq:effAD1AD2}$ with $\epsilon  = 0$, the expected error can be upper bounded as
\begin{equation} \label{eq:AD1AD2err_BIBD}
\mathbb{E}[\text{Err}_{\mathcal{F}}(\mathbf{B})] \le mk -   \frac{m\delta^2(n-s)} {m\delta + (n-s-1)\lambda}.
\end{equation}
\end{cor}
\begin{proof} 
% \aditya{Mention this fact in the Preliminaries}
By \eqref{eq: bibd_prop}, $\mathbf{\Delta}_{\mathcal{F}} = \mathbf{A}_{\mathcal{F}}^T\mathbf{A}_{\mathcal{F}} =  (\delta - \lambda) \mathbf{I}_{n-s} + \lambda \mathbf{J}_{n-s} $. Observe that $\mathbf{\Delta}_{\mathcal{F}}$ has eigenvalue $\delta - \lambda$ with multiplicity $n-s-1$ and eigenvalue $\delta - \lambda + \lambda(n-s)$ with multiplicity 1. Since $\delta  > \lambda$ for a BIBD \cite{StinsonBook}, the eigenvalues of $\mathbf{\Delta}_{\cF}$ are strictly positive. Since $\mathbf{\Delta}_{\mathcal{F}}$ is symmetric and all the eigenvalues are strictly positive, it is positive definite. Since by construction $\tilde{\bfD}_j$ is invertible, $\tilde{\bfD}_j \mathbf{\Delta}_{\cF} \tilde{\bfD}_j$ is positive definite for each $j \in [m]$.  Since the sum of positive definite matrices is positive definite,  it follows that $\mathbf{B}_{\mathcal{F}}^T\mathbf{B}_{\mathcal{F}}  =  \sum_{j = 1}^m \tilde{\bfD}_j\mathbf{\Delta}_{\mathcal{F}}\tilde{\bfD}_j$ is positive definite and hence invertible.
Now, $\mathbf{\Delta}^{(\text{diag})}_{\mathcal{F}} = \delta \mathbf{I}_{n-s}$. %As noted in Remark~\ref{rem:c_equals_one}, when $\epsilon = 0$, we have $c = 1$.
Hence $\mathbf{K} = \mathbf{\Delta}_{\mathcal{F}} + (m-1)\mathbf{\Delta}^{(\text{diag})}_{\mathcal{F}} = (m\delta - \lambda) \mathbf{I}_{n-s} + \lambda \mathbf{J}_{n-s} $. Also, note that $\mathbf{A}_{\mathcal{F}}^T\mathbf{1}_k = \delta \mathbf{1}_{n-s}$ in this case. Observe that $\mathbf{K}$ has eigenvalue $m\delta - \lambda$ with multiplicity $n-s-1$, and an eigenvalue $m\delta - \lambda + \lambda(n-s)$ with multiplicity 1 with corresponding eigenvector $\mathbf{1}_{n-s}$. %Let $\mathbf{K}$ have the following eigen-decomposition, 
Let $\mathbf{K} = \sum_{i = 1}^{n-s} \lambda_i \mathbf{v}_i\mathbf{v}_i^T$ (eigen-decomposition). Therefore,
\begin{align*}
\mathbf{K}^{-1} &= \sum_{i = 1}^{n-s}\frac{1}{\lambda_i}\mathbf{v}_i\mathbf{v}_i^T
= \frac{1}{m\delta + \lambda(n-s-1)} \frac{\mathbf{1}_{n-s}\mathbf{1}^T_{n-s}}{n-s} + \sum_{i = 2}^{n-s}\frac{1}{\lambda_i}\mathbf{v}_i\mathbf{v}_i^T.    
\end{align*}
% The final equality follows by applying matrix inversion lemma \cite{44d50e5e-b5a2-3887-ba7d-b3766ee95319}.

Since $ \mathbf{A}_{\mathcal{F}}^T\mathbf{1}_k = \delta \mathbf{1}_{n-s}$, it follows that $\mathbf{1}_k^T \mathbf{A}_{\mathcal{F}} \mathbf{K}^{-1} \mathbf{A}_{\mathcal{F}}^T\mathbf{1}_k =  \frac{\delta^2(n-s)} {m\delta + (n-s-1)\lambda}$
since $\mathbf{v}_i^T \mathbf{1}_{n-s} = 0$, for $i \ne 1$ and $i \in [n-s]$. The result follows from \eqref{eq:AD1AD2err}.
\end{proof}
\begin{remark} When $\mathbf{A}$ is chosen as above, the work of \cite{kadheKR19} considered the case of $m=1$ and showed that for any set of non-straggling workers $\mathcal{F}$ of size $n-s$, $\text{Err}_{\mathcal{F}}(\mathbf{B}) = \text{Err}_{\mathcal{F}}(\mathbf{A}) = k - \frac{\delta^2(n-s)}{\delta +  (n-s-1)\lambda}$.

%\aditya{CHanged to $\bfA$ below since their scheme only works for $m=1$}
% \begin{equation}\label{eq:ineffbibderr}
% \text{Err}_{\mathcal{F}}(\mathbf{A}) = k - \frac{\delta^2(n-s)}{\delta +  (n-s-1)\lambda}.
% \end{equation}
Our work is a generalization of their result for the communication-efficient case when $m > 1$. Indeed, a naive application of their scheme is to simply construct $\mathbf{B}$ by stacking $\mathbf{A}$ vertically $m$ times, i.e., 
\begin{align}
  \bfB^T &= [ \bfA^T~|~ \dots ~|~ \bfA^T].  \label{eq:trivial_cons}
\end{align}
    


% \begin{equation} \label{eq:trivialeff}
% \mathbf{B} = \begin{bmatrix} \mathbf{A} \\ \mathbf{A} \\ \vdots \\ \mathbf{A}\end{bmatrix}\end{equation}
In this case, their error expression becomes $\text{Err}_{\mathcal{F}}(\mathbf{B}) = mk - \frac{\delta^2(n-s)}{\delta + (n-s-1)\lambda}$ (follows from Theorem \ref{thm:random_diag}).
% \begin{equation}\label{eq:trivialefferr}
%   \text{Err}_{\mathcal{F}}(\mathbf{B}) = mk - \frac{\delta^2(n-s)}{\delta + (n-s-1)\lambda}  .
% \end{equation}
Our construction reduces the expected error to substantially below this.
\end{remark}
% For simplicity, we will refer to $\mathbf{A}$ as the incidence matrix of a $(v,b,p,r,\lambda)$-BIBD.  

% In this work, we wish to design encoding matrix for the case $m \ge 1$ while keeping the data assignment the same as the case $m = 1$. A trivial way is to stack any assignment matrix $\mathbf{A}$ vertically $m$ times to obtain the encoding matrix $\mathbf{B}$, i.e., 
%  We propose a scheme that uses a probabilistic approach and achieves an upper bound on the expected approximation error that is significantly lower than the trivial scheme, when $\mathbf{A} = \mathbf{E}$. Next we discuss the construction of our first scheme.  




% \aditya{Leverage the defn in the structured matrix section and reword below.}
% \sifat{done}
% \begin{defn}
% A strongly regular graph with $n$ vertices is a regular graph with parameters $\lambda$ and $\mu$ such that any pair of adjacent vertices have $\lambda$ common neighbors and any pair of non-adjacent vertices have $\mu$ common neighbors. If the graph has degree $\delta$, we denote the graph by $(n, \delta, \lambda, \mu)$-SRG.
% \end{defn}


\begin{cor}
\label{cor:SRG}
Suppose that the assignment matrix $\mathbf{A}$ is the adjacency matrix of a $(n, \delta, \lambda, \mu)$ SRG with $0 <\delta \ne \mu$. 
% Let $\mathbf{A} = \mathbf{A}_G$. Consequently, $\mathbf{A}$ is an $(n,n, \delta, \delta)$ assignment matrix. 
Then, for the construction as in $\eqref{eq:effAD1AD2}$, with $\epsilon  = 0$, the expected error can be upper bounded as follows.
\begin{equation}
\mathbb{E}[\text{Err}_{\mathcal{F}}(\mathbf{B})] \le mk -   \frac{m\delta^2(n-s)} {(m\delta - \mu)+ \mu(n-s) + (\lambda - \mu)\theta},
\end{equation}
where
\begin{align*}
    \theta = \begin{cases}
        \delta & \text{~if~} \lambda \geq \mu,\\
        \frac{1}{2} \left[ (\lambda - \mu) - \sqrt{(\lambda - \mu)^2 + 4 (\delta - \mu)} \right] & \text{~otherwise.}
    \end{cases}
\end{align*}

\end{cor}

\begin{proof}%[Proof of Corollary \ref{cor:SRG}] 
% Since for an $(n, \delta, \lambda, \mu)$ strongly regular graph any pair of adjacent vertices have $\lambda$ common neighbors and any pair of non-adjacent vertices have $\mu$ common neighbors it can be shown (see\cite{Brouwer_SRG}) that the adjacency matrix $\bfA$ satisfies
% \[ \bfA^2 = \delta \bf{I}_n + \lambda \bfA + \mu (\mathbf{J}_n-\bf{I}_n-\bfA). \]
% Consequently, 
% \begin{align*}
% \mathbf{\Delta}_{\mathcal{F}} &= \bfA_{\mathcal{F}}^T \bfA_{\mathcal{F}} \\
% %&= \delta \bfI_{n-s} + \lambda \bfA(\mathcal{F}, \mathcal{F}) + \mu (\mathbf{J} _{n-s} - \bfI_{n-s} - \bfA(\mathcal{F}, \mathcal{F}))\\
% &= (\delta-\mu) \bfI_{n-s} + \mu \mathbf{J} _{n-s} + (\lambda - \mu) \bfA(\mathcal{F},\mathcal{F}).
% \end{align*}

It is known that $\mathbf{A}$ has only three distinct eigenvalues given by $\delta , \tilde{r}, \tilde{s}$ \cite{Brouwer_SRG} where 
\[ 
\tilde{r} = \frac{1}{2}[ \lambda - \mu + \sqrt{(\lambda - \mu)^2 + 4(\delta - \mu)}],\]
and \[ 
\tilde{s} =  \frac{1}{2}[ \lambda - \mu - \sqrt{(\lambda - \mu)^2 + 4(\delta - \mu)}].\]

Since $0<\delta \ne \mu$, all the eigenvalues of $\mathbf{A}$ are nonzero.
Consequently, $\mathbf{A}^T\mathbf{A} = \mathbf{A}^2$ will have all strictly positive eigenvalues. Since $\mathbf{\Delta}_{\mathcal{F}} = \bfA_{\mathcal{F}}^T \bfA_{\mathcal{F}} = (\mathbf{A}^T\mathbf{A})(\mathcal{F}, \mathcal{F})$, by Cauchy's interlacing theorem (Theorem 4.3.17 in \cite{horn2012matrix}),  $\mathbf{\Delta}_{\mathcal{F}}$ will have all  strictly positive eigenvalues. Since $\mathbf{\Delta}_{\mathcal{F}}$ is symmetric and all the eigenvalues are strictly positive, it is positive definite. Since by construction $\tilde{\bfD}_j$ is invertible, $\tilde{\bfD}_j \mathbf{\Delta}_{\cF} \tilde{\bfD}_j$ is positive definite for each $j \in [m]$. Since the sum of positive definite matrices is positive definite, it follows that $\mathbf{B}_{\mathcal{F}}^T\mathbf{B}_{\mathcal{F}}  =  \sum_{j = 1}^m \tilde{\bfD}_j\mathbf{\Delta}_{\mathcal{F}}\tilde{\bfD}_j$ is positive definite and is hence invertible. By \eqref{eq: srg_prop},
% \cite{Brouwer_SRG}, 
% \begin{align*} \bfA^2 &= \delta \bfI_n + \lambda \bfA + \mu (\mathbf{J}_n-\bfI_n-\bfA)\\
% &= (\delta - \mu) \mathbf{I}_n + \mu \mathbf{J}_n + (\lambda - \mu) \mathbf{A}
% \end{align*} 
\begin{align*}
\mathbf{\Delta}_{\mathcal{F}} &= \mathbf{A}_{\mathcal{F}}^T \mathbf{A}_{\mathcal{F}}\\
& = (\delta - \mu) \mathbf{I}_{n-s} + \mu \mathbf{J}_{n-s} + (\lambda - \mu) \mathbf{A}(\mathcal{F}, \mathcal{F}).
\end{align*}

Therefore, %\aditya{Fix this proof based on the revised notation for the diagonal matrix corresponding to $\Delta_\cF$ I pointed out earlier}
%\sifat{done}
\begin{align*}
\mathbf{\Delta}^{(\text{diag})}_{\mathcal{F}} & = \delta \mathbf{I}_{n-s}, \text{~and}\\
\mathbf{K}& = \mathbf{\Delta}_{\mathcal{F}} + (m-1) \mathbf{\Delta}^{(\text{diag})}_{\mathcal{F}} = (m\delta-\mu) \bfI_{n-s} + \mu \mathbf{J} _{n-s} + (\lambda - \mu) \bfA(\mathcal{F},\mathcal{F}).
\end{align*}

% As noted in Remark~\ref{rem:c_equals_one}, when $\epsilon  = 0$, $c = 1$. Consequently,
% \begin{align*}
% \mathbf{K}& = \mathbf{\Delta}_{\mathcal{F}} + (m-1) \mathbf{\Delta}^{(\text{diag})}_{\mathcal{F}}\\
% &= (m\delta-\mu) \bfI_{n-s} + \mu \mathbf{J} _{n-s} + (\lambda - \mu) \bfA(\mathcal{F},\mathcal{F}).
% \end{align*}

Since both $\mathbf{\Delta}_{\mathcal{F}}$ and $\mathbf{\Delta}^{(\text{diag})}_{\mathcal{F}}$ are positive definite, it follows that $\mathbf{K}$ is positive definite and hence invertible. 
% Here, $\bfA_{\mathcal{F} \times \mathcal{F}}$ is the adjacency matrix of the induced subgraph with vertex set $\mathcal{F}$.
Note that the eigenvalues of $(m\delta-\mu) \bfI_{n-s} + \mu \mathbf{J} _{n-s}$ are $m\delta-\mu + \mu (n-s)$ (with multiplicity 1) and $m\delta-\mu$ (with multiplicity $n-s-1$). Also, by Cauchy's interlacing theorem, the maximum eigenvalue of $\mathbf{A}(\mathcal{F}, \mathcal{F})$, $\lambda_{\text{max}}(\mathbf{A}(\mathcal{F}, \mathcal{F})) \le \delta$ and the minimum eigenvalue of $\mathbf{A}(\mathcal{F}, \mathcal{F})$, $\lambda_{\text{min}}(\mathbf{A}(\mathcal{F}, \mathcal{F})) \ge \tilde{s}$. Now, by applying Weyl's inequality (Theorem 4.3.1 in \cite{horn2012matrix}) to $(m\delta-\mu) \bfI_{n-s} + \mu \mathbf{J} _{n-s}$ and $(\lambda - \mu) \bfA(\mathcal{F},\mathcal{F})$, we get
\begin{align*}
    \lambda_{\max} (\mathbf{K}) &\leq m\delta-\mu + \mu (n-s) + (\lambda - \mu) \theta \label{eq:upper_bd_lmax};
\end{align*}
where $\theta = \delta$ if $\lambda \ge \mu$ and otherwise $\theta = \tilde{s} = \frac{1}{2}[ \lambda - \mu - \sqrt{(\lambda - \mu)^2 + 4(\delta - \mu)}]$

Consequently, 
\[\lambda_{\text{min}}(\mathbf{K}^{-1}) \ge \frac{1}{ m\delta-\mu + \mu (n-s) + (\lambda - \mu) \theta}. \]

Note that, $\mathbf{A}_{\mathcal{F}}^T\mathbf{1}_k = \delta \mathbf{1}_{n-s}$. Let $ \mathbf{u} := \frac{1}{\sqrt{n-s}}\mathbf{1}_{n-s}$ so that $\|\mathbf{u}\|_2 = 1$. Then, 
 $\mathbf{u}^T \mathbf{K}^{-1} \mathbf{u} \ge \lambda_{\text{min}}(\mathbf{K}^{-1})$ by Rayleigh Quotient Theorem (see Theorem 4.2.2 in \cite{horn2012matrix}).
 By Theorem \ref{thm:random_diag}, 

\begin{align*}
\mathbb{E}[\text{Err}_{\mathcal{F}}(\mathbf{B})] 
% &\le  mk - m\mathbf{1}_k^T \bfA_{\cF} \left(\mathbf{\Delta}_{\mathcal{F}} + (m-1)\mathbf{\Delta}^{(\text{diag})}_{\mathcal{F}}\right)^{-1} \bfA_{\cF}^T\mathbf{1}_k\\
& \le  mk - m\delta\mathbf{1}_{n-s}^T (\mathbf{K})^{-1} \delta \mathbf{1}_{n-s}\\
& = mk - m\delta^2(n-s)\mathbf{u}^T(\mathbf{K})^{-1}\mathbf{u}\\
& \le mk - m\delta^2(n-s)\lambda_{\text{min}}(\mathbf{K}^{-1})\\
& \le mk - \frac{m\delta^2(n-s)}{m\delta-\mu + \mu (n-s) + (\lambda - \mu) \theta}. 
\end{align*}
\end{proof}



A desirable property of a scheme is that it recovers the exact gradient when there are no stragglers, i.e., when $s=0$. The construction in  \eqref{eq:effAD1AD2} with assignment matrices from BIBDs and SRGs has this property only when $m=1$. Indeed, exact recovery requires the existence of a decoding matrix $\mathbf{R}$ such that $\mathbf{B}\mathbf{R}=\mathbf{F}$. From the block structures of $\mathbf{B}$ and $\mathbf{F}$, for every $i \in [m]$, exact recovery requires
$
\mathbf{A}\mathbf{D}_i\mathbf{r}_i=\mathbf{1}_k
$ and $
\mathbf{A}\mathbf{D}_j\mathbf{r}_i=\mathbf{0}_k
$ for every $j\neq i$, $j\in [m]$. For the BIBD construction, $\mathbf{A}$ has full column rank, whereas for the SRG construction with $\delta\neq\mu$, $\mathbf{A}$ is invertible. Since each $\mathbf{D}_j$ is invertible, every matrix $\mathbf{A}\mathbf{D}_j$ has a trivial null space. Thus, when $m>1$, the condition $
\mathbf{A}\mathbf{D}_j\mathbf{r}_i=\mathbf{0}_k
$ for some $j\neq i$ implies $\mathbf{r}_i=\mathbf{0}_n$, which contradicts
$
\mathbf{A}\mathbf{D}_i\mathbf{r}_i=\mathbf{1}_k.
$ Thus, although the approximation error is significantly reduced compared to the naive scheme for general values of $s$, it remains strictly positive when $s=0$ and $m>1$.
% PREV

% A desirable property of a scheme is that it recovers the exact gradient
% when there are no stragglers, i.e., when $s=0$. The construction in \eqref{eq:effAD1AD2}
% with assignment matrices from BIBDs and SRGs has this property only when
% $m=1$. Indeed, when $m>1$, although the error is significantly reduced
% compared to the naive scheme for general values of $s$, there remains
% nonzero error even when $s=0$.

We address this issue by considering assignment matrices given by the
bi-adjacency matrix of a $(k,m,\delta)$ coset bipartite graph.
A sufficient condition for exact gradient recovery in the absence of
stragglers is that the encoding matrix $\bfB$ be invertible.
The following proposition shows that this condition holds for coset
bipartite graphs with appropriate parameters.






% A desirable property of a scheme is that it recovers the exact gradient when there are no stragglers, i.e., when $s=0$. The construction as in $\eqref{eq:effAD1AD2}$ with assignment matrices from BIBDs and SRGs has this property only in the situation when $m=1$. Indeed, when $m > 1$, even though we have significantly lower error as compared to the naive scheme for general values of $s$, there is non-zero error even when $s=0$. Our following result addresses this issue where the assignment matrix is  the bi-adjacency matrix of a $(k, m, \delta)$ coset bipartite graph. To guarantee exact gradient recovery in the absence of stragglers,
% it is necessary that the encoding matrix $\bfB$ be invertible. The following proposition shows that this property holds for coset bipartite graphs with specific parameters.

% The constructed encoding matrix in this case is invertible, thereby achieving zero error for $s = 0$.   



% \aditya{Define the random ensemble of directed graph}

% Our next result involves random $\delta$-regular directed graphs. We first define the ensemble of random  $\delta$-regular directed graphs, generated by the configuration model \cite{bender1978asymptotic}. The configuration model generates a random $\delta$-regular directed graph with $n$ vertices as follows: {\it (i)} Associate each vertex $i \in [n]$ with a set $F_i$ of $\delta$ distinct points such that $F_i \cap F_j = \emptyset$ for any $i, j \in [n]$. Consequently, there are $|\cup_{i \in [n]} F_i| = n\delta$ points in total.
% {\it (ii)}  Select a permutation $\sigma$ of the $n\delta$ points uniformly at random. 
% {\it (iii)}  For $i,j \in [n]$, there is a directed edge from $i$ to $j$ if for $u \in F_i$, $\sigma(u) \in F_j$.

% Note that the configuration model can generate directed graphs with multiple edges. In our constructions, we will not use the adjacency matrices of such graphs since that will result in a non-uniform computation load. 
% % \sifat{Configuration model can generate a directed graph that has multiple edges. In this case, the adjacency matrix does not satisfy the properties of an assignment matrix as defined.}
% \begin{cor}
% \label{cor: bi_adj}
% Let $\mathbf{A}_1, \mathbf{A}_2, \dots, \mathbf{A}_m \in \mathbb{R}^{k \times k}$ be the adjacency matrices of random $\delta$-regular directed graphs generated by the configuration model. We construct the assignment matrix $\mathbf{A}$ such that
% \begin{equation}{\label{eq:dir_adj}}
% \mathbf{A} = \begin{bmatrix} \mathbf{A}_1 &\mathbf{A}_2 &\dots &\mathbf{A}_m \end{bmatrix}.
% \end{equation}
% % Consequently $\mathbf{A}$ is a $(mk, k, \delta, m\delta)$ assignment matrix. 
% Then, for the construction as in $\eqref{eq:effAD1AD2}$,
% % For $i \in[m]$, let $\bfD_i \in \mathbb{R}^{n \times n}$ be a diagonal matrix whose entries are i.i.d uniform in the range $[-1-a,-1+a] \cup [1-a, 1+a]$, where $a \in \mathbb{R}$. Consider the construction as in $\eqref{eq:effAD1AD2}$. For a given non-straggler set $\mathcal{F}$ of size $n-s$, suppose that $\mathbf{\Delta}_{\mathcal{F}} := \mathbf{A}_{\mathcal{F}}^T\mathbf{A}_{\mathcal{F}} \in \mathbb{R}^{n-s \times n-s}$ be invertible. 
% the expected error can be upper bounded as
% \begin{equation}
% \mathbb{E}[\text{Err}_{\mathcal{F}}(\mathbf{B}] \le mk - \frac{m\delta^2(n-s)}{m\delta^2 + c(m-1)\delta},
% \end{equation}
% where the expectation is over the randomness of the diagonal matrices $\bfD_i$, for $i \in [m]$ (not over the randomness in the construction of the directed graph).
% % Here, $c = \mathbb{E}[X^2]\mathbb{E}[\frac{1}{X^2}] = \frac{1 + \frac{a^2}{3}}{1-a^2}$, where $X$ is uniformly distributed random variable in the range $[-1-a,-1+a] \cup [1-a, 1+a]$ for some $a \in \mathbb{R}$.
% \end{cor}
% \begin{proof}
% % It is easy to see from the proof of lemma \ref{lemma:random_diag} that, 
% % \[ \text{Err}_{\mathcal{F}}(\mathbf{B}) \le mk - m\mathbf{1}_k^T \bfA_{\cF} \left(\tilde{\mathbf{K}}\right)^{-1} \bfA_{\cF}^T\mathbf{1}_k.
% % \]
% Let $n=mk$ and for each $i \in [m]$, let $\mathbf{X}_i \in \mathbb{R}^{n \times n}$ be a diagonal matrix such that
% \[ \mathbf{X}_i = \begin{bmatrix}
% x^{(i)}_{1} & 0   & \cdots & 0 \\
% 0   & x^{(i)}_{2} & \cdots & 0 \\
% \vdots & \vdots & \ddots & \vdots \\
% 0 & 0 & \cdots & x^{(i)}_{n}
% \end{bmatrix}\]
% Where $x^{(i)}_{j} \in \mathbb{R}$ for each $j \in [n]$. Let $X := (\mathbf{X}_1, \dots \mathbf{X}_m)$. Let $\mathbf{B}_X$  be such that \[ \mathbf{B}_X^T := [ \bfX_1 \bfA^T ~|~ \bfX_2 \bfA^T ~|~ \dots ~|~ \bfX_m \bfA^T].\]
%  Let $\text{det}(\mathbf{B}_{X})$ denote the determinant of $\mathbf{B}_X$. Note that the determinant of $\mathbf{B}_{X} \in \mathbb{R}^{n \times n}$ is a polynomial where the variables are the entries of the diagonal matrices $\mathbf{X}_i$s and with a degree at most $n$. 
%  % \sifat{Define identically zero? Define determinant? How's the degree at most $n$?}
%  % \begin{equation} \label{eq:det}
%  % \text{det}(\mathbf{B}) = \sum_{\sigma \in S_n} sgn(\sigma) \mathbf{B}(1 , \sigma(1))\mathbf{B}(2 , \sigma(2))\dots\mathbf{B}(n , \sigma(n)).
%  % \end{equation}
%  % It is immediate from \eqref{eq:det}, the degree of the determinant is at most $n$.
%  We first prove that $\text{det}(\mathbf{B}_X)$ is not an identically zero polynomial with high probability. For $i \in [m]$, let \[\mathbf{X}_i' = \begin{bmatrix}
% \mathbf{0}_{k \times k} & \cdots & \mathbf{0}_{k \times k} & \cdots & \mathbf{0}_{k \times k} \\
% \vdots  & \ddots & \vdots & \ddots & \vdots \\
% \mathbf{0}_{k \times k} & \cdots & \mathbf{I}_k & \cdots & \mathbf{0}_{k \times k} \\
% \vdots  & \ddots & \vdots & \ddots & \vdots \\
% \mathbf{0}_{k \times k} & \cdots & \mathbf{0}_{k \times k} & \cdots & \mathbf{0}_{k \times k}
% \end{bmatrix}.\] 

% Let $X' = (\mathbf{X}_1', \dots, \mathbf{X}_m')$.
% Consequently,  
% \[
% \mathbf{B}_{X'} = \begin{bmatrix}
% \mathbf{A}_1 & \mathbf{0}_{k \times k}   & \cdots & \mathbf{0}_{k \times k} \\
% \mathbf{0}_{k \times k} & \mathbf{A}_2   & \cdots & \mathbf{0}_{k \times k} \\
% \vdots  & \vdots & \ddots   & \vdots \\
% \mathbf{0}_{k \times k} & \mathbf{0}_{k \times k}   & \cdots & \mathbf{A}_m \\
% \end{bmatrix}.
% \]
% Now, $\text{det}(\mathbf{B}_{X'}) = \text{det}(\mathbf{A}_1) \dots \text{det}(\mathbf{A}_m)$. It is shown in \cite{huang2018invertibilityadjacencymatricesrandom} that for each $i \in [m]$, 
% \[ \mathbb{P}(\text{det}({\mathbf{A}_i}) = 0) = o(1).\]

% Thus, \[\mathbb{P}(\text{det}({\mathbf{B}_{X'}}) = 0) = m.o(1).\]

% Since $m << k$ in practice, $\mathbb{P}(\text{det}(\mathbf{B}_{X'}) = 0) = o(1)$. Consequently, 
% \[\mathbb{P}(\text{det}({\mathbf{B}_X}) \  \text{is identically zero}) \le o(1) .\]

% Note that the entries of the diagonal matrices $\mathbf{D}_i$ are uniformly distributed over the set $S$, and $|S| = 2p$. Also, by Leibniz Formula \cite{horn2012matrix}, we have 


% \begin{equation}\label{eq:det}
% \text{det}(\mathbf{B}_X) = \sum_{ \tau \in S_n} \text{sgn}(\tau) \prod_{i = 1}^n \mathbf{B}_X(i, \tau(i)), \end{equation}

% where $S_n$ is the permutation group of $[n]$, $\text{sgn}$ is the sign function of $S_n$, which returns $+1$ and $-1$ for even and odd functions, respectively. Since each term in the sum in \eqref{eq:det} contains product of at most $n$ variables in $\bfB_X$, the degree of $\text{det}(\mathbf{B}_X)$ is at most $n$.
%  It follows from Schwartz-Zippel lemma \cite{zippel1979probabilistic} that, for any encoding matrix $\mathbf{B}$,   
% \[ \mathbb{P}(\text{det}(\mathbf{B}) \ne 0 | \text{det}(\mathbf{B}_X) \  \text{is not identically zero}) \ge 1 - \frac{n}{2p}.\]
% % \aditya{need p larger than n for it to be useful and n large}
% We have, 
% \begin{align*}
% &\mathbb{P}(\text{det}(\mathbf{B}) \ne 0) \\&= \mathbb{P}(\text{det}(\mathbf{B}) \ne 0 | \text{det}(\mathbf{B}_X) \  \text{is not identically zero}) \mathbb{P}(\text{det}(\mathbf{B}_X) \  \text{is not identically zero}) + \\ &\quad \quad \quad \quad \mathbb{P}(\text{det}(\mathbf{B}) \ne 0 | \text{det}(\mathbf{B}_X) \  \text{is identically zero}) \mathbb{P}(\text{det}(\mathbf{B}_X) \text{is identically zero}) \\
% & \ge (1- \frac{n}{2p})(1-o(1)) + 0\\
% & = 1 - \frac{n}{2p} - o(1) + \frac{n}{2p} o(1).
% \end{align*}

% Consequently, when $n$ is sufficiently large and $p>>n$, then the encoding matrix $\mathbf{B}$ is invertible  with high probability. Consequently, $\mathbf{B}_{\mathcal{F}}^T\mathbf{B}_{\mathcal{F}}$ is invertible and we can apply Lemma $\ref{lemma:random_diag}$. Then, we will further upper bound the expected error by the maximum eigenvalue of $\mathbf{K}$.

% We now calculate the maximum eigenvalue of $\mathbf{\Delta}_{[n]} + c(m-1)\mathbf{\Delta}^{(\text{diag})}_{[n]}$. Note that $\mathbf{\Delta}_{[n]} = \bfA^T \bfA$ by definition and since each column of $\bfA$ has $\delta$ nonzero entries that are $1$, $\bfA^T \mathbf{1}_k = \delta \mathbf{1}_n$ and $\mathbf{\Delta}^{(\text{diag})}_{[n]} = \delta \mathbf{I}_n$. Also, $\bfA \mathbf{1}_n = m\delta \mathbf{1}_k$. We have, 

% \begin{align*}
% (\mathbf{\Delta}_{[n]} + c(m-1)\mathbf{\Delta}^{(\text{diag})}_{[n]})\mathbf{1}_n &= (\mathbf{A}^T\mathbf{A} + c(m-1)\delta \mathbf{I}_n) \mathbf{1}_n\\
% &= \mathbf{A}^T\mathbf{A} \mathbf{1}_n + c(m-1) \delta \mathbf{1}_n\\
% &= \mathbf{A}^T m\delta \mathbf{1}_k + c(m-1) \delta \mathbf{1}_n\\
% & = m\delta^2 \mathbf{1}_n + c(m-1)\delta \mathbf{1}_n\\
% &= (m\delta^2 + c(m-1) \delta) \mathbf{1}_n
% \end{align*}

% So, $\mathbf{1}_n$ is an eigenvector of $\mathbf{\Delta}_{[n]} + c(m-1)\mathbf{\Delta}^{(\text{diag})}_{[n]}$ with eigenvalue $m\delta^2 + c(m-1)\delta$. Since $\mathbf{\Delta}_{[n]} + c(m-1)\mathbf{\Delta}^{(\text{diag})}_{[n]}$ is a matrix with non-negative entries, by Theorem 8.3.4 in \cite{horn2012matrix} it follows that $(m\delta^2 + c(m-1))$ is the maximum eigenvalue of $\mathbf{\Delta}_{[n]} + c(m-1)\mathbf{\Delta}^{(\text{diag})}_{[n]}$, i.e.,
% % Now, ${\mathbf{K}} = \mathbf{\Delta}_{\mathcal{F}} + c(m-1)\mathbf{\Delta}^{(\text{diag})}_{\mathcal{F}}$. 
% % and $\tilde{\mathbf{\Delta}}_{\mathcal{F}}(i,j) = 0$ for $i \ne j$ and 
% % for $i \in [n-s]$.
% $\lambda_{max}\left(\mathbf{\Delta}_{[n]} + c(m-1)\mathbf{\Delta}^{(\text{diag})}_{[n]}\right) = m\delta^2 + c(m-1)\delta$. 
% % \aditya{include proof of it}
% % Since $\mathbf{\Delta}_{\mathcal{F}}$ is a principal submatrix of $\mathbf{\Delta}_{[n]} $, 
% By Cauchy's interlacing theorem (Theorem 4.3.17 in \cite{horn2012matrix}), 
% \begin{align*}
% \lambda_{max}\left(\mathbf{\Delta}_{\mathcal{F}} + c(m-1)\mathbf{\Delta}^{(\text{diag})}_{\mathcal{F}}\right) \le m\delta^2 + c(m-1)\delta.
% \end{align*}

% Consequently, 
% % \aditya{typo}
% \begin{align*}
% \lambda_{min}\left(\left(\mathbf{\Delta}_{\mathcal{F}} + c(m-1)\mathbf{\Delta}^{(\text{diag})}_{\mathcal{F}}\right)^{-1}\right) \ge \frac{1}{m\delta^2 + c(m-1)\delta}.
% \end{align*}

% Now, 
% \begin{align*}
% \text{Err}_{\mathcal{F}}(\mathbf{B}) &\le mk - m\mathbf{1}_k^T \bfA_{\cF} \left({\mathbf{K}}\right)^{-1} \bfA_{\cF}^T\mathbf{1}_k\\
% & = mk - m\delta^2(n-s)\frac{\mathbf{1}_{n-s}^T}{\sqrt{n-s}}\left({\mathbf{K}}\right)^{-1} \frac{\mathbf{1}_{n-s}}{\sqrt{n-s}}\\
% &\le  mk - m\delta^2(n-s)\lambda_{min}\left({\mathbf{K}}\right)^{-1}\\
% & \le mk - \frac{m\delta^2(n-s)}{m\delta^2 + c(m-1)\delta}.
% \end{align*} 

% Here, the second equality follows from  Theorem 4.2.2 in \cite{horn2012matrix}.
% \end{proof}

% \begin{remark}
% Note that since the encoding matrix in Corollary \ref{cor: bi_adj} is invertible with high probability, the error is zero when $s = 0$ in this case. This is not the case when the assignment matrix is the transpose of a BIBD incidence matrix or the adjacency matrix of a strongly regular graph.    
% \end{remark}

\begin{prop}{\label{prop: cos_inv}}
Suppose that the assignment matrix $\bfA$ is the bi-adjacency matrix of a $(k, m, \delta)$ coset bipartite graph such that $k = p^a$, where $p$ is a prime, $a \in \mathbb{Z}_{\ge 1}$ and $p \nmid \delta$. Then, for the construction as in $\eqref{eq:effAD1AD2}$ with $\epsilon > 0$, the encoding matrix $\bfB$ is almost surely invertible. 
\end{prop}


\begin{proof}
By construction, for $i \in \{0, \dots, k-1\}$ and $x \in \{0 , \dots, mk-1\}$, $\bfA(i+1, x+1) = 1$ if and only if $x \in i+ S$. By definition, for $y \in \{0, \dots, k-1\}$, $y \in i+S$ if and only if $y + tk \in i + S$ for any $t \in \{0, \dots, m-1\}$. Consequently, $\bfA(i+1, y+1) = 1$ if and only if $\bfA(i+1, tk + y+1) = 1$. Therefore, $\bfA$ can be written as $m$ equal column blocks such that 
\[ \bfA = \begin{bmatrix} \bfC & \dots & \bfC \end{bmatrix},\]
 where $\bfC \in \mathbb{R}^{k \times k}$. 

Now we show that $\bfC$ is circulant. For $i,j \in \{0, \dots, k-1\}$, by definition $\bfC(i+1, j+1) = 1$ if and only if $j \in i + S$, which implies that $j = i + u \bmod mk$ with $u \in S$. This implies that
\begin{align*}
    j-i &= u \bmod mk\\
    j-i &= u \bmod k \text{~~(since $k ~|~ mk$)}.
\end{align*}
Next, since $u \in S$, there exists $b \in B$ and $h \in H$ such that $u = b+h$ and since $h$ is a multiple of $k$, we have $u = b \bmod k$. We conclude that $(j-i) \bmod k \in B$. The reverse implication is easy to see. We conclude that $\bfC(i+1,j+1)$ only depends on $(j-i) \bmod k$. Therefore, each row of $\bfC$ is obtained from the preceding row by a cyclic shift of one position to the right. Thus, $\bfC$ is a circulant matrix.
% Suppose $\bfA(i+1, x +1) = 1$. So, $ x \in i +S$. So, there exists $u \in S$ such that $ x = (i+u) \bmod mk$. Consequently, $x \equiv i+u \pmod {mk}$. So, $ x +1 \equiv (i+u+1)  \pmod{mk}$. Thus, $(x +1 ) \bmod mk = (i+u+1) \bmod mk$. Therefore, $(x+1) \bmod mk \in (i+1)+S$. Since $ k + S = S$, $(i + 1) +S = ((i+1) \bmod k) + S$. Consequently, $(x+1) \bmod mk \in ((i+1) \bmod k) +S$ and thus $\bfA(((i+1) \bmod k) + 1 , ((x+1) \bmod mk) + 1) = 1$. Similarly, it can be shown that if $\bfA(((i+1) \bmod k) + 1 , ((x+1) \bmod mk) + 1) = 1$, then  $\bfA(i+1, x +1) = 1$. Thus, each row of $\bfA$ is a cyclic shift of the previous row by one position. Since $\bfA = \begin{bmatrix} \bfC & \dots & \bfC \end{bmatrix}$, consequently  $\bfC$ is a circulant matrix. \aditya{this does not follow directly}
%  % Also, note that $x \in i + S$ if and only if $x +1 \in i +1 + S$. Therefore, each row of $\bfC$ is a cyclic shift of the previous row, so $\bfC$ is a circulant matrix. 

Note that in this case $n = mk$. For $t \in \{1,\dots,m \}$, let
\[
\bfX_t := \mathrm{diag}\big(x^{(t)}_1,\dots,x^{(t)}_n\big)
\]
be diagonal matrices whose entries are independent indeterminates. Define $\bfB_{\bfX}$ such that
\begin{equation}\label{eq:poly}
%\mathbf{B} = \begin{bmatrix} \bfA\bfD_1 \\ \bfA\bfD_2 \\ \vdots \\ \bfA\bfD_m \end{bmatrix}
\mathbf{B}_{\bfX}^T := [ \bfX_1 \bfA^T ~|~ \bfX_2 \bfA^T ~|~ \dots ~|~ \bfX_m \bfA^T].
\end{equation}

Let $\text{det}(\bfM)$ denote the determinant of a matrix $\bfM$. We will show that $\text{det}(\bfB_X)$ is not an identically zero polynomial, so that the Schwartz-Zippel lemma \cite{zippel1979probabilistic} can be applied to conclude that the encoding matrix $\mathbf{B}$ is invertible almost surely. For $i \in [m]$,
let $\mathbf{e}_i \in \mathbb{R}^m$ be the $i$-th standard basis vector.
Define
\[
\bfX_i' := (\mathbf{e}_i \mathbf{e}_i^T) \otimes \bfI_k \in \mathbb{R}^{mk\times mk},
\]
where $\otimes$ denotes the Kronecker product. Let 
\begin{equation}
%\mathbf{B} = \begin{bmatrix} \bfA\bfD_1 \\ \bfA\bfD_2 \\ \vdots \\ \bfA\bfD_m \end{bmatrix}
\mathbf{B}_{\bfX'}^T = [ \bfX_1' \bfA^T ~|~ \bfX_2' \bfA^T ~|~ \dots ~|~ \bfX_m' \bfA^T].
\end{equation}

Since $\bfX_i'$ selects the $i$-th $k\times k$ block and each block of $\bfA$
equals $\bfC$, we have $\mathbf{B}_{\bfX'} = \bfI_m \otimes \bfC$. So, the determinant of $\mathbf{B}_{\bfX'}$, $\text{det}(\mathbf{B}_{\bfX'}) = (\text{det} (\bfC))^m$.
We show that the $k\times k$ circulant matrix $\bfC$ is invertible when
$k=p^{a}$ for a prime $p$ and $p \nmid \delta$. Let the first row of $\bfC$ be the
indicator of a set $T\subseteq \{0, \dots, k-1 \}$ with $|T|=\delta$, where index $j \in \{0, \dots,k-1\}$ corresponds to the $(j+1)$-th column of $\bfC$ following our convention on the labeling of the group elements in $\mathbb{Z}_{k}$. Define the mask
polynomial,
\[
p_T(x)\;:=\;\sum_{j\in T} x^{j} \in \mathbb{Z}[x]~~ (\text{ring of polynomials with coefficients in $\mathbb{Z}$}).
\]
Let $\omega:=e^{-2\pi i/k}$. It is well known that for a circulant matrix the eigenvalues are
\begin{equation}\label{eq:eigs-circ}
\lambda_r \;=\; p_T(\omega^{r}) \;=\; \sum_{j\in T}\omega^{rj}, \qquad r=0,1,\dots,k-1.
\end{equation}
In particular, $\lambda_0=p_T(1)=\delta\neq 0$. Assume for contradiction that $\bfC$ is singular. Then $\lambda_r=0$ for some
$r\in\{1,\dots,k-1\}$, i.e.,
\begin{equation}\label{eq:vanish}
\sum_{j\in T} (\omega^{r})^{j}=0.
\end{equation}
Let $q:=\gcd(k,r)$ and set $t=\frac{k}{q}$. Then $\omega^{r}$ is a primitive $t$-th root
of unity. Since $k=p^{a}$ we have $t=p^{b}$ for some $1\le b\le a$. Let $\Phi_t(x)$ denote the $t$-th cyclotomic
polynomial, i.e., the minimal polynomial over $\mathbb{Q}$ of a primitive $t$-th root of unity. For $t=p^b$, we have $\Phi_{t}(x)=1+x^{p^{b-1}}+\cdots+x^{(p-1)p^{b-1}}$, hence
$\Phi_{t}(1)=p$.
Since $\omega^r$ is a primitive $t$-th root and $p_T(\omega^r)=0$, it follows that there exists $g(x) \in \mathbb{Q}[x]$ such that 
\[
\Phi_t(x)g(x) = p_T(x).
\]
Suppose $g(x)$ is a constant. Then, $g(x) = 1$, since $\Phi_t(x)$ and $p_T(x)$ are both monic polynomials. Now, $\delta = p_T(1) =\Phi_t(1)g(1) = p$, contradicting the
assumption $p\nmid \delta$. Thus, $g(x)$ cannot be a constant. Next, suppose that $g(x)$ is not a constant polynomial. Since the coefficients of $\Phi_t(x)$ are coprime, by Gauss's Lemma (see Theorem 17.14 in \cite{judson2021abstract}) there exists  $h(x)\in \mathbb{Z}[x]$ such that $\Phi_t(x)h(x) = p_T(x)$. So, $\delta = p_T(1) = \Phi_t(1) h(1) = p\cdot h(1)$, i.e., $p \mid \delta$ which contradicts the assumption $p \nmid \delta$. Hence $\lambda_r\neq 0$ for all $r$, so $\bfC$ has no zero eigenvalues and is invertible.  

% Since $\Phi_t(x)\in\mathbb{Z}[x]$ is primitive, Gauss's lemma (see \cite{DummitFoote}) implies that the divisibility
% holds in $\mathbb{Z}[x]$, i.e., there exists $h(x)\in\mathbb{Z}[x]$ such that
% \[
% p_T(x)=\Phi_t(x)h(x).
% \]
% Evaluating at $x=1$ yields
% \[
% \delta \;=\; p_T(1) \;=\; \Phi_t(1)\,h(1).
% \]
% For $t=p^b$, we have $\Phi_{p^b}(x)=1+x^{p^{b-1}}+\cdots+x^{(p-1)p^{b-1}}$, hence
% $\Phi_{p^b}(1)=p$. Therefore $\delta=p \cdot h(1)$, i.e., $p\mid \delta$, contradicting the
% assumption $p\nmid \delta$. 




% Equation~\eqref{eq:vanish-sum} is therefore a vanishing sum of $\delta$ (distinct)
% $t$-th roots of unity. By Theorem 3.3 in \cite{LamLeung2000}, if $t$ is a prime power,
% any vanishing sum of distinct $t$-th roots of unity has
% cardinality divisible by $p$. Hence \eqref{eq:vanish-sum} implies $p\mid \delta$, which is a contradiction. Therefore $\lambda_r\neq 0$ for all $r$, so by~\eqref{eq:eigs-circ} the circulant
% matrix $\bfC$ has no zero eigenvalues and is invertible. 
Consequently, $\text{det}(\mathbf{B}_{\bfX'}) \ne 0$ and  $\text{det}(\bfB_X)$ is not an identically zero polynomial. Also, since for the encoding matrix $\bfB$ the diagonal entries of $\bfD_i$ are i.i.d. continuous (as $\epsilon > 0)$, it follows from Schwartz-Zippel lemma \cite{zippel1979probabilistic} that, for the encoding matrix $\mathbf{B}$ generated by this construction, 
\begin{align*}
\mathbb{P}(\text{det}(\mathbf{B}) \ne 0) &= 1.
\end{align*} 


% \begin{equation}
% \mathbb{P}(\text{det}(\mathbf{B}) \ne 0 | \text{det}(\mathbf{B}_X) \not\equiv 0) = 1.
% \end{equation}
% Now, 
% \begin{align*}
% &\mathbb{P}(\text{det}(\mathbf{B}) \ne 0) \\&= \mathbb{P}(\text{det}(\mathbf{B}) \ne 0 | \text{det}(\mathbf{B}_X)  \not\equiv 0 ) \mathbb{P}(\text{det}(\mathbf{B}_X) \not\equiv 0) + \mathbb{P}(\text{det}(\mathbf{B}) \ne 0 | \text{det}(\mathbf{B}_X) \equiv 0 ) \mathbb{P}(\text{det}(\mathbf{B}_X) \equiv 0) \\
% & = 1 \cdot 1 + 0= 1
% \end{align*} 


%Thus, the encoding matrix $\bfB$ is almost surely invertible.
\end{proof}


We now turn to the error analysis. Using the invertibility of the encoding
matrix established in Proposition~\ref{prop: cos_inv}, we derive an
upper bound on the expected error.


%\aditya{So, I think we can use the idea you had for the next section, i.e., taking a m by m orthogonal matrix and choosing diagonal matrices by Kronecker products with all-ones. However, the catch I think is that you will need to use Hadamard matrices to get a bound that can be calculated based on just the cardinality of the surviving set.}



\begin{cor}  \label{cor:bipartite}
Suppose that the assignment matrix $\bfA$ is the bi-adjacency matrix of a $(k,m,\delta)$ coset bipartite graph such that $k = p^a$, where $p$ is a prime, $a \in \mathbb{Z}_{\ge 1}$ and $p \nmid \delta$. Then, for the construction as in $\eqref{eq:effAD1AD2}$ with $\epsilon > 0$,
% For $i \in[m]$, let $\bfD_i \in \mathbb{R}^{n \times n}$ be a diagonal matrix whose entries are i.i.d uniform in the range $[-1-a,-1+a] \cup [1-a, 1+a]$, where $a \in \mathbb{R}$. Consider the construction as in $\eqref{eq:effAD1AD2}$. For a given non-straggler set $\mathcal{F}$ of size $n-s$, suppose that $\mathbf{\Delta}_{\mathcal{F}} := \mathbf{A}_{\mathcal{F}}^T\mathbf{A}_{\mathcal{F}} \in \mathbb{R}^{n-s \times n-s}$ be invertible. 
the expected error can be upper bounded as
\begin{equation} \label{eq:AD1AD2err_Coset}
\mathbb{E}[\text{Err}_{\mathcal{F}}(\mathbf{B})] \le mk - \frac{m\delta^2(mk-s)}{m\delta^2 + c(m-1)\delta}.
\end{equation}
\end{cor}
\begin{proof}



% Let $\text{det}(\mathbf{B}_{\bfX})$ denote the determinant of $\mathbf{B}_X$. By Leibniz Formula \cite{horn2012matrix}, we have 
% \begin{equation}\label{eq:det}
% \text{det}(\mathbf{B}_X) = \sum_{ \tau \in S_n} \text{sgn}(\tau) \prod_{i = 1}^n \mathbf{B}_X(i, \tau(i)), \end{equation}
% where $S_n$ is the permutation group of $[n]$, $\text{sgn}$ is the sign function of $S_n$, which returns $+1$ and $-1$ for even and odd functions, respectively. Since each term in the sum in \eqref{eq:det} contains product of at most $n$ variables in $\bfB_X$, the degree of $\text{det}(\mathbf{B}_X)$ is at most $n$.



% Let $b \in B$. For $t \in [m]$ and $i \in \{0, \dots, k-1 \}$, define $c(t, i+1 ) = (i +b + (t-1)k) \bmod mk$. By construction, vertex $i + H \in L$ is connected to any $x \in i + S$. Note that $i+b+H \subseteq i+S$. Therefore, vertex $i +H \in L$ is connected to any $x \in i+b+H$. Note also that $c(t, i+1) \in i+b+H$. Thus, by definition of bi-adjacency matrix, $\bfA(i, c(t, i+1)) = 1$. It can be observed that all $c(t, i+1)$ are distinct. It follows from \eqref{eq:poly} and \eqref{eq:det} that the determinant expansion contains the monomial $\prod_{t, i+1} x_{c(t, i+1)}^{(t)}$ with coefficient $\pm 1$. Hence, $\text{det}(\mathbf{B}_{\bfX})$ is not identically zero, i.e., $\mathbb{P}(\text{det}(\bfB_{\bfX}) \not \equiv 0) = 1$. Also, since for the encoding matrix $\bfB$ the diagonal entries of $\bfD_i$ are i.i.d continuous, it follows from Schwartz-Zippel lemma \cite{zippel1979probabilistic} that, for any encoding matrix $\mathbf{B}$,  \sifat {this is incorrect.} 

% \begin{equation}
% \mathbb{P}(\text{det}(\mathbf{B}) \ne 0 | \text{det}(\mathbf{B}_X) \not\equiv 0) = 0.
% \end{equation}
% Now, 
% \begin{align*}
% &\mathbb{P}(\text{det}(\mathbf{B}) \ne 0) \\&= \mathbb{P}(\text{det}(\mathbf{B}) \ne 0 | \text{det}(\mathbf{B}_X)  \not\equiv 0 ) \mathbb{P}(\text{det}(\mathbf{B}_X) \not\equiv 0) + \mathbb{P}(\text{det}(\mathbf{B}) \ne 0 | \text{det}(\mathbf{B}_X) \equiv 0 ) \mathbb{P}(\text{det}(\mathbf{B}_X) \equiv 0) \\
% & = 0.1 + 0= 0
% \end{align*}

% \begin{equation}\label{eq:det}
% \text{det}(\mathbf{B}_X) = \sum_{ \tau \in S_n} \text{sgn}(\tau) \prod_{i = 1}^n \mathbf{B}_X(i, \tau(i)), \end{equation}









% If $\bfA$ is the bi-adjacency matrix of a $(k,m,\delta)$ coset bipartite graph, then $\mathbf{A} \mathbf{1}_n = m \delta \mathbf{1}_k$ and $\mathbf{A}^T \mathbf{1}_k = \delta \mathbf{1}_n$. Let $\mathbf{X}_i \in \mathbb{R}^{n \times n}$ be a diagonal matrix such that
% \[ \mathbf{X}_i = \begin{bmatrix}
% x^{(i)}_{1} & 0   & \cdots & 0 \\
% 0   & x^{(i)}_{2} & \cdots & 0 \\
% \vdots & \vdots & \ddots & \vdots \\
% 0 & 0 & \cdots & x^{(i)}_{n}
% \end{bmatrix}\]
% Where $x^{(i)}_{j} \in \mathbb{R}$ for each $j \in [n]$. Let $X := (\mathbf{X}_1, \dots \mathbf{X}_m)$. Let $\mathbf{B}_X$  be such that \[ \mathbf{B}_X^T := [ \bfX_1 \bfA^T ~|~ \bfX_2 \bfA^T ~|~ \dots ~|~ \bfX_m \bfA^T].\]
% By construction, $\text{det}(\bfB_X)$ is a not an identically zero polynomial.
By Proposition \ref{prop: cos_inv}, $\bfB$ is almost surely invertible. Consequently, for any $\cF \subseteq [n]$ of size $n-s$, $\mathbf{B}_{\mathcal{F}}^T\mathbf{B}_{\mathcal{F}}$ is invertible and thus we can apply Theorem $\ref{thm:random_diag}$. 
% By Lemma $\ref{lemma:random_diag}$,  $\mathbb{E}[\text{Err}_{\mathcal{F}}(\mathbf{B})] \le mk - m\mathbf{1}_k^T \bfA_{\cF} \left(\mathbf{K}\right)^{-1} \bfA_{\cF}^T\mathbf{1}_k $. Here $\mathbf{K} = \mathbf{\Delta}_{\mathcal{F}} + c(m-1)\mathbf{\Delta}^{(\text{diag})}_{\mathcal{F}}$, $\mathbf{\Delta}_{\mathcal{F}} = \mathbf{A}_{\mathcal{F}}^T\mathbf{A}_{\mathcal{F}}$ and  $\mathbf{\Delta}^{(\text{diag})}_{\mathcal{F}}$ is a diagonal matrix such that $\mathbf{\Delta}^{(\text{diag})}_{\mathcal{F}}(i,i) := \mathbf{\Delta}_{\mathcal{F}}(i,i)$ 
% % and $\tilde{\mathbf{\Delta}}_{\mathcal{F}}(i,j) = 0$ for $i \ne j$ and 
% for $i \in [n-s]$. 
Note that in this case $n = mk$. We now calculate the maximum eigenvalue of $\mathbf{\Delta}_{[n]} + c(m-1)\mathbf{\Delta}^{(\text{diag})}_{[n]}$. Note that $\mathbf{\Delta}_{[n]} = \bfA^T \bfA$ by definition and since each column of $\bfA$ has $\delta$ nonzero entries that are $1$, $\bfA^T \mathbf{1}_k = \delta \mathbf{1}_n$ and $\mathbf{\Delta}^{(\text{diag})}_{[n]} = \delta \mathbf{I}_n$. Also, $\bfA \mathbf{1}_n = m\delta \mathbf{1}_k$. We have, 

\begin{align*}
(\mathbf{\Delta}_{[n]} + c(m-1)\mathbf{\Delta}^{(\text{diag})}_{[n]})\mathbf{1}_n &= (\mathbf{A}^T\mathbf{A} + c(m-1)\delta \mathbf{I}_n) \mathbf{1}_n\\
&= \mathbf{A}^T\mathbf{A} \mathbf{1}_n + c(m-1) \delta \mathbf{1}_n\\
&= \mathbf{A}^T m\delta \mathbf{1}_k + c(m-1) \delta \mathbf{1}_n\\
& = m\delta^2 \mathbf{1}_n + c(m-1)\delta \mathbf{1}_n\\
&= (m\delta^2 + c(m-1) \delta) \mathbf{1}_n.
\end{align*}

So, $\mathbf{1}_n$ is an eigenvector of $\mathbf{\Delta}_{[n]} + c(m-1)\mathbf{\Delta}^{(\text{diag})}_{[n]}$ with eigenvalue $m\delta^2 + c(m-1)\delta$. Since $\mathbf{\Delta}_{[n]} + c(m-1)\mathbf{\Delta}^{(\text{diag})}_{[n]}$ is a matrix with non-negative entries, by a consequence of the Perron–Frobenius theorem for nonnegative matrices (see Theorem 8.3.4 in \cite{horn2012matrix}) %it follows that $m\delta^2 + c(m-1)\delta$ is the maximum eigenvalue of $\mathbf{\Delta}_{[n]} + c(m-1)\mathbf{\Delta}^{(\text{diag})}_{[n]}$, i.e.,
% Now, ${\mathbf{K}} = \mathbf{\Delta}_{\mathcal{F}} + c(m-1)\mathbf{\Delta}^{(\text{diag})}_{\mathcal{F}}$. 
% and $\tilde{\mathbf{\Delta}}_{\mathcal{F}}(i,j) = 0$ for $i \ne j$ and 
% for $i \in [n-s]$.
$\lambda_{max}\left(\mathbf{\Delta}_{[n]} + c(m-1)\mathbf{\Delta}^{(\text{diag})}_{[n]}\right) = m\delta^2 + c(m-1)\delta$. 
% \aditya{include proof of it}
% Since $\mathbf{\Delta}_{\mathcal{F}}$ is a principal submatrix of $\mathbf{\Delta}_{[n]} $, 
By Cauchy's interlacing theorem (Theorem 4.3.17 in \cite{horn2012matrix}), 
\begin{align*}
\lambda_{max}\left(\bfK = \mathbf{\Delta}_{\mathcal{F}} + c(m-1)\mathbf{\Delta}^{(\text{diag})}_{\mathcal{F}}\right) \le m\delta^2 + c(m-1)\delta.
\end{align*}

For $m>1$, since $\mathbf{\Delta}_{\cF}$ is positive semidefinite and $c(m-1) \mathbf{\Delta}^{(\text{diag})}_{\cF} = c(m-1) \delta \bfI_{n-s}$ is positive definite, $\bfK$ is positive definite. For $m=1$, $\mathbf K=\mathbf{\Delta}_{\mathcal F}$, which is
positive definite since $\mathbf A_{\mathcal F}$ has full column rank.
Hence, $\mathbf K$ is invertible. Consequently, 
% \aditya{typo}
\begin{equation}\label{eq:lambda_min}
\lambda_{min}\left(\bfK^{-1}\right) \ge \frac{1}{m\delta^2 + c(m-1)\delta}.
\end{equation}

By Theorem \ref{thm:random_diag}, 
\begin{align*}
\mathbb{E}[\text{Err}_{\mathcal{F}}(\mathbf{B})] &\le mk - m\mathbf{1}_k^T \bfA_{\cF} \left({\mathbf{K}}\right)^{-1} \bfA_{\cF}^T\mathbf{1}_k\\
& = mk - m\delta^2(n-s)\frac{\mathbf{1}_{n-s}^T}{\sqrt{n-s}}\left({\mathbf{K}}\right)^{-1} \frac{\mathbf{1}_{n-s}}{\sqrt{n-s}}\\
&\le  mk - m\delta^2(n-s)\lambda_{min}\left({\mathbf{K}}^{-1}\right)\\
& \le mk - \frac{m\delta^2(n-s)}{m\delta^2 + c(m-1)\delta} = mk - \frac{m\delta^2(mk-s)}{m\delta^2 + c(m-1)\delta}.
\end{align*} 

Here, the second inequality follows from  Rayleigh Quotient Theorem (Theorem 4.2.2 in \cite{horn2012matrix}). The third inequality follows from $\eqref{eq:lambda_min}$.
\end{proof}



\remark{For the coset bipartite graph construction, $\epsilon>0$ ensures that the diagonal entries of $\mathbf{D}_i$ are continuously distributed, which is needed to establish almost-sure invertibility in Proposition~1; consequently, $c>1$. Under the parameter settings of Proposition~1, the encoding matrix is almost surely invertible, so the approximation error is zero in the absence of stragglers even when $m>1$, and the exact gradient can be recovered.}

% \remark{ 
% For the coset bipartite graph construction, $\epsilon>0$ is needed to establish almost-sure invertibility in Proposition \ref{prop: cos_inv}, and consequently $c>1$. Under the parameter settings of Proposition~1, the encoding matrix is almost surely invertible, so the approximation error is zero in the absence of stragglers even when $m>1$, and the exact gradient can be recovered.
% }

% \remark{For coset bipartite graphs with appropriate parameters, the constructed encoding matrix is almost surely invertible. This ensures that in case of no stragglers, the approximation error is zero even for $m > 1$ and thus the exact gradient can be computed.}


%%\aditya{What is the need for the remark below; why is it needed?}
%\sifat{I included it in order to address a reviewer's comment. Included the paragraph after that for the same reason.}
\remark{The derived upper bounds illustrate the tradeoff between communication efficiency and approximation accuracy. For the BIBD and SRG constructions, with the remaining parameters fixed, each upper bound is the difference between the linearly increasing term $mk$ and a term that converges to the finite limit $\delta(n-s)$ as $m$ increases. Consequently, the upper bounds increase approximately linearly with $m$ for large $m$, indicating a weaker approximation guarantee. This behavior is expected, since increasing $m$ reduces the length of the vector transmitted by each worker ($d/m$), thereby improving communication efficiency at the expense of approximation accuracy. For the coset bipartite graph based construction, the number of workers $n=mk$, so varying $m$ also changes the number of workers. However, for fixed $k$, $s$, $\delta$, and $c$ the corresponding upper bound also grows linearly with $m$ for large $m$, implying a similar communication--accuracy tradeoff.}


A direct comparison of the three upper bounds under identical system
parameters is generally not possible. In particular, the SRG based
construction uses an $n\times n$ assignment matrix and hence requires
$k=n$, whereas the coset bipartite graph based construction requires $n=mk$.
Therefore, for $m>1$, these two constructions cannot be instantiated
with common values of $n$ and $k$. Furthermore, a fair comparison would
require the computation load $\delta$ to be fixed, while the existence
conditions discussed in Section~IV-B may prevent a common admissible
parameter choice across all three structures. Accordingly, a direct
parameter-matched comparison of all three upper bounds is generally not
possible.

\subsection{Discussion about condition number of the recovery system of equations}
For the BIBD and SRG based constructions, explicit upper bounds on $\kappa(\mathbf{B}_{\mathcal{F}}^T\mathbf{B}_{\mathcal{F}})$ are provided in Appendix \ref{sec:appendix_cond_number}. These bounds show that in these cases the condition number is low for many practical cases of interest. Thus, we have numerically stable gradient recovery at the PS. 


For the construction above based on coset bipartite graph, however, a similar condition number bound is hard to obtain when considering post-multiplication by random diagonal matrices. Instead for this case, we consider post-multiplication by diagonal matrices that are obtained from columns of Hadamard matrices.

\begin{defn} {[{\it Hadamard matrix}]}
A Hadamard matrix of order $\ell$ is a matrix
\[
\mathbf{H}_\ell \in \{\pm 1\}^{\ell\times \ell}
\]
whose rows are pairwise orthogonal. Equivalently,
\[
\mathbf{H}_\ell \mathbf{H}_\ell^{T} = \mathbf{H}_\ell^{T} \mathbf{H}_\ell = \ell I_\ell.
\]    
For $\ell>2$, a necessary condition for the existence of a Hadamard matrix
is that $\ell$ be divisible by $4$. Hadamard matrices are known to exist for many
orders, including every order which is a power-of-2.
\end{defn}
To obtain a construction based on coset bipartite graph for which the condition number can be bounded, consider a $(k,m,\delta)$ coset bipartite graph such that $k = p^a$, where $p$ is a prime, $a \in \mathbb{Z}_{\ge 1}$ and $p \nmid \delta$. Its bi-adjacency matrix can be written as $\mathbf{A}=[\mathbf{C}\ \cdots\ \mathbf{C}],$ where $\mathbf{C}\in\mathbb{R}^{k\times k}$ is circulant. Let $\mathbf{H}_m$ be the Hadamard matrix of order $m$, and define the encoding matrix,% \aditya{Define the Hadamard matrix formally first.}
\begin{equation} \label{eq:hada_coset_cons} \mathbf{B}=\mathbf{H}_m^T\otimes\mathbf{C}. \end{equation} 
For this construction, using the Cauchy interlacing theorem, we can conclude that $\kappa(\mathbf{B}_{\mathcal{F}}^T\mathbf{B}_{\mathcal{F}}) \le \kappa(\bfB^T\bfB)$. Furthermore, $\kappa(\bfB^T\bfB) = \kappa(\bfH_m\bfH^T_m) \kappa (\bfC^T\bfC) =  \kappa (\bfC^T\bfC)$.


\remark{For example, consider the $(27,2,5)$ coset bipartite graph with generating set $B = \{0,1,4,9,15\}$. For the corresponding circulant matrix $\mathbf{C}$,
$
\kappa(\mathbf{C}^T\mathbf{C}) \approx 110.49 $. Therefore, for the construction in \eqref{eq:hada_coset_cons},
$
\kappa(\mathbf{B}_{\mathcal{F}}^T\mathbf{B}_{\mathcal{F}})
\leq 110.5
$, for every non-straggler set $\mathcal{F}$.}



% where $\theta = \delta$ if $\lambda \ge \mu$. Otherwise, $\theta = \frac{1}{2} \left[ (\lambda - \mu) - \sqrt{(\lambda - \mu)^2 + 4 (\delta - \mu)} \right]$

% \begin{proof}
% see appendix. \aditya{TBD}
% \end{proof}
%%%%%%
%% An example of a floating figure using the graphicx package.
%% Note that \label must occur AFTER (or within) \caption.
%% For figures, \caption should occur after the \includegraphics.
%%
% \begin{figure}[htbp]
%   \centering
%   \includegraphics[width=0.3\textwidth]{myfigure}
%   % where an .eps filename suffix will be assumed under latex,
%   % and a .pdf suffix will be assumed for pdflatex
%   \caption{Simulation results.}
%   \label{fig:sim}
% \end{figure}
%%%%%%

%%%%%%
%% An example of a double column floating figure using two subfigures.
%% (The subfigure.sty package must be loaded for this to work.)  The
%% subfigure \label commands are set within each subfigure command,
%% the \label for the overall figure must come after \caption.  
%% \hfil must be used as a separator to get equal spacing
%%
% \begin{figure*}[htbp]
%   \centerline{\subfigure[Case I]{\includegraphics[width=2.5in]{subfigcase1}
%       % where an .eps filename suffix will be assumed under latex,
%       % and a .pdf suffix will be assumed for pdflatex
%       \label{fig:first_case}}
%     \hfil
%     \subfigure[Case II]{\includegraphics[width=2.5in]{subfigcase2}
%       % where an .eps filename suffix will be assumed under latex,
%       % and a .pdf suffix will be assumed for pdflatex
%       \label{fig:second_case}}}
%   \caption{Simulation results.}
%   \label{fig:sim}
% \end{figure*}
%%%%%%

%%%%%%
%% An example of a floating table. 
%% Note that, for IEEE style tables, the \caption command should come
%% BEFORE the table. Table text will default to \footnotesize as IEEE
%% normally uses this smaller font for tables. The \label must come
%% after \caption as always.
%%
% \begin{table}[htbp]
%   % increase table row spacing, adjust to taste
%   \renewcommand{\arraystretch}{1.3}
%   \caption{An Example of a Table}
%   \label{tab:table_example}
%   \centering
%   % Some packages, such as MDW tools, offer better commands for making tables
%   % than the plain LaTeX2e tabular which is used here.
%   \begin{tabular}{|c||c|}
%     \hline
%     One & Two\\
%     \hline
%     Three & Four\\
%     \hline
%   \end{tabular}
% \end{table}
%%%%%%


% \section{TBD}
% We present a brief review of block designs and relevant results of \cite{kadheKR19}. 

% \begin{defn}
% A design is a pair $(X, \mathcal{A})$ where $X$ is a set of elements called points, and $\mathcal{A}$ is a collection of nonempty subsets of $X$, each of which is called a block. Let $X = \{ x_1, x_2, \dots, x_v\}$ and $\mathcal{A} = \{ A_1, A_2, \dots, A_b\}$. The incidence matrix $\mathbf{M} \in \mathbb{R}^{v \times b}$ of $(X, \mathcal{A})$ is defined such that $\mathbf{M}(i,j) = 1$ if $x_i \in A_j$ and $\mathbf{M}(i,j) = 0$ if $x_i \notin A_j$.
% \end{defn}
% \begin{defn}
% A $(n,k,\gamma,\delta,\lambda)$-balanced incomplete block design (BIBD) is a design $(X, \mathcal{A})$ with $n$ points and $k$ blocks such that every point is contained in exactly $\delta$ blocks, every block contains exactly $\gamma$ points and any pair of distinct points is contained in exactly $\lambda$ blocks. If furthermore $n = k$ (and hence $\delta = \gamma$)  the BIBD is said to be symmetric. 
% \end{defn}
% For a $(n,k,\gamma,\delta,\lambda)$-BIBD, it is well known that $n\delta = k\gamma$ and $\delta(\gamma-1) = \lambda(n-1)$.

% Let $\mathbf{E} \in \mathbb{R}^{k \times n}$ be the trasnpose of the incidence matrix of a $(n,k,\gamma,\delta,\lambda)$ BIBD. Let $\mathbf{A} = \mathbf{E}$. Consequently, $\mathbf{A}$ is an $(n,k,\delta,\gamma)$ assignment matrix where the workers correspond to points and the data subsets correspond to blocks of the BIBD. In addition, any two workers will have exactly $\lambda$ data subsets in common. For $m =1$, let the encoding matrix $\mathbf{B} = \mathbf{A}_1 = \mathbf{E}$.
% % Consequently we have the number of workers $n = v$, number of data subsets $k = b$, the computation load $\delta = r$ and the replication factor $\gamma = p$. Also, any two workers have exactly $\lambda$ data subsets in common. 




%\section{Second Construction}
\section{Construction based on Hadamard product with null-space constraints}
\label{sec:construction_RHPNSC}

We now propose another construction based on bi-adjacency matrices of bipartite graphs or incidence matrices of BIBDs that achieves zero approximation error in the case of zero stragglers. 

% A desirable property of a scheme is that it recovers the exact gradient when there are no stragglers, i.e., when $s=0$. The scheme presented in Section \ref{sec:construction_RDM}, has this property only in the situation when $m=1$. Indeed, when $m > 1$, even though we have significantly lower error as compared to the naive scheme for general values of $s$, there is non-zero error even when $s=0$.

%For the construction as in $\eqref{eq:effAD1AD2}$, for $s = 0$,
% $\text{Err}_{[n]}(\mathbf{B}) =  mk -   \frac{m\delta^2n} {m\delta + (n-1)\lambda} = mk -   \frac{m\delta k \gamma} {m\delta + \delta (\gamma-1)}$.
% So, $\text{Err}_{[n]}(\mathbf{B}) = 0$ if and only if $\frac{\gamma}{m+\gamma-1} = 1 \implies m = 1$. So if $m >1 $, the approximation error is not zero when there are no stragglers. 

% In this section, we discuss a construction that overcomes this limitation. 


In the following construction, let $\mathbf{A}$ be either the $k\times n$ bi-adjacency matrix of a bi-regular bipartite graph or the incidence matrix of a $(k,n,\gamma,\delta,\lambda)$ BIBD,\footnote{This differs from Section \ref{sec:construction_RDM}, where we considered a $(n,k,\gamma,\delta,\lambda)$ BIBD and took $\mathbf{A}$ to be the transpose of its $n\times k$ incidence matrix.} where $n=mk$. Thus, in the BIBD case, the design has $k$ points and $n$ blocks, and its incidence matrix has dimensions $k\times n$. Note that, for $m>1$, we have $k<n$, so $\mathbf{A}$ has more columns than rows.


%%%PREV
% In the following construction, let $\bfA$ be either the $k \times n$ bi-adjacency matrix of a bi-regular bipartite graph or the incidence matrix of a BIBD\footnote{Note that this differs from Section \ref{sec:construction_RDM}, where we used the transpose of the incidence matrix of a BIBD.}, where $n = mk$. Note that for $m > 1$, we have $k < n$ so $\bfA$ has more columns than rows.

Let $\bfC \circ \bfD$ denote the Hadamard product (element-wise product) of matrices $\bfC$ and $\bfD$ of the same dimension. At a top-level, the construction proceeds by first selecting a vector $\bfv_1 \in \mathbb{R}^n$ with no zero entries. Following this, we first construct a matrix $\bfA_1$ by setting its $i$-th row, $\bfA_1(i,:)$ to be the Hadamard product $\bfA(i,:) \circ \bfv_1$ and then normalizing it by $||\bfA_1(i,:)||_2^2$. Since $\bfA_1$ has a non-trivial null space (since $k < n$ for $m>1$), we then choose $\bfv_2$ from $\text{Null}(\bfA_1)$ with no zero entries and construct $\bfA_2$ in the same manner. Thus, at the $j$-th iteration, $j \leq m$ we choose $\bfv_j$ with no zero entries such that $ \bfv_j \in \text{Null}(\bfA_1) \cap \dots \cap \text{Null}(\bfA_{j-1})$ (the corresponding null space is nontrivial since $k(j-1) < n$) and construct $\bfA_j$ in a similar manner. At the end of this process, we set the encoding matrix $\bfB$ by vertically stacking the $\bfA_i$'s so that

\begin{equation}\label{eq:recons}
\bfB^T = [\bfA_1^T ~|~ \dots ~|~\bfA_m^T].    
\end{equation}
% \begin{equation} \label{eq:reconsscheme}
%  \mathbf{B} = \begin{bmatrix} \mathbf{A}_1 \\ \mathbf{A}_2 \\ \vdots \\ \mathbf{A}_m  \end{bmatrix}   
% \end{equation}


The requirement that every entry of $\mathbf{v}_j$ be nonzero ensures that
$
\left\|\mathbf{v}_j(\mathcal{H}_i)\right\|_2^2 > 0$ for all  $i \in [k]$, where $\mathcal{H}_i =  \text{supp}(\mathbf{A}(i,:))$. Consequently, the normalization used in the construction is well-defined. Moreover, this condition ensures that 
$\text{supp}(\mathbf{A}_i)
=
\text{supp}(\mathbf{A})$ for all $ i \in [m]$. At the $j$-th iteration, the vector $\mathbf{v}_j$ is selected from the null space of the matrix obtained by vertically stacking $\mathbf{A}_1,\ldots,\mathbf{A}_{j-1}$ (see  Algorithm  \ref{alg:example} for a formal description).

We emphasize that although the rank-nullity theorem guarantees that this null space is nontrivial, it does not, in general, guarantee the existence of a null-space vector with all entries nonzero. Accordingly, the construction is stated under the assumption that such $\bfv_i$'s can be found. There are examples of assignment matrices $\mathbf{A}$, where this can be guaranteed; we discuss this in the upcoming Proposition~\ref{prop:coset-full-support}.

%We show in Proposition~\ref{prop:coset-full-support} that this requirement can always be satisfied when $\mathbf{A}$ is the bi-adjacency matrix of a $(k,m,\delta)$ coset bipartite graph. 





\begin{example}
Consider a $(4,2,1,2)$ bi-regular bipartite graph. The bi-adjacency matrix is 
\[
\mathbf{A}=
\begin{bmatrix}
1&0&1&0\\
0&1&0&1
\end{bmatrix}.
\]
Choose $\mathbf{v}_1=[1\ 1\ 1\ 1]^T$. Algorithm~\ref{alg:example} gives
\[
\mathbf{A}_1=\frac{1}{2}
\begin{bmatrix}
1&0&1&0\\
0&1&0&1
\end{bmatrix}.
\]
Next, choose $\mathbf{v}_2=[1\ 1\ {-1}\ {-1}]^T$. Since
$\mathbf{A}_1\mathbf{v}_2=\mathbf{0}_2$, we have
$\mathbf{v}_2\in\operatorname{Null}(\mathbf{A}_1)$, and hence
\[
\mathbf{A}_2=\frac{1}{2}
\begin{bmatrix}
1&0&-1&0\\
0&1&0&-1
\end{bmatrix}.
\]
Therefore,
\[
\mathbf{B}
=
\begin{bmatrix}
\mathbf{A}_1\\
\mathbf{A}_2
\end{bmatrix}
=
\frac{1}{2}
\begin{bmatrix}
1&0&1&0\\
0&1&0&1\\
1&0&-1&0\\
0&1&0&-1
\end{bmatrix},
\]
with
$\text{supp}(\mathbf{A}_1)
=\text{supp}(\mathbf{A}_2)
=\text{supp}(\mathbf{A})$.
\end{example}









Our scheme allows for exact gradient recovery when $s=0$. Furthermore, we provide an upper bound on $\operatorname{Err}_{\mathcal{F}}(\mathbf{B})$ that can be computed numerically and analyzed for certain choices of the vectors $\mathbf{v}_j$. These results are summarized in the following lemma.









% Note that $\text{supp}(\mathbf{A}_i) = \text{supp}(\mathbf{A})$ for all $i \in [m]$. A formal description of the scheme appears in Algorithm  \ref{alg:example}.

% Our scheme allows for exact gradient recovery when $s = 0$. Furthermore, we can provide an upper bound on the $\text{Err}_{\mathcal{F}}(\mathbf{B})$ that can be computed numerically and analyzed for certain choices of the $\bfv_i$'s. These results are summarized in the following lemma. %\aditya{Need to discuss the requirements of invertibility}



%we use 

% \begin{defn}
% Let $G = (\mathcal{X} \cup \mathcal{Y}, \mathcal{E})$ be a bipartite graph such that $\mathcal{X} = [n]$, $\mathcal{Y} = [k]$, each vertex in set $\mathcal{X}$ has degree $\delta$ and each vertex in set $\mathcal{Y}$ has degree $\gamma$. We refer to this as an $(n,k,\delta, \gamma)$-regular bipartite graph. An incidence matrix of an $(n,k,\delta, \gamma)$-regular bipartite graph is a matrix $\tilde{\mathbf{A}} \in \mathbb{R}^{k \times n}$ such that $\tilde{\mathbf{A}}(i,j) = 1$ if $ij \in \mathcal{E}$ where $i \in [k]$,$j \in [n]$ and $\tilde{\mathbf{A}}(i,j) = 0$ otherwise.        
% \end{defn}





% Let $m > 1$ and $\tilde{\mathbf{A}} \in \mathbb{R}^{k \times n}$ be the incidence matrix of an $(n,k,\delta, \gamma)$-regular bipartite graph 
% % or a BIBD
% such that $n = mk$, hence $k < n$. Let $\mathbf{A} = \tilde{\mathbf{A}}$. Consequently, $\mathbf{A}$ is an $(n,k,\delta,\gamma)$ assignment matrix. Let $\mathbf{A}_1, \mathbf{A}_2 ,\dots, \mathbf{A}_m \in \mathbb{R}^{k\times n}$ be output matrices of Algorithm 1. 

% In Algorithm 1, $\text{Null}(\mathbf{H})$ denotes the nullspace of the matrix $\mathbf{H}$. For each $j = [m] \setminus \{1\}$, the size of $\mathbf{H}$  is ${k(j-1) \times n}$. Since $n = mk$, $k(j-1) < n$ and therefore $\text{rank}(\mathbf{H}) \le k(j-1) < n$. By rank-nullity theorem, there exists $\mathbf{v}_j \in \text{Null}(\mathbf{H})$ such that $\mathbf{v}_j \ne \mathbf{0}_n$.

\begin{algorithm}[t!]
\caption{Algorithm to construct $\mathbf{B}$}
\label{alg:example}
\begin{algorithmic}[1] % The [1] numbers the lines
\REQUIRE Assignment matrix $\mathbf{A} \in \mathbb{R}^{k \times n}$ such that $n = mk$, communication reduction factor $m >1$.
\ENSURE Matrices $\mathbf{A}_1, \mathbf{A}_2, \dots, \mathbf{A}_m \in \mathbb{R}^{k \times n}$ such that for any $i \in [m]$, $\text{supp}(\mathbf{A}_i) = \text{supp}(\mathbf{A})$.
\STATE Initialize $\mathbf{A}_1, \dots, \mathbf{A}_m$ as zero matrices.  
% \STATE Generate a random vector $\mathbf{v}_1 \in \mathbb{R}^n$.
\STATE Pick a vector $\mathbf{v}_1 \in \mathbb{R}^n$ such that $\mathbf{v}_1(u) \ne 0$ for every $u \in [n]$.
% \sim \mathcal{N}(\mathbf{0}_n, \mathbf{I}_n)$ %\aditya{This will change since we don't necessarily use Gausssian}
%\STATE Generate random Gaussian vector $\mathbf{v}_1 \sim \mathcal{N}(\mathbf{0}_n, \mathbf{I}_n)$ \aditya{This will change since we don't necessarily use Gausssian}
\FOR{$i = 1$ to $k$}
 \STATE $\mathcal{H} \gets \text{supp}(\mathbf{A}(i,:))$.
 \STATE $\mathbf{A}_1(i, \mathcal{H}) \gets \frac{(\mathbf{v}_1(\mathcal{H}))^T}{\|\mathbf{v}_1(\mathcal{H})\|_2^2}$.
 % \STATE $\mathbf{A}_1(i, \mathcal{H}) \gets \frac{\mathbf{A}_1(i, \mathcal{H})}{\|\mathbf{A}_1(i,\mathcal{H})\|_2^2}$
\ENDFOR
\FOR{$j = 2$ to $m$}
\STATE $\mathbf{H} \gets \begin{bmatrix} \mathbf{A}_1 \\ \vdots \\ \mathbf{A}_{j-1}\end{bmatrix}$.
\STATE Pick a vector $\mathbf{v}_j \in \mathbb{R}^n$ such that $\mathbf{v}_j \in \text{Null}(\mathbf{H})$ and $\mathbf{v}_j(u) \ne 0$ for every $u \in [n]$. 
\FOR{$i = 1$ to $k$}
 \STATE $\mathcal{H} \gets \text{supp}(\mathbf{A}(i,:))$.
 \STATE $\mathbf{A}_j(i, \mathcal{H}) \gets\frac{(\mathbf{v}_j(\mathcal{H}))^T}{\|\mathbf{v}_j(\mathcal{H})\|_2^2}$.
 % \STATE $\mathbf{A}_j(i, \mathcal{H}) \gets \frac{\mathbf{A}_j(i, \mathcal{H})}{\|\mathbf{A}_j(i,\mathcal{H})\|_2^2}$
\ENDFOR
\ENDFOR

\RETURN $\mathbf{A}_1, \mathbf{A}_2, \dots, \mathbf{A}_m$.
\end{algorithmic}
\end{algorithm}

\begin{lemma}
\label{lemma:second_cons}
The construction that uses Algorithm 1 recovers the exact gradient when $s = 0$. For a given set of non-stragglers corresponding to $\mathcal{F}$, let $\mathbf{\Sigma}_{\mathcal{F}} := \mathbf{B}_{\mathcal{F}}^T\mathbf{B}_{\mathcal{F}}$, and $\tilde{\mathbf{\Sigma}}_{\mathcal{F}}$ be a $(n-s) \times (n-s)$ diagonal matrix such that for $u \in [n-s]$, 
 $\tilde{\mathbf{\Sigma}}_{\mathcal{F}}(u,u) := \mathbf{\Sigma}_{\mathcal{F}}(u,u) + \sum_{j \ne u} |\mathbf{\Sigma}_{\mathcal{F}}(u,j)|$. Suppose that $\mathbf{\Sigma}_{\mathcal{F}}$ is invertible. Then,
\begin{equation}\label{eq:bounddiagdom}
\text{Err}_{\mathcal{F}}(\mathbf{B}) \le mk - \sum_{i = 1}^m \mathbf{1}_k^T {\mathbf{A}_i}_{\mathcal{F}} (\tilde{\mathbf{\Sigma}}_{\mathcal{F}})^{-1} {\mathbf{A}_i}_{\mathcal{F}}^T \mathbf{1}_k.  
\end{equation}
\end{lemma}

\begin{proof}%[Proof of Lemma 2]
 Let $\mathbf{v}_1, \mathbf{v}_2, \dots, \mathbf{v}_m \in \mathbb{R}^n$ be vectors as generated by algorithm 1. For any $j \in [m]$ and $u\in [k]$, let $\text{supp}(\mathbf{A}_j(u, :)) = \mathcal{H}_u$. We emphasize that $\mathcal{H}_u$ is the same for all $\bfA_j$. Now, 

\begin{align*}(\mathbf{A}_j\mathbf{v}_j)(u) = \mathbf{A}_j(u,:)\mathbf{v}_j &= \mathbf{A}_j(u,\mathcal{H}_u)  \mathbf{v}_j({\mathcal{H}_u})\\
&= \frac{\mathbf{v}_j(\mathcal{H}_u)^T \mathbf{v}_j(\mathcal{H}_u)}{\|\mathbf{v}_j(\mathcal{H}_u)\|_2^2} = 1.
\end{align*}
So, $\mathbf{A}_j\mathbf{v}_j = \mathbf{1}_k$ for all $j \in [m]$. 
Note that by construction, $\bfv_j \in \cap_{i=1}^{j-1} \text{Null}(\bfA_i)$ (Step 9 of Algorithm \ref{alg:example}). Next we show that $\bfv_j \in \text{Null}(\bfA_{j+1}) \cap \dots \cap \text{Null}(\bfA_{m})$.

Consider $i$ such that $j < i \leq m$. Now,
\begin{align*}(\mathbf{A}_i\mathbf{v}_j)(u) = \mathbf{A}_i(u,:)\mathbf{v}_j &= \mathbf{A}_i(u,\mathcal{H}_u)  \mathbf{v}_j({\mathcal{H}_u})\\
&= \frac{\mathbf{v}_i(\mathcal{H}_u)^T \mathbf{v}_j(\mathcal{H}_u)}{\|\mathbf{v}_i(\mathcal{H}_u)\|_2^2}.
\end{align*}

We note here that by construction $\bfv_i \in \text{Null}(\bfA_j)$ so that  $\mathbf{A}_j\mathbf{v}_i = \mathbf{0}_k$. However, this means that for $u \in [k]$, we have $(\mathbf{A}_j\mathbf{v}_i)(u) = 0 \implies \frac{\mathbf{v}_j(\mathcal{H}_u)^T \mathbf{v}_i(\mathcal{H}_u)}{\|\mathbf{v}_j(\mathcal{H}_u)\|_2^2 } = 0$. This in turn implies that $\mathbf{v}_i(\mathcal{H}_u)^T \mathbf{v}_j(\mathcal{H}_u) = 0$ so that we can conclude that $\bfA_i \bfv_j = 0_k$ for $j < i \leq m$. It follows that $\bfB \bfv_i = \mathbf{f}_i$ for $i \in [m]$, i.e., the exact gradient recovery is possible when $s=0$.




Now we prove the upper bound on $\text{Err}_{\mathcal{F}}(\mathbf{B})$. Let $\mathbf{Q} := \tilde{\mathbf{\Sigma}}_{\mathcal{F}} - \mathbf{\Sigma}_{\mathcal{F}}$. $\bfQ$ is evidently Hermitian. Now, for any $u \in [n-s]$, we have
\begin{align*}
&\mathbf{Q}(u,u) = \tilde{\mathbf{\Sigma}}_{\mathcal{F}}(u,u) - \mathbf{\Sigma}_{\mathcal{F}}(u,u)\\ 
&= \mathbf{\Sigma}_{\mathcal{F}}(u,u) + \sum_{j \ne u} |\mathbf{\Sigma}_{\mathcal{F}}(u,j)| - \mathbf{\Sigma}_{\mathcal{F}}(u,u) \\
&=  \sum_{j \ne u} |\mathbf{\Sigma}_{\mathcal{F}}(u,j)|.
\end{align*}

On the other hand, for $j \ne u$ and $j \in [n-s]$,
\begin{align*}
\mathbf{Q}(u,j) &= -\mathbf{\Sigma}_{\mathcal{F}}(u,j) \text{, so that}\\
\sum_{j \ne u} |\mathbf{Q}(u,j)| &=  \sum_{j \ne u} |\mathbf{\Sigma}_{\mathcal{F}}(u,j)|.
\end{align*}

This shows that $\mathbf{Q}$ is a diagonally dominant matrix and hence positive semidefinite \cite{horn2012matrix}.

%Consequently,  Also, $\mathbf{Q}^T = (\tilde{\mathbf{\Delta}}_{\mathcal{F}} - \mathbf{\Delta}_{\mathcal{F}})^T = \tilde{\mathbf{\Delta}}_{\mathcal{F}} - \mathbf{\Delta}_{\mathcal{F}}$. So, $\mathbf{Q}$ is hermitian. So, all the eigenvalues of $\mathbf{Q}$ are real. Since $\mathbf{Q}$ is diagonally dominant, all the eigenavlues are non-negative. Since $\mathbf{Q}$ is hermitian and all the eigenvalues are non-negative, it is positive semi-definite. 

Thus, 
$\mathbf{Q} \succeq 0 \implies \tilde{\mathbf{\Sigma}}_{\mathcal{F}} - \mathbf{\Sigma}_{\mathcal{F}} \succeq 0 \implies \tilde{\mathbf{\Sigma}}_{\mathcal{F}} \succeq \mathbf{\Sigma}_{\mathcal{F}} \implies (\tilde{\mathbf{\Sigma}}_{\mathcal{F}})^{-1} \preceq (\mathbf{\Sigma}_{\mathcal{F}})^{-1}$ (see \cite{bhatia2009positive}). Finally,
\begin{align*}
\text{Err}_{\mathcal{F}}(\mathbf{B}) 
% & = \min\limits_{\substack{\mathbf{R}\in \mathbb{R}^{n \times m} \\ supp(\mathbf{R}(:,i)) \subseteq \mathcal{F} \\ \forall i \in [m]}}  {\|\mathbf{B} \mathbf{R}- \mathbf{F}\|_F^2} \\
% & = \sum_{i = 1}^m  \min\limits_{\substack{\mathbf{r}_i\in \mathbb{R}^{n} \\ supp(\mathbf{r}_i) \subseteq \mathcal{F} }}   \|\bfB\mathbf{r}_i - \mathbf{f}_i\|_2^2\\
& = \sum_{i =1}^m \mathbf{f}_i^T \mathbf{f}_i - \mathbf{f}_i^T \bfB_{\cF}(\bfB_{\cF}^T\bfB_{\cF})^{-1}\bfB_{\cF}^T\mathbf{f}_i \\
& = mk -  \sum_{i =1}^m \mathbf{1}_k^T {\mathbf{A}_i}_{\mathcal{F}}(\mathbf{\Sigma}_{\mathcal{F}})^{-1}  {\mathbf{A}_i}_{\mathcal{F}}^T\mathbf{1}_k\\
& \le mk -  \sum_{i =1}^m \mathbf{1}_k^T {\mathbf{A}_i}_{\mathcal{F}}(\tilde{\mathbf{\Sigma}}_{\mathcal{F}})^{-1}  {\mathbf{A}_i}_{\mathcal{F}}^T\mathbf{1}_k.
\end{align*}
\end{proof}



% \noindent {\it Sketch of Proof:} 
% For $1 \leq i \leq m$ it can be observed that $\bfA_i \bfv_i = \mathbf{1}_k$ owing to the unit-norm rows of $\bfA_i$. Furthermore, by construction we have $\bfv_i \in \text{Null}(\bfA_1) \cap \dots \cap \text{Null}(\bfA_{i-1})$. It turns out that we can also show that $\bfv_i \in \text{Null}(\bfA_{i+1}) \cap \dots \cap \text{Null}(\bfA_{m})$, from which it holds that $\bfB \bfv_i = \mathbf{f}_i$.
% For the upper bound on $\text{Err}_{\mathcal{F}}(\mathbf{B})$, we show that $\tilde{\mathbf{\Sigma}}_{\mathcal{F}} - \mathbf{\Sigma}_{\mathcal{F}} \succeq 0$, i.e., positive semi-definite, from which it follows that $(\tilde{\mathbf{\Sigma}}_{\mathcal{F}})^{-1} \preceq (\mathbf{\Sigma}_{\mathcal{F}})^{-1}$. Subsequently, the result follows with additional steps. The details appear in Appendix B.

\begin{remark}
    The matrix $\tilde{\mathbf{\Sigma}}_{\mathcal{F}}$ is a diagonal matrix and hence the bound in \eqref{eq:bounddiagdom} is easier to calculate as compared to calculating the actual least-squares error. The presented bound depends on the set $\cF$. However, it may be possible to arrive at weaker upper bounds that hold for all $\cF$ with $|\cF|=n-s$, by considering matrices $\tilde{\mathbf{\Sigma}}'_{\mathcal{F}} \succeq \tilde{\mathbf{\Sigma}}_{\mathcal{F}}$ and using them instead in the RHS of \eqref{eq:bounddiagdom}.
\end{remark}
%\aditya{Include corollary that shows that if the $\bfv_i$'s can be chosen from $\pm 1$ distribution, then the above bound can actually be analyzed.}
%\aditya{Important point: we need to discuss the significance of this upper bound? The fact that it can be calculated simply by knowing the vectors. Can we optimize over choice of the vectors? etc. Can we discuss why the bound is more useful than pure simulation?}
%\sifat{I need to think about this.}

% \begin{proof}
% Let $\mathbf{v}_1, \mathbf{v}_2, \dots, \mathbf{v}_m \in \mathbb{R}^n$ be vectors as generated by algorithm 1. For any $i \in [m]$ and $u\in [k]$, let $\text{supp}(\mathbf{A}_i(u, :)) = \mathcal{H}$. Now, 

% \begin{align*}(\mathbf{A}_i\mathbf{v}_i)(u) = \mathbf{A}_i(u,:)\mathbf{v}_i &= \mathbf{A}_i(u,\mathcal{H})  \mathbf{v}_i({\mathcal{H}})\\
% &= \frac{\mathbf{v}_i(\mathcal{H})^T \mathbf{v}_i(\mathcal{H})}{\|\mathbf{v}_i(\mathcal{H})\|_2^2} = 1
% \end{align*}
% So, $\mathbf{A}_i\mathbf{v}_i = \mathbf{1}_k$ for all $i \in [m]$. 

% Let $j \in  [m] \setminus \{1\}$. Let $\mathbf{H} = \begin{bmatrix} \mathbf{A}_1 \\ \vdots \\ \mathbf{A}_{j-1}\end{bmatrix}$. By construction, $\mathbf{v}_j \in \text{Nul}(\mathbf{H})$. So, $\mathbf{H}\mathbf{v}_j = \mathbf{0}_ {k(j-1)}$ and consequently, for $i \in [m]$ and $i < j$, 
% $\mathbf{A}_i \mathbf{v}_j = \mathbf{0}_k$. Now let $i \in [m]$ and $i >j$. We have,

% \begin{align*}(\mathbf{A}_i\mathbf{v}_j)(u) = \mathbf{A}_i(u,:)\mathbf{v}_j &= \mathbf{A}_i(u,\mathcal{H})  \mathbf{v}_j({\mathcal{H}})\\
% &= \frac{\mathbf{v}_i(\mathcal{H})^T \mathbf{v}_j(\mathcal{H})}{\|\mathbf{v}_i(\mathcal{H})\|_2^2} 
% \end{align*}
% Now since $i > j$, $\mathbf{A}_j\mathbf{v}_i = \mathbf{0}_k$. 

% So, $(\mathbf{A}_j\mathbf{v}_i)(u) = 0 \implies \frac{\mathbf{v}_j(\mathcal{H})^T \mathbf{v}_i(\mathcal{H})}{\|\mathbf{v}_j(\mathcal{H})\|_2^2 } = 0$. Consequently, 
% $\mathbf{v}_j(\mathcal{H})^T \mathbf{v}_i(\mathcal{H}) = 0$ and $(\mathbf{A}_i\mathbf{v}_j)(u) = 0$. So, $\mathbf{A}_i\mathbf{v}_j = \mathbf{0}_k$ for $i > j$.
% Consequently we have that for any $i,j\in [m]$, $\mathbf{A}_i\mathbf{v}_j = \mathbf{1}_k$ for $ i = j$ and $\mathbf{A}_i\mathbf{v}_j = \mathbf{0}_k$ for $i \ne j$. Let $\mathbf{V} = \begin{bmatrix} \mathbf{v}_1 & \dots & \mathbf{v}_m \end{bmatrix}$.
% Now, 
% \begin{align*}
% \mathbf{B}\mathbf{V} &= \begin{bmatrix} \mathbf{A}_1 \\ \vdots \\ \mathbf{A}_{m}\end{bmatrix}\begin{bmatrix} \mathbf{v}_1 & \dots & \mathbf{v}_m \end{bmatrix}\\
% & = \begin{bmatrix}
% \mathbf{1}_{k} & \mathbf{0}_{k} & \cdots & \mathbf{0}_{k} \\
% \mathbf{0}_{k} & \mathbf{1}_{k} & \cdots & \mathbf{0}_{k} \\
% \vdots & \vdots & \ddots & \vdots \\
% \mathbf{0}_{k}  &\mathbf{0}_{k}  & \cdots & \mathbf{1}_{k}
% \end{bmatrix} = \mathbf{F}
% \end{align*}

% Since $\mathbf{B} \mathbf{V} = \mathbf{F}$, it follows that the approximation error is zero in case of no stragglers. Now we prove the upper bound. Let $\mathbf{Q} = \tilde{\mathbf{\Delta}}_{\mathcal{F}} - \mathbf{\Delta}_{\mathcal{F}}$. Now, for any $u \in [n-s]$, 
% $\mathbf{Q}(u,u) = \tilde{\mathbf{\Delta}}_{\mathcal{F}}(u,u) - \mathbf{\Delta}_{\mathcal{F}}(u,u) = \mathbf{\Delta}_{\mathcal{F}}(u,u) + \sum_{j \ne u} |\mathbf{\Delta}_{\mathcal{F}}(u,j)| - \mathbf{\Delta}_{\mathcal{F}}(u,u) =  \sum_{j \ne u} |\mathbf{\Delta}_{\mathcal{F}}(u,j)|$
% Also, for $j \ne u$ and $j \in [n-s]$, 
% $\mathbf{Q}(u,j) = -\mathbf{\Delta}_{\mathcal{F}}(u,j)$. So, $\sum_{j \ne u} |\mathbf{Q}(u,j)| =  \sum_{j \ne u} |\mathbf{\Delta}_{\mathcal{F}}(u,j)|$. Consequently, $\mathbf{Q}$ is a diagonally dominant matrix. Also, $\mathbf{Q}^T = (\tilde{\mathbf{\Delta}}_{\mathcal{F}} - \mathbf{\Delta}_{\mathcal{F}})^T = \tilde{\mathbf{\Delta}}_{\mathcal{F}} - \mathbf{\Delta}_{\mathcal{F}}$. So, $\mathbf{Q}$ is hermitian. So, all the eigenvalues of $\mathbf{Q}$ are real. Since $\mathbf{Q}$ is diagonally dominant, all the eigenavlues are non-negative. Since $\mathbf{Q}$ is hermitian and all the eigenvalues are non-negative, it is positive semi-definite. So, 
% $\mathbf{Q} \ge 0 \implies \tilde{\mathbf{\Delta}}_{\mathcal{F}} - \mathbf{\Delta}_{\mathcal{F}} \ge 0 \implies \tilde{\mathbf{\Delta}}_{\mathcal{F}} \ge \mathbf{\Delta}_{\mathcal{F}} \implies (\tilde{\mathbf{\Delta}}_{\mathcal{F}})^{-1} \le (\mathbf{\Delta}_{\mathcal{F}})^{-1}$.

% Finally,

% \begin{align*}
% \text{Err}_{\mathcal{F}}(\mathbf{B}) 
% % & = \min\limits_{\substack{\mathbf{R}\in \mathbb{R}^{n \times m} \\ supp(\mathbf{R}(:,i)) \subseteq \mathcal{F} \\ \forall i \in [m]}}  {\|\mathbf{B} \mathbf{R}- \mathbf{F}\|_F^2} \\
% % & = \sum_{i = 1}^m  \min\limits_{\substack{\mathbf{r}_i\in \mathbb{R}^{n} \\ supp(\mathbf{r}_i) \subseteq \mathcal{F} }}   \|\bfB\mathbf{r}_i - \mathbf{f}_i\|_2^2\\
% & = \sum_{i =1}^m \mathbf{f}_i^T \mathbf{f}_i - \mathbf{f}_i^T \bfB_{\cF}(\bfB_{\cF}^T\bfB_{\cF})^{-1}\bfB_{\cF}^T\mathbf{f}_i \\
% & = mk -  \sum_{i =1}^m \mathbf{1}^T {\mathbf{A}_i}_{\mathcal{F}}(\mathbf{\Delta}_{\mathcal{F}})^{-1}  {\mathbf{A}_i}_{\mathcal{F}}^T\mathbf{1}\\
% & \le mk -  \sum_{i =1}^m \mathbf{1}^T {\mathbf{A}_i}_{\mathcal{F}}(\tilde{\mathbf{\Delta}}_{\mathcal{F}})^{-1}  {\mathbf{A}_i}_{\mathcal{F}}^T\mathbf{1}
% \end{align*}





% \end{proof}
It turns out that if we choose $\mathbf{A}$ as the bi-adjacency matrix of a $(k,m,\delta)$ coset bipartite graph, there is a choice of the $\mathbf{v}_i$'s so that the above construction is guaranteed to work. In particular, consider pairwise orthogonal vectors $\mathbf{q}_1, \dots, \mathbf{q}_m \in \mathbb{R}^{m}$ such that $\mathbf{q}_1 = \mathbf{1}_m$ and every entry of each $\mathbf{q}_j$ is nonzero. We choose $\mathbf{v}_i = \mathbf{q}_i \otimes \mathbf{1}_k$. This is stated as a proposition below; the proof appears in Appendix \ref{app:coset-full-support}. 
\begin{prop}
\label{prop:coset-full-support}
Suppose that $\mathbf{A}$ is the bi-adjacency matrix of a $(k,m,\delta)$ coset bipartite graph. Then, there exist vectors $\mathbf{v}_1,\ldots,\mathbf{v}_m\in\mathbb{R}^{mk}$ such that for $j \in [m]$ every entry of $\mathbf{v}_j$ is nonzero  and for $j \in \{ 2, \dots, m\}$,
\[
\mathbf{v}_j \in
\text{Null}\left(
\left[
\mathbf{A}_1^{T}\ \cdots\ \mathbf{A}_{j-1}^{T}
\right]^{T}
\right),\] 
where $\mathbf{A}_1,\ldots,\mathbf{A}_{j-1}$ are constructed from $\bfA$ according to Algorithm \ref{alg:example}.
\end{prop}


Moreover, the vectors $\mathbf{q}_1,\ldots,\mathbf{q}_m$ can be scaled such that
$\|\mathbf{q}_j\|_2^2=m$ for all $j\in[m]$, while preserving their orthogonality
and nonzero entries. Since $\mathbf{A}=[\mathbf{C}\ \cdots\ \mathbf{C}]$ and each
row of $\mathbf{C}$ has $\delta$ nonzero entries, for
$\mathbf{v}_j=\mathbf{q}_j\otimes\mathbf{1}_k$,
$
\|\mathbf{v}_j(\mathcal{H}_i)\|_2^2=m\delta,
$
for all $i\in[k]$ and $j\in[m]$. Thus, Algorithm~1 gives
$
\mathbf{A}_j=\frac{1}{m\delta}
\big[q_j(1)\mathbf{C}\ \cdots\ q_j(m)\mathbf{C}\big].
$
Let $\mathbf{Q}=[\mathbf{q}_1\ \cdots\ \mathbf{q}_m]$, then the resulting encoding
matrix can be written as
$
\mathbf{B}=\frac{1}{m\delta}\left(\mathbf{Q}^T\otimes\mathbf{C}\right).
$
When $k=p^a$, where $p$ is prime, $a\in\mathbb{Z}_{\geq 1}$, and
$p\nmid\delta$, the matrix $\mathbf{C}$ is invertible. Since
$\mathbf{Q}\mathbf{Q}^T=m\mathbf{I}_m$, it follows that
$
\kappa(\mathbf{B}_{\mathcal{F}}^T\mathbf{B}_{\mathcal{F}})
\leq
\kappa(\mathbf{B}^T\mathbf{B})
=
\kappa(\mathbf{C}^T\mathbf{C}).
$






\section{Lower Bound}
%We propose the following lower bound for any communication-efficient approximate gradient coding scheme. 
The following lower bound applies to any communication-efficient approximate gradient coding scheme. 
%\aditya{Need to also say that some of the ideas are borrowed from the previous paper}
% \aditya{Check whether argument goes thru when the load is at most delta.}
% \sifat{It does.}
\begin{thm}
\label{thm:lowerbound}
For any communication-efficient approximate gradient coding encoding matrix $\mathbf{B} \in \mathbb{R}^{mk \times n}$ with communication reduction factor $m$, number of stragglers $s$, computation load at most $\delta$, number of workers $n$ and number of partitions $k$, there exists a non-straggler set $\mathcal{C} \subseteq [n]$ of size $n-s$ such that,
\begin{equation}\label{eq:lowerbound}
\text{Err}_{\mathcal{C}}(\mathbf{B})  
% {\|\mathbf{B} \mathbf{R}- \mathbf{F}\|_F^2}
\ge \max_{u \in [m]} \min \left\{k,{\floor{\frac{k(s+m-u)}{n\delta}}} \right \}u.
\end{equation}

% \begin{equation}\label{eq:lowerbound}
% \min\limits_{\substack{\mathbf{R}\in \mathbb{R}^{n \times m} \\\text{supp}(\mathbf{R}(:,i)) \subseteq \mathcal{C} \\ \forall i \in [m]}}  {\|\mathbf{B} \mathbf{R}- \mathbf{F}\|_F^2}  \ge \max_{u \in [m]}{\floor{\frac{k(s+m-u)}{n\delta}}} u.
% \end{equation}

% \aditya{when using supp in an equation it needs to be in a text environment. Make this change throughout}
% \sifat{done}
\end{thm}


\label{sec:lowerbound_proof}
\begin{proof}
Our proof leverages the basic ideas of \cite{dimakis_cyclic_mds}.
 Let $G = (\mathcal{W} \cup \mathcal{D}, \mathcal{E})$ be a bipartite graph where vertex set $\mathcal{W} = \{W_1 , \dots, W_n\}$ corresponds to the set of workers and vertex set $\mathcal{D} = \{ \mathcal{D}_1, \dots , \mathcal{D}_k\}$ corresponds to the set of data subsets. Also, let $(W_i,\mathcal{D}_j) \in \mathcal{E}$ if and only if the worker $W_i$ is assigned the data subset $\mathcal{D}_j$. For a vertex $v$ in graph $G$, let $\text{deg}(v)$ denote its degree. Without loss of generality, let us assume that $\deg(\mathcal{D}_1) \le \deg(\mathcal{D}_2) \le \dots \le \deg(\mathcal{D}_{k})$. Let $d_{j} := \frac{1}{j} \sum_{i = 1}^j \deg(\mathcal{D}_i)$, for any $j \in [k]$. Since each vertex in $\mathcal{W}$ has degree at most $\delta$, it follows that $d_1 \le d_2 \le \dots \le d_{k} \le \frac{n \delta}{k}$. For a vertex set $\mathcal{V}$ in $G$, let $N(\mathcal{V})$ denote the set of vertices that are adjacent to the vertices in $\mathcal{V}$. It follows that for $j \in [k]$, there exists a set of data subsets $\mathcal{P}_{j} \triangleq \{\mathcal{D}_1, \dots, \mathcal{D}_j\}$ of size $j$ for which $|N(\mathcal{P}_j)| \le \sum_{i = 1}^j \deg(\mathcal{D}_i) = jd_j \le {}\frac{jn \delta}{k}$. Fix $u \in [m]$,  and let
$
j_u:=\min\left\{k,
\left\lfloor\frac{k(s+m-u)}{n\delta}\right\rfloor\right\}$. Let $\mathcal{P}_0:=\emptyset$. Also, let $\mathcal{Q}_u:=\mathcal{P}_{j_u}$. Then,
$
|\mathcal{N}(\mathcal{Q}_u)|
\leq \frac{j_u n\delta}{k}
\leq s+m-u.
$
Thus, the data subsets in $\mathcal{Q}_u$ are assigned to at most
$s+m-u$ workers.
% Now, there exists a set of data subsets $\mathcal{Q}_u := \mathcal{P}_{\floor{\frac{k(s+m-u)}{n\delta}}} = \{ \mathcal{D}_1, \dots , \mathcal{D}_{\floor{\frac{k(s+m-u)}{n\delta}}} \}$ of size $\floor{\frac{k(s+m-u)}{n\delta}}$ such that  $|N(\mathcal{Q}_u)| \le \floor{\frac{k(s+m-u)}{n\delta}} \frac{n\delta}{k} \le s+m-u$. So, the data subsets in the set $\mathcal{Q}_u$ are assigned to at most $s+m-u$ workers. 

Let $\mathcal{S}$ be the set of stragglers, so $|\mathcal{S}| = s$. If $|N(\mathcal{Q}_u)| \ge s$, then choose $\mathcal{S}$ such that $\mathcal{S} \subseteq N(\mathcal{Q}_u)$. In this case, the data subsets in the set $\mathcal{Q}_u$ are assigned to at most $m-u$ non-stragglers. Otherwise, choose $\mathcal{S}$ such that $N(\mathcal{Q}_u)$ is contained in $\mathcal{S}$, i.e., $N(\mathcal{Q}_u) \subseteq \mathcal{S}$. In this case, the data subsets in the set $\mathcal{Q}_u$ are assigned to zero non-stragglers. In both cases, the data subsets in the set $\mathcal{Q}_u$ are assigned to at most $m-u$ non-stragglers. Let us denote the set of non-stragglers by $\mathcal{C} \triangleq [n] \setminus \mathcal{S}$.



% \aditya{General remark: Mathematics also needs punctuation. If you're ending a line, you need a period and commas should be used as appropriate.}
% \sifat{Added appropriate punctuations.}

% Denote $\mathbf{B}_{\mathcal{C}} := \mathbf{B}(:, \mathcal{C})$. 
For $i \in [k]$, denote $\mathcal{H}_i \triangleq \{ i, k+i, \dots,  (m-1)k + i \}$.
% Let $\mathcal{C}_i$ be the set of non-stragglers $\mathcal{D}_i$ is assigned to. Let $c := |\mathcal{C}_i| \le m-u$.
Define a matrix $\mathbf{B}^{(i)}$ such that $\mathbf{B}^{(i)} := \mathbf{B}(\mathcal{H}_i, \mathcal{C}) \in \mathbb{R}^{m \times (n-s)}$. 
% Since $u \in [m]$, $c\le m-u < m$. 
%Let $j \in \left[\floor{\frac{k(s+m-u)}{n\delta}}\right]$. 
 For any $\mathcal{D}_j\in\mathcal{Q}_u$, let $\text{rank}(\mathbf{B}^{(j)}) = \rho$. Since $\mathcal{D}_j$ is assigned to at most $m-u$ non-stragglers, at most $m-u < m$ columns of $\mathbf{B}^{(j)}$ will be nonzero. Hence, $\rho \le m-u$. For any $\mathbf{R} \in \mathbb{R}^{(n-s) \times m}$, $\text{rank}(\mathbf{B}^{(j)}\mathbf{R}) \le \min\{\text{rank}(\mathbf{B}^{(j)}), \text{rank}(\mathbf{R}) \} \le \rho$.
% Let $\mathbf{I}_m$ be the $m \times m$ identity matrix. 
Let $\{\sigma_i = 1\}_{i = 1}^m$ be the singular values of $\mathbf{I}_m$. Thus, for every $\mathcal{D}_j\in\mathcal{Q}_u$,
\begin{align*} \min\limits_{\substack{\mathbf{R}\in \mathbb{R}^{(n-s) \times m} }}  {\|\mathbf{B}^{(j)} \mathbf{R}- \mathbf{I}_m\|_F^2} & \ge \min\limits_{\substack{\tilde{\mathbf{I}}_m\in \mathbb{R}^{m \times m} \\ \text{rank}(\tilde{\mathbf{I}}_m) \le \rho }}  {\|\tilde{\mathbf{I}}_m- \mathbf{I}_m\|_F^2} \\
& = \sum_{ i = \rho+1}^m \sigma_i^2 = m-\rho \ge u.
\end{align*}
The equality holds by the Eckart-Young theorem \cite{Eckart_Young_1936}. The second inequality holds since $\rho \le m-u$. 

% In CITE, we show that for any $\mathbf{R} \in \mathbb{R}^{n-s \times m},$
% \begin{equation} \label{eq:claim3claim}
%   {\|\mathbf{B}_{\mathcal{C}} \mathbf{R}- \mathbf{F}\|_F^2}  = \sum_{i = 1}^k {\|\mathbf{B}^{(i)} \mathbf{R}- \mathbf{I}_m\|_F^2}.
% \end{equation}
%See appendix for the proof of $\eqref{eq:claim3claim}$. 
 Finally, for a fixed $u \in [m]$,
\begin{align*}
 %&\min\limits_{\substack{\mathbf{R}\in \mathbb{R}^{n \times m} \\ \text{supp}(\mathbf{R}(:,i)) \subseteq \mathcal{C} \\ \forall i \in [m]}}  {\|\mathbf{B} \mathbf{R}- \mathbf{F}\|_F^2}\\
 %
 \text{Err}_{\mathcal{C}}(\mathbf{B})  & =\min\limits_{\substack{\mathbf{R}\in \mathbb{R}^{(n-s) \times m}}} {\|\mathbf{B}_{\mathcal{C}} \mathbf{R}- \mathbf{F}\|_F^2}\\
 &= \min\limits_{\substack{\mathbf{R}\in \mathbb{R}^{(n-s) \times m} }}\sum_{i = 1}^k   {\|\mathbf{B}^{(i)} \mathbf{R}- \mathbf{I}_m\|_F^2}\\
 &\ge \sum_{i = 1}^k \min\limits_{\substack{\mathbf{R}\in \mathbb{R}^{(n-s) \times m} }}\  {\|\mathbf{B}^{(i)} \mathbf{R}- \mathbf{I}_m\|_F^2}\\
 &\ge \sum_{i:\mathcal{D}_i \in \mathcal{Q}_u} \min\limits_{\substack{\mathbf{R}\in \mathbb{R}^{(n-s) \times m} }}  {\|\mathbf{B}^{(i)} \mathbf{R}- \mathbf{I}_m\|_F^2}\\
 & \ge  j_u u.
\end{align*}
The second equality above holds from the argument in Appendix \ref{sec:appendix_claim_frob}. Since $u \in [m]$ was arbitrary, by maximizing over all $u \in [m]$, we have

\begin{align*}
 \text{Err}_{\mathcal{C}}(\mathbf{B}) \ge \max_{u \in [m]} \min\left\{k,
\left\lfloor\frac{k(s+m-u)}{n\delta}\right\rfloor\right\}u.
\end{align*}

% Since $u \in [m]$, we have
% % \[ \min\limits_{\substack{\mathbf{R}\in \mathbb{R}^{n \times m} \\ \text{supp}(\mathbf{R}(:,i)) \subseteq \mathcal{C} \\ \forall i \in [m]}}  {\|\mathbf{B} \mathbf{R}- \mathbf{F}\|_F^2}  \ge \max_{u \in [m]}{\floor{\frac{k(s+m-u)}{n\delta}}} u.
% % \]
% \[ \text{Err}_{\mathcal{C}}(\mathbf{B})  \ge \max_{u \in [m]}{\floor{\frac{k(s+m-u)}{n\delta}}} u.
% \]
\end{proof}











% \noindent {\it Sketch of Proof:} Our proof leverages the basic ideas of \cite{dimakis_cyclic_mds}. Specifically, we identify a group of subsets $\mathcal{Q}_u$ that are assigned to at most $s+m-u$ workers for $u \in [m]$ such that at most $m-u$ of the workers are non-stragglers. We further arrive at a lower bound on the least-squares error by focusing our attention on the corresponding group of subsets. The details appear in Appendix D.

 




% \aditya{Do we keep remark below etc?}
% \begin{remark}
% If $k = n$, and $s \ge \delta + 1$ then the lower bound reduces to $\floor{\frac{s}{\delta}}m$. \aditya{I'm not sure that previous line is true.}
% \sifat{I'll try to prove this. Previous analysis was incorrect.}
% % If $k = n$, 
% % \begin{align*}

% % \end{align*}
\begin{example}
Suppose that $k=n \ge 5$, $\delta = 4$, $m=2$ and $s = \delta +1 =5$. Then the above bound is achieved at $u=m=2$ and the lower bound is $2$. On the other hand, if $s = \delta-1 = 3$, then the bound is achieved at $u=1$ and takes the value $1$. Thus, the maximum in \eqref{eq:lowerbound} is achieved for different values of $u$ in different ranges of $s$.
\end{example}

\begin{remark}
The lower bound is associated with the worst-case straggler set for a given number of stragglers $s$. Consequently, the approximation error can be lower than this bound for other choices of straggler set $\mathcal{S}$ with $|\mathcal{S}| = s$. 
\end{remark}
%Suppose that $k=n$, $\delta = 4$,$m=2$ and $s \geq \delta +1 =5$. Then the above bound is achieved at $u=m=2$ and the lower bound is $2$. On the other hand, if $s = \delta-1 = 3$, then the bound is achieved at $u=1$ and takes the value $1$.
% \end{remark}
%\sifat{I need to elaborate this discussion. Also, do I need to prove this?}
% Since $P_i \in \mathcal{Q}$, the number of non-straggling workers for the partition $P_i$ is at most $m-u$. Let $\mathcal{S}_i$ be the set of non-straggling workers $P_i$ is assigned to. Let $\bfB_{\mathcal{S}_i}$ be the submatrix of $\bfB$ where the columns are the columns of $\bfB$ corresponding to the workers $\mathcal{S}_i$ and the rows are $i+(j-1)m$-th row of $\bfB$ for each $j \in [m]$. So, the minimum contribution to $\|\bfB \mathbf{r}_i-\mathbf{f}_i\|_2^2$, for each $i \in [m]$ by $P_i$ is the least square solution to $\|\bfB_{\mathcal{S}_i} \mathbf{x} - \mathbf{a}_i\|_2^2$, where $\mathbf{a}_1 \triangleq \begin{bmatrix}0 & 0& \dots& 1& \dots &0 \end{bmatrix}$ (every entry is zero except for the $i$-th entry) and $\mathbf{x}$ is the variable. 


% Now since,
% \[
% \min_{\mathbf{x}} \|\bfB_{\mathcal{S}_i} \mathbf{x} - \mathbf{a}_1\|_2^2 + \dots +  \min_{\mathbf{x}} \|\bfB_{\mathcal{S}_i} \mathbf{x} - \mathbf{a}_m\|_2^2 \ge u 
% \]

% The contribution to $\|\mathbf{BR-F}\|_F^2$ by each $P_i$ such that $P_i \in \mathcal{Q}$ is at least u.




% Since there are $\floor{\frac{s+m-u}{d}}$ partitions in $\mathcal{Q}$, we have 

% $\|\mathbf{BR-F}\|_F^2 \ge \floor{\frac{s+m-u}{d}} u$.

% Since $u \in [m]$, we have 
% \[ 
% \|\mathbf{BR-F}\|_F^2 \ge \max_{u} \floor{\frac{s+m-u}{\delta}} u \]


\section{Convergence Analysis}
\label{sec:convergence}
% \aditya{Too short,. You need a formal statement of SGD convergence for this to make sense. Need to be clear about probabilistic failure pattern}\\
Since our schemes propose approximations of the gradient, it is necessary to show that the gradient descent algorithm converges in these cases. In this section, we establish convergence for the construction as in $\eqref{eq:effAD1AD2}$ with BIBD, vertex-transitive SRG and coset bipartite graph assignment matrices, and for the construction as in $\eqref{eq:hada_coset_cons}$.
% Let $ \mathbf{g}^{(t)}$ be the exact gradient, 

Let $\partial L(\mathbf{w}^{(t)})$ be the set of subgradients of the loss function $L$ at $\mathbf{w}^{(t)}$. In Stochastic Gradient Descent (SGD) (Section 14.3 in \cite{books/daglib/0033642}), the parameter is updated as 

\[ \mathbf{w}^{(t+1)} = \mathbf{w}^{(t)} - \eta^{(t)} {\mathbf{v}}^{(t)},\]

where $\mathbf{v}^{(t)}$ is a random vector such that $\mathbb{E}[ \mathbf{v}^{(t)} | \mathbf{w}^{(t)}] \in \partial L (\mathbf{w}^{(t)})$. When $L$ is differentiable, this condition reduces to $\mathbb{E}\!\left[\mathbf{v}^{(t)} \mid \mathbf{w}^{(t)}\right]
=
\nabla L\!\left(\mathbf{w}^{(t)}\right)$. In what follows, we show that, after an appropriate rescaling, the computed gradient satisfies this unbiasedness condition. Therefore, the resulting update can be viewed as an instance of SGD. We also establish that the variance of the rescaled computed gradient is bounded by a constant (assuming the spectral norm of $\bfZ$ is bounded), independent of the iteration number. Under this property, Corollary~2.2 in~\cite{doi:10.1137/120880811} provides a stationary-point convergence guarantee for SGD when the loss function is bounded below and has a Lipschitz-continuous gradient, with the learning rate chosen as specified, where the guarantee is stated for $\mathbf{w}^{(R_T)}$ with $R_T$ chosen uniformly at random from $\{1,\ldots,T\}$.

%\aditya{what does randomized output condition mean?} 
%\sifat{fixed}

%establishes a stationary-point convergence guarantee 

%Accordingly, in Subsection~VIII-A, we establish the unbiasedness-up-to-scaling property for the BIBD construction. In Subsection~VIII-B, we establish the corresponding property for vertex-transitive SRGs and coset bipartite graphs. In Subsection~VIII-C, we use the expected approximation error bounds derived in Appendix~A to establish uniformly bounded variance of the appropriately rescaled computed gradient. Finally, in Subsection~VIII-D, we state the resulting stationary-point convergence guarantee by applying Corollary~2.2 of \cite{doi:10.1137/120880811}.









% Under certain standard assumptions on the loss function, SGD can be shown to converge (see Theorem 14.8 in \cite{books/daglib/0033642} for convex loss functions and Example 4.2 in \cite{books/lib/BertsekasT96} for non-convex loss functions). We will show that the expected value of the computed gradient for our schemes is the exact gradient. Consequently, our schemes can be considered as a special case of SGD and thus they will converge under standard assumptions.


In this part we assume that the arrival of the computed result from each worker within a fixed time interval is modeled by a Bernoulli random variable \cite{DBLP:journals/corr/abs-1901-09671}. In particular, we will assume that within a fixed time interval, each worker sends its computed result independently with probability $1-q$ for $q \in [0, 1)$. Therefore, each worker straggles independently with probability $q$.





% Let $\tilde{\mathbf{g}}^{(t)}$ be the approximation of the gradient at iteration $t$.  Then, the parameter is updated as 






% In stochastic gradient descent (SGD), the gradient descent algorithm converges if the expected value of the computed gradient is the exact gradient. CITE. 
% For the random diagonal matrix based construction, the computed gradient is a function of the random diagonal matrices $\mathbf{D}_i$. 
For the random diagonal matrix based construction, since in each iteration the encoding matrix varies based on the diagonal matrices $\mathbf{D}_i$, $i \in [m]$, the computed gradient is a function of random diagonal matrices $\mathbf{D}_i$. As before we let $\mathcal{F}$ denote the index set corresponding to non-straggling workers. 
%
%Let $\mathcal{F}$ be a random subset such that for any $i \in [n]$, $i \in \mathcal{F}$ with probability $1-q$ and $i \notin \mathcal{F}$ with probability $q$. Then, $\mathcal{F}$ is the index set corresponding to non-straggling workers. Note that in this case, $\cF$ can be an empty set since all the workers straggle with probability $q^n$. 
As described in Section \ref{sec:prob_form}, for a given encoding matrix $\mathbf{B}$ and a given set of non-stragglers corresponding to $\mathcal{F}$, the computed gradient for our schemes is $\mathbf{Z} \mathbf{B}_{\cF}\mathbf{R}$, where $\bfR \in \mathbb{R}^{(n-s) \times m}$ is the optimal decoding matrix. Let $\mathbf{D} := (\mathbf{D}_1, \mathbf{D}_2, \dots, \mathbf{D}_m)$. %Note that the diagonal entries of $\bfD_i$ are i.i.d. uniformly distributed over the set $S$, as described in Section \ref{sec:construction_RDM}. 
For any function $\phi$ of $\mathbf{D}$, we use $\mathbb{E}_{\mathbf{D}}[\phi(\mathbf{D})]$ to denote the expectation with respect to the joint distribution of the diagonal entries of $\mathbf{D}_1,\ldots,\mathbf{D}_m$. For any function $\psi$ defined on $2^{[n]}$ (codomain can vary), define 
% \[
% \mathbb{E}{\mathbf{D}}[\phi(\mathbf{D})]
% :=
% \int \phi\left(\operatorname{diag}(x{1,1},\ldots,x_{1,n}),\ldots,
% \operatorname{diag}(x_{m,1},\ldots,x_{m,n})\right)
% \prod_{i=1}^{m}\prod_{j=1}^{n}d\mu_{\epsilon}(x_{i,j}).
% \]

% \[
% \mathbb{E}_{\mathbf{D}}[\phi(\mathbf{D})]
% := \sum_{\substack{\mathbf{D}_1^{\star},\dots,\mathbf{D}_m^{\star}\\
% \mathbf{D}_i^{\star}\ \text{diagonal}, i \in [m] \\ \ \mathbf{D}_i^{\star}(j,j)\in S, j \in [n]}}
% \phi(\mathbf{D}_1^{\star},\dots,\mathbf{D}_m^{\star})\,
% \prod_{i=1}^m \mathbb{P}\!\left(\mathbf{D}_i=\mathbf{D}_i^{\star}\right).
% \]





% \[
% \mathbb{E}_{\mathbf{D}}[\phi(\mathbf{D})]
% := \sum_{\mathbf{D}_1^{\star},\dots,\mathbf{D}_m^{\star}}
% \phi(\mathbf{D}_1^{\star},\dots,\mathbf{D}_m^{\star})\,
% \mathbb{P}\!\left(
% \mathbf{D}_1=\mathbf{D}_1^{\star},\dots,
% \mathbf{D}_m=\mathbf{D}_m^{\star}
% \right).
% \]

% \[ \mathbb{E}_{\bfD}[\phi(\bfD)] := \sum_{\substack{\bfD_1,\dots,\bfD_m \ \text{diagonal}\\ 
% \bfD_i(j,j) \in  \{1, -1\} \\
% i \in [m], j \in [n]}}
% \phi(\bfD) \mathbb{P}(\bfD_1, \dots, \bfD_m) .\]

\[
\mathbb{E}_{\mathcal{F}}[\psi(\mathcal{F})]
:= \sum_{F \subseteq [n]} \psi(F)\, \mathbb{P}(\mathcal{F} = F)
= \sum_{i = 0}^n \ \sum_{\substack{|F| = i}}
\psi(F)\, q^{\,n-i}(1-q)^{\,i}.
\]
% $\mathbb{E}_{\mathcal{F}}[\phi(\mathcal{F})]
% := \sum_{\cF \subseteq [n]} \phi(\cF)\, \mathbb{P}(\mathcal{F}) = \sum_{i = 0}^n \sum_{\cF: |\cF| = i} \phi(\cF)  q^{i}(1-q)^{\,n-i}
% $.
In the discussion below, we will show that $\mathbb{E}_{\mathbf{D}}[\mathbb{E}_{\mathcal{F}}[\mathbf{B}_{\cF} \mathbf{R}]]= \zeta \mathbf{F}$, so that the true gradient can be obtained by calculating $\frac{1}{\zeta}\mathbb{E}_{\mathbf{D}}[\mathbb{E}_{\mathcal{F}}[\mathbf{Z}\mathbf{B}_{\cF}\mathbf{R}]]$. We need the following basic definitions.
\begin{defn}{[{\it Permutation}]}
%A permutation of $[n]$ is a bijective function $\sigma : [n] \to [n]$.  
A permutation is a reordering of the elements of $[n]$, denoted by $\sigma : [n] \to [n]$.  
The corresponding permutation matrix
$\bfP_{\sigma} \in \{0,1\}^{n \times n}$ is such that for $i, j \in [n]$, $\bfP_{\sigma}(i,j) = 1$ if $i = \sigma(j)$ and $\bfP_{\sigma}(i,j) = 0$ otherwise. Note that for standard basis vector $\mathbf{e}_i \in \mathbb{R}^n$,  $ \mathbf{e}_{\sigma(i)}^T\bfP_{\sigma} = \mathbf{e}_i ^T$. Also, $\bfP_{\sigma}^{-1} = \bfP_{\sigma}^T$.      
\end{defn}
    

% \[
%   \bfP_{\sigma}(i,j) =
%   \begin{cases}
%     1, & \text{if } i = \sigma(j),\\
%     0, & \text{otherwise}.
%   \end{cases}
% \]








\begin{defn}{[{\it Graph Automorphism}]}
Let $G=([n], E)$ be a graph with adjacency matrix $\bfA_G$. An automorphism of $G$ is a permutation $\sigma$ (with corresponding matrix $\bfP_{\sigma}$) of the vertex set $[n]$ such that for all $i, j \in [n]$, $(i,j) \in E$ if and only if $(\sigma(i),\sigma(j)) \in E$. In this case, $\bfP_{\sigma} \bfA_G \bfP_{\sigma}^T = \bfA_G$. %If $\bfA$ is the adjacency matrix $G$ and $\bfP_{\sigma}$ is a permutation matrix associated with an automorphism of $G$, then $\bfP_{\sigma} \bfA \bfP_{\sigma}^T = \bfA$. 
% We will use this fact to prove convergence for vertex-transitive strongly regular graphs.
\end{defn}

\begin{defn}{[{\it Vertex Transitive Graph}]}
A graph $G = ([n],E)$ is vertex-transitive if for any $i, j \in [n]$, there exists an automorphism $\sigma$ such that $\sigma(i) = j$.
\end{defn}

% \begin{remark}
% If $\bfA$ is the adjacency matrix of a graph $G$ and $\bfP_{\sigma}$ is a permutation matrix associated with an automorphism of $G$, then $\bfP_{\sigma} \bfA \bfP_{\sigma}^T = \bfA$. We will use this fact to prove convergence for vertex-transitive strongly regular graphs.
% \end{remark}

% \begin{defn}{[{\it Distributional Equivalence}]}
% Let $X$ and $Y$ be random variables. %taking values in a measurable space $(\mathcal{X}, \mathcal{F})$. 
% Then $X$ and $Y$ are said to be distributionally equivalent, denoted by $X \overset{d} = Y$ if for every event $A$, $\mathbb{P}(X \in A) = \mathbb{P}(Y \in A)$. If $X$ and $Y$ are distributionally equivalent, then it follows that $\mathbb{E}[X] = \mathbb{E}[Y]$.
% \end{defn}




% \sifat{define permutation matrix, distributional equivalence, automorphism, vertex transitivity}\\
\subsection{Unbiasedness for BIBDs}
% \aditya{typo}\\
% The following lemma proves convergence for the random diagonal matrix based construction with BIBD assignment matrices.

The following lemma establishes that the computed gradient for the random diagonal matrix-based construction with BIBD assignment matrices is unbiased up to a nonzero scaling factor.


\begin{lemma}
\label{lemma:BIBD_unbiasedness}
Suppose that the assignment matrix $\mathbf{A}$ is the transpose of the incidence matrix of a $(n,k,\gamma,\delta,\lambda)$ BIBD. %and let $\mathbf{A} = \mathbf{E}$ so that $\mathbf{A}$ is an $(n,k,\delta, \gamma)$ assignment matrix. 
Then, for the construction as in $\eqref{eq:effAD1AD2}$ with $\epsilon  = 0$, we have 

\[ \mathbb{E}_{\mathbf{D}}[\mathbb{E}_{\mathcal{F}}[\mathbf{B}_{\cF} \mathbf{R}]]= \alpha \mathbf{F},\]
where $\alpha$ is a strictly positive constant. %nonzero constant.


\end{lemma}

\begin{proof}

% Let $\bf{A} \in \mathbb{R   }^\mathbf{k \times n}$ be the transpose of the incidence matrix of a $(n,k, \gamma, \delta, \lambda)$ BIBD. Let $\bfD_1, \bfD_2, \dots, \bfD_m \in \mathbb{R}^{n \times n}$ be diagonal matrices where the diagonal entries are i.i.d with $\mathbb{P} (1) = \mathbb{P} (-1) = .5$ ($a = 0, p =1$ in this case). We construct the encoding matrix $\bfB$ as follows: 
% % \[ \bfB = \begin{bmatrix} \mathbf{A} \bfD_1 \\ \bfA \bfD_2\\ \vdots \\ \bfA\bfD_m\end{bmatrix}\]
% \[\mathbf{B}^T = [ \bfD_1 \bfA^T ~|~ \bfD_2 \bfA^T ~|~ \dots ~|~ \bfD_m \bfA^T]\]


% \[ \bfB_{\cF} = \begin{bmatrix} \mathbf{A}_{\cF} \tilde{\bfD}_1 \\ \bfA_{\cF} \tilde{\bfD}_2\\ \vdots \\  \mathbf{A}_{\cF} \tilde{\bfD}_m \end{bmatrix}\] 


Fix $i \in [m]$. Let $\bfr_i* = \text{argmin}_{\mathbf{r}} \|\bfB_{\cF} \mathbf{r} - \mathbf{f}_i \|_2^2$. When $\cF$ is empty, we assume $\mathbf{B}_{\cF} = \mathbf{0}_{mk}$. Thus, $\mathbf{B}_{\cF}\mathbf{r}_i* = \mathbf{0}_{mk}$. For non-empty $\cF$, we have $\mathbf{B}_{\mathcal{F}}^T = [ \tilde{\bfD}_1 \bfA_{\mathcal{F}}^T ~|~ \tilde{\bfD}_2 \bfA_{\mathcal{F}}^T ~|~ \dots ~|~ \tilde{\bfD}_m\bfA_{\mathcal{F}}^T]$. Here, for $t \in [m]$, $\tilde{\mathbf{D}}_t := \mathbf{D}_t(\mathcal{F}, \mathcal{F})$. Let $\mathbf{\Delta}_{\mathcal{F}} := \mathbf{A}_{\mathcal{F}}^T \bfA_{\cF}$. Then, $\mathbf{B}_{\mathcal{F}}^T\mathbf{B}_{\mathcal{F}} = \sum_{t = 1}^m\tilde{\bf{D}}_t  \mathbf{\Delta}_{\cF} \tilde{\bf{D}}_t$. So, for non-empty $\mathcal{F}$ $(s \ne n)$ of size $n-s$,  
% For $i \in [m]$, let . 
% So, $\bfr_i* = (\bfB_{\cF}^T\bfB_{\cF})^{-1} \bfB_{\cF}^T  \mathbf{f}_i$. 
% We have, 
% \begin{align*} 
% \bfr_i* &=  (\bfB_{\cF}^T\bfB_{\cF})^{-1} \bfB_{\cF}^T  \mathbf{f}_i\\
% & =  \left(\sum_{t = 1}^m\tilde{\bf{D}}_t  \mathbf{\Delta}_{\cF} \tilde{\bf{D}}_t\right)^{-1} \tilde{\bf{D}}_i \mathbf{A}_{\cF}^T \mathbf{1}_k\\
% % & = \tilde{\bf{D}}_i^{-1}\left( \mathbf{\Delta}_{\cF}  + \sum_{t = 1, t\ne i}^m\tilde{\bf{D}}_i^{-1}\tilde{\bf{D}}_t \mathbf{\Delta}_{\cF} \tilde{\bf{D}}_t\tilde{\bf{D}}_i^{-1}\right)^{-1} \tilde{\bf{D}}_i^{-1}\tilde{\bf{D}}_i \mathbf{A}_{\cF}^T \mathbf{1}_k\\
%  & \overset{(a)}=  \tilde{\mathbf{D}}_i^{-1} \left(\mathbf{\Delta}_{\cF}  + \sum_{t = 1, t\ne i}^m\tilde{\bf{D}}_i^{-1}\tilde{\bf{D}}_t \mathbf{\Delta}_{\cF} \tilde{\bf{D}}_t\tilde{\bf{D}}_i^{-1}\right)^{-1}  \mathbf{A}_{\cF}^T \mathbf{1}_k.
%  % &= \tilde{\mathbf{D}}_i^{-1}\bfM_i^{-1}\mathbf{A}_{\cF}^T \mathbf{1}_k.
% \end{align*}
\begin{align*} 
\bfr_i* &=  (\bfB_{\cF}^T\bfB_{\cF})^{-1} \bfB_{\cF}^T  \mathbf{f}_i \overset{(a)}=  \tilde{\mathbf{D}}_i^{-1} \left(\mathbf{\Delta}_{\cF}  + \sum_{t = 1, t\ne i}^m\tilde{\bf{D}}_i^{-1}\tilde{\bf{D}}_t \mathbf{\Delta}_{\cF} \tilde{\bf{D}}_t\tilde{\bf{D}}_i^{-1}\right)^{-1}  \mathbf{A}_{\cF}^T \mathbf{1}_k,
\end{align*}
where
$(a)$ holds since $\tilde{\bfD}_i^{-1}\tilde{\bfD}_i = \mathbf{I}_{n-s}$. Let $\bfM_i := \mathbf{\Delta}_{\cF}  + \sum_{t = 1, t\ne i}^m\tilde{\bf{D}}_i^{-1}\tilde{\bf{D}}_t \mathbf{\Delta}_{\cF} \tilde{\bf{D}}_t\tilde{\bf{D}}_i^{-1}$. From the discussion in Corollary \ref{cor:BIBD}, we have that $\bfM_i$ is positive-definite. Now,
%Note that for a BIBD, $\mathbf{\Delta}_{\cF} =  (\delta - \lambda) \mathbf{I}_{n-s} + \lambda \mathbf{J}_{n-s} =  (\delta - \lambda) \mathbf{I}_{n-s} + \lambda \mathbf{1}_{n-s} \mathbf{1}_{n-s}^T$. 
%Observe that $\mathbf{\Delta}_{\mathcal{F}}$ has eigenvalue $\delta - \lambda$ with multiplicity $n-s-1$ and eigenvalue $\delta - \lambda + \lambda(n-s)$ with multiplicity 1. Since $\delta  > \lambda$ for a BIBD \cite{StinsonBook}, the eigenvalues of $\mathbf{\Delta}_{\cF}$ are strictly positive. Since $\mathbf{\Delta}_{\mathcal{F}}$ is symmetric and all the eigenvalues are strictly positive, it is positive definite. Also, since the diagonal entries of $\tilde{\bf{D}}_t$ are nonzero for $t \in [m]$, each term in the sum 
%$\sum_{t = 1, t\ne i}^m\tilde{\bf{D}}_i^{-1}\tilde{\bf{D}}_t \mathbf{\Delta}_{\cF} \tilde{\bf{D}}_t\tilde{\bf{D}}_i^{-1}$ is positive definite. So, $\bfM_i$ is positive definite. Now, 

\begin{align*}
\mathbf{B}_{\mathcal{F}}\mathbf{r}_i* = \begin{bmatrix} \mathbf{A}_{\cF} \tilde{\bfD}_1\tilde{\mathbf{D}}_i^{-1} \\ \vdots\\ \bfA_{\cF}\\ \vdots \\  \mathbf{A}_{\cF} \tilde{\bfD}_m\tilde{\mathbf{D}}_i^{-1} \end{bmatrix} \bfM_i^{-1}\mathbf{A}_{\cF}^T \mathbf{1}_k.
\end{align*}
For $j \ne i$, $j \in [m]$, let 

\begin{equation}
f_j(\mathbf{D}_i)  := \mathbb{E}_{\mathbf{D}_t, t \in [m], t \ne i}[ \mathbf{A}_{\cF} \tilde{\bfD}_j\tilde{\bf{D}}_i^{-1}\bfM_i^{-1}\mathbf{A}_{\cF}^T \mathbf{1}_k].
\end{equation}

Note that  
$f_j(-\mathbf{D}_i) = - f_j(\mathbf{D}_i)$.
% \sifat{needs to be more rigorous}\\
So, $f_j(\mathbf{D}_i)$ is an odd function of $\mathbf{D}_i$. Also, note that the distribution of $\mathbf{D}_i$ is symmetric, i.e., $\bfD_i \overset{d}= - \bfD_i$ (~$\overset{d}=$ denotes distributional equivalence). Consequently, for $j \ne i$, $j \in [m]$, we have $\mathbb{E}_{\mathbf{D}_i}[f_j(\mathbf{D}_i)] = \mathbf{0}_k$. Therefore, 
% \sifat{needs to be more rigorous}
\begin{align*}\mathbb{E}_{\mathbf{D}} [ \mathbf{A}_{\cF} \tilde{\bfD}_j\tilde{\bf{D}}_i^{-1}\bfM_i^{-1}\mathbf{A}_{\cF}^T \mathbf{1}_k] &= \mathbb{E}_{\mathbf{D}_i}[\mathbb{E}_{\mathbf{D}_t, t \in [m], t \ne i}[ \mathbf{A}_{\cF} \tilde{\bfD}_j\tilde{\bf{D}}_i^{-1}\bfM_i^{-1}\mathbf{A}_{\cF}^T \mathbf{1}_k]]\\
&= \mathbb{E}_{\mathbf{D}_i}[f_j(\mathbf{D}_i)] \\
&=  \mathbf{0}_{k}. 
\end{align*}

Next we calculate $\mathbb{E}_{\bfD} [\mathbf{A}_{\cF}\bfM_i^{-1}\mathbf{A}_{\cF}^T \mathbf{1}_k]$. 
% Since for a BIBD, $\mathbf{\Delta}_{\cF} =  (\delta - \lambda) \mathbf{I}_{n-s} + \lambda \mathbf{J}_{n-s} =  (\delta - \lambda) \mathbf{I}_{n-s} + \lambda \mathbf{1}_{n-s} \mathbf{1}_{n-s}^T$, 
We have,
\begin{align*}
\mathbf{M}_i &= \mathbf{\Delta}_{\cF} + \sum_{t = 1, t\ne i}^m\tilde{\bf{D}}_i^{-1}\tilde{\bf{D}}_t \mathbf{\Delta}_{\cF} \tilde{\bf{D}}_t\tilde{\bf{D}}_i^{-1}\\
&= (\delta - \lambda) \mathbf{I}_{n-s} + \lambda \mathbf{1}_{n-s} \mathbf{1}_{n-s}^T + \sum_{t = 1, t\ne i}^m \tilde{\bf{D}}_i^{-1}\tilde{\bf{D}}_t  \left[ (\delta - \lambda) \mathbf{I}_{n-s} + \lambda \mathbf{1}_{n-s} \mathbf{1}_{n-s}^T \right]\tilde{\bf{D}}_t\tilde{\bf{D}}_i^{-1}\\
& = m(\delta - \lambda) \bfI_{n-s} + \lambda \mathbf{1}_{n-s} \mathbf{1}_{n-s}^T+ \lambda \sum_{t = 1, t\ne i}^m\mathbf{u}_{it} \mathbf{u}_{it}^T,
% &=  \sum_{k = 1, k\ne i}^m \left[(\delta - \lambda) \tilde{\bf{D}}_i^{-1}\tilde{\bf{D}}_k \tilde{\bf{D}}_k\tilde{\bf{D}}_i^{-1} + \lambda \tilde{\bf{D}}_i^{-1}\tilde{\bf{D}}_k \mathbf{1}_{n-s}\mathbf{1}_{n-s}^T \tilde{\bf{D}}_k \tilde{\bf{D}}_i^{-1} \right]\\
%& = (m-1)(\delta - \lambda) \bfI_{n-s} + \lambda \sum_{t = 1, t\ne i}^m\mathbf{u}_{it} \mathbf{u}_{it}^T. 
\end{align*}
where $\mathbf{u}_{it}:= \tilde{\bf{D}}_i^{-1}\tilde{\bf{D}}_t \mathbf{1}_{n-s}$. The final equality holds since for $\epsilon = 0$, the entries of $\tilde{\bfD_{j}}$ are $\pm 1$ and therefore $\tilde{\mathbf{D}}_j^{-1} =\tilde{\mathbf{D}}_j$ for all $j \in [m]$. %Consequently, 
%\[ \mathbf{M}_i = m(\delta - \lambda) \bfI_{n-s} + \lambda \mathbf{1}_{n-s} \mathbf{1}_{n-s}^T+ \lambda \sum_{t = 1, t\ne i}^m\mathbf{u}_{it} \mathbf{u}_{it}^T.\]
% Let,  $\mathbf{P}_i := m(\delta - \lambda) \mathbf{I}_{n-s} + \lambda\sum_{t = 1, t\ne i}^m\mathbf{u}_{it} \mathbf{u}_{it}^T$. So, $\mathbf{M}_i = \mathbf{P}_i + \lambda \mathbf{1}_{n-s}\mathbf{1}_{n-s}^T$. Applying the Woodbury Matrix Identity \cite{woodbury1950}, we get \[ \mathbf{M}_i^{-1} = \mathbf{P}_i^{-1} - \frac{\lambda \mathbf{P}_i^{-1} \mathbf{1}_{n-s}\mathbf{1}_{n-s}^T \mathbf{P}_i^{-1}}{1 + \lambda \mathbf{1}_{n-s}^T \mathbf{P}_i^{-1} \mathbf{1}_{n-s}}.\]
% Now, 
% \begin{align*}
% \mathbf{M}_i^{-1} \mathbf{1}_{n-s} &= \mathbf{P}_i^{-1}\mathbf{1}_{n-s} - \frac{\lambda \mathbf{P}_i^{-1} \mathbf{1}_{n-s}\mathbf{1}_{n-s}^T \mathbf{P}_i^{-1} \mathbf{1}_{n-s}}{1 + \lambda \mathbf{1}_{n-s}^T \mathbf{P}_i^{-1} \mathbf{1}_{n-s}}\\
% &= \mathbf{P}_i^{-1}\mathbf{1}_{n-s}\left(1- \frac{\lambda \mathbf{1}_{n-s}^T \mathbf{P}_i^{-1} \mathbf{1}_{n-s}}{1 + \lambda \mathbf{1}_{n-s}^T \mathbf{P}_i^{-1} \mathbf{1}_{n-s}}\right)\\
% & = \frac{\mathbf{P}_i^{-1}\mathbf{1}_{n-s}}{1 + \lambda \mathbf{1}_{n-s}^T \mathbf{P}_i^{-1} \mathbf{1}_{n-s}}.
% \end{align*}
Let $\mathbf{\Pi} \in \mathbb{R}^{(n-s) \times (n-s)}$ be a permutation matrix such that for $u, v \in [n-s]$,  $\mathbf{\Pi}^T \mathbf{e}_u = \mathbf{e}_v$. %Note that since $\mathbf{\Pi}$ is a permutation matrix, $\mathbf{\Pi} \mathbf{u}_{it}$ is a reordering of the entries of  $\mathbf{u}_{it}$. Since the entries of $\mathbf{u}_{it}$ are i.i.d. and take values $\pm 1$ with equal probability, the joint distribution of $\mathbf{u}_{it}$ is invariant under permutations of its coordinates. Consequently, $\mathbf{u}_{it} \;\stackrel{d}{=}\; \mathbf{\Pi}\mathbf{u}_{it}$. Also, since $\mathbf{\Pi}$ is a permutation matrix, $\mathbf{\Pi} \mathbf{1}_{n-s} = \mathbf{1}_{n-s}$. 
% \begin{align*}
% \mathbf{P}_i &= m(\delta - \lambda) \mathbf{I}_{n-s} + \lambda\sum_{k = 1, k\ne i}^m\mathbf{u}_{ik} \mathbf{u}_{ik}^T\\
% & \overset{d}= m(\delta - \lambda) \mathbf{I}_{n-s} + \lambda\sum_{k = 1, k\ne i}^m\mathbf{\Pi}\mathbf{u}_{ik} \mathbf{u}_{ik}^T\mathbf{\Pi}^T\\
% &= m(\delta - \lambda) \mathbf{\Pi} \mathbf{\Pi}^T + \lambda\sum_{k = 1, k\ne i}^m\mathbf{\Pi}\mathbf{u}_{ik} \mathbf{u}_{ik}^T\mathbf{\Pi}^T\\
% & = \mathbf{\Pi} \mathbf{P}_i \mathbf{\Pi}^T
% \end{align*}
We have,
\begin{align*}
\mathbf{\Pi}\mathbf{M}_i\mathbf{\Pi}^T &= m(\delta - \lambda) \mathbf{\Pi}\mathbf{\Pi}^T + \lambda \mathbf{1}_{n-s} \mathbf{1}_{n-s}^T+ \lambda \sum_{t = 1, t\ne i}^m \mathbf{\Pi}\mathbf{u}_{it} \mathbf{u}_{it}^T \mathbf{\Pi}^T\\
%& \overset{d} = m(\delta - \lambda) \mathbf{I}_{n-s} + \lambda \mathbf{1}_{n-s} \mathbf{1}_{n-s}^T+ \lambda \sum_{t = 1, t\ne i}^m \mathbf{u}_{it} \mathbf{u}_{it}^T 
&\overset{d}=\mathbf{M}_i.
% &= \mathbf{\Pi}\mathbf{P}_i\mathbf{\Pi}^T + \lambda \mathbf{\Pi}\mathbf{1}_{n-s}\mathbf{1}_{n-s}^T\mathbf{\Pi}^T\\
% & = \mathbf{\Pi}\mathbf{P}_i\mathbf{\Pi}^T + \lambda \mathbf{1}_{n-s}\mathbf{1}_{n-s}^T\\
% & = \mathbf{P}_i +   \lambda \mathbf{1}_{n-s}\mathbf{1}_{n-s}^T = 
%&=\mathbf{M}_i.
\end{align*}
where the second equality above holds by $\mathbf{\Pi}\mathbf{\Pi}^T = \mathbf{I}_{n-s}$, $\mathbf{\Pi} \mathbf{1}_{n-s} = \mathbf{1}_{n-s}$, and since $\mathbf{u}_{it} \;\stackrel{d}{=}\; \mathbf{\Pi}\mathbf{u}_{it}$ because the entries of $\mathbf{u}_{it}$ are i.i.d. and take values $\pm 1$ with equal probability. Therefore, % Thus, $\mathbf{u}_{it} \;\stackrel{d}{=}\; \mathbf{\Pi}\mathbf{u}_{it}$.%  the joint distribution of $\mathbf{u}_{it}$ is invariant under permutations of its coordinates. Consequently, $\mathbf{u}_{it} \;\stackrel{d}{=}\; \mathbf{\Pi}\mathbf{u}_{it}$.
\begin{align*}
\mathbf{M}_i^{-1} \mathbf{1}_{n-s} &\overset{d}= (\mathbf{\Pi} \mathbf{M}_i \mathbf{\Pi}^T)^{-1} \mathbf{1}_{n-s} \\
%&= \mathbf{\Pi} \mathbf{M}_i^{-1} \mathbf{\Pi}^T \mathbf{1}_{n-s}\\
& = \mathbf{\Pi} \mathbf{M}_i^{-1} \mathbf{1}_{n-s}.
\end{align*}
%So, $\mathbf{\Pi}\mathbf{M}_i^{-1} \mathbf{1}_{n-s} \overset{d} = \mathbf{M}_i^{-1} \mathbf{1}_{n-s}$ and thus $\mathbb{E}_{\bfD}[\mathbf{\Pi}\mathbf{M}_i^{-1} \mathbf{1}_{n-s}] = \mathbb{E}_{\bfD}[\mathbf{M}_i^{-1} \mathbf{1}_{n-s}]$. 
Next,
\begin{align*}
\mathbb{E}_{\mathbf{D}} [ (\mathbf{M}_i^{-1}\mathbf{1}_{n-s})(u)] & = \mathbb{E}_{\mathbf{D}} [\mathbf{e}_u^T \mathbf{M}_i^{-1}\mathbf{1}_{n-s}]\\
%&= \mathbf{e}_u^T \mathbb{E}_{\mathbf{D}} [ \mathbf{M}_i^{-1}\mathbf{1}_{n-s}]\\
& \overset{(a)}{=} \mathbf{e}_u^T \mathbb{E}_{\mathbf{D}} [\mathbf{\Pi} \mathbf{M}_i^{-1}\mathbf{1}_{n-s}]\\
% & = \mathbf{e}_u^T \mathbf{\Pi}\mathbb{E}_{\mathbf{D}} [ \mathbf{M}_i^{-1}\mathbf{1}_{n-s}]\\
& \overset{(b)}{=} \mathbf{e}_v^T \mathbb{E}_{\mathbf{D}} [ \mathbf{M}_i^{-1}\mathbf{1}_{n-s}]\\
%& = \mathbb{E}_{\mathbf{D}} [\mathbf{e}_v^T \mathbf{M}_i^{-1}\mathbf{1}_{n-s}]\\
& = \mathbb{E}_{\mathbf{D}} [ (\mathbf{M}_i^{-1}\mathbf{1}_{n-s})(v)],
\end{align*}
where $(a)$ holds since $\mathbf{\Pi}\mathbf{M}_i^{-1} \mathbf{1}_{n-s} \overset{d} = \mathbf{M}_i^{-1} \mathbf{1}_{n-s}$ and $(b)$ holds since $\mathbf{\Pi}^T \mathbf{e}_u = \mathbf{e}_v$.
Since $u,v \in [n-s]$ are arbitrary, it follows that all the entries of $\mathbb{E}_{\mathbf{D}}[\mathbf{M}_i^{-1} \mathbf{1}_{n-s}]$ are equal. Also, since $\mathbf{M}_i$ is positive definite, $\mathbf{M}_i^{-1}$ is positive definite. Hence, $\mathbb{E}_{\mathbf{D}}[\mathbf{M}_i^{-1}]$ is positive definite. So, $\mathbf{1}_{n-s}^T\mathbb{E}_{\mathbf{D}}[\mathbf{M}_i^{-1}] \mathbf{1}_{n-s} = \mathbf{1}_{n-s}^T\mathbb{E}_{\mathbf{D}}[\mathbf{M}_i^{-1} \mathbf{1}_{n-s}] > 0$. So, $\mathbb{E}_{\mathbf{D}}[\mathbf{M}_i^{-1} \mathbf{1}_{n-s}]$ is a nonzero vector. Since $\bfM_i$ is a function of $\mathbf{\Delta}_{\cF}$ and for a BIBD $\mathbf{\Delta}_{\cF}$ depends only on the size of $\cF$, $\mathbb{E}_{\mathbf{D}} [ \mathbf{M}_i^{-1}\mathbf{1}_{n-s}] = \alpha_{i,|\cF|} \mathbf{1}_{n-s}$, where $\alpha_{i,|\cF|} > 0$. So, when $\mathcal{F}$ is non-empty,
\begin{align*}
\mathbb{E}_{\mathbf{D}}\left[\mathbf{B}_{\mathcal{F}}\mathbf{r}_i*\right] & = \mathbb{E}_{\mathbf{D}}\left[\begin{bmatrix} \mathbf{A}_{\cF} \tilde{\bfD}_1\tilde{\bf{D}}_i^{-1} \\ \vdots \\\mathbf{A}_{\cF} \\ \vdots \\  \mathbf{A}_{\cF} \tilde{\bfD}_m\tilde{\bf{D}}_i^{-1}\end{bmatrix} \bfM_i^{-1}\mathbf{A}_{\cF}^T \mathbf{1}_k \right]\\
&=\mathbb{E}_{\mathbf{D}}\left[\begin{bmatrix} \mathbf{A}_{\cF} \tilde{\bfD}_1\tilde{\bf{D}}_i^{-1} \bfM_i^{-1}\mathbf{A}_{\cF}^T \mathbf{1}_k \\ \vdots \\\mathbf{A}_{\cF}\bfM_i^{-1}\mathbf{A}_{\cF}^T \mathbf{1}_k  \\ \vdots \\  \mathbf{A}_{\cF} \tilde{\bfD}_m\tilde{\bf{D}}_i^{-1}\bfM_i^{-1}\mathbf{A}_{\cF}^T \mathbf{1}_k \end{bmatrix} \right]\\
&= \begin{bmatrix} \mathbf{0}_k\\ \vdots \\ \delta \mathbf{A}_{\mathcal{F}}\mathbb{E}_{\mathbf{D}}[ \mathbf{M}_i^{-1}\mathbf{1}_{n-s}] \\ \vdots\\ \mathbf{0}_k \end{bmatrix} =  \begin{bmatrix} \mathbf{0}_k\\ \vdots \\ \delta \alpha_{i, |\cF|} \mathbf{A}_{\mathcal{F}} \mathbf{1}_{n-s} \\ \vdots\\ \mathbf{0}_k \end{bmatrix}.
\end{align*}
For non-empty $\cF$, let $\mathbf{P}_{\mathcal{F}} = \mathbf{I}_n(:, \mathcal{F}) \in \mathbb{R}^{n \times (n-s)}$ be the selection matrix that selects the columns of a matrix corresponding to the set $\mathcal{F}$. So, $\mathbf{A}_{\cF} = \bfA \bfP_{\cF}$. 
% Note that when $\mathcal{F}$ is empty, $\mathbb{E}_{\bfD}[\bfB_{\mathcal{F}} \mathbf{r}_i*]=  \mathbb{E}_{\bfD}[\bfB \bfP_{\mathcal{F}} \mathbf{r}_i*] = \mathbf{0}_{mk} = \begin{bmatrix} \mathbf{0}_k\\ \vdots \\ \delta c_i \mathbf{A} \bfP_{\mathcal{F}} \mathbf{1}_{n-s} \\ \vdots\\ \mathbf{0}_k \end{bmatrix}$
Also, let $\mathbf{1}_{n}^{(\cF)} := \bfP_{\cF} \mathbf{1}_{n-s}$, where $\mathbf{1}_{n}^{(\cF)}$ is a vector of length $n$ whose $i$-th entry is $1$ if $i \in \cF$ and zero otherwise. So, $\bfA_{\cF} \mathbf{1}_{n-s} = \bfA \bfP_{\cF} \mathbf{1}_{n-s} = \bfA\mathbf{1}_{n}^{(\cF)}$. Now, if $\cF$ is empty, $\mathbf{1}_{n}^{(\cF)} = \mathbf{0}_n$. So, $\bfA\mathbf{1}_{n}^{(\cF)} = \bfA \mathbf{0}_n = \mathbf{0}_k$. Since for $\cF = \emptyset$, $\mathbb{E}_{\bfD}[\bfB_{\cF} \mathbf{r}_i*] = \mathbf{0}_{mk}$,  we have that for all $\cF$,  
\[ \mathbb{E}_{\bfD}[\mathbf{B}_{\cF}  \mathbf{r}_i*] = \begin{bmatrix} \mathbf{0}_k\\ \vdots \\ \delta \alpha_{i,|\cF|} \mathbf{A} \mathbf{1}_{n}^{(\cF)} \\ \vdots\\ \mathbf{0}_k \end{bmatrix}.\]
Now we compute $\mathbb{E}_{\cF} [\alpha_{i,|\cF|} \mathbf{1}_{n}^{(\cF)}]$. Since the worker failures are i.i.d., all coordinates of $\mathbb{E}_{\cF} [\alpha_{i,|\cF|} \mathbf{1}_{n}^{(\cF)}]$ are equal. So, $\mathbb{E}_{\cF} [\alpha_{i,|\cF|} \mathbf{1}_{n}^{(\cF)}] = \alpha_i \mathbf{1}_n$. Moreover, $\alpha_i>0$ since $\alpha_{i,|\mathcal{F}|}>0$ for every nonempty
$\mathcal{F}$ and each worker belongs to $\mathcal{F}$ with probability
$1-q>0$.
% \begin{align*}
% \mathbb{E}_{\cF}[\bfA_{\cF} \mathbf{1}_{n-s}] &= \mathbb{E}[\bfA \bfP_{\cF} \mathbf{1}_{n-s}]\\
% & = \mathbb{E}_{\cF}[\bfA \mathbf{1}_{n}^{(\cF)}] = \bfA \mathbb{E}_{\cF}[\mathbf{1}_{n}^{(\cF)}]. 
% \end{align*}
% Since $i \in {\cF}$ with probability $1-q$ and $i \notin  {\cF}$ with probability $q$, 
% \[ \mathbb{E} [\mathbf{1}_{n}^{(\cF)}(i)] = 1\cdot(1-q) + 0\cdot q = 1-q.\]
% Thus, $\mathbb{E}_{\cF} [\mathbf{1}_{n}^{(\cF)}] = (1-q) \mathbf{1}_n$. Therefore, 
% \begin{align*}
% \mathbb{E}_{\cF}[\bfA\mathbf{1}_{n}^{(\cF)}] = (1-q) \bfA \mathbf{1}_n = \gamma (1-q) \mathbf{1}_k.
% \end{align*}
% \sifat{Proof required}
Thus,
\[ 
\mathbb{E}_{\mathcal{F}}[\mathbb{E}_{\mathbf{D}}\left[\mathbf{B}_{\mathcal{F}}\mathbf{r}_i*\right]] = \begin{bmatrix} \mathbf{0}_k\\ \vdots \\  \delta \alpha_i \gamma \mathbf{1}_{k} \\ \vdots\\ \mathbf{0}_k \end{bmatrix} =  \delta \alpha_i \gamma  \mathbf{f}_i
.\]

Since the diagonal matrices $\mathbf{D}_i$ are i.i.d., $\mathbb{E}_{\bfD}[\mathbf{M}_i^{-1} \mathbf{1}_{n-s}] = \mathbb{E}_{\bfD}[\mathbf{M}_j^{-1} \mathbf{1}_{n-s}]$ for all $i,j \in [m]$. Therefore, $\alpha_i = \alpha_ j = \alpha'$ (let) for all $i ,j \in [m]$. So, for each $i \in [m]$,  
\[ 
\mathbb{E}_{\mathcal{F}}[\mathbb{E}_{\mathbf{D}}\left[\mathbf{B}_{\mathcal{F}}\mathbf{r}_i*\right]]=  \delta \alpha' \gamma \mathbf{f}_i
.\]
Since $\bfR = \begin{bmatrix}\mathbf{r}_1* & \mathbf{r}_2* & \dots & \mathbf{r}_m* \end{bmatrix}$, and $\bfF = \begin{bmatrix}\mathbf{f}_1 & \mathbf{f}_2 & \dots & \mathbf{f}_m \end{bmatrix}$, we have 
\[ \mathbb{E}_{\mathcal{F}}[\mathbb{E}_{\mathbf{D}}\left[\mathbf{B}_{\mathcal{F}}\mathbf{R}\right]] = \delta \alpha' \gamma \mathbf{F} = \alpha \mathbf{F},\] 
where $\alpha = \delta \alpha' \gamma $. Since $\alpha' \ne 0$ and $\delta, \gamma > 0$, $\alpha$ is a nonzero constant.

\end{proof}
\subsection{Unbiasedness for Vertex-Transitive SRGs and Coset Bipartite Graphs}

The following lemma establishes that the computed gradient for the random diagonal matrix-based construction is unbiased up to a nonzero scaling factor under an appropriate symmetry condition on the assignment matrix.







\begin{lemma}
{\label{lemma:SRG_COSET_unbiasedness}}
Let $\bfA \in \mathbb{R}^{k \times n}$ be an assignment matrix such that $\mathbf{A}^T\mathbf{1}_k\neq \mathbf{0}_n$ and for any $u,v \in [k]$ there exist permutations $\sigma$ of $[k]$ and $\pi$ of $[n]$ such that for the corresponding permutation matrices $\bfP_{\sigma}$ and $\bfQ_{\pi}$, $\mathbf{e}_{u}^T\bfP_{\sigma} = \mathbf{e}_{v}^T$  and $\bfP_{\sigma} \mathbf{A}\bfQ_{\pi}^T = \bfA$. Suppose the encoding matrix $\bfB$ is constructed as in $\eqref{eq:effAD1AD2}$ and that $\bfB_{\cF}^T \bfB_{\cF}$ is invertible for any non-empty $\cF \subseteq [n]$. Then, 

\[ \mathbb{E}_{\mathbf{D}}[\mathbb{E}_{\mathcal{F}}[\mathbf{B}_{\cF} \mathbf{R}]]= \beta\mathbf{F},\]
where $\beta$ is a nonzero constant.

\end{lemma}


% \begin{lemma}
% Suppose that the assignment matrix $\mathbf{A}$ is the adjacency matrix of a $(n, \delta, \lambda, \mu)$ vertex-transitive strongly regular graph $G$. %and let $\mathbf{A} = \mathbf{E}$ so that $\mathbf{A}$ is an $(n,k,\delta, \gamma)$ assignment matrix. 
% Then, for the construction as in $\eqref{eq:effAD1AD2}$ with $a = 0$ and $p = 1$, we have 

% \[ \mathbb{E}_{\mathbf{D}}[\mathbb{E}_{\mathcal{F}}[\mathbf{B} \mathbf{R}]]= c\mathbf{F},\]
% where $c$ is a non-zero constant.
% \end{lemma}
% In this section, we will show convergence for the random diagonal matrix based construction where the assignment matrix is the adjacency matrix of a vertex transitive strongly regular graph. Let $\bf{A} \in \mathbb{R   }^\mathbf{n \times n}$ be the adjacency matrix of a $(n, \delta, \lambda, \mu)$ strongly regular graph. Let $\bfD_1, \bfD_2, \dots, \bfD_m \in \mathbb{R}^{n \times n}$ be diagonal matrices as in the random diagonal matrix based construction. We construct the encoding matrix $\bfB$ as follows: 

% % \[ \bfB = \begin{bmatrix} \mathbf{A} \bfD_1 \\ \bfA \bfD_2\\ \vdots \\ \bfA\bfD_m\end{bmatrix}\]
% \[\mathbf{B}^T = [ \bfD_1 \bfA^T ~|~ \bfD_2 \bfA^T ~|~ \dots ~|~ \bfD_m \bfA^T]\]

% We have, 
% % \[ \bfB_{\cF} = \begin{bmatrix} \mathbf{A}_{\cF} \tilde{\bfD}_1 \\ \bfA_{\cF} \tilde{\bfD}_2\\ \vdots \\  \mathbf{A}_{\cF} \tilde{\bfD}_m \end{bmatrix}\] 
% \[\mathbf{B}_{\mathcal{F}}^T = [ \tilde{\bfD}_1 \bfA_{\mathcal{F}}^T ~|~ \tilde{\bfD}_2 \bfA_{\mathcal{F}}^T ~|~ \dots ~|~ \tilde{\bfD}_m\bfA_{\mathcal{F}}^T]\]


% \[ \bfB = \begin{bmatrix} \mathbf{A} \bfD_1 \\ \bfA \bfD_2\\ \vdots \\ \bfA\bfD_m\end{bmatrix}\]

\begin{proof}
Fix $i \in [m]$. Let $\bfr_i* = \text{argmin}_{\mathbf{r}} \|\bfB_{\cF} \mathbf{r} - \mathbf{f}_i \|_2^2$. When $\cF$ is empty, we assume $\mathbf{B}_{\cF} = \mathbf{0}_{mk}$. Thus, $\mathbf{B}_{\cF}\mathbf{r}_i* = \mathbf{0}_{mk}$. For non-empty $\cF$, we have
\[ \bfB_{\cF} =\begin{bmatrix} \mathbf{A} \bfD_1 \\ \bfA \bfD_2\\ \vdots \\ \bfA\bfD_m\end{bmatrix} \mathbf{P}_{\mathcal{F}},\] 
where $\mathbf{P}_{\mathcal{F}} = \mathbf{I}_n(:, \mathcal{F})$ is the selection matrix that selects the columns of $\mathbf{B}$ corresponding to the set $\mathcal{F}$. Also, for non-empty $\cF$, $\bfr_i* = (\bfB_{\cF}^T\bfB_{\cF})^{-1} \bfB_{\cF}^T \mathbf{f}_i$. We have,  
\begin{align*} 
\bfB_{\cF} \bfr_i* &=  \bfB_{\cF}(\bfB_{\cF}^T\bfB_{\cF})^{-1} \bfB_{\cF}^T  \mathbf{f}_i\\ 
& = \begin{bmatrix} \mathbf{A} \bfD_1 \\ \bfA \bfD_2\\ \vdots \\ \bfA\bfD_m\end{bmatrix} \mathbf{P}_{\mathcal{F}} \left(\sum_{t=1}^m\mathbf{P}_{\mathcal{F}}^T \mathbf{D}_t\mathbf{A}^T\mathbf{A}\mathbf{D}_t \mathbf{P}_{\mathcal{F}}\right)^{-1}\mathbf{P}_{\mathcal{F}}^T \mathbf{D}_i \mathbf{A}^T \mathbf{1}_k.\\
% & = \begin{bmatrix} h_{1i}(\mathbf{A}, \mathcal{F}, \mathbf{D}_1, \dots, \mathbf{D}_m)\\ \vdots \\ h_{mi}(\mathbf{A}, \mathcal{F}, \mathbf{D}_1, \dots, \mathbf{D}_m)\end{bmatrix} \mathbf{1}.
\end{align*}

For $j \in [m]$, let 

\[
h_{ji}(\mathbf{A}, \mathcal{F}, \mathbf{D}) :=
\begin{cases}
\mathbf{A}\mathbf{D}_j\mathbf{P}_{\mathcal{F}}
\left(\sum_{t = 1}^m \mathbf{P}_{\mathcal{F}}^T \mathbf{D}_t \mathbf{A}^T \mathbf{A} \mathbf{D}_t \mathbf{P}_{\mathcal{F}}\right)^{-1}
\mathbf{P}_{\mathcal{F}}^T \mathbf{D}_i \mathbf{A}^T, & \mathcal{F} \neq \emptyset, \\[0.5em]
\mathbf{0}_{k \times k}, & \mathcal{F} = \emptyset.
\end{cases}
\]
% \[h_{ji}(\mathbf{A}, \mathcal{F}, \mathbf{D}) := \mathbf{A}\mathbf{D}_j\mathbf{P}_{\mathcal{F}} \left(\sum_{k = 1}^m\mathbf{P}_{\mathcal{F}}^T \mathbf{D}_k \mathbf{A}^T\mathbf{A}\mathbf{D}_k \mathbf{P}_{\mathcal{F}}\right)^{-1}\mathbf{P}_{\mathcal{F}}^T  \mathbf{D}_i \mathbf{A}^T.\]

Also, let 
\[ f_{ji}( \mathbf{A}, \mathbf{D}) := \mathbb{E}_{\mathcal{F}}\left[h_{ji}(\mathbf{A}, \mathcal{F}, \mathbf{D})\right].\]

% By assumption $\bfB_{\cF}^T \bfB_{\cF}$ is invertible for non-empty $\cF$, so it is positive definite. Hence $(\bfB_{\cF}^T \bfB_{\cF})^{-1}$ is positive definite and consequently $h_{ii}(\bfA, \cF, \bfD)$ is positive definite for non-empty $\cF$. Since $f_{ii}(\bfA, \bfD)$ is a sum of positive definite matrices, it is positive definite.
% Note that for each non-empty $\cF$, since $\bfB_{\cF}^T \bfB_{\cF}$ is invertible, $\mathbf{A}\mathbf{D}_j\mathbf{P}_{\mathcal{F}}
% (\bfB_{\cF}^T \bfB_{\cF})^{-1}
% \mathbf{P}_{\mathcal{F}}^T \mathbf{D}_i \mathbf{A}^T$ is positive definite. Consequently, $f_{ji}(\bfA, \bfD)$ is positive definite. 
Now,

\[ \mathbb{E}_{\mathbf{D}} [ \mathbb{E}_{\mathcal{F}}[\mathbf{B}_{\mathcal{F}} \mathbf{r}_i*]] = \mathbb{E}_{\mathbf{D}} \left[ \begin{bmatrix}f_{1i}(\mathbf{A}, \mathbf{D}) \\ f_{2i}(\mathbf{A}, \mathbf{D})\\ \vdots \\ f_{mi}(\mathbf{A}, \mathbf{D}) \end{bmatrix} \mathbf{1}_k\right] .\]
For $j \in [m]$, let $f_j(\mathbf{D}_i)  := \mathbb{E}_{\mathbf{D}_t, t\in [m], t \ne i}[f_{ji}(\mathbf{A}, \mathbf{D}) \mathbf{1}_k] $. Note that if $i \ne j$, 
\[f_j(-\mathbf{D}_i) = - f_j(\mathbf{D}_i).\]
So, $f_j(\mathbf{D}_i)$ is an odd function of $\mathbf{D}_i$. Also, the distribution of $\mathbf{D}_i$ is symmetric, i.e., $\bfD_i \overset{d}= - \bfD_i$. Thus, for $i \ne j$, we have, $\mathbb{E}_{\mathbf{D}_i}[f_j(\mathbf{D}_i)] =  \mathbf{0}_k$.
Therefore, for $i \ne j$, 
\begin{align*}\mathbb{E}_{\mathbf{D}} [f_{ji}(\mathbf{A}, \mathbf{D}) \mathbf{1}_k] &= \mathbb{E}_{\mathbf{D}_i}[\mathbb{E}_{\mathbf{D}_t,   t \in [m], t \ne i}[f_{ji} (\mathbf{A}, \mathbf{D})\mathbf{1}_k]]\\&= \mathbb{E}_{\mathbf{D}_i}[f_j(\mathbf{D}_i)] =  \mathbf{0}_k.\end{align*}

Next we find $\mathbb{E}_{\bfD}[f_{ii}(\bfA, \bfD) \mathbf{1}_k]$. Suppose $u, v \in [k]$. So, by assumption there exist permutations $\sigma$ of $[k]$ and $\pi$ of $[n]$ such that for the corresponding permutation matrices $\bfP_{\sigma}$ and $\bfQ_{\pi}$, $\mathbf{e}_u ^T\bfP_{\sigma}= \mathbf{e}_{v}^T$ and $\bfP_{\sigma} \mathbf{A}\bfQ_{\pi}^T = \bfA$.
% where $\bfQ_{\pi}$ is some permutation matrix corresponding to some permutation $\pi$ of $[n]$. 
We will show that $\mathbf{P}_{\sigma} f_{ii}( \mathbf{A}, \mathbf{D}) \mathbf{P}_{\sigma}^T \overset{d}{=} f_{ii}( \mathbf{A}, \mathbf{D})$.


% Since $\sigma$ is an automorphism, $\bfQ_{\sigma} \bfA \bfQ_{\sigma}^T = \bfA$.

% $f_{ji}(\mathbf{A}, \mathbf{D}_1, \dots, -\mathbf{D}_i, \dots, \mathbf{D}_m) =  -f_{ji}(\mathbf{A}, \mathbf{D}_1, \dots, \mathbf{D}_i, \dots, \mathbf{D}_m)$.





% We have


% Since we assume each worker straggles randomly with probability $q$, we have

% \begin{align*}
% f_{ji}( \mathbf{A}, \mathbf{D}_1, \dots, \mathbf{D}_m) = \sum_{i = 0}^n \sum_{\mathcal{F}: |\mathcal{F}| = i} (1-q)^i q^{n -i} h_{ji}(\mathbf{A}, \mathcal{F}, \mathbf{D}_1, \dots, \mathbf{D}_m).
% \end{align*}


% Since $\sigma$ is an automorphism, $\bfQ_{\sigma} \bfA \bfQ_{\sigma}^T = \bfA$. 
Let $\mathcal{H}_r$ be the set of subsets of $[n]$ of size $r$, $r \in [n]$. Let $\pi_{\mathcal{H}_r}: \mathcal{H}_r \to \mathcal{H}_r$ be such that for $\cF \in \mathcal{H}_r$, $\pi_{\mathcal{H}_r}(\mathcal{F}) = \{\pi(u)| u \in \cF\}$. Since $\pi$ is a bijection on $[n]$, $\pi_{\mathcal{H}_r}$ is also a bijection on $\mathcal{H}_r$. Also, since $\bfQ_{\pi} \mathbf{e}_w = \mathbf{e}_{\pi(w)}$ for $w \in [n]$, $\bfQ_{\pi} \bfP_{\cF} = \bfP_{\pi_{\mathcal{H}_r}(\cF)}$. Let $\mathbf{Y}(\bfP_{\cF}, \bfD) := \sum_{t=1}^m 
\mathbf{P}_{\cF}^T \mathbf{D}_t \mathbf{A}^T\mathbf{A}\mathbf{D}_t \mathbf{P}_{\cF}$. Let $c_r = q^{n-r} (1-q)^{r}$. Since we assume each worker straggles independently with probability $q$, we have
% \sifat{ Above argument requires proof.}\\


\begin{align*}
f_{ii}(\mathbf{A}, \mathbf{D})
& \overset{(a)}= q^n  \mathbf{0}_{k \times k} + \sum_{r = 1}^n 
    \sum_{\substack{F \subseteq [n] \\ |F| = r}} c_r \, h_{ii}(\mathbf{A}, F, \mathbf{D})\\
& \overset{(b)}= \sum_{r = 1}^n 
    \sum_{\substack{F \subseteq [n] \\ |F| = r}}
    c_r \, h_{ii}(\mathbf{A}, F, \mathbf{D})\\
& \overset{(c)}= \sum_{r = 1}^n 
    \sum_{\substack{F \subseteq [n] \\ |F| = r}}
    c_r \,
    h_{ii}(\mathbf{A}, \pi_{\mathcal{H}_r}(F), \mathbf{D})\\
& \overset{(d)}= \sum_{r = 1}^n 
    \sum_{\substack{F \subseteq [n] \\ |F| = r}}
    c_r \,
   \mathbf{A}\mathbf{D}_i\mathbf{P}_{\pi_{\mathcal{H}_r}(F)} (\mathbf{Y}(\bfP_{\pi_{\mathcal{H}_r}(F)}, \bfD))^{-1}
    \mathbf{P}_{\pi_{\mathcal{H}_r}(F)}^T  \mathbf{D}_i\mathbf{A}^T\\ 
& \overset{(e)}= \sum_{r = 1}^n 
    \sum_{\substack{F \subseteq [n] \\ |F| = r}}
   c_r \,
   \mathbf{A}\mathbf{D}_i\bfQ_{\pi}\mathbf{P}_{F} (\mathbf{Y}(\bfQ_{\pi}\bfP_{F}, \bfD))^{-1}
    \mathbf{P}_{F}^T \bfQ_{\pi}^T \mathbf{D}_i\mathbf{A}^T.
%     & \overset{(e)}= \sum_{r = 1}^n 
%     \sum_{\substack{F \subseteq [n] \\ |F| = r}}
%     (1-q)^r q^{\,n-r} \,
%     \mathbf{A}\mathbf{D}_i\mathbf{Q}_{\pi}\mathbf{P}_{F} \\
% &\qquad\qquad\qquad\qquad\qquad\quad
% \left(\sum_{t = 1}^m 
%     \mathbf{P}_{F}^T \mathbf{Q}_{\pi}^T 
%     \mathbf{D}_t \mathbf{A}^T \mathbf{A} \mathbf{D}_t 
%     \mathbf{Q}_{\pi} \mathbf{P}_{F}\right)^{-1}
%     \mathbf{P}_{F}^T \mathbf{Q}_{\pi}^T \mathbf{D}_i \mathbf{A}^T.
\end{align*}




% \begin{align*}
% &  f_{ii}(\mathbf{A}, \mathbf{D})\\
% & \overset{(a)}= \sum_{r = 0}^n \sum_{F: |F| = r} (1-q)^r q^{\,n-r} \, h_{ii}(\mathbf{A}, F, \mathbf{D})\\
% & \overset{(b)}= \sum_{r = 1}^n \sum_{F: |F| = r} (1-q)^r q^{\,n-r} \, h_{ii}(\mathbf{A}, F, \mathbf{D})\\
% & \overset{(c)}= \sum_{r = 1}^n \sum_{F: |F| = r} (1-q)^r q^{\,n-r} \,
%     h_{ii}(\mathbf{A}, \pi_{\mathcal{H}_r}(F), \mathbf{D})\\
% & \overset{(d)}= \sum_{r = 1}^n \sum_{F: |F| = r} (1-q)^r q^{\,n-r} \,
%     \mathbf{A}\mathbf{D}_i\mathbf{Q}_{\pi}\mathbf{P}_{F} \\
% &\qquad\qquad\qquad\qquad\qquad\quad
% \left(\sum_{t = 1}^m \mathbf{P}_{F}^T \mathbf{Q}_{\pi}^T 
%     \mathbf{D}_t \mathbf{A}^T \mathbf{A} \mathbf{D}_t 
%     \mathbf{Q}_{\pi} \mathbf{P}_{F}\right)^{-1}
%     \mathbf{P}_{F}^T \mathbf{Q}_{\pi}^T \mathbf{D}_i \mathbf{A}^T.
% \end{align*}




% \begin{align*}
% &  f_{ii}(\bfA, \bfD)\\
% & \overset{(a)}= \sum_{r = 0}^n \sum_{\mathcal{F}: |\mathcal{F}| = r} (1-q)^r q^{n -r} h_{ii}(\mathbf{A}, \mathcal{F}, \mathbf{D})\\
% & \overset{(b)}= \sum_{r = 1}^n \sum_{\mathcal{F}: |\mathcal{F}| = r} (1-q)^r q^{n -r} h_{ii}(\mathbf{A}, \mathcal{F}, \mathbf{D})\\
% & \overset {(c)}= \sum_{r = 1}^n \sum_{\mathcal{F}: |\mathcal{F}| = r} (1-q)^r q^{n -r} h_{ii}(\mathbf{A}, \pi_{\mathcal{H}_r}(\mathcal{F}), \mathbf{D})\\
% & \overset{(d)} = \sum_{r = 1}^n \sum_{\mathcal{F}: |\mathcal{F}| = r} (1-q)^r q^{n -r} \mathbf{A}\mathbf{D}_i\mathbf{Q}_{\pi}\mathbf{P}_{\mathcal{F}} \\
% &\quad \quad \quad \quad \quad \quad \quad  \left(\sum_{t = 1}^m\mathbf{P}_{\mathcal{F}}^T\mathbf{Q}_{\pi}^T \mathbf{D}_t \mathbf{A}^T\mathbf{A}\mathbf{D}_t \mathbf{Q}_{\pi}\mathbf{P}_{\mathcal{F}}\right)^{-1}\mathbf{P}_{\mathcal{F}}^T \mathbf{Q}_{\pi}^T  \mathbf{D}_i\mathbf{A}^T.
% % &\overset{(d)} = \sum_{i = 0}^n \sum_{\mathcal{F}: |\mathcal{F}| = i} (1-q)^i q^{n -i} \mathbf{A}\mathbf{D}_i\mathbf{Q}_{\sigma}\mathbf{P}_{\mathcal{F}} \left(\sum_{k = 1}^m\mathbf{P}_{\mathcal{F}}^T\mathbf{Q}_{\sigma}^T \mathbf{D}_k \mathbf{A}^T\mathbf{A}\mathbf{D}_k \mathbf{Q}_{\sigma}\mathbf{P}_{\mathcal{F}}\right)^{-1}\mathbf{P}_{\mathcal{F}}^T \mathbf{Q}_{\sigma}^T  \mathbf{D}_i\mathbf{A}^T.
% \end{align*}

Here, $(a)$ holds by definition of $f_{ii}(\bfA, \bfD)$. $(b)$ holds since $h_{ii}(\bfA, \cF, \bfD)$ is a $k \times k$ zero matrix for $|\cF| = 0$. $(c)$ holds since $\pi_{\mathcal{H}_r}$ is a bijection. $(d)$ holds by definition of $h_{ii}(\mathbf{A}, \pi_{\mathcal{H}_r}(F), \mathbf{D})$. $(e)$ holds since  $\bfQ_{\pi} \bfP_{F} = \bfP_{\pi_{\mathcal{H}_r}(F)}$, for each $r \in [n]$. 

Recall that $\mathbf{D} = (\mathbf{D}_1, \dots, \mathbf{D}_m)$. In the expression below, $\mathbf{Q}_{\pi} \mathbf{D} \mathbf{Q}_{\pi}^T$ is used to denote $(\mathbf{Q}_{\pi} \mathbf{D}_1 \mathbf{Q}_{\pi}^T, \dots, \mathbf{Q}_{\pi} \mathbf{D}_m \mathbf{Q}_{\pi}^T)$. We have,
\begin{align*}
&\mathbf{P}_{\sigma} f_{ii}( \mathbf{A}, \mathbf{D}) \mathbf{P}_{\sigma}^T\\
& \overset{(f)}= q^n \mathbf{0}_{k \times k} + \sum_{r = 1}^n 
    \sum_{\substack{F \subseteq [n] \\ |F| = r}} c_r \, \mathbf{P}_{\sigma} h_{ii}(\mathbf{A}, F, \mathbf{D}) \mathbf{P}_{\sigma}^T\\
& \overset{(g)}=  \sum_{r = 1}^n 
    \sum_{\substack{F \subseteq [n]\\ |F| = r}} 
    c_r \,
    \mathbf{P}_{\sigma} \mathbf{A}\mathbf{D}_i\mathbf{P}_{F} (\bfY(\bfP_{F}, \bfD))^{-1}
    \mathbf{P}_{F}^T  \mathbf{D}_i\mathbf{A}^T\mathbf{P}_{\sigma}^T\\[0.5em]
& \overset{(h)}=  \sum_{r = 1}^n 
    \sum_{\substack{F \subseteq [n]\\ |F| = r}} 
    c_r \,
    \mathbf{A}\mathbf{Q}_{\pi}\mathbf{D}_i\mathbf{Q}_{\pi}^T\mathbf{Q}_{\pi}\mathbf{P}_{F} (\bfY(\bfQ_{\pi}\bfP_{\cF}, \bfQ_{\pi}\bfD \bfQ_{\pi}^T))^{-1}
    \mathbf{P}_{F}^T\mathbf{Q}_{\pi}^T\mathbf{Q}_{\pi}  
    \mathbf{D}_i \mathbf{Q}_{\pi}^T\mathbf{A}^T\\[0.5em]
& \overset{(i)}=  f_{ii}(\mathbf{A}, \mathbf{Q}_{\pi}\mathbf{D} \mathbf{Q}_{\pi}^T).
\end{align*}



% \begin{align*}
% &\mathbf{P}_{\sigma} f_{ii}( \mathbf{A}, \mathbf{D}) \mathbf{P}_{\sigma}^T\\
% & \overset{(e)}=  \sum_{r = 1}^n \sum_{\mathcal{F}: |\mathcal{F}| = r} (1-q)^r q^{n -r} \mathbf{P}_{\sigma} \mathbf{A}\mathbf{D}_i\mathbf{P}_{\mathcal{F}}\\ 
% & \quad \quad  \quad \quad \quad \quad \quad  \left( \sum_ {t= 1}^m\mathbf{P}_{\mathcal{F}}^T \mathbf{D}_t \mathbf{A}^T\mathbf{A}\mathbf{D}_t \mathbf{P}_{\mathcal{F}}\right)^{-1}\mathbf{P}_{\mathcal{F}}^T  \mathbf{D}_i\mathbf{A}^T\mathbf{P}_{\sigma}^T\\
% % & =  \sum_{i = 0}^n \sum_{\mathcal{F}: |\mathcal{F}| = i} (1-q)^i q^{n -i} \mathbf{Q}_{\sigma}\mathbf{A}\mathbf{Q}_{\sigma}^T\mathbf{Q}_{\sigma}\mathbf{D}_j\mathbf{Q}_{\sigma}^T\mathbf{Q}_{\sigma}\mathbf{P}_{\mathcal{F}} \left(\sum_{i=1}^m\mathbf{P}_{\mathcal{F}}^T\mathbf{Q}_{\sigma}^T\mathbf{Q}_{\sigma} \mathbf{D}_i \mathbf{Q}_{\sigma}^T\mathbf{Q}_{\sigma}\mathbf{A}^T\mathbf{Q}_{\sigma}^T\mathbf{Q}_{\sigma}\mathbf{A}\mathbf{Q}_{\sigma}^T\mathbf{Q}_{\sigma}\mathbf{D}_i\mathbf{Q}_{\sigma}^T\mathbf{Q}_{\sigma}\mathbf{P}_{\mathcal{F}}\right)^{-1}\\ 
% % & \quad \quad  \quad \quad \quad \quad \quad \mathbf{P}_{\mathcal{F}}^T\mathbf{Q}_{\sigma}^T\mathbf{Q}_{\sigma}  \mathbf{D}_i \mathbf{Q}_{\sigma}^T\mathbf{Q}_{\sigma}\mathbf{A}^T\mathbf{Q}_{\sigma}^T\\
% & \overset{(f)}=  \sum_{r = 1}^n \sum_{\mathcal{F}: |\mathcal{F}| = r} (1-q)^r q^{n -r} \mathbf{A}\mathbf{Q}_{\pi}\mathbf{D}_i\mathbf{Q}_{\pi}^T\mathbf{Q}_{\pi}\mathbf{P}_{\mathcal{F}}\\ 
% & \quad \quad  \quad \left(\sum_{t=1}^m\mathbf{P}_{\mathcal{F}}^T\mathbf{Q}_{\pi}^T\mathbf{Q}_{\pi} \mathbf{D}_t \mathbf{Q}_{\pi}^T\mathbf{A}^T\mathbf{A}\mathbf{Q}_{\pi}\mathbf{D}_t\mathbf{Q}_{\pi}^T\mathbf{Q}_{\pi}\mathbf{P}_{\mathcal{F}}\right)^{-1}\mathbf{P}_{\mathcal{F}}^T\mathbf{Q}_{\pi}^T\mathbf{Q}_{\pi}  \mathbf{D}_i \mathbf{Q}_{\pi}^T\mathbf{A}^T\\
% & \overset{(g)}=  f_{ii}(\mathbf{A}, \mathbf{Q}_{\pi}\mathbf{D} \mathbf{Q}_{\pi}^T) .
% \end{align*}








% \begin{align*}
% &\mathbf{Q}_{\sigma} f_{ii}( \mathbf{A}, \mathbf{D}) \mathbf{Q}_{\sigma}^T\\
% & \overset{(e)}=  \sum_{i = 0}^n \sum_{\mathcal{F}: |\mathcal{F}| = i} (1-q)^i q^{n -i} \mathbf{Q}_{\sigma} \mathbf{A}\mathbf{D}_j\mathbf{P}_{\mathcal{F}} \left( \sum_ {i= 1}^m\mathbf{P}_{\mathcal{F}}^T \mathbf{D}_i \mathbf{A}^T\mathbf{A}\mathbf{D}_i \mathbf{P}_{\mathcal{F}}\right)^{-1}\mathbf{P}_{\mathcal{F}}^T  \mathbf{D}_i\mathbf{A}^T\mathbf{Q}^T\\
% % & =  \sum_{i = 0}^n \sum_{\mathcal{F}: |\mathcal{F}| = i} (1-q)^i q^{n -i} \mathbf{Q}\mathbf{A}\mathbf{Q}^T\mathbf{Q}\mathbf{D}_j\mathbf{Q}^T\mathbf{Q}\mathbf{P}_{\mathcal{F}} \left(\sum_{i=1}^m\mathbf{P}_{\mathcal{F}}^T\mathbf{Q}^T\mathbf{Q} \mathbf{D}_i \mathbf{Q}^T\mathbf{Q}\mathbf{A}^T\mathbf{Q}^T\mathbf{Q}\mathbf{A}\mathbf{Q}^T\mathbf{Q}\mathbf{D}_i\mathbf{Q}^T\mathbf{Q}\mathbf{P}_{\mathcal{F}}\right)^{-1}\\ & \quad \quad  \quad \quad \quad \quad \quad \mathbf{P}_{\mathcal{F}}^T\mathbf{Q}^T\mathbf{Q}  \mathbf{D}_i \mathbf{Q}^T\mathbf{Q}\mathbf{A}^T\mathbf{Q}^T\\
% & \overset{(f)}=  \sum_{i = 0}^n \sum_{\mathcal{F}: |\mathcal{F}| = i} (1-q)^i q^{n -i} \mathbf{A}\mathbf{Q}\mathbf{D}_j\mathbf{Q}^T\mathbf{Q}\mathbf{P}_{\mathcal{F}} \left(\sum_{i=1}^m\mathbf{P}_{\mathcal{F}}^T\mathbf{Q}^T\mathbf{Q} \mathbf{D}_i \mathbf{Q}^T\mathbf{A}^T\mathbf{A}\mathbf{Q}\mathbf{D}_i\mathbf{Q}^T\mathbf{Q}\mathbf{P}_{\mathcal{F}}\right)^{-1}\\ & \quad \quad  \quad \quad \quad \quad \quad \mathbf{P}_{\mathcal{F}}^T\mathbf{Q}^T\mathbf{Q}  \mathbf{D}_i \mathbf{Q}^T\mathbf{A}^T\\
% & \overset{(g)}=  f_{ii}(\mathbf{A}, \mathbf{Q}\mathbf{D} \mathbf{Q}^T) 
% \end{align*}
% \sifat{I will have to find a way to make this look better.}\\


Here, $(f)$ holds by definition of $f_{ii}(\bfA, \bfD)$. $(g)$ holds by definition of $h_{ii}(\bfA, F, \bfD)$. $(h)$ holds since $\mathbf{Q}_{\pi}^T \mathbf{Q}_{\pi} = \mathbf{I}_n$ and $\bfP_{\sigma} \bfA \bfQ_{\pi}^T = \bfA$. $(i)$ holds as a consequence of $(e)$. Since for each $t \in [m]$, we have $\mathbf{D}_t \overset{d}{=} \mathbf{Q}_{\pi}\mathbf{D}_t \mathbf{Q}_{\pi}^T$, 
$ f_{ii}(\mathbf{A}, \mathbf{Q}_{\pi}\mathbf{D} \mathbf{Q}_{\pi}^T) \overset{d}{=}  f_{ii}(\mathbf{A}, \mathbf{D}) 
.$
Consequently, 
\[ \mathbf{P}_{\sigma} f_{ii}( \mathbf{A}, \mathbf{D}) \mathbf{P}_{\sigma}^T \overset{d}{=} f_{ii}( \mathbf{A}, \mathbf{D}),\]

and thus $\mathbb{E}_{\bfD}[\mathbf{P}_{\sigma} f_{ii}( \mathbf{A}, \mathbf{D}) \mathbf{P}_{\sigma}^T] = \mathbb{E}_{\bfD}[f_{ii}( \mathbf{A}, \mathbf{D})]$. Using this distributional equivalence and the properties of $\mathbf{P}_\sigma$ we have,
% For $i, j \in [n]$, let $\mathbf{Q}$ be a permutation such that $\mathbf{Q}^T \mathbf{e}_i = \mathbf{e}_j$.
\begin{align*}
(\mathbb{E}_{\mathbf{D}} [f_{ii}( \mathbf{A}, \mathbf{D}) \mathbf{1}_k ])(u) 
%&=\mathbb{E}_{\mathbf{D}} [(f_{ii}( \mathbf{A}, \mathbf{D}) \mathbf{1}_k)(u) ]\\ 
& = \mathbb{E}_{\mathbf{D}}[\mathbf{e}_u^T f_{ii}( \mathbf{A}, \mathbf{D}) \mathbf{1}_k]\\
%&= \mathbf{e}_u^T  \mathbb{E}_{\mathbf{D}}[f_{ii}( \mathbf{A}, \mathbf{D})] \mathbf{1}_k\\
& = \mathbf{e}_u^T  \mathbb{E}_{\mathbf{D}}[\mathbf{P}_{\sigma}f_{ii}( \mathbf{A}, \mathbf{D}) \mathbf{P}_{\sigma}^T] \mathbf{1}_k\\
%& =  \mathbb{E}_{\mathbf{D}}[ \mathbf{e}_u^T\mathbf{P}_{\sigma}f_{ii}( \mathbf{A}, \mathbf{D}) \mathbf{P}_{\sigma}^T\mathbf{1}_k] \\
& = \mathbb{E}_{\mathbf{D}}[ \mathbf{e}_v^Tf_{ii}( \mathbf{A}, \mathbf{D}) \mathbf{1}_k]\\
%& = \mathbb{E}_{\mathbf{D}} [(f_{ii}( \mathbf{A}, \mathbf{D}) \mathbf{1}_k)(v) ]\\
& = (\mathbb{E}_{\mathbf{D}} [f_{ii}( \mathbf{A}, \mathbf{D}) \mathbf{1}_k ])(v).
\end{align*}


% \begin{align*}
% (\mathbb{E}_{\mathbf{D}} [f_{ii}( \mathbf{A}, \mathbf{D}) \mathbf{1}_n ])(u) &=\mathbb{E}_{\mathbf{D}} [(f_{ii}( \mathbf{A}, \mathbf{D}) \mathbf{1}_n)(u) ]\\ 
% & = \mathbb{E}_{\mathbf{D}}[\mathbf{e}_u^T f_{ii}( \mathbf{A}, \mathbf{D}) \mathbf{1}_n]\\
% &= \mathbf{e}_u^T  \mathbb{E}_{\mathbf{D}}[f_{ii}( \mathbf{A}, \mathbf{D}] \mathbf{1}_n\\
% & = \mathbf{e}_u^T  \mathbb{E}_{\mathbf{D}}[\mathbf{P}_{\sigma}f_{ii}( \mathbf{A}, \mathbf{D}) \mathbf{P}_{\sigma}^T] \mathbf{1}_n\\
% & =  \mathbb{E}_{\mathbf{D}}[ \mathbf{e}_u^T\mathbf{P}_{\sigma}f_{ii}( \mathbf{A}, \mathbf{D}) \mathbf{P}_{\sigma}^T\mathbf{1}_n] \\
% & = \mathbb{E}_{\mathbf{D}}[ \mathbf{e}_v^Tf_{ii}( \mathbf{A}, \mathbf{D}) \mathbf{1}_n]\\
% & = \mathbb{E}_{\mathbf{D}} [(f_{ii}( \mathbf{A}, \mathbf{D}) \mathbf{1}_n)(v) ]\\
% & = (\mathbb{E}_{\mathbf{D}} [f_{ii}( \mathbf{A}, \mathbf{D}) \mathbf{1}_n ])(v).
% \end{align*}




% \begin{align*}
% (\mathbb{E}_{\mathbf{D}} [f_{ii}( \mathbf{A}, \mathbf{D}) \mathbf{1}_n ])_u &=\mathbb{E}_{\mathbf{D}} [(f_{ii}( \mathbf{A}, \mathbf{D}) \mathbf{1}_n)_u ]\\ & = \mathbb{E}_{\mathbf{D}}[\mathbf{e}_u^T f_{ii}( \mathbf{A}, \mathbf{D}) \mathbf{1}_n]\\
% &= \mathbf{e}_u^T  \mathbb{E}_{\mathbf{D}}[f_{ii}( \mathbf{A}, \mathbf{D}] \mathbf{1}_n\\
% & = \mathbf{e}_u^T  \mathbb{E}_{\mathbf{D}}[\mathbf{Q}f_{ii}( \mathbf{A}, \mathbf{D}) \mathbf{Q}^T] \mathbf{1}_n\\
% & =  \mathbb{E}_{\mathbf{D}}[ \mathbf{e}_u^T\mathbf{Q}f_{ii}( \mathbf{A}, \mathbf{D}) \mathbf{Q}^T\mathbf{1}_n] \\
% & = \mathbb{E}_{\mathbf{D}}[ \mathbf{e}_v^Tf_{ii}( \mathbf{A}, \mathbf{D}) \mathbf{1}_n]\\
% & = \mathbb{E}_{\mathbf{D}} [(f_{ii}( \mathbf{A}, \mathbf{D}) \mathbf{1}_n)_v ]\\
% & = (\mathbb{E}_{\mathbf{D}} [f_{ii}( \mathbf{A}, \mathbf{D}) \mathbf{1}_n ])_v.
% \end{align*}

By our assumptions $u,v \in [k]$ are arbitrary. Thus, it follows that all the entries of $\mathbb{E}_{\mathbf{D}} [f_{ii}( \mathbf{A}, \mathbf{D}) \mathbf{1}_k ]$ are equal. Moreover, $h_{ii}(\bfA, \cF, \bfD)$ is positive semidefinite for every non-empty $\cF$. For $\mathcal F=[n]$, since $\mathbf{A}^T\mathbf{1}_k \ne \mathbf{0}_n$ and $\mathbf D_i$ is invertible,
\[
\mathbf 1_k^T h_{ii}(\mathbf A,[n],\mathbf D)\mathbf 1_k
=
(\mathbf D_i\mathbf A^T\mathbf 1_k)^T
(\mathbf B^T\mathbf B)^{-1}
(\mathbf D_i\mathbf A^T\mathbf 1_k)>0.
\]
Since $q < 1$, $\cF = [n]$ has nonzero probability. It follows that
$
\mathbf 1_k^T\mathbb{E}_{\mathbf D}
[f_{ii}(\mathbf A,\mathbf D)]\mathbf 1_k>0
$.
Hence, $\mathbb{E}_{\mathbf D}[f_{ii}(\mathbf A,\mathbf D)\mathbf 1_k]$
is a nonzero vector. Consequently,
$
\mathbb{E}_{\mathbf D}[f_{ii}(\mathbf A,\mathbf D)\mathbf 1_k]
=\beta_i\mathbf 1_k
$, where $\beta_i$ is a nonzero constant.  
% Also, since ${f}_{ii}(\bfA,\bfD)$ is positive definite, $\mathbb{E}_{\mathbf{D}}[{f}_{ii}(\bfA,\bfD)]$ is positive definite. Thus, $\mathbf{1}_{k}^T\mathbb{E}_{\mathbf{D}}[{f}_{ii}(\bfA,\bfD)] \mathbf{1}_{k} = \mathbf{1}_{k}^T\mathbb{E}_{\mathbf{D}}[{f}_{ii}(\bfA,\bfD) \mathbf{1}_{k}] > 0$. So, $\mathbb{E}_{\mathbf{D}}[{f}_{ii}(\bfA,\bfD) \mathbf{1}_{k}]$ is a non-zero vector. Consequently, $\mathbb{E}_{\mathbf{D}} [ {f}_{ii}(\bfA,\bfD)\mathbf{1}_{k}] = \beta_i \mathbf{1}_k$, where $\beta_i$ is a non-zero constant. 
Therefore, for $i \in [m]$,
\[\mathbb{E}_{\mathbf{D}} [\mathbb{E}_{\mathcal{F}}[ \mathbf{B}_{\mathcal{F}} \mathbf{r}_i*] ] = \begin{bmatrix}\mathbf{0}_k \\ \vdots \\ \beta_i\mathbf{1}_k\\ \vdots \\ \mathbf{0}_k \end{bmatrix} = \beta_i\mathbf{f}_i.\]
Since the diagonal matrices $\mathbf{D}_i$ are i.i.d., $\mathbb{E}_{\mathbf{D}}[ {f}_{ii}(\bfA,\bfD) \mathbf{1}_{k}] = \mathbb{E}_{\mathbf{D}}[{f}_{jj}(\bfA,\bfD) \mathbf{1}_{k}]$ for all $i,j \in [m]$. Therefore, $\beta_i = \beta_ j = \beta$ (let) for all $i ,j \in [m]$.
Since $\bfR = \begin{bmatrix}\mathbf{r}_1* & \mathbf{r}_2* & \dots & \mathbf{r}_m* \end{bmatrix}$, and $\bfF = \begin{bmatrix}\mathbf{f}_1 & \mathbf{f}_2 & \dots & \mathbf{f}_m \end{bmatrix}$, we have  \[\mathbb{E}_{\mathbf{D}} [\mathbb{E}_{\mathcal{F}}[ \mathbf{B}_{\mathcal{F}} \mathbf{R}] ] = \beta\mathbf{F},\]
where $\beta$ is a nonzero constant.
\end{proof}

\begin{cor}
\label{cor:SRG_unbiased}
 Suppose that the assignment matrix $\mathbf{A}$ is the adjacency matrix of a $(n, \delta, \lambda, \mu)$ vertex-transitive strongly regular graph $G = ([n], E)$. 
 %and let $\mathbf{A} = \mathbf{E}$ so that $\mathbf{A}$ is an $(n,k,\delta, \gamma)$ assignment matrix. 
Then, for the construction as in $\eqref{eq:effAD1AD2}$,

 \[ \mathbb{E}_{\mathbf{D}}[\mathbb{E}_{\mathcal{F}}[\mathbf{B}_{\cF} \mathbf{R}]]= \beta\mathbf{F},\]
where $\beta$ is a nonzero constant.

\end{cor}

\begin{proof}
By definition, for a vertex transitive strongly regular graph, for any $u, v \in [n]$ there exists an automorphism $\sigma$ such that $\sigma(v) = u$. So, for                    the corresponding permutation matrix $\bfP_{\sigma}$, $\mathbf{e}_{u}^T\bfP_{\sigma} = \mathbf{e}_{v}^T$ and $\bfP_{\sigma}\bfA \bfP_{\sigma}^T = \bfA$. Also, as proved in Corollary $\ref{cor:SRG}$, $\bfB_{\cF}^T \bfB_{\cF}$ is invertible for non-empty $\cF \subseteq [n]$. Since $\bfA^T \mathbf{1}_k = \delta \mathbf{1}_n$, the conclusion follows from Lemma $\ref{lemma:SRG_COSET_unbiasedness}$.
\end{proof}


\begin{cor}
\label{cor:COSET_unbiased}
Suppose that the assignment matrix $\mathbf{A}$ is the bi-adjacency matrix of a $(k,m,\delta)$ coset bipartite graph such that $k = p^a$, where $p$ is a prime, $a \in \mathbb{Z}_{\ge 1}$ and $p \nmid \delta$. Then, for the construction as in $\eqref{eq:effAD1AD2}$ with $\epsilon > 0$,
\[ \mathbb{E}_{\mathbf{D}}[\mathbb{E}_{\mathcal{F}}[\mathbf{B}_{\cF} \mathbf{R}]]= \beta\mathbf{F},\]
where $\beta$ is a nonzero constant.
\end{cor}

\begin{proof}
% As proved in Corollary $\ref{cor:bipartite}$, $\bfB_{\cF}^T \bfB_{\cF}$ is invertible for non-empty $\cF \subseteq [n]$.
Let $G = (L \cup R, E)$ be a  $(k, m, \delta)$ coset bipartite graph. So, $R = \mathbb{Z}_{mk}$ and $L = \{\, i+H : i=0,1,\dots,k-1 \,\}$ where $H$ is an order $m$ subgroup of $\mathbb{Z}_{mk}$. By construction, $i + H \in L$ is adjacent to $x \in R = \mathbb{Z}_{mk}$ if and only if $x \in i + S$, for $i \in \{0, \dots, k-1\}$ and $x \in \mathbb{Z}_{mk}$. Consequently, for its bi-adjacency matrix $\bfA$, we have $\bfA(i+1, x+1) = 1$ if and only if $x \in i + S $ (note that the sets of row and column indices of $\bfA$ are $[k]$ and $[mk]$ respectively).
For $g \in \{ 0, \dots, k-1\}$, define the row permutation $\sigma_g:[k]\to[k]$ by
\[
\sigma_g(i) =  ((i-1+g) \bmod{k}) + 1,
\]
and the column permutation $\tau_g:[mk]\to [mk]$ by
\[
\tau_g(x) =  ((x-1+g) \bmod{mk}) +1 .
\]
Let $\mathbf{P}_{\sigma_g}\in\{0,1\}^{k\times k}$ and $\mathbf{Q}_{\tau_g}\in\{0,1\}^{(mk)\times(mk)}$ denote the corresponding permutation matrices.
% \[
% \mathbf{P}_g \mathbf{e}_i = \mathbf{e}_{\sigma_g(i)},\qquad \mathbf{Q}_g \mathbf{e}_x = \mathbf{e}_{\tau_g(x)}.
% \]
We will first show that for every $g\in \{0 , \dots, k-1\}$,
% \begin{equation}\label{eq:PAQt}
$\bfP_{\sigma_g}\bfA \mathbf{Q}_{\tau_g}^T= \bfA.$
% \end{equation}
We have,
\[
(\bfP_{\sigma_g}\bfA \mathbf{Q}_{\tau_g}^T )(i+1,x+1)
= \bfA(\sigma_g^{-1}(i+1),\,\tau_g^{-1}(x+1))
% = \bfA(i+1-g ,\;x+1-g).
\]
Moreover, by the definition of $\sigma_g$ and $\tau_g$,
\[
\sigma_g^{-1}(y) = ((y-1-g)\bmod{k})+1,
\qquad
\tau_g^{-1}(z) = ((z-1-g)\bmod{mk})+1.
\]
Thus,
\[
(\bfP_{\sigma_g}\bfA \mathbf{Q}_{\tau_g}^T )(i+1,x+1)
= \bfA\!\Big( ((i-g)\bmod{k})+1,\; ((x-g)\bmod{mk})+1 \Big).
\]


But, by definition, $\bfA\big((i-g)\bmod k + 1,\,(x-g)\bmod mk + 1\big)=1$
\text{if and only if }
$(x-g)\bmod mk \in \big((i-g)\bmod k\big)+S$. 
% \begin{align*}
% &\bfA\!\Big( (i-g)\bmod{k}+1,\; (x-g)\bmod{mk}+1 \Big)=1 \\
% &\iff \big((x-g)\bmod{mk}\big)\in \big((i-g)\bmod{k}\big)+S.
% \end{align*}
Let $i':=(i-g)\bmod{k}\in\{0,\dots,k-1\}$ and $x':=(x-g)\bmod{mk}\in\{0,\dots,mk-1\}$.
% By definition of $\bfA$, the above is equivalent to $x'\in i'+S$, i.e.,
% \[
% x'-i'\in S.
% \]
We can view $i',x'$ as elements of $\mathbb{Z}_{mk}$ (via the natural embedding of $\{0,\dots,k-1\}$ into $\mathbb{Z}_{mk}$). 
Since $i' = i-g \pmod{k}$, there exists an integer $t$ such that
\[
i' = i-g + tk.
\]
Also, since $x' = x-g \pmod{mk}$, there exists an integer $u$ such that
\[
x' = x-g + umk.
\]
So, $x-g+umk \in i -g + tk + S$. By definition of $S$, $tk +S = S$. So, $x-g+umk \in i -g + S$. Therefore, by definition of $S$ there exists $b \in S$ such that for some integer $v$, 
\[ x-g + umk = i-g+b + vmk.\] So, $x = i+b + (v-u) mk$. Since $b \in S$, $b + (v-u)mk \in S$. Now, $ i + b + (v-u)mk = (i +  b  + (v-u)mk) \bmod mk$, since $x < mk$. Thus, $i +  b  + (v-u)mk \in i +S$ and therefore, $x \in i +S$. Similarly we can show that if $x \in i +S$, then $x' \in i' +S$. Consequently, 
$\bfA\!\Big( ((i-g)\bmod{k})+1,\; ((x-g)\bmod{mk})+1 \Big)=1$
if and only if  $\bfA(i+1,x+1)=1$. Since the entries of $\bfA$ are either $0$ or $1$,  $(\bfP_{\sigma_g} \bfA \mathbf{Q}_{\tau_g}^{T})(i+1,x+1)= \bfA(i+1,x+1)$ for all $i \in \{0, \dots, k-1\}$ and $x \in \mathbb{Z}_{mk}$. So, $\bfP_{\sigma_g} \bfA \mathbf{Q}_{\tau_g}^{T} = \bfA$. Now, given any pair of rows $i,j \in [k]$, we can choose $g =  (i-j) \bmod k$. Then, $\sigma_g(j) = i$. Thus,
$
\mathbf{e}_i^T\bfP_{\sigma_g}  = \mathbf{e}_j^T.
$
So, for any $i,j \in [k]$, there exist permutation matrices $\bfP_{\sigma_g}$ and $\bfQ_{\tau_g}$ such that $\bfP_{\sigma_g} \bfA \bfQ_{\tau_g}^T = \bfA$ and $\mathbf{e}_i^T \bfP_{\sigma_g} = \mathbf{e}_j^T$. Also, as proved in Proposition $\ref{prop: cos_inv}$, for $k = p^a$ where $p$ is a prime, $a \in \mathbb{Z}_{\ge 1}$ and $p \nmid \delta$, $\bfB$ is almost surely invertible and thus $\bfB_{\cF}^T \bfB_{\cF}$ is invertible for non-empty $\cF \subseteq [n]$. Since $\bfA^T \mathbf{1}_k = \delta \mathbf{1}_n$, the conclusion follows from Lemma $\ref{lemma:SRG_COSET_unbiasedness}$.
% so $\bfP_g$ sends row $i$ to row $j$; by \eqref{eq:PAQt} the corresponding $Q_g$ satisfies $\bfP_g \bfA \mathbf{Q}_g^{\top}= \bfA$. Hence for \emph{any} pair of rows $(i,j)$ there exist permutation matrices $(\bfP, \mathbf{Q})$ with $\bfP \mathbf{e}_i= \mathbf{e}_j$ and $\bfP \bfA \mathbf{Q}^{\top}= \bfA$.                                                                            
\end{proof}









\subsection{Bounded Variance of the Computed Gradient}

For the construction in \eqref{eq:hada_coset_cons}, the corresponding unbiasedness result is established in Lemma \ref{lemma:Cost_Hadamard_Convergence} in Appendix \ref{app:hada_cos_unbiased}.


\begin{thm}
\label{thm:Bound_variance}
Suppose that the spectral norm of $\mathbf{Z}$ is bounded, i.e.,
there exists a constant $c_z>0$ such that $\|\mathbf{Z}\|_2\leq c_z$. Then, for the construction as in $\eqref{eq:effAD1AD2}$ with BIBD, vertex-transitive SRG and coset bipartite graph assignment matrices, and for the construction as in $\eqref{eq:hada_coset_cons}$, the variance of the rescaled computed gradient is bounded.
\end{thm}

\begin{proof}
% \aditya{why do you assume a bound on each of the individual gi's?}
%   \sifat{That assumption is not necessary. Fixed accordingly.}
% \begin{equation}
% \|\mathbf Z\|_2^2
% \leq
% \|\mathbf Z\|_F^2
% =
% \sum_{i=1}^{k}\|\mathbf g_i\|_2^2
% \leq kG^2.
% \end{equation}
By Lemma \ref{lemma:BIBD_unbiasedness} for the BIBD construction, Corollaries~\ref{cor:SRG_unbiased} and \ref{cor:COSET_unbiased} for the
vertex-transitive SRG and coset bipartite graph constructions as in
\eqref{eq:effAD1AD2}, respectively, the rescaled computed gradient is unbiased. Thus,
\begin{equation}
\mathbb E_{\mathbf D}
\left[
\mathbb E_{\mathcal F}
\left[
\mathbf B_{\mathcal F}\mathbf R
\right]
\right]
=
\zeta\mathbf F.
\end{equation}
Therefore,
\begin{align}
\mathbb E_{\mathbf D}
\left[
\mathbb E_{\mathcal F}
\left[
\left\|
\mathbf B_{\mathcal F}\mathbf R-\mathbf F
\right\|_F^2
\right]
\right] =
\mathbb E_{\mathbf D}
\left[
\mathbb E_{\mathcal F}
\left[
\left\|
\mathbf B_{\mathcal F}\mathbf R-\zeta\mathbf F
\right\|_F^2
\right]
\right]
+
\left\|
(\zeta-1)\mathbf F
\right\|_F^2.
\end{align}
Consequently,
\begin{align}
\mathbb E_{\mathbf D}
\left[
\mathbb E_{\mathcal F}
\left[
\left\|
\mathbf B_{\mathcal F}\mathbf R-\zeta\mathbf F
\right\|_F^2
\right]
\right]
\leq
\mathbb E_{\mathbf D}
\left[
\mathbb E_{\mathcal F}
\left[
\operatorname{Err}_{\mathcal F}(\mathbf B)
\right]
\right],
\end{align}
where
\begin{equation}
\left\|
\mathbf B_{\mathcal F}\mathbf R-\mathbf F
\right\|_F^2
=
\operatorname{Err}_{\mathcal F}(\mathbf B)
\end{equation}
since $\mathbf R$ is the optimal decoding matrix. Hence,
\begin{align}
\mathbb E_{\mathbf D}
\left[
\mathbb E_{\mathcal F}
\left[
\left\|
\frac{1}{\zeta}
\mathbf Z\mathbf B_{\mathcal F}\mathbf R
-
\mathbf Z\mathbf F
\right\|_F^2
\right]
\right]
&\leq
\frac{\|\mathbf Z\|_2^2}{\zeta^2}
\mathbb E_{\mathbf D}
\left[
\mathbb E_{\mathcal F}
\left[
\left\|
\mathbf B_{\mathcal F}\mathbf R-\zeta\mathbf F
\right\|_F^2
\right]
\right]
\\
&\leq
\frac{c_z^2}{\zeta^2}
\mathbb E_{\mathbf D}
\left[
\mathbb E_{\mathcal F}
\left[
\operatorname{Err}_{\mathcal F}(\mathbf B)
\right]
\right].
\end{align}

Corollaries~\ref{cor:expected_error_BIBD_strr}, \ref{cor:expected_error_SRG_strr}, and \ref{cor:expected_error_coset_strr} in Appendix \ref{sec:exp_app_err_rand_strr} show that for each of the corresponding
constructions, there exists a finite constant $U$ depending only on
the construction and straggler model parameters such that
\begin{equation}
\mathbb E_{\mathbf D}
\left[
\mathbb E_{\mathcal F}
\left[
\operatorname{Err}_{\mathcal F}(\mathbf B)
\right]
\right]
\leq U.
\end{equation}
Therefore,
\begin{equation}
\mathbb E_{\mathbf D}
\left[
\mathbb E_{\mathcal F}
\left[
\left\|
\frac{1}{\zeta}
\mathbf Z\mathbf B_{\mathcal F}\mathbf R
-
\mathbf Z\mathbf F
\right\|_F^2
\right]
\right]
\leq
\frac{c_z^2U}{\zeta^2}.
\end{equation}

Also, for the construction as in \eqref{eq:hada_coset_cons}, by Lemma \ref{lemma:Cost_Hadamard_Convergence}, $\mathbb E_{\mathcal F}
\left[
\mathbf B_{\mathcal F}\mathbf R
\right]
=
\zeta\mathbf F$, and by Lemma \ref{lemma:hada_coset_bound}, $
\mathbb E_{\mathcal F}
\left[
\operatorname{Err}_{\mathcal F}(\mathbf B)
\right] 
\leq  nq$. Therefore, by the same argument as above, $\mathbb E_{\mathcal F}
\left[
\left\|
\frac{1}{\zeta}
\mathbf Z\mathbf B_{\mathcal F}\mathbf R
-
\mathbf Z\mathbf F
\right\|_F^2
\right]
\leq
\frac{c_z^2 nq}{\zeta^2}.$ So the variance bound holds with $U = nq$.

\end{proof}



\subsection{Stationary-Point Convergence}

\begin{cor}
Consider the construction as in $\eqref{eq:effAD1AD2}$ with BIBD, vertex-transitive SRG and coset bipartite graph assignment matrices, and the construction as in $\eqref{eq:hada_coset_cons}$,
and suppose that the spectral norm of $\mathbf{Z}$ is bounded, i.e.,
there exists a constant $c_z>0$ such that $\|\mathbf{Z}\|_2\leq c_z$. Furthermore, assume that $L$ is bounded below by $L_{\inf}$ and has
an $L_s$-Lipschitz-continuous gradient. Let
\begin{equation}
\sigma^2=\frac{c_z^2U}{\zeta^2}
\end{equation}
and define
\begin{equation}
D_L
=
\left(
\frac{2\left(L(\mathbf w^{(1)})-L_{\inf}\right)}
{L_s}
\right)^{1/2}.
\end{equation}
For a given number of iterations $T$, let $\mathbf{v}^{(t)}$ be the computed gradient obtained from the columns of $\frac{1}{\zeta}\mathbf{Z}^{(t)}\mathbf{B}_{\cF}^{(t)}\mathbf{R}^{(t)}$ at iteration $t$.
\begin{equation}
\mathbf w^{(t+1)}
=
\mathbf w^{(t)}
-
\eta
\mathbf{v}^{(t)},
\qquad t=1,\ldots,T,
\end{equation}
where
\begin{equation}
\eta
=
\min\left\{
\frac{1}{L_s},
\frac{\widetilde D}{\sigma\sqrt{T}}
\right\}
\end{equation}
for some $\widetilde D>0$. Let $R_T$ be uniformly distributed over
$\{1,\ldots,T\}$. Then,
\begin{equation}
\frac{1}{L_s}
\mathbb E
\left[
\left\|
\nabla L\left(\mathbf w^{(R_T)}\right)
\right\|_2^2
\right]
\leq
\frac{L_sD_L^2}{T}
+
\left(
\widetilde D+\frac{D_L^2}{\widetilde D}
\right)
\frac{\sigma}{\sqrt{T}}.
\end{equation}
Consequently,
\begin{equation}
\lim_{T\rightarrow\infty}
\mathbb E
\left[
\left\|
\nabla L\left(\mathbf w^{(R_T)}\right)
\right\|_2^2
\right]
=0.
\end{equation}
\end{cor}

\begin{proof}
By Lemmas \ref{lemma:BIBD_unbiasedness} and $\ref{lemma:Cost_Hadamard_Convergence}$ and Corollaries \ref{cor:SRG_unbiased} and \ref{cor:COSET_unbiased}, the
rescaled computed gradient is unbiased. Theorem \ref{thm:Bound_variance} establishes that its variance is bounded. Therefore, under the stated assumptions, the result follows from
Corollary~2.2 in~\cite{doi:10.1137/120880811}.
\end{proof}

% For the Hadamard matrix based coset bipartite graph construction in \eqref{eq:hada_coset_cons}, the corresponding unbiasedness result is established in Lemma \ref{lemma:Cost_Hadamard_Convergence} in Appendix \ref{app:hada_cos_unbiased}.




















\section{Numerical Experiments}
\label{sec:numerical_exp}

% \sifat{I need to think about what SRG to choose. Bound seems to be pretty loose. Also, getting infeasibility for m = 3 for the reconstruction scheme with $\pm 1$ distribution.}
In what follows, we refer to the construction obtained by stacking the underlying assignment matrix vertically $m$ times as the baseline construction (similar to  \eqref{eq:trivial_cons}). We present the results of numerical experiments to demonstrate the performance of our constructions and compare them with the baseline construction, along with the lower bound derived in Theorem \ref{thm:lowerbound} (see \cite{cE_agc_repo} for the source code used for the experiments). 

% \begin{figure}[t!]
% %\vspace{.4 cm}
%     \centering
%     \includegraphics[width=0.45\textwidth]{Journal/Figure_3_v2.pdf}
%     \caption{{\small Performance of the construction as in \eqref{eq:effAD1AD2} for $\epsilon = 0$,  and  $m = 2$ with a $(91,91,10,10,1)$ BIBD.}}
%     \label{fig:AD1bound}
% \end{figure}

\begin{figure*}[t!]
    \centering

    \begin{subfigure}[t]{0.45\textwidth}
        \centering
        \includegraphics[
            width=\linewidth,
            ,
            clip
        ]{Journal/Fig_3_a.pdf}
        \caption{Comparison with the exact communication-efficient gradient coding scheme in \cite{tayyebehM21} for $m =2$.}
        \label{fig:bibd_31_exact}
        
    \end{subfigure}
    \hfill
    \begin{subfigure}[t]{0.45\textwidth}
        \centering
        \includegraphics[
            width=\linewidth,
          ,
            clip
        ]{Journal/Fig_3_b.pdf}
        \caption{Performance for different values of $m$.}
        \label{fig:different_m}
        
    \end{subfigure}

    \caption{Performance of the construction as in \eqref{eq:effAD1AD2} with
    $\epsilon=0$ and a $(31,31,6,6,1)$ BIBD.}
    \label{fig:bibd_comparison}
\end{figure*}



Firstly, we chose a $(31,31,6,6,1)$ BIBD as the assignment matrix and $m = 2$ to evaluate the construction discussed in Section \ref{sec:construction_RDM} (see \eqref{eq:effAD1AD2}). To compare with the upper bound \eqref{eq:AD1AD2err_BIBD}, the empirical average error was calculated as follows. For a given number of stragglers $s$, $100$ different realizations of the encoding matrix $\mathbf{B}$ were generated by generating the matrices $\mathbf{D}_1, \mathbf{D}_2$ each time from the underlying distribution with $\epsilon = 0$. For each of these realizations, the approximation error \eqref{eq:ApproxError} was calculated for each of $1000$ random choices of stragglers. Finally, the average of the errors $(100 \times 1000)$ was calculated (see \autoref{fig:bibd_31_exact}). Notice that this average error lies close to the upper bound \eqref{eq:AD1AD2err_BIBD}. Also, the upper bound is significantly lower than the error for the baseline construction. We additionally compare this construction against the exact communication-efficient gradient coding scheme in~\cite{tayyebehM21}. For the same set of parameters, the exact scheme is 
designed to tolerate $s=\gamma-m=4$ stragglers. Since its polynomial encoder 
is linear, we represent it in encoding matrix form and evaluate its error using 
the same optimal least-squares decoder. For each value of $s$, we report the maximum approximation error of the exact scheme over the considered straggler sets. %We note that the exact scheme can exhibit numerical instability in this setting. 
We point out that the condition number of the corresponding matrix is very high and can lead to numerical instability in scenarios with higher number of workers.
As shown in \autoref{fig:bibd_31_exact}, 
it achieves zero error within its designed tolerance, but its error increases 
sharply beyond this threshold, whereas the proposed construction exhibits a 
more gradual increase.

We further evaluate this construction for $m\in\{2,3,4,5\}$ using the same procedure. As shown in \autoref{fig:different_m}, the empirical average error remains close to the corresponding upper bound for all the considered values of $m$, while the lower bound is also shown for comparison.



% Firstly, we chose a $(31,31,6,10,1)$ BIBD as the assignment matrix and $m =2$ to evaluate the construction discussed in Section \ref{sec:construction_RDM} (see \eqref{eq:effAD1AD2}). To compare with the upper bound \eqref{eq:AD1AD2err_BIBD}, the empirical average error was calculated as follows. For a given number of stragglers $s$, $100$ different realizations of the encoding matrix $\mathbf{B}$ were generated by generating the matrices $\mathbf{D}_1, \mathbf{D}_2$ each time from the underlying distribution with $\epsilon = 0$. For each of these realizations, the approximation error \eqref{eq:ApproxError} was calculated for each of $1000$ random choices of stragglers. Finally, the average of the errors $(100 \times 1000)$ was calculated (see \autoref{fig:AD1bound}). Notice that this average error lies close to the upper bound \eqref{eq:AD1AD2err_BIBD}. Also, the upper bound is significantly lower than the error for the baseline construction.
%\sifat{Use the word construction in place of scheme in the plots.}

% \begin{figure}[t!]
% %\vspace{.4 cm}
%     \centering
%     \includegraphics[width=0.45\textwidth]{Journal/Figure_5_v2.pdf}
%     \caption{{\small Performance of the construction as in \eqref{eq:effAD1AD2} for $\epsilon  = .1$, and $m = 2$ with a $(27, 2, 5) $ coset bipartite graph.}}
%     \label{fig:bi-rand}
% \end{figure}

\begin{figure*}[t!]
    \centering

    \begin{subfigure}[t]{0.45\textwidth}
        \centering
        \includegraphics[
            width=\linewidth,
            ,
            clip
        ]{Journal/Fig_4_a.pdf}
        \caption{Empirical average error and corresponding bounds.}
        \label{fig:rdm_cos_bound}
        
    \end{subfigure}
    \hfill
    \begin{subfigure}[t]{0.45\textwidth}
        \centering
        \includegraphics[
            width=\linewidth,
          ,
            clip
        ]{Journal/Fig_4_b.pdf}
        \caption{Comparison of empirical average error and relative squared gradient error.}
        \label{fig:err_gradd_err}
        
    \end{subfigure}

    \caption{{\small Performance of the construction as in \eqref{eq:effAD1AD2} for $\epsilon  = 0.1$, and $m = 2$ with a $(27, 2, 5) $ coset bipartite graph.}}
    \label{fig:coset_rdm}
\end{figure*}




\begin{figure*}[t!]
    \centering

    \begin{subfigure}[t]{0.45\textwidth}
        \centering
        \includegraphics[
            width=\linewidth,
            ,
            clip
        ]{Journal/Fig_5_a.pdf}
        \caption{ Performance of the construction as in  $\eqref{eq:recons}$ for a $(40,20,3,6)$ bi-regular bipartite graph.}
        \label{fig:DiagDom}
        
    \end{subfigure}
    \hfill
    \begin{subfigure}[t]{0.45\textwidth}
        \centering
        \includegraphics[
            width=\linewidth,
          ,
            clip
        ]{Journal/Fig_5_b.pdf}
        \caption{Performance of the construction as in  \eqref{eq:hada_coset_cons} for a $(27,2,5)$ coset bipartite graph.}
        \label{fig:hada_coset_cons}
        
    \end{subfigure}

    \caption{{\small Performance of the constructions in \eqref{eq:recons} and \eqref{eq:hada_coset_cons}.}}
    \label{fig:coset_cons}
   
\end{figure*}






% \begin{figure}[t!]
% %\vspace{.4 cm}
%     \centering
%     \includegraphics[width=0.45\textwidth]{Journal/Figure_4_v2.pdf}
%     \caption{{\small Performance of the construction as in  $\eqref{eq:recons}$ for a $(40,20,3,6)$ bi-regular bipartite graph and $m = 2$. }}
%     \label{fig:DiagDom}
    
% \end{figure}






% \begin{figure}[t!]
% %\vspace{.4 cm}
%     \centering
%     \includegraphics[width=0.45\textwidth]{Journal/Training_loss_2.pdf}
%     \caption{{\small Training loss for a small neural network with a $(20, 2, 3)$ coset bipartite graph and $m = 2$.}}
%     \label{fig:loss_coset}
% \end{figure}

% 


Next, we chose a $(27, 2, 5)$ coset bipartite graph to  evaluate the construction as in \eqref{eq:effAD1AD2} with $\epsilon = 0.1$. In this case, $ m = 2$, $n = 54$, $k = 27$ and $\delta = 5$. The empirical average was calculated the same way as in \autoref{fig:bibd_31_exact}. Note that in this case the bound in \eqref{eq:AD1AD2err} depends on the specific sets of stragglers and the bound in \eqref{eq:AD1AD2err_Coset} depends on the number of stragglers only. Although the bound appears loose for zero stragglers, the approximation error is zero because the encoding matrix in this case is invertible (see \autoref{fig:rdm_cos_bound}). To examine how the approximation error translates to the actual
gradient computation, we also evaluate the relative squared gradient error,
defined as
$
\frac{\left\|\mathbf{Z}\mathbf{B}_{\mathcal{F}}\mathbf{R}
-\mathbf{Z}\mathbf{F}\right\|_F^2}
{\left\|\mathbf{Z}\mathbf{F}\right\|_F^2},
$
where $\mathbf{Z}\mathbf{F}$ is the true gradient and
$\mathbf{Z}\mathbf{B}_{\mathcal{F}}\mathbf{R}$ is the decoded gradient. 
This evaluation was performed using the MNIST dataset and a fully connected neural network with input dimension 784, one hidden layer with 128 neurons and ReLU activation, and an output layer with 10 neurons and softmax activation. We calculate the empirical average of this
relative squared gradient error over the sampled realizations. \autoref{fig:err_gradd_err} compares this empirical average relative squared gradient error with the empirical average approximation error. The two
quantities increase with the straggler fraction and exhibit a similar
overall trend.

We also evaluate the construction in \eqref{eq:hada_coset_cons} using the $(27,2,5)$
coset bipartite graph with $m=2$. For each number of stragglers, the
approximation error is averaged over 1000 random choices of the straggler set and
compared with the upper bound in Lemma \ref{lemma:hada_coset_bound} and the lower bound.
As shown in \autoref{fig:hada_coset_cons}, the empirical average error remains well below the
upper bound.


Next, we consider the construction in Section \ref{sec:construction_RHPNSC} (see Algorithm \ref{alg:example}). We chose the bi-adjacency matrix of a $(40,20,3,6)$ bi-regular bipartite graph as the assignment matrix and $m = 2$.
% We constructed the encoding matrix in two ways based on the choice of $\mathbf{v}_1$ and $\mathbf{v}_2$. Firstly, 
In Algorithm 1, we picked $\mathbf{v}_1 = \mathbf{1}_n$ and $\mathbf{v}_2$ such that the entries of $\mathbf{v}_2$ are $\pm 1$. Note that in this particular case, we were successful in finding a $\bfv_2$ that satisfied the null-space constraints. However, this is not always guaranteed. If no restrictions are placed on $\bfv_2$ other than the null-space constraints, then the existence of a nonzero null-space vector is guaranteed. %\aditya{What about the Gaussian simulations?}
%\sifat{Adding}
For a given number of stragglers, the upper bound \eqref{eq:bounddiagdom} was calculated for each of 1000 random choices of stragglers and the maximum of them is shown (see \autoref{fig:DiagDom}). The average approximation error $\eqref{eq:ApproxError}$ was calculated for the same 1000 choices of stragglers, for this construction and the baseline construction. The upper bound in this case is significantly lower than the baseline scheme error in this case as well.
Next, we picked $\mathbf{v}_1$ from a Gaussian $\mathcal{N}(\mathbf{0}_n,\bfI_n)$ distribution and $\bfv_2$ according to the null-space constraints of Algorithm \ref{alg:example}. In this case, the simulation results show that the error is quite low. However, the upper bound based on \eqref{eq:bounddiagdom} is quite loose. For both cases, as expected, the error is zero when there are no stragglers. We note that in \autoref{fig:coset_rdm} and \autoref{fig:coset_cons}, the average errors are lower than the lower bound (obtained from \eqref{eq:lowerbound}) at slightly higher straggler fractions. This is because the lower bound corresponds to a constructed worst-case non-straggler set of size $n-s$ that has the corresponding approximation error, whereas the average error curves are generated by random choice of the non-stragglers. 

% We additionally compare the random diagonal construction using a BIBD assignment matrix against the exact communication-efficient gradient coding scheme in~\cite{tayyebehM21}. 
% We consider a $(31,31,6,6,1)$ BIBD assignment matrix with $m=2$, for which the exact scheme is 
% designed to tolerate $s=\gamma-m=4$ stragglers. Since its polynomial encoder 
% is linear, we represent it in encoding matrix form and evaluate its error using 
% the same optimal least-squares decoder. As shown in Fig.~\ref{fig:bibd_31_exact}, 
% it achieves zero error within its designed tolerance, but its error increases 
% sharply beyond this threshold, whereas the proposed construction exhibits a 
% more gradual increase.






% We note that in \autoref{fig:DiagDom, fig:bi-rand}, the average errors for the $\pm 1$ and the Gaussian constructions are lower than the lower bound (obtained from \eqref{eq:lowerbound}) at slightly higher straggler fractions. This is because the lower bound corresponds to a constructed worst case non-straggler set of size $n-s$ that has the corresponding approximation error, whereas the average error curves are generated by random choice of the non-stragglers.

% \aditya{Need to provide a clear reason why the lower bound is higher than experimental curve at some places.}
Finally, we compare the convergence of our scheme with the baseline construction. For training, we considered the MNIST dataset, and the fully connected neural network discussed above.
%
%For training, we use the MNIST \cite{6296535} dataset and a fully connected neural network with input dimension 784, one hidden layer with 128 neurons and ReLU activation, and an output layer with 10 neurons with softmax activation. 
For comparison, we chose a  $(7, 7, 3,3, 1)$ BIBD and a $(27, 2, 5)$ coset bipartite graph with $m = 2$ for both cases. We set the straggling probability to $q = .25$. The encoding matrix follows construction  \eqref{eq:effAD1AD2}, with $\epsilon  = 0$ for the BIBD case and $\epsilon  = .1$ for the coset bipartite graph case, compared against the corresponding baseline construction, i.e., $\bfB^T = \begin{bmatrix} \bfA^T & \bfA^T\end{bmatrix}$. The results are averaged over 20 random initializations. It can be observed that in both cases, our construction converges faster (see  \autoref{fig:loss_coset} and \autoref{fig:loss_bibd}).

% We compare the convergence for the construction as in \eqref{eq:effAD1AD2} with the corresponding baseline construction, i.e., $\bfB^T = \begin{bmatrix} \bfA^T & \bfA^T\end{bmatrix}$. 
%In the case when $\mathbf{v}_1$ and $\mathbf{v}_2$ are picked from Gaussian distribution, $\mathbf{v}_2$ satisfying the null-space constraint, the upper bound is loose, while the average error is relatively small.
%Furthermore, as expected, the error is zero when there are no stragglers. 


\begin{figure}[t!]
%\vspace{.4 cm}
    \centering
     \begin{subfigure}[t]{0.48\linewidth}
        \centering
        \includegraphics[width= .9\linewidth]{Journal/Fig_6_a.pdf}
    \caption{{\small $(27, 2, 5)$ coset bipartite graph}}
    \label{fig:loss_coset}
    
 \end{subfigure}
    \hfill
    \begin{subfigure}[t]{0.48\linewidth}
        \centering
        \includegraphics[width= .9\linewidth]{Journal/Fig_6_b.pdf}
    \caption{{\small  $(7, 7, 3,3, 1)$ BIBD }}
    \label{fig:loss_bibd}
    \end{subfigure}
    \caption{Training loss for a small neural network with different constructions, $m = 2$.}
    \label{fig:training_loss_comparison}



    
\end{figure}










\section{Conclusions and Future Work}
In this work, we propose approximate gradient coding schemes that are communication-efficient. Our constructions are based on post multiplication by diagonal matrices and Hadamard products with null-space constraints, applied to structured assignment matrices from BIBDs, strongly regular graphs, and coset bipartite graphs. Moreover, we prove convergence for some of our constructions with specific assignment matrices. We validate our constructions through numerical experiments, which show improved approximation error and faster convergence relative to the baseline constructions. 

There is ample scope for future work. We expect to develop constructions with tighter approximation error bounds and lower empirical error across a wide range of straggler sets. Additionally, schemes that incorporate partial computations from straggling workers (rather than ignoring the partial results entirely) are of particular interest. Extending our schemes to heterogeneous straggler settings while maintaining communication efficiency is also an important open problem.


% \aditya{Need to put in.}


%\aditya{Sometimes I remove conclusion from conference paper}
% \section{Conclusion}
% We proposed two constructions of communication-efficient approximate gradient coding schemes. Also, the error for these schemes can be upper bounded in terms the number of stragglers and communication reduction factor. While the previous works only consider the effect of stragglers, our schemes allow a trade-off between the error and the communication reduction factor. We also provided a universal lower bound for any communication efficient scheme in terms of $m$ and $s$. 

%%%%%%
%% Appendix:
%% If needed a single appendix is created by
%%
%\appendix
%%
%% If several appendices are needed, then the command
%%
% \appendices
%%
%% in combination with further \section commands can be used.
%%%%%%







% \section*{Acknowledgment}



%%%%%%
%% To balance the columns at the last page of the paper use this
%% command:
%%
%\enlargethispage{-1.2cm} 
%%
%% If the balancing should occur in the middle of the references, use
%% the following trigger:
%%
%\IEEEtriggeratref{7}
%%
%% which triggers a \newpage (i.e., new column) just before the given
%% reference number. Note that you need to adapt this if you modify
%% the paper.  The "triggered" command can be changed if desired:
%%
%\IEEEtriggercmd{\enlargethispage{-20cm}}
%%
%%%%%%


%%%%%%
%% References:
%% We recommend the usage of BibTeX:
%%
%\bibliographystyle{IEEEtran}
%\bibliography{definitions,bibliofile}
%%
%% where we here have assumed the existence of the files
%% definitions.bib and bibliofile.bib.
%% BibTeX documentation can be obtained at:
%% http://www.ctan.org/tex-archive/biblio/bibtex/contrib/doc/
%%%%%%



%% Or you use manual references (pay attention to consistency and the
%% formatting style!):
% \begin{thebibliography}{9}

% \bibitem{Laport:LaTeX}
% L.~Lamport,
%   \emph{\LaTeX: A Document Preparation System,} 
%   Addison-Wesley, Reading, Massachusetts, USA, 2nd~ed., 1994. 

% \bibitem{GMS:LaTeXComp}
% F.~Mittelbach, M,~Goossens, J.~Braams, D.~Carlisle, and
% C.~Rowley, \emph{The {\LaTeX} Companion,} Addison-Wesley,
% Reading, Massachusetts, USA, 2nd~ed., 2004.

% \bibitem{oetiker_latex}
% T.~Oetiker, H.~Partl, I.~Hyna, and E.~Schlegl, \emph{The Not So Short
%   Introduction to {\LaTeX2e}}, version 5.06, Jun.~20, 2016. [Online].
%   Available: \url{https://tobi.oetiker.ch/lshort/}

% \bibitem{typesetmoser}
% S.~M. Moser, \emph{How to Typeset Equations in {\LaTeX}}, version 4.6,
%   Sep. 29, 2017. [Online]. Available:
%   \url{http://moser-isi.ethz.ch/manuals.html#eqlatex}

% \bibitem{IEEE:pdfsettings}
% IEEE, \emph{Preparing Conference Content for the IEEE Xplore Digital
%   Library.} [Online]. Available:
%   \url{http://www.ieee.org/conferences_events/conferences/organizers/pubs/preparing_content.html}

% \bibitem{IEEE:AuthorToolbox}
% IEEE, \emph{Author Digital Toolbox.} [Online.] Available:
%   \url{http://www.ieee.org/publications_standards/publications/authors/authors_journals.html}

% \end{thebibliography}
\bibliographystyle{IEEEtran}
\bibliography{BG, merged_clean}
% \bibliography{BG, citations, tipnewpfmt_kfcsfullpap, citations_v2, tipnewpfmt_kfcsfullpap_v2}


%\clearpage

% \appendix




\appendices
\section{Expected Approximation Error Under Random Straggling}
\label{sec:exp_app_err_rand_strr}
We first average the fixed non-straggler set error bounds derived in Corollaries \ref{cor:BIBD}, \ref{cor:SRG}, and \ref{cor:bipartite} over the random set of non-straggling workers. The following lemma provides a common bound that is subsequently applied to the BIBD, SRG, and coset bipartite graph constructions.





\begin{lemma}
\label{lemma:expected_error_random_straggling}
Suppose that each worker straggles independently with probability $q$, and let $\mathcal{F}\subseteq[n]$ be the random set of non-straggling workers, and let $f = |\mathcal{F}|$. Suppose that there exist constants $C_0>0$, $a_0>0$, and $b_0\geq 0$ such that, for every nonempty $\cF$,
\[
\mathbb{E}_{\mathbf{D}}
\left[
\text{Err}_{\mathcal{F}}(\mathbf{B})
\right]
\leq
mk-\frac{C_0f}{a_0+b_0(f-1)}.
\]
Then,
\[
\mathbb{E}_{\mathcal{F}}
\left[
\mathbb{E}_{\mathbf{D}}
\left[
\text{Err}_{\mathcal{F}}(\mathbf{B})
\right]
\right]
\leq
mk-
\frac{C_0n(1-q)}
{a_0+b_0(n-1)}.
\]
\end{lemma}

\begin{proof}

When $\mathcal{F}=\emptyset$, no worker returns a message. So,
\[
\text{Err}_{\emptyset}(\mathbf{B})
=
\|\mathbf{F}\|_F^2
=
mk.
\]
Therefore,

\begin{align*}
&\mathbb{E}_{\mathcal{F}}
\left[
\mathbb{E}_{\mathbf{D}}
\left[
\operatorname{Err}_{\mathcal{F}}(\mathbf{B})
\right]
\right]
\\
&=
\sum_{{F}\subseteq[n]} \mathbb{P}(\mathcal{F} = F)
\mathbb{E}_{\mathbf{D}}
\left[
\operatorname{Err}_{{F}}(\mathbf{B})
\right]
\\
&=
q^n mk
+
\sum_{f=1}^{n}
\sum_{\substack{{F}\subseteq[n]\\|{F}|=f}} (1-q)^f q^{n-f}
\mathbb{E}_{\mathbf{D}}
\left[
\operatorname{Err}_{{F}}(\mathbf{B})
\right]
\\
% &\leq
% q^n mk
% +
% \sum_{f=1}^{n}
% \sum_{\substack{\mathcal{F}\subseteq[n]\\|\mathcal{F}|=f}} (1-q)^f q^{n-f}
% \left(
% mk-\frac{Cf}{a+b(f-1)}
% \right)
% \\
&\le
q^n mk
+
\sum_{f=1}^{n}
\binom{n}{f} (1-q)^f q^{n-f}
\left(
mk-\frac{C_0f}{a_0+b_0(f-1)}
\right)
.
\end{align*}




For every $f\in\{1,\ldots,n\}$, since $b_0\geq0$,
\[
a_0+b_0(f-1)
\leq
a_0+b_0(n-1).
\]
Consequently,
\[
\frac{f}{a_0+b_0(f-1)}
\geq
\frac{f}{a_0+b_0(n-1)}.
\]
It follows that
\[
\begin{aligned}
&\mathbb{E}_{\mathcal{F}}
\left[
\mathbb{E}_{\mathbf{D}}
\left[
\text{Err}_{\mathcal{F}}(\mathbf{B})
\right]
\right]
\\
&\leq
q^nmk
+
\sum_{f=1}^{n}
\binom{n}{f} (1-q)^f q^{n-f}
\left(
mk-\frac{C_0f}{a_0+b_0(n-1)}
\right)
\\
&=
mk-
\frac{C_0}{a_0+b_0(n-1)}
\sum_{f=1}^{n}f\binom{n}{f} (1-q)^f q^{n-f}
\\
&=
mk-
\frac{C_0n(1-q) ((1-q)+q)^{n-1}}
{a_0+b_0(n-1)}.
\end{aligned}
\]

So,
\[
\mathbb{E}_{\mathcal{F}}
\left[
\mathbb{E}_{\mathbf{D}}
\left[
\text{Err}_{\mathcal{F}}(\mathbf{B})
\right]
\right]
\leq
mk-
\frac{C_0n(1-q)}
{a_0+b_0(n-1)}.
\]
\end{proof}

We now apply Lemma~\ref{lemma:expected_error_random_straggling} to the BIBD, SRG, and coset bipartite graph constructions.

\begin{cor}
\label{cor:expected_error_BIBD_strr}
Suppose that the assignment matrix $\mathbf{A}$ is the transpose of the incidence matrix of a $(n,k,\gamma,\delta,\lambda)$ BIBD. Then, for the construction as in \eqref{eq:effAD1AD2} with $\epsilon=0$,
\[
\mathbb{E}_{\mathcal{F}}
\left[
\mathbb{E}_{\mathbf{D}}
\left[
\text{Err}_{\mathcal{F}}(\mathbf{B})
\right]
\right]
\leq
mk-
\frac{m\delta^2n(1-q)}
{m\delta+\lambda(n-1)}.
\]
\end{cor}

\begin{proof}
For a fixed non-empty set $\mathcal{F}$ with $|\mathcal{F}|=f$, Corollary~\ref{cor:BIBD} gives
\[
\mathbb{E}_{\mathbf{D}}
\left[
\text{Err}_{\mathcal{F}}(\mathbf{B})
\right]
\leq
mk-
\frac{m\delta^2f}
{m\delta+\lambda(f-1)}.
\]
The result follows from Lemma~\ref{lemma:expected_error_random_straggling} by setting
\[
C_0=m\delta^2,
\qquad
a_0=m\delta,
\qquad
b_0=\lambda.
\]
For a BIBD, $\delta>0$ and $\lambda\geq0$. Hence, $a_0>0$, $b_0\geq0$ and $C_0 > 0$.
\end{proof}

\begin{cor}
\label{cor:expected_error_SRG_strr}
Suppose that the assignment matrix $\mathbf{A}$ is the adjacency matrix of a $(n,\delta,\lambda,\mu)$ SRG with $0<\delta\neq\mu$. Then, for the construction as in \eqref{eq:effAD1AD2} with $\epsilon=0$,
\[
\mathbb{E}_{\mathcal{F}}
\left[
\mathbb{E}_{\mathbf{D}}
\left[
\text{Err}_{\mathcal{F}}(\mathbf{B})
\right]
\right]
\leq
mk-
\frac{m\delta^2n(1-q)}
{m\delta+(\lambda-\mu)\theta+\mu(n-1)},
\]
where $\theta$ is defined as in Corollary~\ref{cor:SRG}.
\end{cor}

\begin{proof}
For a fixed non-empty set $\mathcal{F}$ with $|\mathcal{F}|=f$, Corollary~\ref{cor:SRG} gives
\[
\mathbb{E}_{\mathbf{D}}
\left[
\text{Err}_{\mathcal{F}}(\mathbf{B})
\right]
\leq
mk-
\frac{m\delta^2f}
{m\delta+(\lambda-\mu)\theta+\mu(f-1)}.
\]
The result follows from Lemma~\ref{lemma:expected_error_random_straggling} by setting
\[
C_0=m\delta^2,
\qquad
a_0=m\delta+(\lambda-\mu)\theta,
\qquad
b_0=\mu.
\]

By definition, $\mu\geq0$. We next verify that $a_0>0$. If $\lambda\geq\mu$, then $\theta=\delta$, and hence
\[
a_0
=
m\delta+(\lambda-\mu)\delta
>0.
\]
If $\lambda<\mu$, then
\[
\theta
=
\frac{1}{2}
\left[
(\lambda-\mu)
-
\sqrt{(\lambda-\mu)^2+4(\delta-\mu)}
\right]
<0.
\]
Since $\lambda-\mu<0$, we have
\[
(\lambda-\mu)\theta>0,
\]
and therefore
\[
a_0=m\delta+(\lambda-\mu)\theta>0.
\]
Thus, the conditions of Lemma~\ref{lemma:expected_error_random_straggling} are satisfied.
\end{proof}

\begin{cor}
\label{cor:expected_error_coset_strr}
Suppose that the assignment matrix $\mathbf{A}$ is the bi-adjacency matrix of a $(k,m,\delta)$ coset bipartite graph such that $k=p^a$, where $p$ is a prime, $a\in\mathbb{Z}_{\geq1}$, and $p\nmid\delta$. Then, for the construction as in \eqref{eq:effAD1AD2} with $\epsilon > 0$,
\[
\mathbb{E}_{\mathcal{F}}
\left[
\mathbb{E}_{\mathbf{D}}
\left[
\text{Err}_{\mathcal{F}}(\mathbf{B})
\right]
\right]
\leq
mk-
\frac{m\delta^2n(1-q)}
{m\delta^2+c(m-1)\delta},
\]
where $n=mk$.
\end{cor}

\begin{proof}
For a fixed non-empty set $\mathcal{F}$ with $|\mathcal{F}|=f$, Corollary~\ref{cor:bipartite} gives
\[
\mathbb{E}_{\mathbf{D}}
\left[
\text{Err}_{\mathcal{F}}(\mathbf{B})
\right]
\leq
mk-
\frac{m\delta^2f}
{m\delta^2+c(m-1)\delta}.
\]
The result follows from Lemma~\ref{lemma:expected_error_random_straggling} by setting
\[
C_0=m\delta^2,
\qquad
a_0=m\delta^2+c(m-1)\delta,
\qquad
b_0=0.
\]
Since $m,\delta,c>0$, we have $a_0>0$. Thus, the conditions of Lemma~\ref{lemma:expected_error_random_straggling} are satisfied.
\end{proof}

\begin{lemma}
\label{lemma:hada_coset_bound}
Suppose that the assignment matrix $\bfA$ is the bi-adjacency matrix of a
$(k,m,\delta)$ coset bipartite graph such that $k=p^a$, where $p$ is a prime,
$a\in\mathbb{Z}_{\geq 1}$, and $p\nmid\delta$. Then, for the construction as
in \eqref{eq:hada_coset_cons}, for any non-straggler set $\cF$ with
$|\cF|=n-s$,
\begin{equation}
\label{eq:Hada_Coset_err}
\text{Err}_{\cF}(\bfB)\leq s.
\end{equation}
Consequently,
\begin{equation}
\label{eq:Hada_Coset_exp_err}
\mathbb{E}_{\cF}\left[\text{Err}_{\cF}(\bfB)\right]\leq nq.
\end{equation}
\end{lemma}

\begin{proof}
First consider a non-empty set $\cF$. For $i\in[m]$, observe that
$\mathbf{f}_i^T\mathbf{f}_i=k$. Since $\bfB$ is invertible,
$\bfB_{\cF}^T\bfB_{\cF}$ is invertible.  Let
$\bfP_{\cF}:=\bfI_n(:,\cF)\in\mathbb{R}^{n\times(n-s)}$
be the selection matrix that selects the columns corresponding to $\cF$.
Thus, $\bfB_{\cF}=\bfB\bfP_{\cF}$. We have
\begin{align*}
&\mathbf{f}_i^T\mathbf{f}_i
-\mathbf{f}_i^T\bfB\bfP_{\cF}
(\bfB_{\cF}^T\bfB_{\cF})^{-1}
\bfP_{\cF}^T\bfB^T\mathbf{f}_i\\
&=k-\|\bfP_{\cF}^T\bfB^T\mathbf{f}_i\|_2^2
\frac{\mathbf{f}_i^T\bfB\bfP_{\cF}}
{\|\bfP_{\cF}^T\bfB^T\mathbf{f}_i\|_2}
(\bfB_{\cF}^T\bfB_{\cF})^{-1}
\frac{\bfP_{\cF}^T\bfB^T\mathbf{f}_i}
{\|\bfP_{\cF}^T\bfB^T\mathbf{f}_i\|_2}\\
&\leq
k-\|\bfP_{\cF}^T\bfB^T\mathbf{f}_i\|_2^2
\lambda_{\min}\left(
(\bfB_{\cF}^T\bfB_{\cF})^{-1}
\right).
\end{align*}

Recall that $\mathbf{f}_i=\mathbf{e}_i\otimes\mathbf{1}_k$. Since
$\bfB=\bfH^T\otimes\bfC$,
\begin{align*}
\bfB^T\mathbf{f}_i
&=(\bfH\otimes\bfC^T)
(\mathbf{e}_i\otimes\mathbf{1}_k)\\
&=(\bfH\mathbf{e}_i)\otimes
(\bfC^T\mathbf{1}_k)\\
&=\delta(\bfH\mathbf{e}_i)\otimes\mathbf{1}_k.
\end{align*}
Since every entry of $\bfH\mathbf{e}_i$ has magnitude one,
every entry of $\bfB^T\mathbf{f}_i$ has magnitude $\delta$. Therefore,
\begin{equation}
\label{eq:Hada_Coset_f_norm}
\|\bfP_{\cF}^T\bfB^T\mathbf{f}_i\|_2^2
=\delta^2(n-s).
\end{equation}

Moreover,
$
\bfB^T\bfB
=
\bfH\bfH^T\otimes\bfC^T\bfC
=
m\bfI_m\otimes\bfC^T\bfC.
$ Since $\bfC$ is $\delta$-regular,
$
\lambda_{\max}(\bfC^T\bfC)=\delta^2,
$
and hence
$
\lambda_{\max}(\bfB^T\bfB)=m\delta^2.
$
Also,
$
\bfB_{\cF}^T\bfB_{\cF}
=
\bfP_{\cF}^T\bfB^T\bfB\bfP_{\cF}.
$
Thus, by Cauchy's interlacing theorem,
$
\lambda_{\max}(\bfB_{\cF}^T\bfB_{\cF})
\leq m\delta^2
$. Consequently,
\[
\lambda_{\min}\left(
(\bfB_{\cF}^T\bfB_{\cF})^{-1}
\right)
=
\frac{1}{
\lambda_{\max}(\bfB_{\cF}^T\bfB_{\cF})
}
\geq \frac{1}{m\delta^2}.
\]

Using \eqref{eq:Hada_Coset_f_norm}, for each $i\in[m]$,
\begin{align*}
&\mathbf{f}_i^T\mathbf{f}_i
-\mathbf{f}_i^T\bfB\bfP_{\cF}
(\bfB_{\cF}^T\bfB_{\cF})^{-1}
\bfP_{\cF}^T\bfB^T\mathbf{f}_i\\
&\leq
k-\delta^2(n-s)\frac{1}{m\delta^2}\\
&=
k-\frac{n-s}{m}.
\end{align*}
Summing over $i\in[m]$ gives
\[
\text{Err}_{\cF}(\bfB)
\leq
mk-(n-s).
\]
Since $n=mk$ for a $(k,m,\delta)$ coset bipartite graph,
$
\text{Err}_{\cF}(\bfB)\leq s.
$ For $\cF=\emptyset$, we have
$\text{Err}_{\emptyset}(\bfB)=\|\bfF\|_F^2=mk=n$. So \eqref{eq:Hada_Coset_err} also holds since in this case $s=n$. Since each worker straggles independently with probability $q$, 
$
\mathbb{E}_{\cF}\left[\text{Err}_{\cF}(\bfB)\right]
\leq
\mathbb{E}_{\cF}[s]
=
nq,
$ which proves \eqref{eq:Hada_Coset_exp_err}.
\end{proof}



% \begin{lemma}
% Suppose that the assignment matrix $\bfA$ is the bi-adjacency matrix of a
% $(k,m,\delta)$ coset bipartite graph such that $k=p^a$, where $p$ is a
% prime, $a\in\mathbb{Z}_{\geq 1}$, and $p\nmid\delta$. Then, for the
% construction as in \eqref{eq:Hada_Coset}, for any non-empty non-straggler
% set $\cF$ of size $n-s$,
% \begin{equation}
% \label{eq:Hada_Coset_err}
% \text{Err}_{\cF}(\bfB) \leq s.
% \end{equation}
% \end{lemma}

% \begin{proof}
% For $i\in[m]$, observe that $\mathbf{f}_i^T\mathbf{f}_i=k$. Since
% $\bfB$ is invertible, $\bfB_{\cF}^T\bfB_{\cF}$ is invertible for every
% non-empty $\cF$. We have
% \begin{align*}
% &\mathbf{f}_i^T\mathbf{f}_i
% -\mathbf{f}_i^T\bfB\bfP_{\cF}
% (\bfB_{\cF}^T\bfB_{\cF})^{-1}
% \bfP_{\cF}^T\bfB^T\mathbf{f}_i\\
% &=k-\|\bfP_{\cF}^T\bfB^T\mathbf{f}_i\|_2^2
% \frac{\mathbf{f}_i^T\bfB\bfP_{\cF}}
% {\|\bfP_{\cF}^T\bfB^T\mathbf{f}_i\|_2}
% (\bfB_{\cF}^T\bfB_{\cF})^{-1}
% \frac{\bfP_{\cF}^T\bfB^T\mathbf{f}_i}
% {\|\bfP_{\cF}^T\bfB^T\mathbf{f}_i\|_2}\\
% &\leq k-\|\bfP_{\cF}^T\bfB^T\mathbf{f}_i\|_2^2
% \lambda_{\min}\left(
% (\bfB_{\cF}^T\bfB_{\cF})^{-1}
% \right).
% \end{align*}

% Recall that $\mathbf{f}_i=\mathbf{e}_i\otimes\mathbf{1}_k$. Since
% $\bfB=\bfH^T\otimes\bfC$,
% \begin{align*}
% \bfB^T\mathbf{f}_i
% &=(\bfH\otimes\bfC^T)
% (\mathbf{e}_i\otimes\mathbf{1}_k)\\
% &=(\bfH\mathbf{e}_i)\otimes
% (\bfC^T\mathbf{1}_k)\\
% &=\delta(\bfH\mathbf{e}_i)\otimes\mathbf{1}_k.
% \end{align*}
% Since every entry of $\bfH\mathbf{e}_i$ has magnitude one, every entry
% of $\bfB^T\mathbf{f}_i$ has magnitude $\delta$. Therefore,
% \begin{equation}
% \label{eq:Hada_Coset_f_norm}
% \|\bfP_{\cF}^T\bfB^T\mathbf{f}_i\|_2^2
% =\delta^2(n-s).
% \end{equation}

% Moreover,
% \[
% \bfB^T\bfB
% =
% \bfH\bfH^T\otimes\bfC^T\bfC
% =
% m\bfI_m\otimes\bfC^T\bfC.
% \]
% Since $\bfC$ is $\delta$-regular,
% $\lambda_{\max}(\bfC^T\bfC)=\delta^2$. Hence,
% \[
% \lambda_{\max}(\bfB^T\bfB)=m\delta^2.
% \]
% Also,
% \[
% \bfB_{\cF}^T\bfB_{\cF}
% =
% \bfP_{\cF}^T\bfB^T\bfB\bfP_{\cF},
% \]
% and therefore, by Cauchy's interlacing theorem,
% \[
% \lambda_{\max}(\bfB_{\cF}^T\bfB_{\cF})
% \leq m\delta^2.
% \]
% Consequently,
% \[
% \lambda_{\min}\left(
% (\bfB_{\cF}^T\bfB_{\cF})^{-1}
% \right)
% =
% \frac{1}{
% \lambda_{\max}(\bfB_{\cF}^T\bfB_{\cF})
% }
% \geq \frac{1}{m\delta^2}.
% \]

% Using \eqref{eq:Hada_Coset_f_norm}, for each $i\in[m]$,
% \[
% \mathbf{f}_i^T\mathbf{f}_i
% -\mathbf{f}_i^T\bfB\bfP_{\cF}
% (\bfB_{\cF}^T\bfB_{\cF})^{-1}
% \bfP_{\cF}^T\bfB^T\mathbf{f}_i
% \leq
% k-\frac{n-s}{m}.
% \]
% Summing over $i\in[m]$ gives
% \[
% \text{Err}_{\cF}(\bfB)
% \leq
% mk-(n-s).
% \]
% Since $n=mk$ for a $(k,m,\delta)$ coset bipartite graph,
% \[
% \text{Err}_{\cF}(\bfB)\leq s.
% \]
% \end{proof}








\section{Condition Number of the Decoding Step}
\label{sec:appendix_cond_number}

In this appendix, we analyze the numerical stability of the least-squares decoding step for the construction in $\eqref{eq:effAD1AD2}$ with BIBD and SRG assignment matrices. For a fixed non-straggler set $\mathcal{F}$, we derive upper bounds on the condition number $\kappa(\mathbf{B}_{\mathcal{F}}^T\mathbf{B}_{\mathcal{F}})$. The following claim first bounds this condition number in terms of $\boldsymbol{\Delta}_{\mathcal{F}}=\mathbf{A}_{\mathcal{F}}^{T}\mathbf{A}_{\mathcal{F}}$.




% In this appendix, we analyze the numerical stability of the least-squares problem solved by the parameter server for a fixed non-straggler set $\mathcal{F}$. For the random diagonal matrix based construction with BIBD and SRG assignment matrices, we derive upper bounds on the condition number of the matrix $\mathbf{B}_{\mathcal{F}}^{T}\mathbf{B}_{\mathcal{F}}$. The following claim first bounds this condition number in terms of $\boldsymbol{\Delta}_{\mathcal{F}}=\mathbf{A}_{\mathcal{F}}^{T}\mathbf{A}_{\mathcal{F}}$.

\begin{claim}
\label{claim:cond_num}
Suppose that $\mathbf{\Delta}_{\mathcal{F}}$ is invertible. Then, for the construction in~(7),
$$
\kappa\!\left(\mathbf{B}_{\mathcal{F}}^{T}\mathbf{B}_{\mathcal{F}}\right)
\leq
\kappa\!\left(\mathbf{\Delta}_{\mathcal{F}}\right)
\frac{\sum_{j=1}^{m}\left\|\tilde{\mathbf{D}}_{j}\right\|_{2}^{2}}
{\sum_{j=1}^{m}\sigma_{\min}^{2}\!\left(\tilde{\mathbf{D}}_{j}\right)}
\leq
\left(\frac{1+\epsilon}{1-\epsilon}\right)^{2}
\kappa\!\left(\mathbf{\Delta}_{\mathcal{F}}\right).
$$
\end{claim}

% \begin{claim}
% The condition number of the communication efficient scheme is bounded for any $\epsilon$ given $\mathbf{\Delta}_{\mathcal{F}}$ is invertible.
% \end{claim}

\begin{proof}
Recall that
\[
\mathbf{B}_{\mathcal F}^{T}\mathbf{B}_{\mathcal F}
=
\sum_{j=1}^{m}
\tilde{\mathbf{D}}_j
\mathbf{\Delta}_{\mathcal F}
\tilde{\mathbf{D}}_j,
\]
where each $\tilde{\mathbf{D}}_j$ is an invertible diagonal matrix. Since $\mathbf{\Delta}_{\mathcal{F}}=\mathbf{A}_{\mathcal{F}}^T\mathbf{A}_{\mathcal{F}}$ is positive semidefinite and invertible, it is positive definite. For any vector $\mathbf{y}$,
\[
\lambda_{\min}(\mathbf{\Delta}_{\mathcal F})\|\mathbf{y}\|_2^2
\le
\mathbf{y}^{T}\mathbf{\Delta}_{\mathcal F}\mathbf{y}
\le
\lambda_{\max}(\mathbf{\Delta}_{\mathcal F})\|\mathbf{y}\|_2^2.
\]
Let $\mathbf{y}= \tilde{\mathbf{D}}_j\mathbf{x}$. Then, for any unit-norm vector $\mathbf{x}$,
\[
\mathbf{x}^{T}
\tilde{\mathbf{D}}_j
\mathbf{\Delta}_{\mathcal F}
\tilde{\mathbf{D}}_j
\mathbf{x}
=
\mathbf{y}^{T}{\mathbf{\Delta}}_{\mathcal F}\mathbf{y}
\le
\lambda_{\max}({\mathbf{\Delta}}_{\mathcal F})
\|\mathbf{y}\|_2^2.
\]
Since
\[
\|\mathbf{y}\|_2
=
\|\tilde{\mathbf{D}}_j\mathbf{x}\|_2
\le
\|\tilde{\mathbf{D}}_j\|_2\|\mathbf{x}\|_2
=
\|\tilde{\mathbf{D}}_j\|_2,
\]
we get
\[
\mathbf{x}^{T}
\tilde{\mathbf{D}}_j
\mathbf{\Delta}_{\mathcal F}
\tilde{\mathbf{D}}_j
\mathbf{x}
\le
\lambda_{\max}(\mathbf{\Delta}_{\mathcal F})
\|\tilde{\mathbf{D}}_j\|_2^2.
\]
Taking the maximum over all unit-norm vectors $\mathbf{x}$ gives
\[
\lambda_{\max}
\!\left(
\tilde{\mathbf{D}}_j
\mathbf{\Delta}_{\mathcal F}
\tilde{\mathbf{D}}_j
\right)
\le
\|\tilde{\mathbf{D}}_j\|_2^2
\lambda_{\max}(\mathbf{\Delta}_{\mathcal F}).
\]

Similarly,
\[
\mathbf{x}^{T}
\tilde{\mathbf{D}}_j
\mathbf{\Delta}_{\mathcal F}
\tilde{\mathbf{D}}_j
\mathbf{x}
=
\mathbf{y}^{T}\mathbf{\Delta}_{\mathcal F}\mathbf{y}
\ge
\lambda_{\min}(\mathbf{\Delta}_{\mathcal F})
\|\mathbf{y}\|_2^2.
\]
Since
\[
\|\mathbf{y}\|_2
=
\|\tilde{\mathbf{D}}_j\mathbf{x}\|_2
\ge
\sigma_{\min}(\tilde{\mathbf{D}}_j)\|\mathbf{x}\|_2
=
\sigma_{\min}(\tilde{\mathbf{D}}_j),
\]
we get
\[
\mathbf{x}^{T}
\tilde{\mathbf{D}}_j
\mathbf{\Delta}_{\mathcal F}
\tilde{\mathbf{D}}_j
\mathbf{x}
\ge
\lambda_{\min}({\mathbf{\Delta}}_{\mathcal F})
\sigma_{\min}^2(\tilde{\mathbf{D}}_j).
\]
Taking the minimum over all unit vectors $\mathbf{x}$ gives
\[
\lambda_{\min}
\!\left(
\tilde{\mathbf{D}}_j
\mathbf{\Delta}_{\mathcal F}
\tilde{\mathbf{D}}_j
\right)
\ge
\sigma_{\min}^2(\tilde{\mathbf{D}}_j)
\lambda_{\min}(\mathbf{\Delta}_{\mathcal F}).
\]

Now, by Weyl's inequality,
\[
\lambda_{\max}
\!\left(
\mathbf{B}_{\mathcal F}^{T}
\mathbf{B}_{\mathcal F}
\right)
\le
\sum_{j=1}^{m}
\lambda_{\max}
\!\left(
\tilde{\mathbf{D}}_j
\mathbf{\Delta}_{\mathcal F}
\tilde{\mathbf{D}}_j
\right),
\]
and therefore
\[
\lambda_{\max}
\!\left(
\mathbf{B}_{\mathcal F}^{T}
\mathbf{B}_{\mathcal F}
\right)
\le
\lambda_{\max}(\mathbf{\Delta}_{\mathcal F})
\sum_{j=1}^{m}\|\tilde{\mathbf{D}}_j\|_2^2.
\]
Similarly,
\[
\lambda_{\min}
\!\left(
\mathbf{B}_{\mathcal F}^{T}
\mathbf{B}_{\mathcal F}
\right)
\ge
\sum_{j=1}^{m}
\lambda_{\min}
\!\left(
\tilde{\mathbf{D}}_j
\mathbf{\Delta}_{\mathcal F}
\tilde{\mathbf{D}}_j
\right),
\]
and hence
\[
\lambda_{\min}
\!\left(
\mathbf{B}_{\mathcal F}^{T}
\mathbf{B}_{\mathcal F}
\right)
\ge
\lambda_{\min}(\mathbf{\Delta}_{\mathcal F})
\sum_{j=1}^{m}\sigma_{\min}^2(\tilde{\mathbf{D}}_j).
\]

Thus,
\[
\kappa
\!\left(
\mathbf{B}_{\mathcal F}^{T}
\mathbf{B}_{\mathcal F}
\right)
=
\frac{
\lambda_{\max}
\!\left(
\mathbf{B}_{\mathcal F}^{T}
\mathbf{B}_{\mathcal F}
\right)
}{
\lambda_{\min}
\!\left(
\mathbf{B}_{\mathcal F}^{T}
\mathbf{B}_{\mathcal F}
\right)
}
\le
\kappa(\mathbf{\Delta}_{\mathcal F})
\frac{
\sum_{j=1}^{m}\|\tilde{\mathbf{D}}_j\|_2^2
}{
\sum_{j=1}^{m}\sigma_{\min}^2(\tilde{\mathbf{D}}_j)
}.
\]

Since the absolute value of every diagonal entry of $\tilde{\mathbf{D}}_j$ is between $1-\epsilon$ and $1+\epsilon$, we have
$$
\|\tilde{\mathbf{D}}_j\|_2\leq 1+\epsilon
$$
and
$$
\sigma_{\min}(\tilde{\mathbf{D}}_j)\geq 1-\epsilon.
$$
Therefore,
$$
\sum_{j=1}^{m}\|\tilde{\mathbf{D}}_j\|_2^2
\leq
m(1+\epsilon)^2
$$
and
$$
\sum_{j=1}^{m}\sigma_{\min}^2(\tilde{\mathbf{D}}_j)
\geq
m(1-\epsilon)^2.
$$
Consequently,
$$
\begin{aligned}
\kappa\left(\mathbf{B}_{\mathcal{F}}^T
\mathbf{B}_{\mathcal{F}}\right)
&\leq
\kappa\left(\mathbf{\Delta}_{\mathcal{F}}\right)
\frac{
\sum_{j=1}^{m}\|\tilde{\mathbf{D}}_j\|_2^2
}{
\sum_{j=1}^{m}\sigma_{\min}^2(\tilde{\mathbf{D}}_j)
}\\
&\leq
\left(\frac{1+\epsilon}{1-\epsilon}\right)^2
\kappa\left(\mathbf{\Delta}_{\mathcal{F}}\right).
\end{aligned}
$$
This completes the proof.




% Now, by Weyl's inequality,
% \[
% \lambda_{\max}
% \!\left(
% \mathbf{B}_{\mathcal F}^{T}
% \mathbf{B}_{\mathcal F}
% \right)
% \le
% \sum_{j=1}^{m}
% \lambda_{\max}
% \!\left(
% \mathbf{D}_j
% \mathbf{\Delta}_{\mathcal F}
% \mathbf{D}_j
% \right),
% \]
% and therefore
% \[
% \lambda_{\max}
% \!\left(
% \mathbf{B}_{\mathcal F}^{T}
% \mathbf{B}_{\mathcal F}
% \right)
% \le
% \lambda_{\max}(\mathbf{\Delta}_{\mathcal F})
% \sum_{j=1}^{m}\|\mathbf{D}_j\|_2^2.
% \]
% Similarly,
% \[
% \lambda_{\min}
% \!\left(
% \mathbf{B}_{\mathcal F}^{T}
% \mathbf{B}_{\mathcal F}
% \right)
% \ge
% \sum_{j=1}^{m}
% \lambda_{\min}
% \!\left(
% \mathbf{D}_j
% \mathbf{\Delta}_{\mathcal F}
% \mathbf{D}_j
% \right),
% \]
% and hence
% \[
% \lambda_{\min}
% \!\left(
% \mathbf{B}_{\mathcal F}^{T}
% \mathbf{B}_{\mathcal F}
% \right)
% \ge
% \lambda_{\min}(\mathbf{\Delta}_{\mathcal F})
% \sum_{j=1}^{m}\sigma_{\min}^2(\mathbf{D}_j).
% \]

% Thus,
% \[
% \kappa
% \!\left(
% \mathbf{B}_{\mathcal F}^{T}
% \mathbf{B}_{\mathcal F}
% \right)
% =
% \frac{
% \lambda_{\max}
% \!\left(
% \mathbf{B}_{\mathcal F}^{T}
% \mathbf{B}_{\mathcal F}
% \right)
% }{
% \lambda_{\min}
% \!\left(
% \mathbf{B}_{\mathcal F}^{T}
% \mathbf{B}_{\mathcal F}
% \right)
% }
% \le
% \kappa(\mathbf{\Delta}_{\mathcal F})
% \frac{
% \sum_{j=1}^{m}\|\mathbf{D}_j\|_2^2
% }{
% \sum_{j=1}^{m}\sigma_{\min}^2(\mathbf{D}_j)
% }.
% \]


% Since the absolute value of every diagonal entry of $\mathbf{D}_j$ is between $1-\epsilon$ and $1+\epsilon$, we have
% $$
% \|\mathbf{D}_j\|_2\leq 1+\epsilon
% $$
% and
% $$
% \sigma_{\min}(\mathbf{D}_j)\geq 1-\epsilon.
% $$
% Therefore,
% $$
% \sum_{j=1}^{m}\|\mathbf{D}_j\|_2^2
% \leq
% m(1+\epsilon)^2
% $$
% and
% $$
% \sum_{j=1}^{m}\sigma_{\min}^2(\mathbf{D}_j)
% \geq
% m(1-\epsilon)^2.
% $$
% Consequently,
% $$
% \begin{aligned}
% \kappa\left(\mathbf{B}_{\mathcal{F}}^T
% \mathbf{B}_{\mathcal{F}}\right)
% &\leq
% \kappa\left(\mathbf{\Delta}_{\mathcal{F}}\right)
% \frac{
% \sum_{j=1}^{m}\|\mathbf{D}_j\|_2^2
% }{
% \sum_{j=1}^{m}\sigma_{\min}^2(\mathbf{D}_j)
% }
% \
% &\leq
% \left(\frac{1+\epsilon}{1-\epsilon}\right)^2
% \kappa\left(\mathbf{\Delta}_{\mathcal{F}}\right).
% \end{aligned}
% $$
% This completes the proof.
\end{proof}
\begin{cor}
Suppose that the assignment matrix $\mathbf{A}$ is the transpose of the incidence matrix of a $(n,k,\gamma,\delta,\lambda)$ BIBD. Then, for the construction as in (7), for every nonempty $\mathcal{F}\subseteq\{1,\ldots,n\}$ of size $n-s$,
$$
\kappa\left(\mathbf{B}_{\mathcal{F}}^T
\mathbf{B}_{\mathcal{F}}\right)
\leq
\left(\frac{1+\epsilon}{1-\epsilon}\right)^2
\frac{\delta+(n-s-1)\lambda}{\delta-\lambda}.
$$
% Consequently,
% $$
% \kappa\left(\mathbf{B}_{\mathcal{F}}^T
% \mathbf{B}_{\mathcal{F}}\right)
% \leq
% \left(\frac{1+\epsilon}{1-\epsilon}\right)^2
% \frac{\delta+(n-1)\lambda}{\delta-\lambda}.
% $$
\end{cor}

\begin{proof}
Recall that
$$
\mathbf{\Delta}_{\mathcal{F}}
= (\delta-\lambda)\mathbf{I}_{n-s}
+\lambda\mathbf{J}_{n-s}.$$
The largest eigenvalue of $\mathbf{\Delta}_{\mathcal{F}}$ is
$$
\lambda_{\max} \left(\mathbf{\Delta}_{\mathcal{F}}\right)
= \delta+(n-s-1)\lambda.
$$
Moreover,
$$
\lambda_{\min}\left(\mathbf{\Delta}_{\mathcal{F}}\right)
\geq
\delta-\lambda.
$$
Since $\delta>\lambda$ for a BIBD, $\mathbf{\Delta}_{\mathcal{F}}$ is positive definite. Therefore,
$$
\begin{aligned}
\kappa\left(\mathbf{\Delta}_{\mathcal{F}}\right)
&=
\frac{
\lambda_{\max}\left(\mathbf{\Delta}_{\mathcal{F}}\right)
}{
\lambda_{\min}\left(\mathbf{\Delta}_{\mathcal{F}}\right)
}
\
&\leq
\frac{\delta+(n-s-1)\lambda}{\delta-\lambda}.
\end{aligned}
$$

The result follows from Claim \ref{claim:cond_num}.
\end{proof}

\remark{For a $(91, 91, 10,10,1)$ BIBD with $\epsilon = 0$, we have $\kappa(\bfB_{\cF}^T\bfB_{\cF}) \le \frac{100}{9} \approx 11.11$.}





\begin{cor}
Suppose that the assignment matrix $\mathbf{A}$ is the adjacency matrix of an $(n,\delta,\lambda,\mu)$ SRG with $\delta\neq\mu$. Let
$$
\theta_{+}
= \frac{1}{2}
\left(
(\lambda-\mu)
+
\sqrt{(\lambda-\mu)^2+4(\delta-\mu)}
\right)
$$
and
$$
\theta_{-}
= \frac{1}{2}
\left(
(\lambda-\mu)
- \sqrt{(\lambda-\mu)^2+4(\delta-\mu)}
\right).
$$
Also, let
$$
\rho
= \min\left\{
|\theta_{+}|,
|\theta_{-}|
\right\}.
$$
Then, for the construction as in (7), for every nonempty $\mathcal{F}\subseteq\{1,\ldots,n\}$,
$$
\kappa\left(\mathbf{B}_{\mathcal{F}}^T
\mathbf{B}_{\mathcal{F}}\right)
\leq
\left(\frac{1+\epsilon}{1-\epsilon}\right)^2
\frac{\delta^2}{\rho^2}.
$$
\end{cor}

\begin{proof}
We have,
$$
\mathbf{\Delta}_{\mathcal{F}}
=\mathbf{A}_{\mathcal{F}}^T\mathbf{A}_{\mathcal{F}} = \mathbf{A}^2(\mathcal{F},\mathcal{F}).
$$
% Thus, $\mathbf{\Delta}_{\mathcal{F}}$ is a principal submatrix of $\mathbf{A}^2$.

The eigenvalues of $\mathbf{A}$ are $\delta$, $\theta_{+}$, and $\theta_{-}$. Therefore, the eigenvalues of $\mathbf{A}^2$ are $\delta^2$, $\theta_{+}^2$, and $\theta_{-}^2$. Moreover,
$$
\theta_{+}\theta_{-}
= \mu-\delta.
$$
Since $\delta\neq\mu$, both $\theta_{+}$ and $\theta_{-}$ are nonzero. Hence,
$$
\rho>0.
$$
By the Cauchy interlacing theorem,
$$
\lambda_{\max}\left(\mathbf{\Delta}_{\mathcal{F}}\right)
\leq
\delta^2
$$
and
$$
\lambda_{\min}\left(\mathbf{\Delta}_{\mathcal{F}}\right)
\geq
\rho^2.
$$
Therefore, $\mathbf{\Delta}_{\mathcal{F}}$ is positive definite, and
$$
\begin{aligned}
\kappa\left(\mathbf{\Delta}_{\mathcal{F}}\right)
&=
\frac{
\lambda_{\max}\left(\mathbf{\Delta}_{\mathcal{F}}\right)
}{
\lambda_{\min}\left(\mathbf{\Delta}_{\mathcal{F}}\right)
}
\
&\leq
\frac{\delta^2}{\rho^2}.
\end{aligned}
$$
The result follows from Claim \ref{claim:cond_num}.
\end{proof}

\remark{For a $(126, 25, 8,4)$ SRG with $\epsilon = 0$, we have $\kappa(\bfB_{\cF}^T\bfB_{\cF}) \le \frac{625}{9}  \approx 69.44 $.}





































% \subsection{Proof of Lemma \ref{lemma:random_diag}.}







% \subsection{Proof of Lemma \ref{lemma:second_cons}.}





% \begin{proof}[Proof of Claim 2]
%  Let $\mathbf{v}_1, \mathbf{v}_2, \dots, \mathbf{v}_m \in \mathbb{R}^n$ be vectors as generated by Algorithm 1. For any $i \in [m]$ and $u\in [k]$, let $\text{supp}(\mathbf{A}_i(u, :)) = \mathcal{H}$. Now, 

% \begin{align*}(\mathbf{A}_i\mathbf{v}_i)(u) = \mathbf{A}_i(u,:)\mathbf{v}_i &= \mathbf{A}_i(u,\mathcal{H})  \mathbf{v}_i({\mathcal{H}})\\
% &= \frac{\mathbf{v}_i(\mathcal{H})^T \mathbf{v}_i(\mathcal{H})}{\|\mathbf{v}_i(\mathcal{H})\|_2^2} = 1
% \end{align*}
% So, $\mathbf{A}_i\mathbf{v}_i = \mathbf{1}_k$ for all $i \in [m]$. 

% Let $j \in  [m] \setminus \{1\}$. Let $\mathbf{H} = \begin{bmatrix} \mathbf{A}_1 \\ \vdots \\ \mathbf{A}_{j-1}\end{bmatrix}$. By construction, $\mathbf{v}_j \in \text{Nul}(\mathbf{H})$. So, $\mathbf{H}\mathbf{v}_j = \mathbf{0}_ {k(j-1)}$ and consequently, for $i \in [m]$ and $i < j$, 
% $\mathbf{A}_i \mathbf{v}_j = \mathbf{0}_k$. Now let $i \in [m]$ and $i >j$. We have,

% \begin{align*}(\mathbf{A}_i\mathbf{v}_j)(u) = \mathbf{A}_i(u,:)\mathbf{v}_j &= \mathbf{A}_i(u,\mathcal{H})  \mathbf{v}_j({\mathcal{H}})\\
% &= \frac{\mathbf{v}_i(\mathcal{H})^T \mathbf{v}_j(\mathcal{H})}{\|\mathbf{v}_i(\mathcal{H})\|_2^2} 
% \end{align*}
% Now since $i > j$, $\mathbf{A}_j\mathbf{v}_i = \mathbf{0}_k$. 

% So, $(\mathbf{A}_j\mathbf{v}_i)(u) = 0 \implies \frac{\mathbf{v}_j(\mathcal{H})^T \mathbf{v}_i(\mathcal{H})}{\|\mathbf{v}_j(\mathcal{H})\|_2^2 } = 0$. Consequently, 
% $\mathbf{v}_j(\mathcal{H})^T \mathbf{v}_i(\mathcal{H}) = 0$ and $(\mathbf{A}_i\mathbf{v}_j)(u) = 0$. So, $\mathbf{A}_i\mathbf{v}_j = \mathbf{0}_k$ for $i > j$.
% So, we have that for any $i,j\in [m]$, $\mathbf{A}_i\mathbf{v}_j = \mathbf{1}_k$ for $ i = j$ and $\mathbf{A}_i\mathbf{v}_j = \mathbf{0}_k$ for $i \ne j$. Let $\mathbf{V} = \begin{bmatrix} \mathbf{v}_1 & \dots & \mathbf{v}_m \end{bmatrix}$.
% Now, 
% \begin{align*}
% \mathbf{B}\mathbf{V} &= \begin{bmatrix} \mathbf{A}_1 \\ \vdots \\ \mathbf{A}_{m}\end{bmatrix}\begin{bmatrix} \mathbf{v}_1 & \dots & \mathbf{v}_m \end{bmatrix}\\
% & = \begin{bmatrix}
% \mathbf{1}_{k} & \mathbf{0}_{k} & \cdots & \mathbf{0}_{k} \\
% \mathbf{0}_{k} & \mathbf{1}_{k} & \cdots & \mathbf{0}_{k} \\
% \vdots & \vdots & \ddots & \vdots \\
% \mathbf{0}_{k}  &\mathbf{0}_{k}  & \cdots & \mathbf{1}_{k}
% \end{bmatrix} = \mathbf{F}
% \end{align*}

% Since $\mathbf{B} \mathbf{V} = \mathbf{F}$, it follows that the approximation error is zero in case of no stragglers. Now we prove the upper bound. Let $\mathbf{Q} = \tilde{\mathbf{\Delta}}_{\mathcal{F}} - \mathbf{\Delta}_{\mathcal{F}}$. Now, for any $u \in [n-s]$, 
% $\mathbf{Q}(u,u) = \tilde{\mathbf{\Delta}}_{\mathcal{F}}(u,u) - \mathbf{\Delta}_{\mathcal{F}}(u,u) = \mathbf{\Delta}_{\mathcal{F}}(u,u) + \sum_{j \ne u} |\mathbf{\Delta}_{\mathcal{F}}(u,j)| - \mathbf{\Delta}_{\mathcal{F}}(u,u) =  \sum_{j \ne u} |\mathbf{\Delta}_{\mathcal{F}}(u,j)|$
% Also, for $j \ne u$ and $j \in [n-s]$, 
% $\mathbf{Q}(u,j) = -\mathbf{\Delta}_{\mathcal{F}}(u,j)$. So, $\sum_{j \ne u} |\mathbf{Q}(u,j)| =  \sum_{j \ne u} |\mathbf{\Delta}_{\mathcal{F}}(u,j)|$. Consequently, $\mathbf{Q}$ is a diagonally dominant matrix. Also, $\mathbf{Q}^T = (\tilde{\mathbf{\Delta}}_{\mathcal{F}} - \mathbf{\Delta}_{\mathcal{F}})^T = \tilde{\mathbf{\Delta}}_{\mathcal{F}} - \mathbf{\Delta}_{\mathcal{F}}$. So, $\mathbf{Q}$ is hermitian. So, all the eigenvalues of $\mathbf{Q}$ are real. Since $\mathbf{Q}$ is diagonally dominant, all the eigenavlues are non-negative. Since $\mathbf{Q}$ is hermitian and all the eigenvalues are non-negative, it is positive semi-definite. So, 
% $\mathbf{Q} \ge 0 \implies \tilde{\mathbf{\Delta}}_{\mathcal{F}} - \mathbf{\Delta}_{\mathcal{F}} \ge 0 \implies \tilde{\mathbf{\Delta}}_{\mathcal{F}} \ge \mathbf{\Delta}_{\mathcal{F}} \implies (\tilde{\mathbf{\Delta}}_{\mathcal{F}})^{-1} \le (\mathbf{\Delta}_{\mathcal{F}})^{-1}$.

% Finally,

% \begin{align*}
% \text{Err}_{\mathcal{F}}(\mathbf{B}) 
% % & = \min\limits_{\substack{\mathbf{R}\in \mathbb{R}^{n \times m} \\ supp(\mathbf{R}(:,i)) \subseteq \mathcal{F} \\ \forall i \in [m]}}  {\|\mathbf{B} \mathbf{R}- \mathbf{F}\|_F^2} \\
% % & = \sum_{i = 1}^m  \min\limits_{\substack{\mathbf{r}_i\in \mathbb{R}^{n} \\ supp(\mathbf{r}_i) \subseteq \mathcal{F} }}   \|\bfB\mathbf{r}_i - \mathbf{f}_i\|_2^2\\
% & = \sum_{i =1}^m \mathbf{f}_i^T \mathbf{f}_i - \mathbf{f}_i^T \bfB_{\cF}(\bfB_{\cF}^T\bfB_{\cF})^{-1}\bfB_{\cF}^T\mathbf{f}_i \\
% & = mk -  \sum_{i =1}^m \mathbf{1}_k^T {\mathbf{A}_i}_{\mathcal{F}}(\mathbf{\Delta}_{\mathcal{F}})^{-1}  {\mathbf{A}_i}_{\mathcal{F}}^T\mathbf{1}_k\\
% & \le mk -  \sum_{i =1}^m \mathbf{1}_k^T {\mathbf{A}_i}_{\mathcal{F}}(\tilde{\mathbf{\Delta}}_{\mathcal{F}})^{-1}  {\mathbf{A}_i}_{\mathcal{F}}^T\mathbf{1}_k
% \end{align*}
% \end{proof}



% \subsection{Proof of Corollary \ref{cor:SRG}.}



% \subsection{Proof of Claim \ref{claim:lowerbound}}



% \subsection{\texorpdfstring
% {Proof that ${\|\mathbf{B}_{\mathcal{C}} \mathbf{R}- \mathbf{F}\|_F^2}  = \sum_{i = 1}^k {\|\mathbf{B}^{(i)} \mathbf{R}- \mathbf{I}_m\|_F^2}$.}
% {Proof that ||B_C R - F||_F^2 = sum_{i=1}^k ||B^(i) R - I_m||_F^2.}}




\section{Proof of Proposition~\ref{prop:coset-full-support}}
\label{app:coset-full-support}

\begin{proof}
Since $\mathbf{A}$ is the bi-adjacency matrix of a $(k,m,\delta)$ coset bipartite graph, its columns can be ordered such that
\[
\mathbf{A}
=
\left[
\mathbf{C}\ \cdots\ \mathbf{C}
\right],
\]
where $\mathbf{C}\in\mathbb{R}^{k\times k}$. Choose pairwise orthogonal vectors
$\mathbf{q}_1,\ldots,\mathbf{q}_m\in\mathbb{R}^m$ such that
$\mathbf{q}_1=\mathbf{1}_m$ and every entry of each $\mathbf{q}_j$ is nonzero. Such a collection can be obtained, for example, from a suitably scaled Householder matrix \cite{10.1145/320941.320947}. For each $j\in[m]$, let $
\mathbf{v}_j
=
\mathbf{q}_j\otimes\mathbf{1}_k.$ Since every entry of $\mathbf{q}_j$ is nonzero, every entry of $\mathbf{v}_j$ is also nonzero. For each $\ell\in[m]$, let
$
\mathbf{D}_{\mathbf{v}_\ell}
=
\operatorname{diag}(\mathbf{v}_\ell),
$
and let $\mathbf{R}_\ell$ denote the diagonal row-normalization matrix whose $i$-th diagonal entry is
\[
\mathbf{R}_\ell(i,i)
=
\frac{1}
{\left\|\mathbf{v}_\ell(\mathcal{H}_i)\right\|_2^2}.
\]
The matrix constructed at the $\ell$-th iteration of Algorithm \ref{alg:example} can then be written as
\[
\mathbf{A}_\ell
=
\mathbf{R}_\ell
\mathbf{A}
\mathbf{D}_{\mathbf{v}_\ell}.
\]
It suffices to show that,
$
\mathbf{A}_\ell\mathbf{v}_j=\mathbf{0}_k$, for every $\ell < j$.
For any $\ell<j$, we have
\[
\begin{aligned}
\mathbf{A}_\ell\mathbf{v}_j
&=
\mathbf{R}_\ell
\mathbf{A}
\mathbf{D}_{\mathbf{v}_\ell}
\mathbf{v}_j\\
&=
\mathbf{R}_\ell
\mathbf{A}
\left(
\mathbf{v}_\ell\circ\mathbf{v}_j
\right).
\end{aligned}
\]
By the definition of $\mathbf{v}_\ell$ and $\mathbf{v}_j$,
\[
\mathbf{v}_\ell\circ\mathbf{v}_j
=
\begin{bmatrix}
\mathbf{q}_\ell(1)\mathbf{q}_j(1)\mathbf{1}_k\\
\vdots\\
\mathbf{q}_\ell(m)\mathbf{q}_j(m)\mathbf{1}_k
\end{bmatrix}.
\]
Therefore, 
\[
\begin{aligned}
\mathbf{A}
\left(
\mathbf{v}_\ell\circ\mathbf{v}_j
\right)
&=
\left[
\mathbf{C}\ \cdots\ \mathbf{C}
\right]
\begin{bmatrix}
\mathbf{q}_\ell(1)\mathbf{q}_j(1)\mathbf{1}_k\\
\vdots\\
\mathbf{q}_\ell(m)\mathbf{q}_j(m)\mathbf{1}_k
\end{bmatrix}\\
&=
\sum_{t=1}^{m}
\mathbf{q}_\ell(t)\mathbf{q}_j(t)\mathbf{C}\mathbf{1}_k\\
&=
\left(
\sum_{t=1}^{m}\mathbf{q}_\ell(t)\mathbf{q}_j(t)
\right)
\mathbf{C}\mathbf{1}_k\\
&=
\left(
\mathbf{q}_\ell^T\mathbf{q}_j
\right)
\mathbf{C}\mathbf{1}_k.
\end{aligned}
\]
Since $\mathbf{q}_\ell$ and $\mathbf{q}_j$ are orthogonal,
\[
\mathbf{q}_\ell^T\mathbf{q}_j=0.
\]
It follows that
\[
\mathbf{A}_\ell\mathbf{v}_j
=
\mathbf{0}_k
\qquad
\text{for every }\ell<j.
\]
Consequently,
\[
\mathbf{v}_j
\in
\bigcap_{\ell=1}^{j-1}
\text{Null}(\mathbf{A}_\ell)
=
\text{Null}\!\left(
\left[
\mathbf{A}_1^T\ \cdots\ \mathbf{A}_{j-1}^T
\right]^T
\right).
\]

Thus, there exist vectors
$\mathbf{v}_1,\ldots,\mathbf{v}_m$ satisfying the required null-space conditions, with every entry of each $\mathbf{v}_j$ nonzero.
\end{proof}

% \aditya{Note that if m=2, then there is only one such 2 by 2 orthogonal matrix, if you want all the vectors to have the same norm}

% In particular, after fixing
% $\mathbf{q}_1=\mathbf{1}_m$, the remaining vectors can be chosen as an orthogonal basis of
% $\mathbf{1}_m^{\perp}$. A generic choice of such an orthogonal basis has no zero entries, since the choices for which a specified entry is zero form a proper lower-dimensional subset, and there are only finitely many such entries.














\section{Proof that  ${\|\mathbf{B}_{\mathcal{C}} \mathbf{R}- \mathbf{F}\|_F^2}  = \sum_{i = 1}^k {\|\mathbf{B}^{(i)} \mathbf{R}- \mathbf{I}_m\|_F^2}$.}
\label{sec:appendix_claim_frob}
\begin{proof}
%[Proof that    ${\|\mathbf{B}_{\mathcal{C}} \mathbf{R}- \mathbf{F}\|_F^2}  = \sum_{i = 1}^k {\|\mathbf{B}^{(i)} \mathbf{R}- \mathbf{I}_m\|_F^2}$]

Let $\tilde{\mathbf{b}}_i^T$ and $\tilde{\mathbf{f}_i}^T$ be the $i^{th}$ row of $\mathbf{B}_C$ and $\mathbf{F}$ respectively, for $i \in [mk]$. Since the squared Frobenius norm of a matrix is the sum of the squared Frobenius norm of each row of that matrix, 

\begin{align*}
\|\mathbf{B}_{\mathcal{C}} \mathbf{R}- \mathbf{F}\|_F^2 &= \sum_{i = 1}^{mk} \|\tilde{\mathbf{b}}_i^T\mathbf{R}- \tilde{\mathbf{f}}_i^T\|_F^2\\
&= \sum_{i = 1}^k \sum_{j = 1}^m \|\tilde{\mathbf{b}}_{i+ k(j-1)}^T\mathbf{R}- \tilde{\mathbf{f}}_{i+ k(j-1)}^T\|_F^2.
\end{align*}

By definition, for $i \in [k]$, 
\begin{align*}
\mathbf{B}^{(i)}  = \begin{bmatrix}\tilde{\mathbf{b}}_{i}^T \\ \tilde{\mathbf{b}}_{i+ k}^T\\ \vdots \\ \tilde{\mathbf{b}}_{i+ k(m-1)}^T \end{bmatrix}.
\end{align*}

% \aditya{Indexing above looks wrong. Fix it}
% \sifat{done}
Also, note that, 
\begin{align*}
\begin{bmatrix}
\tilde{\mathbf{f}}_{i}^T \\ \tilde{\mathbf{f}}_{i+ k}^T\\ \vdots \\ \tilde{\mathbf{f}}_{i+ k(m-1)}^T \end{bmatrix} = \mathbf{I}_m.
\end{align*}

Consequently, 
\begin{align*}
\sum_{j = 1}^m \|\tilde{\mathbf{b}}_{i+ k(j-1)}^T\mathbf{R}- \tilde{\mathbf{f}}_{i+ k(j-1)}^T\|_F^2 = \|\mathbf{B}^{(i)}\mathbf{R} - \mathbf{I}_m\|_F^2.
\end{align*}
\end{proof}



\section{Hadamard Matrix Based Construction with Coset Bipartite Graph}
\label{app:hada_cos_unbiased}
\normalfont
In this section, we establish the unbiasedness property of the construction in \eqref{eq:hada_coset_cons}.

% Suppose that the assignment matrix $\mathbf{A} \in \mathbb{R}^{k \times n}$ is the bi-adjacency matrix of a $(k,m,\delta)$ coset bipartite graph. Then, $ \mathbf{A}=[\mathbf{C}\ \cdots\ \mathbf{C}],$ where $\mathbf{C}\in\mathbb{R}^{k\times k}$ is a circulant matrix. Suppose $\mathbf{H}$ is a Hadamard matrix of order $m$. Construct the encoding matrix $\mathbf{B} \in \mathbb{R}^{n \times n}$ as follows: 

% \begin{equation}
% \label{eq:Hada_Coset}
% \mathbf{B} = \mathbf{H}^T \otimes \mathbf{C}
% \end{equation}

\begin{lemma}
\label{lemma:Cost_Hadamard_Convergence}

Suppose that the assignment matrix $\mathbf{A}$ is the bi-adjacency matrix of a $(k, m, \delta)$ coset bipartite graph such that $k = p^a$, where $p$ is a prime, $a \in \mathbb{Z}_{\ge 1}$ and $p \nmid \delta$. Then, for the construction as in $\eqref{eq:hada_coset_cons}$, we have 

\[\mathbb{E}_{\mathcal{F}}[\mathbf{B}_{\cF} \mathbf{R}]= \alpha_H \mathbf{F},\]

where $\alpha_H$ is a nonzero constant.
\end{lemma}
\begin{proof}

Here, 
$\mathbf{B}^T \mathbf{B} = (\mathbf{H}_m^T \otimes \mathbf{C})^T  (\mathbf{H}_m^T \otimes \mathbf{C}) = \mathbf{H}_m\mathbf{H}_m^T \otimes \mathbf{C}^T \mathbf{C} = m \mathbf{I}_m \otimes \mathbf{C}^T \bfC$. For $\cF =\emptyset$, we assume $\bfB_{\cF} \bfR = \mathbf{0}_{mk \times m}$. For non-empty $\cF$,  let $\bfP_{\cF} = \bfI_n(:, \cF)$. We have, 
\begin{align*}
\mathbf{B}_{\mathcal{F}}\mathbf{R} &= \bfB_{\cF}(\bfB_{\cF}^T\bfB_{\cF})^{-1} \bfB_{\cF}^T \mathbf{F}\\
&=\mathbf{B} \mathbf{P}_{\mathcal{F}}(\mathbf{P}_\mathcal{F}^T \mathbf{B}^T\mathbf{B}\mathbf{P}_\mathcal{F})^{-1}\mathbf{P}_{\mathcal{F}}^T \bfB^T\mathbf{F}\\
\end{align*}

For each $i\in[m]$, let
\[
\mathcal{F}_i
:=
\left\{
j\in[k] : (i-1)k+j\in\mathcal{F}
\right\}.
\]
Thus, $\mathcal{F}_i$ denotes the indices of the non-straggling workers within the $i$-th block of $k$ workers. For non-empty $\cF_i$, let $\bfP'_{\cF_i} = \bfI_k(:, \cF_i)$. Thus, when $\cF_i \ne \emptyset$ for all $i \in [m]$,

\[
\mathbf{P}_{\mathcal{F}}
=
\begin{bmatrix}
\mathbf{P}'_{\mathcal{F}_1}
    & \mathbf{0}_{k\times |\mathcal{F}_2|}
    & \cdots
    & \mathbf{0}_{k\times |\mathcal{F}_m|}\\
\mathbf{0}_{k\times |\mathcal{F}_1|}
    & \mathbf{P}'_{\mathcal{F}_2}
    & \cdots
    & \mathbf{0}_{k\times |\mathcal{F}_m|}\\
\vdots
    & \vdots
    & \ddots
    & \vdots\\
\mathbf{0}_{k\times |\mathcal{F}_1|}
    & \mathbf{0}_{k\times |\mathcal{F}_2|}
    & \cdots
    & \mathbf{P}'_{\mathcal{F}_m}
\end{bmatrix}.
\]
% \sifat{Does it handle the case when $\cF_i$ is empty?}

Let $\mathbf{G}:=\mathbf{C}^T\mathbf{C}.$ Since $\bfB^T\bfB = m \bfI_m \otimes \bfC^T \bfC$, it follows that
\begin{align*}
\mathbf{P}_{\mathcal{F}}^T\mathbf{B}^T\mathbf{B}
\mathbf{P}_{\mathcal{F}}
& = m
\begin{bmatrix}
\mathbf{P}'^{T}_{\mathcal{F}_1} \mathbf{G}\mathbf{P}'_{\mathcal{F}_1}
& & \mathbf{0}\\
&
\ddots
&\\
\mathbf{0}
& &
\mathbf{P}'^{T}_{\mathcal{F}_m}\mathbf{G}\mathbf{P}'_{\mathcal{F}_m}
\end{bmatrix}
\end{align*}
 Therefore,
\[
\mathbf{P}_{\mathcal{F}}
(\mathbf{P}_{\mathcal{F}}^T\mathbf{B}^T\mathbf{B}
\mathbf{P}_{\mathcal{F}})^{-1}
\mathbf{P}_{\mathcal{F}}^T
=
\frac{1}{m}
\begin{bmatrix}
\mathbf{P}'_{\mathcal{F}_1}
\left(
\mathbf{P}'^{T}_{\mathcal{F}_1}
\mathbf{G}
\mathbf{P}'_{\mathcal{F}_1}
\right)^{-1}
\mathbf{P}'^{T}_{\mathcal{F}_1} & & \mathbf{0}\\
& \ddots &\\
\mathbf{0} & & \mathbf{P}'_{\mathcal{F}_m}
\left(
\mathbf{P}'^{T}_{\mathcal{F}_m}
\mathbf{G}
\mathbf{P}'_{\mathcal{F}_m}
\right)^{-1}
\mathbf{P}'^{T}_{\mathcal{F}_m}
\end{bmatrix}.
\]


For $i\in[m]$, let
% \[
% \mathbf{T}_{\mathcal{F}_i}
% :=
% \mathbf{P}_{\mathcal{F}_i}
% \left(
% \mathbf{P}_{\mathcal{F}_i}^T
% \mathbf{G}
% \mathbf{P}_{\mathcal{F}_i}
% \right)^{-1}
% \mathbf{P}_{\mathcal{F}_i}^T,
% \]
% where $\mathbf{T}_{\mathcal{F}_i}:=\mathbf{0}_{k\times k}$
% when $\mathcal{F}_i=\emptyset$. 



\[
\mathbf{T}_{\mathcal{F}_i} :=
\begin{cases}
\mathbf{P}'_{\mathcal{F}_i}
\left(
\mathbf{P}'^{T}_{\mathcal{F}_i}
\mathbf{G}
\mathbf{P}'_{\mathcal{F}_i}
\right)^{-1}
\mathbf{P}'^{T}_{\mathcal{F}_i}, & \mathcal{F}_i \neq \emptyset, \\[0.5em]
\mathbf{0}_{k \times k}, & \mathcal{F}_i = \emptyset.
\end{cases}
\]


% \sifat{say that this defn work for any non-empty and also empty F}

Since for $\cF_i = \emptyset$, $\mathbf{T}_{\cF_i} = \mathbf{0}_{k \times k}$, for any $\cF \subseteq [n]$, 

\[\mathbf{B}_{\cF} \bfR = \frac{1}{m}\mathbf{B}\begin{bmatrix}
\mathbf{T}_{\cF_1} & & \mathbf{0}\\
& \ddots &\\
\mathbf{0} & & \mathbf{T}_{\cF_m} 
\end{bmatrix} \bfB^T \bfF\]


Since each worker is independently non-straggling with probability
$1-q$, the random sets $\mathcal{F}_1,\ldots,\mathcal{F}_m$ are
independent. Let
$
\overline{\mathbf{T}}
:=
\mathbb{E}_{\mathcal{F}}
[\mathbf{T}_{\mathcal{F}_i}]
.$ Thus,
\[
\mathbb{E}_{\mathcal{F}}
\left[
\bfB_{\cF}\bfR
\right]
=
\frac{1}{m} \bfB (\mathbf{I}_m\otimes\overline{\mathbf{T}}) \mathbf{B}^T \bfF.
\]

Next, we show that $\overline{\mathbf{T}}$ is circulant. Let
$\mathbf{Q}\in\{0,1\}^{k\times k}$ be the permutation matrix
corresponding to a cyclic shift by one position. Since $\mathbf{C}$
is circulant,
$
\mathbf{Q}^T\mathbf{C}\mathbf{Q}=\mathbf{Q}\mathbf{C}\mathbf{Q}^T =\mathbf{C}.
$ From this we can infer that $\mathbf{Q}^T\mathbf{G}\mathbf{Q} = \mathbf{Q}\mathbf{G}\mathbf{Q}^T
= \mathbf{G}$.


For $\mathcal{F}_i\subseteq[k]$, let $\mathcal{F}_i'$
denote the set obtained by cyclically shifting the elements of
$\mathcal{F}_i$ by one position. Then,
$\mathbf{Q}\mathbf{P}'_{\mathcal{F}_i}$ is an equivalent selection
matrix for $\mathcal{F}_i'$. 
% Hence,
% \[
% \mathbf{T}_{\mathcal{S}'}
% =
% \mathbf{Q}\mathbf{T}_{\mathcal{S}}\mathbf{Q}^T.
% \]
% Since the workers straggle independently, 
% \[
% \mathcal{S}'\overset{d}{=}\mathcal{S}.
% \]
Now,
\begin{align*}
\mathbf{Q} \bfT_{\mathcal{F}_i} \mathbf{Q}^T &= \mathbf{Q} \mathbf{P}'_{\mathcal{F}_i}
\left(
\mathbf{P}'^{T}_{\mathcal{F}_i}
\mathbf{G}
\mathbf{P}'_{\mathcal{F}_i}
\right)^{-1}
\mathbf{P}'^{T}_{\mathcal{F}_i} \mathbf{Q}^T\\
&= \mathbf{Q} \mathbf{P}'_{\mathcal{F}_i}
\left(
\mathbf{P}'^{T}_{\mathcal{F}_i}\mathbf{Q}^T\mathbf{Q}
\mathbf{G}\mathbf{Q}^T\mathbf{Q}
\mathbf{P}'_{\mathcal{F}_i}
\right)^{-1}
\mathbf{P}'^{T}_{\mathcal{F}_i} \mathbf{Q}^T\\
&= \mathbf{Q} \mathbf{P}'_{\mathcal{F}_i}
\left(
\mathbf{P}'^{T}_{\mathcal{F}_i}\mathbf{Q}^T
\mathbf{G}\mathbf{Q}
\mathbf{P}'_{\mathcal{F}_i}
\right)^{-1}
\mathbf{P}'^{T}_{\mathcal{F}_i} \mathbf{Q}^T\\
& = \bfT_{\cF_i'}.
\end{align*}
% \begin{align*}

% \end{align*}
% \[
% \mathbf{Q}\mathbf{T}_{\mathcal{S}}\mathbf{Q}^T
% \overset{d}{=}
% \mathbf{T}_{\mathcal{S}}.
% \]

Since $\cF_i  \overset{d}{=} \cF_i'$, we have $\bfT_{\cF_i} \overset{d} {=} \bfT_{\cF_i'}$.   Thus, $\mathbb{E}_{\cF}\left[\bfT_{\cF_i}\right] = \mathbb{E}_{\cF}\left[\bfT_{\cF_i'}\right]$. Consequently, $
\mathbf{Q}\overline{\mathbf{T}}\mathbf{Q}^T
=
\overline{\mathbf{T}}.
$
So, by noting that $\mathbf{Q}$ corresponds to a cyclic shift by one position, we can conclude that $\overline{\mathbf{T}}$ is circulant. So,
\begin{align*}
\mathbb{E}_{\mathcal{F}}
[\mathbf{B}_{\mathcal{F}}\mathbf{R}]
&=
\frac{1}{m} \bfB (\mathbf{I}_m\otimes\overline{\mathbf{T}}) \mathbf{B}^T \bfF\\
&=
\frac{1}{m}(\mathbf{H}_m^T\otimes\mathbf{C})
\left(\mathbf{I}_m\otimes\overline{\mathbf{T}}
\right)
(\mathbf{H}_m\otimes\mathbf{C}^T)(\bfI_m \otimes \mathbf{1}_k)\\
&=\frac{1}{m} \bfH_m^T\bfH_m \otimes \bfC \overline{\mathbf{T}} \bfC^T\mathbf{1}_k = \bfI_m \otimes  \bfC \overline{\mathbf{T}} \bfC^T\mathbf{1}_k .
\end{align*}

Let
$
\tilde{\mathbf{T}}:=
\mathbf{C}\overline{\mathbf{T}}\mathbf{C}^T.$
Since $\mathbf{C}$ and $\overline{\mathbf{T}}$ are circulant,
$\tilde{\mathbf{T}}$ is also circulant. Therefore, $\mathbf{1}_k$ is an
eigenvector of $\tilde{\mathbf{T}}$, and there exists a constant $\alpha_H$
such that $\tilde{\mathbf{T}}\mathbf{1}_k=\alpha_H\mathbf{1}_k.$


It remains to show that $\alpha_H$ is nonzero. Since $\mathbf{C}$ is
invertible, $\mathbf{G}=\mathbf{C}^T\mathbf{C}$ is positive definite.
Thus, for any non-empty $\mathcal{F}_i$,
$
\mathbf{1}_k^T\mathbf{T}_{\mathcal{F}_i}\mathbf{1}_k
=
\mathbf{1}_{|\mathcal{F}_i|}^T
\left(
\mathbf{P}'^{T}_{\mathcal{F}_i}
\mathbf{G}
\mathbf{P}'_{\mathcal{F}_i}
\right)^{-1}
\mathbf{1}_{|\mathcal{F}_i|}
>0.
$
Since $q<1$, $\mathcal{F}_i$ is non-empty with positive probability.
Therefore,
$
\mathbf{1}_k^T\overline{\mathbf{T}}\mathbf{1}_k>0
.$ Moreover, since $\mathbf{C}$ is $\delta$-regular,
$\mathbf{C}\mathbf{1}_k=\mathbf{C}^T\mathbf{1}_k
=\delta\mathbf{1}_k$. Hence,
$
\alpha_H k
=
\mathbf{1}_k^T\widetilde{\mathbf{T}}\mathbf{1}_k
=
\delta^2
\mathbf{1}_k^T\overline{\mathbf{T}}\mathbf{1}_k
>0.
$ Consequently, $\alpha_H$ is nonzero.






% Recall that
% \[
% \mathbf{f}_i
% =
% \mathbf{e}_i\otimes\mathbf{1}_k,
% \qquad i\in[m].
% \]
% Thus,
% \[
% (\mathbf{I}_m\otimes\mathbf{M})\mathbf{f}_i
% =
% \mathbf{e}_i\otimes
% \mathbf{M}\mathbf{1}_k
% =
% \alpha\mathbf{f}_i.
% \]
% Since
% $\mathbf{F}=[\mathbf{f}_1|\cdots|\mathbf{f}_m]$, it follows that
% \[
% \mathbb{E}_{\mathcal{F}}
% [\mathbf{B}_{\mathcal{F}}\mathbf{R}]
% =
% \alpha\mathbf{F}.
% \]

% Finally, $\alpha$ is non-zero. Since $\mathbf{C}$ is invertible,
% $\mathbf{G}=\mathbf{C}^T\mathbf{C}$ is positive definite, and thus,
% for every non-empty $\mathcal{S}$,
% \[
% \mathbf{1}_k^T\mathbf{T}_{\mathcal{S}}\mathbf{1}_k
% =
% \mathbf{1}_{|\mathcal{S}|}^T
% \left(
% \mathbf{P}_{\mathcal{S}}^T
% \mathbf{G}
% \mathbf{P}_{\mathcal{S}}
% \right)^{-1}
% \mathbf{1}_{|\mathcal{S}|}
% >0.
% \]
% Since $q<1$, a non-empty set of non-stragglers occurs with positive
% probability. Hence,
% \[
% \mathbf{1}_k^T\overline{\mathbf{T}}\mathbf{1}_k>0.
% \]
% Also, $\mathbf{C}\mathbf{1}_k=\delta\mathbf{1}_k$ and
% $\mathbf{C}^T\mathbf{1}_k=\delta\mathbf{1}_k$. Consequently,
% \[
% \mathbf{1}_k^T\mathbf{M}\mathbf{1}_k
% =
% \delta^2
% \mathbf{1}_k^T\overline{\mathbf{T}}\mathbf{1}_k
% >0,
% \]
% which implies $\alpha>0$. Therefore,
% \[
% \mathbb{E}_{\mathcal{F}}
% [\mathbf{B}_{\mathcal{F}}\mathbf{R}]
% =
% \alpha\mathbf{F},
% \]
% where $\alpha$ is a non-zero constant.


















\end{proof}








\end{document}


%%%%%%
%% Some comments about useful packages
%% (extract from bare_conf.tex by Michael Shell)
%%

% *** MISC UTILITY PACKAGES ***
%
%\usepackage{ifpdf}
% Heiko Oberdiek's ifpdf.sty is very useful if you need conditional
% compilation based on whether the output is pdf or dvi.
% usage:
% \ifpdf
%   % pdf code
% \else
%   % dvi code
% \fi
% The latest version of ifpdf.sty can be obtained from:
% http://www.ctan.org/pkg/ifpdf
% Also, note that IEEEtran.cls V1.7 and later provides a builtin
% \ifCLASSINFOpdf conditional that works the same way.
% When switching from latex to pdflatex and vice-versa, the compiler may
% have to be run twice to clear warning/error messages.


% *** CITATION PACKAGES ***
%
%\usepackage{cite}
% cite.sty was written by Donald Arseneau
% V1.6 and later of IEEEtran pre-defines the format of the cite.sty package
% \cite{} output to follow that of the IEEE. Loading the cite package will
% result in citation numbers being automatically sorted and properly
% "compressed/ranged". e.g., [1], [9], [2], [7], [5], [6] without using
% cite.sty will become [1], [2], [5]--[7], [9] using cite.sty. cite.sty's
% \cite will automatically add leading space, if needed. Use cite.sty's
% noadjust option (cite.sty V3.8 and later) if you want to turn this off
% such as if a citation ever needs to be enclosed in parenthesis.
% cite.sty is already installed on most LaTeX systems. Be sure and use
% version 5.0 (2009-03-20) and later if using hyperref.sty.
% The latest version can be obtained at:
% http://www.ctan.org/pkg/cite
% The documentation is contained in the cite.sty file itself.


% *** GRAPHICS RELATED PACKAGES ***
%
\ifCLASSINFOpdf
  % \usepackage[pdftex]{graphicx}
  % declare the path(s) where your graphic files are
  % \graphicspath{{../pdf/}{../jpeg/}}
  % and their extensions so you won't have to specify these with
  % every instance of \includegraphics
  % \DeclareGraphicsExtensions{.pdf,.jpeg,.png}
\else
  % or other class option (dvipsone, dvipdf, if not using dvips). graphicx
  % will default to the driver specified in the system graphics.cfg if no
  % driver is specified.
  % \usepackage[dvips]{graphicx}
  % declare the path(s) where your graphic files are
  % \graphicspath{{../eps/}}
  % and their extensions so you won't have to specify these with
  % every instance of \includegraphics
  % \DeclareGraphicsExtensions{.eps}
\fi
% graphicx was written by David Carlisle and Sebastian Rahtz. It is
% required if you want graphics, photos, etc. graphicx.sty is already
% installed on most LaTeX systems. The latest version and documentation
% can be obtained at: 
% http://www.ctan.org/pkg/graphicx
% Another good source of documentation is "Using Imported Graphics in
% LaTeX2e" by Keith Reckdahl which can be found at:
% http://www.ctan.org/pkg/epslatex
%
% latex, and pdflatex in dvi mode, support graphics in encapsulated
% postscript (.eps) format. pdflatex in pdf mode supports graphics
% in .pdf, .jpeg, .png and .mps (metapost) formats. Users should ensure
% that all non-photo figures use a vector format (.eps, .pdf, .mps) and
% not a bitmapped formats (.jpeg, .png). The IEEE frowns on bitmapped formats
% which can result in "jaggedy"/blurry rendering of lines and letters as
% well as large increases in file sizes.
%
% You can find documentation about the pdfTeX application at:
% http://www.tug.org/applications/pdftex


% *** MATH PACKAGES ***
%
%\usepackage{amsmath}
% A popular package from the American Mathematical Society that provides
% many useful and powerful commands for dealing with mathematics.
%
% Note that the amsmath package sets \interdisplaylinepenalty to 10000
% thus preventing page breaks from occurring within multiline equations. Use:
%\interdisplaylinepenalty=2500
% after loading amsmath to restore such page breaks as IEEEtran.cls normally
% does. amsmath.sty is already installed on most LaTeX systems. The latest
% version and documentation can be obtained at:
% http://www.ctan.org/pkg/amsmath


% *** SPECIALIZED LIST PACKAGES ***
%
%\usepackage{algorithmic}
% algorithmic.sty was written by Peter Williams and Rogerio Brito.
% This package provides an algorithmic environment for describing algorithms.
% You can use the algorithmic environment in-text or within a figure
% environment to provide for a floating algorithm. Do NOT use the algorithm
% floating environment provided by algorithm.sty (by the same authors) or
% algorithm2e.sty (by Christophe Fiorio) as the IEEE does not use dedicated
% algorithm float types and packages that provide these will not provide
% correct IEEE style captions. The latest version and documentation of
% algorithmic.sty can be obtained at:
% http://www.ctan.org/pkg/algorithms
% Also of interest may be the (relatively newer and more customizable)
% algorithmicx.sty package by Szasz Janos:
% http://www.ctan.org/pkg/algorithmicx


% *** ALIGNMENT PACKAGES ***
%
%\usepackage{array}
% Frank Mittelbach's and David Carlisle's array.sty patches and improves
% the standard LaTeX2e array and tabular environments to provide better
% appearance and additional user controls. As the default LaTeX2e table
% generation code is lacking to the point of almost being broken with
% respect to the quality of the end results, all users are strongly
% advised to use an enhanced (at the very least that provided by array.sty)
% set of table tools. array.sty is already installed on most systems. The
% latest version and documentation can be obtained at:
% http://www.ctan.org/pkg/array

% IEEEtran contains the IEEEeqnarray family of commands that can be used to
% generate multiline equations as well as matrices, tables, etc., of high
% quality.


% *** SUBFIGURE PACKAGES ***
%\ifCLASSOPTIONcompsoc
%  \usepackage[caption=false,font=normalsize,labelfont=sf,textfont=sf]{subfig}
%\else
%  \usepackage[caption=false,font=footnotesize]{subfig}
%\fi
% subfig.sty, written by Steven Douglas Cochran, is the modern replacement
% for subfigure.sty, the latter of which is no longer maintained and is
% incompatible with some LaTeX packages including fixltx2e. However,
% subfig.sty requires and automatically loads Axel Sommerfeldt's caption.sty
% which will override IEEEtran.cls' handling of captions and this will result
% in non-IEEE style figure/table captions. To prevent this problem, be sure
% and invoke subfig.sty's "caption=false" package option (available since
% subfig.sty version 1.3, 2005/06/28) as this is will preserve IEEEtran.cls
% handling of captions.
% Note that the Computer Society format requires a larger sans serif font
% than the serif footnote size font used in traditional IEEE formatting
% and thus the need to invoke different subfig.sty package options depending
% on whether compsoc mode has been enabled.
%
% The latest version and documentation of subfig.sty can be obtained at:
% http://www.ctan.org/pkg/subfig


% *** FLOAT PACKAGES ***
%
%\usepackage{fixltx2e}
% fixltx2e, the successor to the earlier fix2col.sty, was written by
% Frank Mittelbach and David Carlisle. This package corrects a few problems
% in the LaTeX2e kernel, the most notable of which is that in current
% LaTeX2e releases, the ordering of single and double column floats is not
% guaranteed to be preserved. Thus, an unpatched LaTeX2e can allow a
% single column figure to be placed prior to an earlier double column
% figure.
% Be aware that LaTeX2e kernels dated 2015 and later have fixltx2e.sty's
% corrections already built into the system in which case a warning will
% be issued if an attempt is made to load fixltx2e.sty as it is no longer
% needed.
% The latest version and documentation can be found at:
% http://www.ctan.org/pkg/fixltx2e


%\usepackage{stfloats}
% stfloats.sty was written by Sigitas Tolusis. This package gives LaTeX2e
% the ability to do double column floats at the bottom of the page as well
% as the top. (e.g., "\begin{figure*}[!b]" is not normally possible in
% LaTeX2e). It also provides a command:
%\fnbelowfloat
% to enable the placement of footnotes below bottom floats (the standard
% LaTeX2e kernel puts them above bottom floats). This is an invasive package
% which rewrites many portions of the LaTeX2e float routines. It may not work
% with other packages that modify the LaTeX2e float routines. The latest
% version and documentation can be obtained at:
% http://www.ctan.org/pkg/stfloats
% Do not use the stfloats baselinefloat ability as the IEEE does not allow
% \baselineskip to stretch. Authors submitting work to the IEEE should note
% that the IEEE rarely uses double column equations and that authors should try
% to avoid such use. Do not be tempted to use the cuted.sty or midfloat.sty
% packages (also by Sigitas Tolusis) as the IEEE does not format its papers in
% such ways.
% Do not attempt to use stfloats with fixltx2e as they are incompatible.
% Instead, use Morten Hogholm's dblfloatfix which combines the features
% of both fixltx2e and stfloats:
%
% \usepackage{dblfloatfix}
% The latest version can be found at:
% http://www.ctan.org/pkg/dblfloatfix


% *** PDF and URL PACKAGES ***
%
%\usepackage{url}
% url.sty was written by Donald Arseneau. It provides better support for
% handling and breaking URLs. url.sty is already installed on most LaTeX
% systems. The latest version and documentation can be obtained at:
% http://www.ctan.org/pkg/url
% Basically, \url{my_url_here}.



% *** Do not adjust lengths that control margins, column widths, etc. ***
% *** Do not use packages that alter fonts (such as pslatex).         ***
%%%%%%


%%% Local Variables:
%%% mode: latex
%%% TeX-master: t
%%% End:
